\section{安全见证与双重值初始化}\label{construct:safety}
\subsection*{\texorpdfstring{{C1.2} {算术条件与安全性的基本事实}}{C1.2 算术条件与安全性的基本事实}}\label{C1-h1}
\addcontentsline{toc}{subsection}{{C1.2} {算术条件与安全性的基本事实}}

{所有平均都与整体平移、整体非零有理缩放相容。} {对非零本原整数零和状态，\ensuremath{G}} {总整除} {\ensuremath{n}：\ensuremath{G}} {整除} {\ensuremath{n*x_{1}}，而} {\ensuremath{\operatorname{gcd}(G,x_{1})=1}。} {若} {\ensuremath{q\ne p}} {是} {\ensuremath{G}} {的素因子，所有初值模} {\ensuremath{q}} {都等于一个非零余数；\ensuremath{Z}{[}1/p{]}} {中的平均保持这个余数，所以不可能归零。} {这证明必要性。}

{称当前整数零和配置} {\ensuremath{E}} \textbf{{合法}}{，如果对每个} {\ensuremath{q}} {属于} {\ensuremath{Q}，其坐标模} {\ensuremath{q}} {不全同余。} {对本原状态，这恰好等价于} {\ensuremath{G}} {为} {\ensuremath{p}} {幂。构造开始以后不做本原约分。} {称一次整数} {\ensuremath{p}} {平均}\textbf{{安全}}{，如果输出仍合法。}

\paragraph*{\texorpdfstring{{引理} {1：危险支持集}}{引理 1：危险支持集}}\label{C1-h2}
\addcontentsline{toc}{paragraph}{{引理} {1：危险支持集}}

{假设} {\ensuremath{n>2p}。若某个模} {\ensuremath{q}} {余类至少包含} {n-p} {个位置，称它为主类，其补集为} {\ensuremath{F_{q}}。} {因为} {\ensuremath{n-p>n/2}，主类唯一，而且}

\[
2\le |F_q|\le p.
\]

{补集非空来自合法性。若补集只有一个位置，记其相对主类余数差为} {\ensuremath{t}；总和模} {\ensuremath{q}} {为} {\ensuremath{n*c+t=t}，迫使} {\ensuremath{t=0}，矛盾。}

{若不存在主类，删除任意} {\ensuremath{p}} {个位置后仍模} {\ensuremath{q}} {不全同余；若存在主类，只要所选位置不包含整个} {\ensuremath{F_{q}}，也保留不同余的未选位置。删除} {\ensuremath{p}} {个位置后主类至少剩} {\ensuremath{n-2p>0}} {个位置。}

{因此，不包含任何完整} {\ensuremath{F_{q}}} {的整数} {\ensuremath{p}} {平均一定安全。}

\paragraph*{\texorpdfstring{{引理} {2：保护池}}{引理 2：保护池}}\label{C1-h3}
\addcontentsline{toc}{paragraph}{{引理} {2：保护池}}

{设一个位置池中的所有值同余模} {\ensuremath{p}，每个实际值在池中最多出现} {p-1} {次。}

\begin{itemize}
\tightlist
\item
  {池大小至少} {d+p} {时，可以选取安全、非恒等的整数} {\ensuremath{p}} {平均。}
\item
  {池大小至少} {d+p+1} {时，对任意指定值} {\ensuremath{a}，还可以要求新均值不等于} {\ensuremath{a}。}
\end{itemize}

\textbf{{证明。}} {逐个处理危险} {\ensuremath{q}：若当前可用池仍包含整个} {\ensuremath{F_{q}}，留出其中一个位置。至多留出} {\ensuremath{d}} {个位置。此后任取} {\ensuremath{p}} {个位置都安全，且因重数小于} {\ensuremath{p}} {而非恒等。} {第二种情形，保护后任取} {p+1} {个位置；如果删去任意一个位置后，剩下} {\ensuremath{p}} {个数的均值总是} {\ensuremath{a}，比较这些和便知} {p+1} {个值全部相等，与重数上界矛盾。证毕。}

\paragraph*{\texorpdfstring{{引理} {3：重复见证}}{引理 3：重复见证}}\label{C1-h4}
\addcontentsline{toc}{paragraph}{{引理} {3：重复见证}}

{以下整数操作安全：}

\begin{enumerate}
\def\labelenumi{\arabic{enumi}.}
\tightlist
\item
  {平均} {(\ensuremath{a^{p-1},x)}，且原来至少有} {\ensuremath{p}} {份} {\ensuremath{a}。}
\item
  {平均由} {\ensuremath{a}、\ensuremath{b}} {的若干份及一个第三值} {\ensuremath{x}} {组成的} {\ensuremath{p}} {元组，且操作后仍各留有至少一份} {\ensuremath{a}、\ensuremath{b}。}
\end{enumerate}

\textbf{{证明。}} {若输出对某个} {\ensuremath{q}} {属于} {\ensuremath{Q}} {全同余为} {\ensuremath{c}，第一种中未选的} {\ensuremath{a}} {迫使} {\ensuremath{a=c}} {\ensuremath{mod}} {\ensuremath{q}，新均值再迫使} {\ensuremath{x=c}。} {第二种中未选的} {\ensuremath{a}、\ensuremath{b}} {都为} {\ensuremath{c}，新均值再迫使唯一的第三值} {\ensuremath{x}} {也为} {\ensuremath{c}。} {所有原输入于是全同余，矛盾。此证明对} {\ensuremath{q=2}} {同样有效，没有除以} {p-1} {或其他可能不可逆的系数。证毕。}

\subsection*{\texorpdfstring{{C1.3} {不变量及初始化}}{C1.3 不变量及初始化}}\label{C1-h5}
\addcontentsline{toc}{subsection}{{C1.3} {不变量及初始化}}

{设} {\ensuremath{f_{E}(a)}} {为实际值} {\ensuremath{a}} {的重数。称} {\ensuremath{f_{E}(a)\ge p}} {的值为} \textbf{{\ensuremath{p}} {重值}}{，令}

\[
H(E)=\{a:f_E(a)\ge p\}.
\]

{保持不变量}

\[
E\in\mathbb Z^n,\qquad \sum E=0,\qquad E\text{ 合法},\qquad |H(E)|\ge2.
\tag{C1.2}
\]

\begin{enumerate}
\def\labelenumi{(\arabic{enumi})}
\tightlist
\item
  {蕴含} {\ensuremath{n\ge d+3p}} {和} {\ensuremath{n>2p}。} {从合法非零输入的} {pX} {出发，所有初值均被} {\ensuremath{p}} {整除。至多两次安全整数操作可建立} {(2)：}
\end{enumerate}

\begin{itemize}
\tightlist
\item
  {已有两个} {\ensuremath{p}} {重值时无需操作。}
\item
  {恰有一个} {\ensuremath{p}} {重值} {\ensuremath{a}，且} {\ensuremath{f(a)\ge 2p-1}} {时，取不同值} {\ensuremath{b}，平均} {(\ensuremath{a^{p-1},b)}。留下至少} {\ensuremath{p}} {份} {\ensuremath{a}，并产生} {\ensuremath{p}} {份不同于} {\ensuremath{a}} {的新均值。}
\item
  {恰有一个} {\ensuremath{p}} {重值} {\ensuremath{a}，且} {\ensuremath{f(a)\le 2p-2}} {时，\ensuremath{a}} {之外至少有} {\ensuremath{n-2p+2\ge d+p+2}} {个位置，每个值重数至多} {p-1。保护池引理给出均值不等于} {\ensuremath{a}} {的安全操作。}
\item
  {没有} {\ensuremath{p}} {重值时，先用整个初态的保护池制造一个。新均值} {\ensuremath{a}} {是唯一可能的} {\ensuremath{p}} {重值，重数至多} {2p-1。\ensuremath{a}} {之外至少有} {\ensuremath{n-2p+1\ge d+p+1}} {个位置，都是尚未改变的初值，仍被} {\ensuremath{p}} {整除，且每值最多} {p-1} {份。再用保护池制造不同于} {\ensuremath{a}} {的新均值。}
\end{itemize}

{第二次操作不必接触第一次新生的均值，所以全程保持整数。}


\section{零和选择与背景值的整数下降}\label{construct:packing}
\subsection*{\texorpdfstring{{C2.3} {用} {EGZ} {定理消除模} {\ensuremath{p}} {分池损失}}{C2.3 用 EGZ 定理消除模  分池损失}}\label{C2-h6}
\addcontentsline{toc}{subsection}{{C2.3} {用} {EGZ} {定理消除模} {\ensuremath{p}} {分池损失}}

\textbf{{零和引理（Erdos--Ginzburg--Ziv} {的素数情形）。}} {任意} {2p-1} {个整数中，可以选择恰好} {\ensuremath{p}} {个位置，使其和被} {\ensuremath{p}} {整除。}

{一个简短证明使用} {Chevalley--Warning} {定理：在} {\ensuremath{F_{p}}} {的} {2p-1} {个变量上考虑}

\[
\sum_i z_i^{p-1}=0,\qquad \sum_i a_i z_i^{p-1}=0.
\]

{两个方程的总次数} {\ensuremath{2(p-1)<2p-1}，所以共同零点数被} {\ensuremath{p}} {整除。} {非零解的支持大小只能是} {\ensuremath{p}；设和为零的} {\ensuremath{p}} {元位置子集数为} {\ensuremath{M}，则零点数为} {\ensuremath{1+(p-1)^{p}}} {\ensuremath{M}，模} {\ensuremath{p}} {等于} {1-M（\ensuremath{p=2}} {时相同）。因此} {\ensuremath{M=1}} {\ensuremath{mod}} {\ensuremath{p}，特别是} {\ensuremath{M>0}。}

{这是一条标准有限域定理的直接应用，不是根据有限素数样本外推。}

{记所有实际值重数小于} {\ensuremath{p}} {的位置数为} {\ensuremath{r}。如果}

\[
r\ge d+2p-1,
\tag{C2.4}
\]

{先保护至多} {\ensuremath{d}} {个位置，再在其余低重数位置中取} {2p-1} {个，应用零和引理。} {得到安全整数} {\ensuremath{p}} {元组。由于任何实际值都不足} {\ensuremath{p}} {份，它不可能恒等。} {原有全部} {\ensuremath{p}} {重值未参与，\ensuremath{I}} {保持。}

\subsection*{\texorpdfstring{{C2.4} {四个重值已经足够继续}}{C2.4 四个重值已经足够继续}}\label{C2-h7}
\addcontentsline{toc}{subsection}{{C2.4} {四个重值已经足够继续}}

\textbf{{引理。}} {如果有至少四个不同的} {\ensuremath{p}} {重值，则存在保持} {\ensuremath{I}} {的安全非恒等整数操作。}

\textbf{{证明。}} {按实际大小选择四个不同重值}

\[
a_1<a_2<a_3<b.
\]

{若前三个值模} {\ensuremath{p}} {全相同，平均} {(\ensuremath{a_{1}^{p-1},a_{2})}，其平均为整数且安全，\ensuremath{a_{3}}、\ensuremath{b}} {都保留为} {\ensuremath{p}} {重值。}

{否则，从前三个中选模} {\ensuremath{p}} {不同的} {a,c，第三个记为} {\ensuremath{x}。} {选择唯一} {\ensuremath{0\le i<p}，使}

\[
i(a-c)+x-c=0\pmod p.
\]

{平均} {(\ensuremath{a^{i},c^{p-1-i},x)}。所选值都来自前三个，每个被选重值还留下副本，所以安全；三元类型的总位置数恰为} {\ensuremath{p}，且操作非恒等。} {新均值严格小于} {\ensuremath{b}，因此它与完整保留的} {\ensuremath{b}} {构成两个} {\ensuremath{p}} {重值。证毕。}

{注意这里利用实际数值的顺序，使新均值自动避开未参与的极端值；不需要对所有重值余类分别分类。}

\subsection*{\texorpdfstring{{C2.5} {线性容量提供背景值}}{C2.5 线性容量提供背景值}}\label{C2-h8}
\addcontentsline{toc}{subsection}{{C2.5} {线性容量提供背景值}}

{现在假设已经存在一个背景值} \(b\){，其重数至少为} \(2p-1\){；另取一个不同的} \(p\) {重值} \(a\){。}

{若} {\ensuremath{a=b}} {\ensuremath{mod}} {\ensuremath{p}，用} {(\ensuremath{b^{p-1},a)}} {继续。} {否则令} {\ensuremath{\delta =a-b}，它在模} {\ensuremath{p}} {下可逆。} {对每个第三值} {\ensuremath{x}，唯一选择} {\ensuremath{0\le i<p}，使}

\[
i\delta+x-b=0\pmod p.
\]

{平均} {(\ensuremath{a^{i},b^{p-1-i},x)}。它是原位置安全整数操作，且操作后} {\ensuremath{b}} {仍至少有} {\ensuremath{p}} {份。} {若新均值不同于} {\ensuremath{b}，\ensuremath{I}} {保持；否则} {\ensuremath{x=b-i*\delta }。}

{因此，在尚未找到保持操作时，全部实际值被迫属于}

\[
\{b+\delta,b,b-\delta,\ldots,b-(p-1)\delta\}.
\tag{C2.5}
\]

{只有} {a,b} {两种实际值不可能：合法性与} {\ensuremath{p}} {不整除} {\ensuremath{\delta }} {给} {\ensuremath{\operatorname{gcd}(\delta ,n)=1}，零和会迫使} {\ensuremath{a}} {的重数被} {\ensuremath{n}} {整除。}

{如果任何非} {\ensuremath{b}} {值至少有} {2p-1} {份，将它当作第二个储备重值，再调节另一个值，就能保持两个旧} {\ensuremath{p}} {重值。} {对于} {b-j*delta，\ensuremath{1\le j<p}，模} {\ensuremath{p}} {的可逆性仍成立。因此剩余情形中每个非} {\ensuremath{b}} {值至多有} {2p-2} {份。}

\subsection*{\texorpdfstring{{C2.6} {系数跨越与背景值终端}}{C2.6 系数跨越与背景值终端}}\label{C2-h9}
\addcontentsline{toc}{subsection}{{C2.6} {系数跨越与背景值终端}}

{写} {(5)} {中的状态为}

\[
((b+\delta)^A,b^B,(b-\delta)^{c_1},\ldots,(b-(p-1)\delta)^{c_{p-1}}),
\]

{其中} {\ensuremath{p\le A\le 2p-2}，\ensuremath{B\ge 2p-1}，\ensuremath{0\le c_{j}\le 2p-2}。} {令} {\ensuremath{S=sum}} {\ensuremath{j*c_{j}}。}

\paragraph*{\texorpdfstring{{6.1} {如果} {\ensuremath{S\ge p}，且某个} {\ensuremath{j\ge 2}} {的例外出现}}{6.1 如果 ，且某个  的例外出现}}\label{C2-h10}
\addcontentsline{toc}{paragraph}{{6.1} {如果} {\ensuremath{S\ge p}，且某个} {\ensuremath{j\ge 2}} {的例外出现}}

{把所有负系数} {\ensuremath{j}} {按从大到小排成序列，取最短前缀} {\ensuremath{j_{1},}\ldots,\ensuremath{j_{s}}，使}

\[
T=j_1+\cdots+j_s\ge p.
\]

{有} {\ensuremath{2\le s\le p-1}：所有} {\ensuremath{j<p}，而至少一个} {\ensuremath{j\ge 2}，所以只要有足够项，至多} {p-1} {项便跨越} {\ensuremath{p}；若总项数更小，则全部项已跨越} {\ensuremath{p}。}

{前} {s-1} {项之和小于} {\ensuremath{p}，故} {\ensuremath{T\le p-1+j_{s}}，且} {(\ensuremath{s-1)j_{s}\le p-1}。} {对} {\ensuremath{2\le s\le p}，由} {(\ensuremath{s-1)(p+1-s)-(p-1)=(s-2)(p-s)\ge 0}，得到} {\ensuremath{j_{s}+s\le p+1}，因此}

\[
T+s\le2p.
\]

{设置} {\ensuremath{t=T-p\ge 0}、\ensuremath{z=p-s-t\ge 0}。选取}

\[
\boxed{(b-j_1\delta,\ldots,b-j_s\delta,
(b+\delta)^t,b^z).}
\tag{C2.6}
\]

{其长度恰为} {\ensuremath{p}，均值恰为} {b-delta。} {又} {\ensuremath{t\le p-s\le p-2}，所以} {\ensuremath{a=b+\delta }} {仍有副本；\ensuremath{b}} {至少留下} {\ensuremath{p}} {份。} {由于} {\ensuremath{\operatorname{gcd}(\delta ,n)=1}，这两个未选值在每个} {\ensuremath{q|n}} {下都不同余，因此操作安全。} {操作非恒等：\ensuremath{s<p}，且除负系数值外还选了} {\ensuremath{a}} {或} {\ensuremath{b}。} {新均值不同于} {\ensuremath{b}，所以} {\ensuremath{I}} {保持。}

{这里} {\ensuremath{\operatorname{gcd}(\delta ,n)=1}} {来自完整等差结构：任何非} {\ensuremath{p}} {素因子若整除} {\ensuremath{\delta }，全部值会同余；\ensuremath{p}} {本身也不整除} {\ensuremath{\delta }。}

\paragraph*{\texorpdfstring{{6.2} {如果} {\ensuremath{S<p}}}{6.2 如果 }}\label{C2-h11}
\addcontentsline{toc}{paragraph}{{6.2} {如果} {\ensuremath{S<p}}}

{零和给}

\[
nb+\delta(A-S)=0.
\]

{因为} {\ensuremath{\operatorname{gcd}(\delta ,n)=1}，有} {\ensuremath{n|(A-S)}。但} {\ensuremath{0<A-S\le 2p-2<n}，矛盾。} {所以该情形根本不会发生。}

\paragraph*{\texorpdfstring{{6.3} {唯一剩余情形}}{6.3 唯一剩余情形}}\label{C2-h12}
\addcontentsline{toc}{paragraph}{{6.3} {唯一剩余情形}}

{只有} {\ensuremath{j=1}} {的例外出现，即实际值只有} {b+delta、\ensuremath{b}、b-delta。} {它们的非} {\ensuremath{b}} {重数分别为} {A,C，均至多} {2p-2。} {零和与} {\ensuremath{\operatorname{gcd}(\delta ,n)=1}} {给} {\ensuremath{n|(A-C)}；因} {\textbar{}\ensuremath{A-C|<n}，得到}

\[
b=0,\qquad A=C.
\]

{于是状态为}

\[
\boxed{(\delta^A,(-\delta)^A,0^B),\qquad B\ge2p-1.}
\tag{C2.7}
\]

{反复平均} {(\ensuremath{\delta ,-\delta ,0^{p-2})}} {即全部归零。每次真实选} {\ensuremath{p}} {个原位置，零位置数不会减少。}

{因此终端分析只需} {\ensuremath{n>2p}，而不再要求} {\ensuremath{n>p(p-1)^{2}}。这正是三次方损失被消除的原因。}


\section{三个重值的顺序选择}\label{construct:three-heavy}
{本节先证明} \(n\ge4p-1+d_p(n)\) {的整数下降；下一构造将其改善为} \(n\ge4p-2+d_p(n)\){。两者均在原位置进行。}

\subsection*{\texorpdfstring{{C3.4} {三个} {\ensuremath{p}} {重值的统一顺序引理}}{C3.4 三个  重值的统一顺序引理}}\label{C3-h13}
\addcontentsline{toc}{subsection}{{C3.4} {三个} {\ensuremath{p}} {重值的统一顺序引理}}

\textbf{{引理。}} {设} {\ensuremath{p\ge 3}} {为奇素数，\ensuremath{n\ge 4p-1}。若合法整数零和状态恰有三个不同} {\ensuremath{p}} {重值，则它有一次保持双} {\ensuremath{p}} {重值不变量的安全非恒等整数平均，或者可由背景值引理直接收尾。}

{将三个实际重值排列为}

\[
a<b<c.
\]

\paragraph*{\texorpdfstring{{情形一：有两个重值同余模} {\ensuremath{p}}}{情形一：有两个重值同余模 }}\label{C3-h14}
\addcontentsline{toc}{paragraph}{{情形一：有两个重值同余模} {\ensuremath{p}}}

{记这两个值为} {\ensuremath{u}、\ensuremath{v}，第三个为} {\ensuremath{w}。} {两种操作}

\[
(u^{p-1},v),\qquad(v^{p-1},u)
\]

{均安全、非恒等且为整数。其均值之差为}

\[
\frac{(p-2)(u-v)}p\ne0.
\]

{因此至少一个均值不等于完整保留的} {\ensuremath{w}，新均值与} {\ensuremath{w}} {保持不变量。} {这里} {\ensuremath{p\ge 3}} {实际被使用；\ensuremath{p=2}} {的两个均值相同，不能套用这一论证。}

\paragraph*{\texorpdfstring{{情形二：三个重值分属不同模} {\ensuremath{p}} {余类，并有低重数值} {\ensuremath{x}}}{情形二：三个重值分属不同模  余类，并有低重数值 }}\label{C3-h15}
\addcontentsline{toc}{paragraph}{{情形二：三个重值分属不同模} {\ensuremath{p}} {余类，并有低重数值} {\ensuremath{x}}}

{\ensuremath{x}} {不等于} {\ensuremath{a}、\ensuremath{b}、\ensuremath{c}。}

\begin{itemize}
\tightlist
\item
  {若} {\ensuremath{x<c}，使用} {\ensuremath{a}、\ensuremath{b}} {作两个模逆调节锚点，选取} {(\ensuremath{a^{i},b^{p-1-i},x)}} {使其平均为整数。所有参与值都小于未参与的} {\ensuremath{c}，故新均值严格小于} {\ensuremath{c}。}
\item
  {若} {\ensuremath{x>c}，使用} {\ensuremath{b}、\ensuremath{c}} {作锚点，选取} {(\ensuremath{b^{i},c^{p-1-i},x)}。所有参与值都大于未参与的} {\ensuremath{a}，故新均值严格大于} {\ensuremath{a}。}
\end{itemize}

{两种操作都安全：所用锚点各至多消耗} {p-1} {份，仍留下副本，第三值只选一份。} {它们非恒等，因为} {\ensuremath{x}} {与锚点是不同的实际值。} {新重值与完整保留的极端重值不同，所以不变量保持。}

{这个论证与低重数值的具体余数无关，不要求它与某个重值同余。}

\paragraph*{\texorpdfstring{{情形三：只有三个实际值，至少一个极端重值超过} {\ensuremath{p}} {份}}{情形三：只有三个实际值，至少一个极端重值超过  份}}\label{C3-h16}
\addcontentsline{toc}{paragraph}{{情形三：只有三个实际值，至少一个极端重值超过} {\ensuremath{p}} {份}}

{假设} {\ensuremath{f(a)\ge p+1}。把} {\ensuremath{a}} {只选一份，并用} {\ensuremath{b}、\ensuremath{c}} {的模逆调节选择其余} {p-1} {份。} {\ensuremath{a}} {仍有至少} {\ensuremath{p}} {份，且新均值严格大于} {\ensuremath{a}，所以保持两个重值。} {若} {\ensuremath{f(c)\ge p+1}，对称处理，新均值严格小于} {\ensuremath{c}。} {安全性仍由两个调节锚点的未选副本保证。}

\paragraph*{\texorpdfstring{{情形四：只有三个实际值，两个极端重值都恰为} {\ensuremath{p}} {份}}{情形四：只有三个实际值，两个极端重值都恰为  份}}\label{C3-h17}
\addcontentsline{toc}{paragraph}{{情形四：只有三个实际值，两个极端重值都恰为} {\ensuremath{p}} {份}}

{此时}

\[
f(b)=n-2p\ge2p-1.
\]

{所以直接应用已证的背景值引理。它给出保持不变量的安全操作或明确的相反值终端。}

{四种情形穷尽所有可能，引理得证。}

\subsection*{\texorpdfstring{{C3.5} {两个} {\ensuremath{p}} {重值的安全预留池}}{C3.5 两个  重值的安全预留池}}\label{C3-h18}
\addcontentsline{toc}{subsection}{{C3.5} {两个} {\ensuremath{p}} {重值的安全预留池}}

{现在恰有两个} {\ensuremath{p}} {重值} {\ensuremath{a}、\ensuremath{b}。} {若某个重数至少为} {2p-1，使用背景值引理。} {否则} {\ensuremath{a}、\ensuremath{b}} {的重数均不超过} {2p-2，其他实际值重数均小于} {\ensuremath{p}。}

{分别保留} {\ensuremath{p}} {份} {\ensuremath{a}、\ensuremath{p}} {份} {\ensuremath{b}。余下池大小为}

\[
n-2p\ge d+2p-1.
\]

{池中} {\ensuremath{a}、\ensuremath{b}} {分别至多剩} {p-2} {份，其他值至多} {p-1} {份。} {再保护至多} {\ensuremath{d}} {个危险位置后，应用} {EGZ} {选出整数} {\ensuremath{p}} {平均。} {它非恒等，因为每个实际值不足} {\ensuremath{p}} {份；它安全，且保留的两组} {\ensuremath{p}} {份完全不动。} {因此不变量保持。}

\subsection*{\texorpdfstring{{C3.6} {充分性与统一上界}}{C3.6 充分性与统一上界}}\label{C3-h19}
\addcontentsline{toc}{subsection}{{C3.6} {充分性与统一上界}}

{从} {pX} {初始化双} {\ensuremath{p}} {重值不变量。至少四个重值用已有顺序引理；三个重值用第} {4} {节；两个重值用第} {5} {节。} {每个非终端都能作一个安全非恒等整数平均，使平方能量至少减少} {2。} {所以路径有限终止于已经给出真实收尾的相反值族。} {在原} {\ensuremath{X}} {上执行同样的位置序列，所有中间值均在} {(1/p)Z。} {非本原整数输入先除去公因子再乘回即可。必要性仍是非} {\ensuremath{p}} {共同非零余数障碍。}

{为证明} {(2)，令} {\ensuremath{\ell =\operatorname{ceil}(log_{2}}} {p)、\ensuremath{n_{0}=4p+\ell +1}。} {因为} {\ensuremath{\ell \le p-1}，有}

\[
n_0\le5p<8\cdot2^\ell,
\]

{从而} {\ensuremath{\operatorname{floor}(log_{2}}} {\ensuremath{n_{0})\le \ell +2}，故}

\[
n_0-\lfloor\log_2n_0\rfloor\ge4p-1.
\]

{函数} {\ensuremath{n-\operatorname{floor}(log_{2}}} {n)} {在整数上单调不减，而} {\ensuremath{d_{p}(n)\le \operatorname{floor}(log_{2}}} {n)。} {因此所有} {\ensuremath{n\ge n_{0}}} {满足} {(1)。}


\section{逆零和引理与宽维数下降}\label{construct:inverse-egz}
\subsection*{\texorpdfstring{{C4.2} {三重值分支只需要两个额外位置}}{C4.2 三重值分支只需要两个额外位置}}\label{C4-h20}
\addcontentsline{toc}{subsection}{{C4.2} {三重值分支只需要两个额外位置}}

{沿用}第\ref{construct:three-heavy}节{的顺序分类：有低重数值、某个极端重值多于} {\ensuremath{p}} {份，或两个重值同余} {\ensuremath{mod}} {\ensuremath{p}} {时，原操作已能保持不变量。} {剩下只有三个实际值}

\[
(a^p,b^{p+r},c^p),\quad a<b<c,\quad r\ge2,
\]

{三个余数模} {\ensuremath{p}} {不同，且} {\ensuremath{pa+(p+r)b+pc=0}。} {唯一选取} {\ensuremath{0\le i<p}，使}

\[
(a^i,c^{p-1-i},b)
\]

{的平均为整数。\ensuremath{b}} {还剩至少} {\ensuremath{p}} {份；若新均值不等于} {\ensuremath{b}，就已完成。}

{若新均值等于} {\ensuremath{b}，得到}

\[
i(a-b)+(p-1-i)(c-b)=0,
\tag{C4.3}
\]

{其中} {\ensuremath{1\le i\le p-2}。} {改选两份} {\ensuremath{b}，并令} {\ensuremath{j}} {为} {2i} {\ensuremath{mod}} {\ensuremath{p}} {在} {{[}0,p-1{]}} {的代表。} {若} {\ensuremath{j\le p-2}，则}

\[
(a^j,c^{p-2-j},b^2)
\tag{C4.4}
\]

{是整数平均，仍完整保留} {\ensuremath{p}} {份} {\ensuremath{b}。} {其均值不可能再次等于} {\ensuremath{b}：否则由} {(3)} {得}

\[
(p-1)j=(p-2)i.
\]

{写} {\ensuremath{j=2i-ep}，\ensuremath{e}} {为} {0} {或} {1，即推出} {\ensuremath{i=e(p-1)}，与} {\ensuremath{1\le i\le p-2}} {矛盾。}

{唯一不能选} {(4)} {的情形是} {\ensuremath{j=p-1}，即} {\ensuremath{i=(p-1)/2}。} {此时} {(3)} {给} {\ensuremath{a+c=2b}；全局零和给} {\ensuremath{nb=0}，所以} {\ensuremath{b=0}、\ensuremath{a=-c}，是明确相反值终端。}

{所有操作均为整数，且每个参与值都保留未选副本，因此自动保持所有非} {\ensuremath{p}} {同余见证。} {这证明三重值分支的门槛是} {3p+2，而不是此前使用的} {4p-1。}

\subsection*{\texorpdfstring{{C4.3} {两个重值可以共享预留容量}}{C4.3 两个重值可以共享预留容量}}\label{C4-h21}
\addcontentsline{toc}{subsection}{{C4.3} {两个重值可以共享预留容量}}

\textbf{{引理。}} {设合法整数零和状态恰有两个} {\ensuremath{p}} {重值} {\ensuremath{a}、\ensuremath{b}，\ensuremath{a\ne b}} {\ensuremath{mod}} {\ensuremath{p}，}

\[
f(a)=p+\alpha,\quad f(b)=p+\beta,\quad
0\le\alpha,\beta\le p-2,\quad\alpha+\beta\ge p-2.
\tag{C4.5}
\]

{则有安全非恒等整数} {\ensuremath{p}} {平均，保持两个不同} {\ensuremath{p}} {重值。}

{令} {\ensuremath{\delta =a-b}。对每个其他实际值} {\ensuremath{x}，作唯一整数调节}

\[
(a^i,b^{p-1-i},x),\qquad0\le i<p.
\tag{C4.6}
\]

{若操作已保持双重值，证明完成。以下假设所有这些操作均失败。}

{因} {\ensuremath{\alpha +\beta \ge p-2}，(6)} {不可能同时把} {\ensuremath{a}、\ensuremath{b}} {的重数都降到} {\ensuremath{p}} {以下。} {所以失败必须是均值恰等于唯一保留的重值，并且} {\ensuremath{x}} {被迫满足下列之一：}

\[
x=b-i\delta,\quad\alpha+1\le i\le p-1;
\qquad
x=a+j\delta,\quad\beta+1\le j\le p-1.
\tag{C4.7}
\]

{两类分别称为下侧与上侧，只指} {\ensuremath{\delta }} {系数的方向；不要求} {\ensuremath{\delta }} {正。} {若} {(6)} {同时保留两个旧重值，本来就不会失败。}

{于是全部值都在} {b+delta} {\ensuremath{Z}} {中。合法性和} {\ensuremath{p}} {不整除} {\ensuremath{\delta }} {给} {\ensuremath{\operatorname{gcd}(\delta ,n)=1}。} {设下侧系数总和为} {\ensuremath{L}，上侧系数总和为} {\ensuremath{H}，上侧位置数为} {\ensuremath{C}。} {零和关系为}

\[
nb+\delta\bigl(f(a)+C+H-L\bigr)=0.
\tag{C4.8}
\]

{若} {L,H} {都小于} {\ensuremath{p}，则括号严格介于} {0} {与} {\ensuremath{n}} {之间：}

\[
f(a)+C+H-L\ge f(a)-(p-1)>0,
\]

\[
f(a)+C+H-L\le n-f(b)+p-1<n.
\]

{这与} {(8)} {及} {\ensuremath{\operatorname{gcd}(\delta ,n)=1}} {矛盾。所以至少一侧的系数总和达到} {\ensuremath{p}。}

{不妨下侧达到} {\ensuremath{p}。将其系数降序排列，取第一次跨越} {\ensuremath{p}} {的前缀} {\ensuremath{i_{1},}\ldots,\ensuremath{i_{s}}，总和} {\ensuremath{T}。} {每个} {\ensuremath{i_{t}\le p-1}，故} {\ensuremath{s\ge 2}。所有这些值是低重数值；若系数全为} {1，总和达到} {\ensuremath{p}} {就会出现第三个} {\ensuremath{p}} {重值，矛盾。} {因此可用既有跨越估计得到} {\ensuremath{s\le p-1}、\ensuremath{T+s\le 2p}。}

{平均}

\[
\boxed{(b-i_1\delta,\ldots,b-i_s\delta,
a^{T-p},b^{2p-s-T})}
\tag{C4.9}
\]

{产生} {\ensuremath{p}} {份} {b-delta。两个锚点各至多消耗} {p-2} {份，所以保留副本，且在所有危险素数下提供差异，操作安全。} {锚点总消耗为} {\ensuremath{p-s\le p-2\le \alpha +\beta }，不可能同时超出各自的富余数量，故至少一个旧锚点仍是} {\ensuremath{p}} {重值。} {新均值既不是} {\ensuremath{a}} {也不是} {\ensuremath{b}，因此保持双重值。} {上侧情形交换} {\ensuremath{a}、\ensuremath{b}} {并将} {\ensuremath{\delta }} {变号即可。引理得证。}

{这里总重数条件是} {\ensuremath{f(a)+f(b)\ge 3p-2}。它允许两个重值共同承担原先由一个} {2p-1} {重背景承担的工作。}

\subsection*{\texorpdfstring{{C4.4} {少一个位置的逆} {EGZ，并给出初等证明}}{C4.4 少一个位置的逆 EGZ，并给出初等证明}}\label{C4-h22}
\addcontentsline{toc}{subsection}{{C4.4} {少一个位置的逆} {EGZ，并给出初等证明}}

\textbf{{引理。}} {对素数} {\ensuremath{p\ge 5}，长度} {2p-2} {的} {\ensuremath{F_{p}}} {序列没有} {\ensuremath{p}} {项零和，当且仅当它恰为}

\[
u^{p-1},v^{p-1},\qquad u\ne v.
\tag{C4.10}
\]

{这是循环素数群的逆} {EGZ} {边界。以下给出本构造所需的独立初等证明，不依赖未检验的外部黑箱。}

{若某余数出现} {\ensuremath{p}} {次，已经有} {\ensuremath{p}} {项零和。若只出现两个余数且每个少于} {\ensuremath{p}} {次，因总长} {2p-2，只能是} {(10)；此时任何} {\ensuremath{p}} {项都各含} {1} {至} {p-1} {份两种余数，不可能零和。}

{剩下有至少三个余数，各重数至多} {p-1。可以取出两对互不重叠的位置，使每对余数不同，并且两对的差不互为正负：}

\begin{itemize}
\tightlist
\item
  {取重数最大的余数} {\ensuremath{a}，其重数至少为} {2。若能找到} {\ensuremath{b}、\ensuremath{c}，使} {b-a} {不等于正负(c-a)，用两份} {\ensuremath{a}} {分别与} {\ensuremath{b}、\ensuremath{c}} {配对。}
\item
  {否则恰有} {\ensuremath{a}、\ensuremath{b}、\ensuremath{c}} {三种余数，且} {\ensuremath{b}、\ensuremath{c}} {关于} {\ensuremath{a}} {对称。\ensuremath{b}、\ensuremath{c}} {至少有一种出现两次，否则总长至多} {\ensuremath{p+1<2p-2}。假设} {\ensuremath{b}} {出现两次，取} {(a,b)、(b,c)，其差相差倍率} {2，在} {\ensuremath{p\ge 5}} {时不互为正负。}
\end{itemize}

{这两对一定消耗至少一份唯一可能具有} {p-1} {重数的余数。因此剩余} {2p-6} {项的最大余数重数至多} {p-2。} {它们能再组成} {p-4} {对余数不同的位置，留下两个单点：在完全多部图中，不同余数位置之间相连，其最大匹配大小为} {\ensuremath{min(\operatorname{floor}(M/2),M-f_{max})}，这里至少为} {p-4。}

{每对可选一个元素，给出一个二元集合。前两对的和集含四个不同余数；其余} {p-4} {对依次由} {Cauchy--Davenport} {使和集至少增长一个余数。} {最后和集为整个} {\ensuremath{F_{p}}。加上两个固定单点，就能选出} {\ensuremath{p-2+2=p}} {项零和，矛盾。} {引理得证。}

{这个证明也给出高效构造：余数桶、最大重数匹配，以及每对一次循环位集更新。没有枚举} {\ensuremath{p}} {元位置子集。}

\subsection*{\texorpdfstring{{C4.5} {用逆定理补上} {4p-2+d} {的双重值分支}}{C4.5 用逆定理补上 4p-2+d 的双重值分支}}\label{C4-h23}
\addcontentsline{toc}{subsection}{{C4.5} {用逆定理补上} {4p-2+d} {的双重值分支}}

{现在} {\ensuremath{n}} {满足} {(1)，恰有两个重值} {\ensuremath{a}、\ensuremath{b}。} {若某重数达到} {2p-1，调用旧背景值引理；若} {\ensuremath{a=b}} {\ensuremath{mod}} {\ensuremath{p}，调用已证的} {3p+d} {同余双重值引理。} {所以只需} {\ensuremath{a\ne b}} {\ensuremath{mod}} {\ensuremath{p}，且两者重数至多} {2p-2。}

\paragraph*{\texorpdfstring{{5.1} {先消除容易的同余低重数值}}{5.1 先消除容易的同余低重数值}}\label{C4-h24}
\addcontentsline{toc}{paragraph}{{5.1} {先消除容易的同余低重数值}}

{若低重数值} {\ensuremath{x=a}} {\ensuremath{mod}} {\ensuremath{p}，平均} {(\ensuremath{a^{p-1},x)}。它安全，\ensuremath{b}} {完整保留。} {唯一碰撞情形为} {\ensuremath{x=pb-(p-1)a}。} {若} {\ensuremath{x}} {至少有两份，改平均} {(\ensuremath{a^{p-2},x^{2})}，新均值为} {2b-a，不等于} {\ensuremath{b}。} {它也安全：若输出全同余，未选} {\ensuremath{a}、\ensuremath{b}} {及该仿射关系会迫使} {\ensuremath{x}} {同余，矛盾。}

{所以若仍未完成，\ensuremath{a}} {的余类中最多有一个低重数单点；\ensuremath{b}} {的余类同理。}

{若} {\ensuremath{f(a)+f(b)\ge 3p-2}，直接使用第} {3} {节。} {以下设}

\[
\alpha+\beta=f(a)+f(b)-2p\le p-3.
\tag{C4.11}
\]

\paragraph*{\texorpdfstring{{5.2} {保护后的} {2p-2} {池}}{5.2 保护后的 2p-2 池}}\label{C4-h25}
\addcontentsline{toc}{paragraph}{{5.2} {保护后的} {2p-2} {池}}

{预留} {\ensuremath{p}} {份} {\ensuremath{a}、\ensuremath{p}} {份} {\ensuremath{b}。} {对每个危险素数，若其整个例外支持仍落在余池中，额外预留其中一个位置。} {这些额外预留都可取低重数值，因为} {\ensuremath{a}、\ensuremath{b}} {已有冻结副本。} {至多预留} {\ensuremath{d}} {个位置，余池仍至少有} {2p-2} {个位置。}

{取其中} {2p-2} {个，必须包括全部} {\ensuremath{\alpha }} {份} {\ensuremath{a}} {和} {\ensuremath{\beta }} {份} {\ensuremath{b}} {的富余副本。由} {(11)} {及容量可做到。} {池中每个实际值都不足} {\ensuremath{p}} {份。} {若逆} {EGZ} {构造给出} {\ensuremath{p}} {项零和模} {\ensuremath{p}，则平均非恒等、安全，并保留两个原} {\ensuremath{p}} {重块，已完成。}

{否则池的余数必须是} {(10)。} {若} {\ensuremath{\alpha >0}，\ensuremath{a}} {余类在池内须有} {p-1} {项；第} {5.1} {节已使其中低重数项至多一个，迫使} {\ensuremath{\alpha \ge p-2}，违背} {(11)。} {\ensuremath{\beta >0}} {同理。} {因此}

\[
f(a)=f(b)=p,
\]

{池中两类都只由低重数值组成，且都不同于} {\ensuremath{a}、\ensuremath{b}} {的模} {\ensuremath{p}} {余数。}

\paragraph*{\texorpdfstring{{5.3} {借一份锚点，碰撞后借两份}}{5.3 借一份锚点，碰撞后借两份}}\label{C4-h26}
\addcontentsline{toc}{paragraph}{{5.3} {借一份锚点，碰撞后借两份}}

{记两个余数类为} {\ensuremath{u}、\ensuremath{v}，各有} {p-1} {个位置。} {唯一选择} {\ensuremath{i}，使一份} {\ensuremath{a}、\ensuremath{i}} {份} {\ensuremath{u}} {类、p-1-i} {份} {\ensuremath{v}} {类的总和被} {\ensuremath{p}} {整除。} {因为} {\ensuremath{a}} {的余数不同于} {\ensuremath{u}、\ensuremath{v}，必有} {\ensuremath{1\le i\le p-2}。}

{若任何一个这样的选择均值不等于} {\ensuremath{b}，就保持} {\ensuremath{b}} {并产生新重值。} {操作安全：借走一份} {\ensuremath{a}} {仍留副本，原来保护的低重数位置也未选。}

{若所有选择都碰撞到} {\ensuremath{b}，交换同一个余类内的一个选中位置与未选位置可知，该余类必须是同一个实际值。} {故两类都实际常值，记为} {\ensuremath{u}、\ensuremath{v}，并满足}

\[
a+i u+(p-1-i)v=pb.
\tag{C4.12}
\]

{若} {\ensuremath{i\ne (p-1)/2}，令} {\ensuremath{j=2i}} {\ensuremath{mod}} {\ensuremath{p}，则} {\ensuremath{j\le p-2}。改平均}

\[
(a^2,u^j,v^{p-2-j}).
\tag{C4.13}
\]

{其均值为} {2b-v} {或} {2b-u；\ensuremath{u}、\ensuremath{v}} {都不是} {\ensuremath{b}，故新均值不等于} {\ensuremath{b}。} {每种参与值都还有未选副本，自动安全。}

{若} {\ensuremath{i=(p-1)/2}，就改从} {\ensuremath{b}} {端作同样构造。} {两端的} {\ensuremath{i}} {不可能相同，因为} {\ensuremath{a\ne b}} {\ensuremath{mod}} {\ensuremath{p}。所以从} {\ensuremath{b}} {端或第一选择已成功，或两份锚点的选择成功。}

{双重值分支至此完整闭合。}

\subsection*{\texorpdfstring{{C4.6} {全输入与有限终止}}{C4.6 全输入与有限终止}}\label{C4-h27}
\addcontentsline{toc}{subsection}{{C4.6} {全输入与有限终止}}

{条件} {(C4.1)} {蕴含} {\ensuremath{n\ge 3p+d}，所以从} {pX} {出发的旧两步初始化有效。} {同时} {\ensuremath{n\ge 3p+2}，所以第} {2} {节与旧顺序分支完整覆盖三个重值。} {至少四个重值使用旧极端锚点引理，两个重值使用第} {5} {节。}

{每次非终端操作都是安全非恒等整数平均，保持至少两个不同} {\ensuremath{p}} {重值，平方能量严格下降。} {唯一单列终端是} {(\ensuremath{t^A,(-t)^A,0^B)}，其中} {\ensuremath{B\ge p-2}，可逐对消去。} {因此有限终止，且缩回} {\ensuremath{X}} {后只需一位额外} {\ensuremath{p}} {进精度。}

{令} {\ensuremath{\ell =\operatorname{ceil}(log_{2}}} {p)，\ensuremath{n_{0}=4p+\ell }。因为} {\ensuremath{n_{0}<2^{\ell +3}}，有} {\ensuremath{\operatorname{floor}(log_{2}}} {\ensuremath{n_{0})\le \ell +2}，故}

\[
n_0-\lfloor\log_2 n_0\rfloor\ge4p-2.
\]

{\ensuremath{n-\operatorname{floor}(log_{2}}} {n)} {单调不减，且} {\ensuremath{d_{p}(n)\le \operatorname{floor}(log_{2}}} {n)，得到} {(2)。} {对} {\ensuremath{p=5}，\ensuremath{n=19}} {的} {\ensuremath{d=1}；对} {\ensuremath{p=7}，\ensuremath{n=27}} {的} {\ensuremath{d=1}。接原} {\ensuremath{n\ge 20}、\ensuremath{n\ge 28}} {定理，分别得到} {\ensuremath{N(5)\le 19}、\ensuremath{N(7)\le 27}。}


\section{三个等大块的完整证明}\label{construct:three-blocks}
\subsection*{\texorpdfstring{{C5.2} {\ensuremath{n=3p}} {的资源与端点不同}}{C5.2  的资源与端点不同}}\label{C5-h28}
\addcontentsline{toc}{subsection}{{C5.2} {\ensuremath{n=3p}} {的资源与端点不同}}

{三个} {\ensuremath{p}} {重块写成}

\[
\mathcal C_p(a,b,c)=(a^p,b^p,c^p),\qquad a+b+c=0.
\]

{每个输入值恰有} {\ensuremath{p}} {个实际副本。} {任意} {3x3} {非负整数矩阵} {\ensuremath{K}，只要每行、每列之和均为} {\ensuremath{p}，就可以按以下方式实现} {K/p：}

\begin{itemize}
\tightlist
\item
  {第} {\ensuremath{i}} {个新块从第} {\ensuremath{j}} {个旧块取} {\ensuremath{K_{ij}}} {个尚未使用的位置；}
\item
  {每列和为} {\ensuremath{p}} {保证恰好用完每个旧块；}
\item
  {每行和为} {\ensuremath{p}} {保证每个新块可用一次真实} {\ensuremath{p}} {平均产生。}
\end{itemize}

{三个新块互不相交，共用原来的全部} {3p} {个位置，恰执行三次} {\ensuremath{p}} {平均。} {这是一层真实的分组操作，不是允许任意双随机矩阵。} {我们下面只使用两个明确的} {\ensuremath{K}，逐项证明其非负性。}

\subsection*{\texorpdfstring{{C5.3} {一对统一的正向剪切运输表}}{C5.3 一对统一的正向剪切运输表}}\label{C5-h29}
\addcontentsline{toc}{subsection}{{C5.3} {一对统一的正向剪切运输表}}

{以下先设} {\ensuremath{p\ge 5}。取} {\ensuremath{\epsilon }} {属于} {\{1,-1\}，使}

\[
p=3m+\epsilon.
\]

{因为} {\ensuremath{p}} {是不等于} {3} {的奇素数，\ensuremath{m\ge 2}。} {令} {\ensuremath{J_{3}}} {是} {3x3} {全一矩阵，置}

\[
D=\begin{pmatrix}
1&-1&0\\
1&-1&0\\
-2&2&0
\end{pmatrix}
=\begin{pmatrix}1\\1\\-2\end{pmatrix}(1,-1,0).
\]

{有}

\[
D\mathbf1=0,\qquad\mathbf1^TD=0,\qquad D^2=0.
\]

{定义两个整数运输表}

\[
\boxed{K_\pm=mJ_3+\epsilon(I\pm D).}
\tag{C5.1}
\]

{I+D} {和} {I-D} {的每个元素都位于} {{[}-2,2{]}。} {因此} {\ensuremath{K_+}、\ensuremath{K_-}} {的每个元素至少为} {\ensuremath{m-2\ge 0}。} {两者每行、每列之和都是} {\ensuremath{3m+\epsilon =p}，所以第} {2} {节给出真实实现。}

{对零和三参数} {\ensuremath{v=(a,b,c)^T}，\ensuremath{J_{3v}=0}，故}

\[
\frac1pK_\pm v=\frac\epsilon p(I\pm D)v.
\tag{C5.2}
\]

{改用整数幺模参数坐标}

\[
x=a,\qquad y=a-b,
\]

{I+D} {的作用是}

\[
(a,b,c)\mapsto(a+(a-b),\ b+(a-b),\ c-2(a-b)),
\]

{所以在} {(x,y)} {上恰为} {U(1)；I-D} {恰为} {U(-1)。} {因此一次三平均层分别实现}

\[
\boxed{\frac\epsilon pU(1),\qquad\frac\epsilon pU(-1).}
\tag{C5.3}
\]

{两层复合在零和参数空间上是} {\ensuremath{p^{-2}I}。} {这是明确的正向逆关系，不需要额外的周期存在性假设。}

\paragraph*{\texorpdfstring{{五平均的两个实际表}}{五平均的两个实际表}}\label{C5-h30}
\addcontentsline{toc}{paragraph}{{五平均的两个实际表}}

{\ensuremath{p=5}} {时，\ensuremath{m=2}、\ensuremath{\epsilon =-1}，两个表为}

\[
K_+=\begin{pmatrix}0&3&2\\1&2&2\\4&0&1\end{pmatrix},
\qquad
K_-=\begin{pmatrix}2&1&2\\3&0&2\\0&4&1\end{pmatrix}.
\]

{例如} {\ensuremath{K_+}} {的第一行就是从} {\ensuremath{b}} {块取三个位置、从} {\ensuremath{c}} {块取两个位置进行五平均。} {每张表只需三次五平均。}

{这说明} {3p} {的正向实现来自足够的副本和非负运输余量，不依赖精细的} {\ensuremath{p}} {参数周期。}

\subsection*{\texorpdfstring{{C5.4} {同一个} {\ensuremath{\Gamma _{0}(3)}} {覆盖所有} {\ensuremath{p}} {的合法三块}}{C5.4 同一个  覆盖所有  的合法三块}}\label{C5-h31}
\addcontentsline{toc}{subsection}{{C5.4} {同一个} {\ensuremath{\Gamma _{0}(3)}} {覆盖所有} {\ensuremath{p}} {的合法三块}}

{交换} {\ensuremath{a}、\ensuremath{b}} {两个块的名称，参数作用为}

\[
\Sigma=\begin{pmatrix}1&-1\\0&-1\end{pmatrix}.
\]

{交换} {\ensuremath{a}、\ensuremath{c}} {两个块的名称，参数作用为}

\[
P=\begin{pmatrix}-2&1\\-3&2\end{pmatrix}.
\]

{这些交换都只是更新完整} {\ensuremath{p}} {元块的原位置账本。} {直接计算}

\[
P\Sigma=U(1)L(-3),\qquad
\boxed{L(-3)=U(-1)P\Sigma.}
\tag{C5.4}
\]

{所以} {U(1)、L(3)} {及其逆都有正向实现。}

{完整} {\ensuremath{\Gamma _{0}(3)}} {包含可以用固定的四个右陪集证明：}

\[
I,\qquad \mathsf S\mathsf T^{-1},\quad\mathsf S,
\quad\mathsf S\mathsf T,
\qquad
\mathsf S=\begin{pmatrix}0&-1\\1&0\end{pmatrix},\quad\mathsf T=U(1).
\]

{非平凡} {Schreier} {回路只有} {U(1)、L(-3)} {及其逆。} {因此上述正向操作包含整个} {\ensuremath{\Gamma _{0}(3)}，而不是仅在模} {3} {上有足够的像。}

{对本原整数三参数，\ensuremath{a+b+c=0}，}

\[
G(a,b,c)=\gcd(a-b,3).
\]

{证明：在} {\ensuremath{y=a-b}} {下，另一差} {\ensuremath{b-c=3a-2y}；因} {\ensuremath{\operatorname{gcd}(a,y)=1}，两个差的} {\ensuremath{gcd}} {等于} {\ensuremath{\operatorname{gcd}(y,3)}。} {复制} {\ensuremath{p}} {次不改变坐标} {\ensuremath{gcd}} {或差分} {\ensuremath{gcd}。} {于是合法三块恰对应本原} {(x,y)} {且} {3} {不整除} {\ensuremath{y}。}

{取} {\ensuremath{\alpha }}\emph{{x+beta}}{\ensuremath{y=1}，并沿} {Bezout} {通解调整使} {3} {整除} {\ensuremath{\alpha }。则}

\[
H=\begin{pmatrix}y&-x\\\alpha&\beta\end{pmatrix}\in\Gamma_0(3),
\qquad H(x,y)^T=(0,1)^T.
\]

{由群包含与正向逆，将} {\ensuremath{H}} {编译成有限个真实运输层，得到}

\[
(0^p,s^p,(-s)^p).
\]

{令} {\ensuremath{h=(p-1)/2}。依次平均}

\[
(s^h,(-s)^h,0),\quad
(s^h,(-s)^h,0),\quad
(s,-s,0^{p-2}),
\]

{三次实际} {\ensuremath{p}} {平均后全部归零。}

{标准模群矩阵} {\ensuremath{S}} {只参与代数分解，不被当成自由物理操作。} {整数} {Euclid} {分解和四陪集重写保证每个输入都有有限操作词；不需逐输入搜索路径。}

\subsection*{\texorpdfstring{{C5.5} {任意} {3p} {元合法输入的三步入口}}{C5.5 任意 3p 元合法输入的三步入口}}\label{C5-h32}
\addcontentsline{toc}{subsection}{{C5.5} {任意} {3p} {元合法输入的三步入口}}

{设} {\ensuremath{X}} {是非零本原整数零和输入，G(X)} {为} {\ensuremath{p}} {幂，\ensuremath{p\ge 5}。} {因为} {\ensuremath{3|3p}} {而} {\ensuremath{3\ne p}，原输入模} {3} {不全同余。}

{任意将位置分成} {\ensuremath{A}、\ensuremath{B}、\ensuremath{C}} {三组，每组} {\ensuremath{p}} {个，记各组和为} {\ensuremath{S_{A}}、\ensuremath{S_{B}}、\ensuremath{S_{C}}。} {如果三个组和模} {3} {不全同余，直接分别平均即可。}

{否则三个组和同余于} {\ensuremath{r}。由于原输入模} {3} {不全同余，必有不同组中的两个值} {\ensuremath{x}、\ensuremath{y}} {模} {3} {不同。} {交换它们在分组中的归属，令} {\ensuremath{\delta =y-x\ne 0}} {\ensuremath{mod}} {3，则新的组和余数是}

\[
r+\delta,\qquad r-\delta,\qquad r,
\]

{它们两两不同。} {按修复后的分组执行三次} {\ensuremath{p}} {平均，进入合法三块状态。} {这一步使用} {\ensuremath{p}} {在模} {3} {下可逆；没有调用更小维数的} {\ensuremath{p}} {平均定理。}

{必要性仍是非} {\ensuremath{p}} {共同非零模余数障碍。} {\ensuremath{p=3}} {时} {\ensuremath{n=9}，用两层三元蝶形网络即可把所有输入平均。} {至此定理} {\ensuremath{A}} {对每个奇素数} {\ensuremath{p}} {都已证明。}

\subsection*{\texorpdfstring{{C5.7} {3p} {以上的同余双重值不再是缺口}}{C5.7 3p 以上的同余双重值不再是缺口}}\label{C5-h33}
\addcontentsline{toc}{subsection}{{C5.7} {3p} {以上的同余双重值不再是缺口}}

{证明引理} {\ensuremath{B}。设恰有两个不同} {\ensuremath{p}} {重值} {\ensuremath{a}、\ensuremath{b}，且} {\ensuremath{a=b}} {\ensuremath{mod}} {\ensuremath{p}，重数为} {\ensuremath{A}、\ensuremath{B}。}

{若} {\ensuremath{B\ge p+1}，平均} {(\ensuremath{a^{p-1},b)}。它是非恒等整数平均，均值严格介于} {\ensuremath{a}、\ensuremath{b}} {之间，而} {\ensuremath{b}} {仍有至少} {\ensuremath{p}} {份。} {\ensuremath{a}} {和} {\ensuremath{b}} {都有未选副本，所以操作安全，并继续有两个不同} {\ensuremath{p}} {重值。} {\ensuremath{A\ge p+1}} {的情况对称。}

{只剩} {\ensuremath{A=B=p}。低重数位置池} {\ensuremath{L}} {有}

\[
|L|=n-2p\ge p+d_p(n).
\]

{对每个危险素数，若其例外支持全部落在} {\ensuremath{L}} {中，就保留其中一个位置不选。} {至多保护} {\ensuremath{d_{p}(n)}} {个位置，仍能取出} {\ensuremath{p}} {个低重数位置} {\ensuremath{x_{1},}\ldots,\ensuremath{x_{p}}。} {若危险支持包含整个} {\ensuremath{a}} {或} {\ensuremath{b}} {重值，它至少占} {\ensuremath{p}} {个位置，下面每次只选至多} {p-1} {份锚点，故也不会完整选中它。}

{对模} {\ensuremath{p}} {的} {\ensuremath{p}} {个差值} {\ensuremath{x_{i}-a}，使用部分和}

\[
0,\quad x_1-a,\quad (x_1-a)+(x_2-a),\quad\ldots
\]

{必有两项相等。两项之间给出非空片段} {\ensuremath{T}，长度} {\ensuremath{s\le p}，且}

\[
\sum_{x\in T}(x-a)=0\pmod p.
\]

{可以用} {\ensuremath{a}} {或} {\ensuremath{b}} {把} {\ensuremath{T}} {补足} {\ensuremath{p}} {个位置。两种均值为}

\[
c_a=\frac{\sum T+(p-s)a}{p},\qquad
c_b=\frac{\sum T+(p-s)b}{p}.
\]

{它们都是整数，且}

\[
c_a-c_b=\frac{p-s}{p}(a-b).
\]

{不可能同时} {\ensuremath{c_{a}=b}、\ensuremath{c_{b}=a}，否则因} {\ensuremath{a\ne b}} {会推出} {\ensuremath{s=2p}，与} {\ensuremath{s\le p}} {矛盾。} {所以至少一种补法产生不同于完整保留重值的新} {\ensuremath{p}} {重值。} {所选片段都是不同于} {\ensuremath{a}、\ensuremath{b}} {的低重数值，保证非恒等；若} {\ensuremath{s=p}，则低重数限制同样排除恒等。} {危险见证全部保留，所以操作安全。引理} {\ensuremath{B}} {得证。}

{这里需要的低重数池长度是} {p+d，不再是} {EGZ} {的} {2p-1+d。} {构造使用一次线性扫描的模} {\ensuremath{p}} {部分和碰撞，无位置子集枚举。}


\section{含零支持的有限收尾}\label{construct:zero-tail}
\subsection*{\texorpdfstring{{C6.2} {支持数传播：子系统保留一个零，合法性自动成立}}{C6.2 支持数传播：子系统保留一个零，合法性自动成立}}\label{C6-h34}
\addcontentsline{toc}{subsection}{{C6.2} {支持数传播：子系统保留一个零，合法性自动成立}}

\textbf{{引理。}} {假设} {\ensuremath{A_{p}(m)}} {已证，\ensuremath{n\ge m}。若} {\ensuremath{n}} {元有理零和状态至少有} {n-m+1} {个零，则该状态可达零。}

\textbf{{证明。}} {非零支持数至多} {m-1。选取} {\ensuremath{m}} {个原位置，包含全部非零坐标，并至少包含一个零。} {其余位置全为零，所以这个子块自身的总和也是零。} {对子块清分母并除坐标公因子后，因为它包含零，差分} {\ensuremath{gcd}} {等于坐标} {\ensuremath{gcd}，即为} {1。} {故可调用} {\ensuremath{A_{p}(m)}。所有操作只在选出的原位置中进行，块外零值保持不动。证毕。}

{这里确实消除了局部均值与全局均值之间的差别：所选子系统包含全部非零坐标，局部均值严格等于全局目标零。} {这不是任意嵌入一个已解窗口。}

{如果只知道``\ensuremath{m}} {元零和状态含一个零就能归零''，同一证明仍成立，不需要} {\ensuremath{A_{p}(m)}} {的全部输入结论。}

\subsection*{\texorpdfstring{{C6.3} {一个对所有} {\ensuremath{n\ge 2p+1}} {的统一造零收尾定理}}{C6.3 一个对所有  的统一造零收尾定理}}\label{C6-h35}
\addcontentsline{toc}{subsection}{{C6.3} {一个对所有} {\ensuremath{n\ge 2p+1}} {的统一造零收尾定理}}

\textbf{{定理（零平均触发收尾）。}} {对每个奇素数} {\ensuremath{p}、每个} {\ensuremath{n\ge 2p+1}，只要当前状态中有一个真实} {\ensuremath{p}} {元组的和为零，就能经有限次} {\ensuremath{p}} {平均把整个状态归零。}

{先平均该} {\ensuremath{p}} {元组，产生} {\ensuremath{p}} {个零。然后按} {\ensuremath{n}} {分四段。}

\paragraph*{\texorpdfstring{{3.1} {\ensuremath{2p+1\le n\le 3p}}}{3.1 }}\label{C6-h36}
\addcontentsline{toc}{paragraph}{{3.1} {\ensuremath{2p+1\le n\le 3p}}}

{已有不依赖任何端点猜想的事实：2p+1} {元零和状态只要有一个零，就能至多五步归零。} {证明是保留这个零，把另外} {2p} {个位置分成两个} {\ensuremath{p}} {元块并平均，得到}

\[
(a^p,(-a)^p,0),
\]

{再用两次各消去} {(p-1)/2} {对相反值，以及最后一对，三步收尾。}

{产生} {\ensuremath{p}} {个零后，当前非零支持数至多} {\ensuremath{n-p\le 2p}。} {所以可选出} {2p+1} {个位置，包含全部非零坐标和至少一个零，调用这个显式尾部。} {从第一次零平均开始，合计至多六次原子操作完成。}

\textbf{{这里不需要证明} {\ensuremath{A_{p}(2p+1)}，也不要求} {2p+1} {为素数。}}

\paragraph*{\texorpdfstring{{3.2} {\ensuremath{3p\le n\le 4p-1}}}{3.2 }}\label{C6-h37}
\addcontentsline{toc}{paragraph}{{3.2} {\ensuremath{3p\le n\le 4p-1}}}

{非零支持数至多} {\ensuremath{n-p\le 3p-1}。} {使用第} {2} {节引理以及已经证明的} {\ensuremath{A_{p}(3p)}。} {这是新的三块运输控制器对整个上半区间的实际用途。}

\paragraph*{\texorpdfstring{{3.3} {\ensuremath{4p\le n\le 5p-1}}}{3.3 }}\label{C6-h38}
\addcontentsline{toc}{paragraph}{{3.3} {\ensuremath{4p\le n\le 5p-1}}}

{已有} {\ensuremath{A_{p}(4p)}：因为} {\ensuremath{d_{p}(4p)=1}，通用条件}

\[
4p\ge4p-1+d_p(4p)
\]

{恰好满足。非零支持数至多} {\ensuremath{n-p\le 4p-1}，再用第} {2} {节引理。}

\paragraph*{\texorpdfstring{{3.4} {\ensuremath{n\ge 5p}}}{3.4 }}\label{C6-h39}
\addcontentsline{toc}{paragraph}{{3.4} {\ensuremath{n\ge 5p}}}

{通用上界} {\ensuremath{N(p)\le 4p+\operatorname{ceil}(log_{2}}} {\ensuremath{p)+1\le 5p}} {已覆盖该区间。} {当前有零坐标，中心化本原后的} {\ensuremath{G=1}，故直接调用该定理。}

{四段覆盖所有} {\ensuremath{n\ge 2p+1}，证明完成。}


\section{安全的双块入口}\label{construct:entry}
\subsection*{\texorpdfstring{{C7.2} {只要求在有限环上取逆}}{C7.2 只要求在有限环上取逆}}\label{C7-h40}
\addcontentsline{toc}{subsection}{{C7.2} {只要求在有限环上取逆}}

{先证明一个稍一般的入口引理。设整数} \(p\ge2\){，} \[
n=2p+r,\qquad1\le r<p,\qquad\gcd(p,n)=1.
\] {记} \(\mathcal P\) {为} \(n\) {的不同素因子集合，}\(d=|\mathcal P|\){，以及} \[
M=\operatorname{rad}(n)=\prod_{\ell\mid n}\ell,\qquad R=\mathbb Z[1/p].
\] {计算始终允许真实} \(p\) {平均产生的} \(R\) {分母。称状态有完整见证，如果它在每个} \(\ell\mid n\) {下不全同余。合法本原整数输入具有此性质。}

{对已具有一个等值} \(p\) {块的状态，写成} \[
(a^p,c_1,\ldots,c_m),\qquad m=n-p=p+r.
\] {选单点角色} \(i\){，平均该单点与块中的} \(p-1\) {个位置，留下最后一份旧块值作为新的第} \(i\) {单点。真实变化是} \[
(a,c_i)\longmapsto\left(\frac{(p-1)a+c_i}{p},a\right).\tag{C7.2}
\] {所有其他单点位置不变。被留下的旧块位置接替单点角色，故每一步都有原位置账本，不需要搬动数值或增加副本。}

{式} {(C7.2)} {的矩阵及形式逆为} \[
E=\begin{pmatrix}(p-1)/p&1/p\\1&0\end{pmatrix},
\qquad E^{-1}=\begin{pmatrix}0&1\\p&-(p-1)\end{pmatrix}.
\] {二者在每个} \(\ell\mid n\) {下可逆并固定常值向量。因此每个真实交换自动保持完整见证。形式逆在这里首先只用于证明安全性。}

{令} \(\lambda=-1/p\){。对每个} \(\ell\mid n\){，该交换模} \(\ell\) {的阶可取} \[
o_\ell=
\begin{cases}
\operatorname{ord}_\ell(\lambda),&p\not\equiv-1\pmod\ell,\\
\ell,&p\equiv-1\pmod\ell.
\end{cases}
\] {第一种情况有两个不同特征值} \(1,\lambda\){。第二种情况} \(E=I+N\){、}\(N^2=0\){，故} \(E^\ell=I\){。这包括} \(\ell=2\) {的幺幂情况。}

{取} \[
L=\operatorname{lcm}_{\ell\mid n}o_\ell.
\] {则同一角色连续交换} \(L\) {次在模} \(M\) {上是恒等，交换} \(L-1\) {次在模} \(M\) {上是逆。全部都是正向真实平均。}

{这不是有理数值上的逆。例如} \(p=3\){、}\(M=7\){，两次交换将} \((a,c)=(0,1)\) {变成} \((2/9,1/3)\){，而形式逆给} \((1,-2)\){。两者模7相同，实际数值不同。入口只要求最终同余见证正确，因此这种逆恰好足够。}

\subsection*{\texorpdfstring{{C7.3} {三个单点给出显式见证转移}}{C7.3 三个单点给出显式见证转移}}\label{C7-h41}
\addcontentsline{toc}{subsection}{{C7.3} {三个单点给出显式见证转移}}

{令} \[
\delta_i=c_i-a,\qquad \delta=(\delta_1,\ldots,\delta_m)^T.
\] {在} \(R^m\) {上，式} {(C7.2)} {的实际差分矩阵为} \[
A_i=I+v_i e_i^T,\qquad v_i=\lambda\mathbf1-e_i.
\tag{C7.3}
\] {其第} \(i\) {差分乘} \(\lambda\){，其他差分加} \(\lambda\delta_i\){。有} \[
e_i^Tv_i=\lambda-1,\qquad
A_i^{-1}=I-\lambda^{-1}v_ie_i^T.
\]

{取三个互异角色} \(i,j,k\){。由于} \(m=p+r\ge3\){，这些角色始终存在。令} \[
T=A_k^{-1}A_jA_k,\quad w=e_k-e_j,\quad h=e_j^T+\lambda e_k^T.
\] {直接计算：} \[
A_k^{-1}v_j=v_j-v_k=w,\qquad e_j^TA_k=h,\qquad T=I+wh.
\] {又} \(e_i^Tw=0\){、}\(hv_i=\lambda(1+\lambda)\){。因此} \[
\boxed{[T,A_i]:=TA_iT^{-1}A_i^{-1}
=I+(1+\lambda)(e_k-e_j)e_i^T.}\tag{C7.4}
\] {其中} \(1+\lambda=(p-1)/p\){。证明可在整个} \(R^m\) {上完成：共轭将} \(v_i\) {改成} \(v_i+\lambda(1+\lambda)w\){，再右乘} \(A_i^{-1}\){，系数恰除去} \(\lambda\){。}

{把} {(4)} {中每个形式逆换成第2节的正向} \(L-1\) {次幂，得到实际词} \(W_{ijk}\){。按时间顺序执行的八段为} \[
A_i^{L-1},\ A_k,\ A_j^{L-1},\ A_k^{L-1},
\ A_i,\ A_k,\ A_j,\ A_k^{L-1}.\tag{C7.5}
\] {总计} \(4L\) {个原子交换。这个词满足模} \(M\) {的精确恒等式} \[
W_{ijk}\equiv I+c(e_k-e_j)e_i^T,\qquad c=(p-1)/p.
\tag{C7.6}
\] {没有宣称} {(5)} {在有理数上等于} {(4)。}

{因} \(((e_k-e_j)e_i^T)^2=0\){，重复该真实词} \(t\) {次给出} \[
\delta_j\mapsto\delta_j-ct\delta_i,\qquad
\delta_k\mapsto\delta_k+ct\delta_i,\qquad
\delta_h\mapsto\delta_h\ (h\ne j,k)
\pmod M.\tag{C7.7}
\] {这是正向路径的单个分支，不是矩阵线性组合或分支求和。}

\subsection*{\texorpdfstring{{C7.4} {失效素数恰好能由保留位置承接}}{C7.4 失效素数恰好能由保留位置承接}}\label{C7-h42}
\addcontentsline{toc}{subsection}{{C7.4} {失效素数恰好能由保留位置承接}}

{定义} \[
\mathcal B=\{\ell\mid n:\ell\mid p-1\},\qquad b=|\mathcal B|.
\] {对} \(\ell\notin\mathcal B\){，}\(c\) {为模} \(\ell\) {单位，式} {(C7.7)} {可以转移一个非零见证。对} \(\ell\in\mathcal B\){，式} {(C7.7)} {在全部差分上为恒等。}

{由于} \[
\gcd(p-1,n)=\gcd(p-1,r+2),
\] {所有失效素数都整除} \(r+2\){。特别地} \[
\boxed{b\le r.}\tag{C7.8}
\] {否则} \(r+2\ge\prod_{\ell\in\mathcal B}\ell\ge2^{r+1}>r+2\){，矛盾。}

{这正是容量与算术的配合：可以为每个失效素数保留一个单点见证，其总数不超过实际剩余位置} \(r\){。其他素数由} {(7)} {在这些保护不变的前提下逐一处理。}

{在端点} \(r=1\){，失效素数至多只有3，因为} \[
\gcd(p-1,2p+1)=\gcd(p-1,3).
\] {这解释了为什么一个保留位置就足够；不涉及} \(2p+1\) {有多少个其他素因子。}

{更具体地，端点模} \(M\) {有} \(\lambda=2\){，所以交换子的系数恰为3，式} {(C7.7)} {是在两个目标单点之间转移} \(3t\delta_i\){。原来依赖} \(2^k\) {的同步避让由一个明确的加法剪切取代。}

{这也解释该入口机制在} \(2p\) {与} \(2p+1\) {的区别。对奇素数} \(p\){，若形式上取} \(r=0,n=2p\){，则失效素数为2，而保留位置为零，式} {(C7.8)} {不再成立；非零相反双块} \((a^p,(-a)^p)\) {正有不可消除的模2障碍。增加到} \(r=1\) {后，失效素数至多为3且已有一个位置承接。这个说明与已证的} \(2p\) {下界相容，不是说所有} \(2p\) {元输入都不可达。}

\subsection*{\texorpdfstring{{C7.5} {无条件建立两个合法} \(p\) {块}}{C7.5 无条件建立两个合法 p 块}}\label{C7-h43}
\addcontentsline{toc}{subsection}{{C7.5} {无条件建立两个合法} \(p\) {块}}

\textbf{{两块入口引理。}} {在第2节的假设下，每个合法输入都能有限步变成具有完整见证的} \[
(a^p,b^p,c_1,\ldots,c_r).
\tag{C7.9}
\]

{先固定原位置0。对每个} \(\ell\mid n\){，保护位置0及另一个与它不同余的位置，总计至多} \(d+1\) {个位置。保护集外足够选出} \(p\) {个位置：若} \(d\ge p+r\){，则} \[
n\ge2^d\ge2^{p+r}>2(p+r)\ge2p+r=n,
\] {矛盾。所以} \(d+1\le p+r=n-p\){。在保护集外平均} \(p\) {个位置，得到一个} \(p\) {块与} \(m=p+r\) {个单点，完整见证由未选位置保留。}

{对每个} \(\ell\in\mathcal B\){，选一个} \(\delta_i\ne0\pmod\ell\) {的单点角色，并把它列入保留集合} \(J\){。完整见证保证这样的单点存在。由} {(8)，可补足至恰好} \(|J|=r\){。从} \(J\) {中固定一个目标角色} \(j\){。}

{假设目标在某个} \(\ell\notin\mathcal B\) {处缺失见证。选择} \(\delta_i\ne0\pmod\ell\) {的角色} \(i\ne j\){，并选任意} \(k\notin\{i,j\}\){。令} \(Q=\prod_{q\in\mathcal P\setminus\mathcal B}q\){。CRT} {给出正整数} \(t<Q\){，使} \[
t\equiv1\pmod\ell,\qquad
t\equiv0\pmod q\quad(q\in\mathcal P\setminus(\mathcal B\cup\{\ell\})).
\] {执行} {(5)} {共} \(t\) {次。式} {(C7.7)} {给目标新增非零差分} \(-c\delta_i\pmod\ell\){，其他非失效素数处的目标差分不变；所有失效素数处的全部差分都不变。此前保留在} \(J\) {中的见证因而仍在。}

{至多} \(d-b\) {个阶段后，目标} \(j\) {在全部非失效素数处都非零，而} \(J\) {在每个失效素数处仍有见证。}

{最后保留} \(J\) {的} \(r\) {个单点，平均其他恰好} \(p\) {个单点。未选的旧} \(p\) {块和} \(J\) {在每个} \(\ell\mid n\) {下都不全同余，所以最后一步安全，得到} {(9)。全过程在原位置执行，至此引理得证。}

{每个中间交换都自动安全；保留的是单点角色，原位置随} {(2)} {更新。不能把角色} \(J\) {误解为始终冻结同一批物理位置。}

\paragraph*{\texorpdfstring{{算法可以比逐素数CRT幂短得多}}{算法可以比逐素数CRT幂短得多}}\label{C7-h44}
\addcontentsline{toc}{paragraph}{{算法可以比逐素数CRT幂短得多}}

{CRT只用于给最直接的存在证明。实际选择} \(t\) {时，可令} \[
\delta_j-ct\delta_i\ne0\pmod q
\] {同时覆盖} \(i,j\) {在所有非失效素数处的支持并集。每个相关素数只禁止一个} \(t\) {类，因此可调用}第\ref{construct:short-choice}节{。}

{对素数} \(p\ge5\){、}\(n<3p\){，筛选在} \(1\le t\le p-1\) {内完成。每阶段严格增加目标支持，得到两块入口的显式上界} \[
2+4L(p-1)(d-b).\tag{C7.10}
\] {这只是有限性与实际算法改进，不是位长多项式或最短路径保证。程序对大参数只指数化固定四角色矩阵或保存重复词，不展开巨大重复数组。}


\section{短配比与上半带三值压缩}\label{construct:short-choice}
\subsection*{\texorpdfstring{{C8.1} {合法性与原位置约定}}{C8.1 合法性与原位置约定}}\label{C8-h45}
\addcontentsline{toc}{subsection}{{C8.1} {合法性与原位置约定}}

{取本原整数零和输入} \(X=(x_1,\ldots,x_n)\){，并记} \[
G(X)=\gcd_{i,j}(x_i-x_j).
\] {因为} \(G(X)\mid n\){，而} \(p\nmid n\){，本区间的允许条件就是} \(G(X)=1\){。全零输入直接终止。}

{令} \(\mathcal P\) {为} \(n\) {的不同素因子集合，}\(d=|\mathcal P|\){。在} \(\mathbb Z[1/p]\) {中，称一个状态具有完整见证，如果对每个} \(\ell\in\mathcal P\){，其坐标在} \(\mathbb F_\ell\) {中不全相等。初始输入具有此性质；下述构造逐步保持这个略强于归一化合法性的性质。因此在任何一步清分母、取本原代表后，仍满足} \(G=1\){。}

{两个基本事实反复使用：}

\begin{enumerate}
\def\labelenumi{\arabic{enumi}.}
\tightlist
\item
  {若未选位置已经在某个} \(\ell\) {处不同余，该次平均不能破坏此处见证。}
\item
  {一旦某状态在} \(\ell\ne p\) {处全同余，之后任何} \(p\) {平均都保持该同余。因此，若一个完整的真实操作段的输出具有完整见证，那么中间状态也都具有完整见证。}
\end{enumerate}

{块均为带原位置标签的等值集合。不同标签块允许恰好取相同数值，不要求它们在实数上互异。取本原代表只用于分析，不是额外执行的缩放。}

\subsection*{\texorpdfstring{{C8.2} {足够短的避同余参数}}{C8.2 足够短的避同余参数}}\label{C8-h46}
\addcontentsline{toc}{subsection}{{C8.2} {足够短的避同余参数}}

\paragraph*{\texorpdfstring{{引理} {1：原容量内的筛选}}{引理 1：原容量内的筛选}}\label{C8-h47}
\addcontentsline{toc}{paragraph}{{引理} {1：原容量内的筛选}}

{设每个} \(\ell\in\mathcal Q\subseteq\mathcal P\) {给定至多一个禁止余数类。则存在} \[
1\le k\le p-1
\tag{C8.2}
\] {同时避开所有禁止类。}

\textbf{{证明。}} {写} \(h=|\mathcal Q|\){、}\(M=\prod_{\ell\in\mathcal Q}\ell\){、}\(\rho=\varphi(M)/M\){。}\(h=0\) {取} \(k=1\){。}

{先处理} \(1\le h\le4\){。长度} \(2^h\) {的区间不能被这些类覆盖。如果素数} {2} {在其中，先保留未被它禁止的一个奇偶类，剩下} \(2^{h-1}\) {项。其余奇素数在该等差数列中分别最多删除} \(\lceil 2^{h-1}/\ell\rceil\) {项。若} {2} {不在其中，直接使用} \(\lceil2^h/\ell\rceil\){。最坏的删除数如下：}

{ % do not increment counter
\begin{longtable}[]{@{}lll@{}}
\toprule\noalign{}
\(h\) & {含} {2：剩余项数／其余类至多删除} & {不含} {2：区间长度／至多删除} \\
\midrule\noalign{}
\endhead
\bottomrule\noalign{}
\endlastfoot
{1} & {1／0} & {2／1} \\
{2} & {2／1} & {4／}\(2+1=3\) \\
{3} & {4／}\(2+1=3\) & {8／}\(3+2+2=7\) \\
{4} & {8／}\(3+2+2=7\) & {16／}\(6+4+3+2=15\) \\
\end{longtable}
}

{这些长度均不超过} \(p-1\){：}\(h\le2\) {使用} \(p\ge5\){；}\(h=3\) {使用} \(30\le M\le n<4p\){，故素数} \(p\ge11\){；}\(h=4\) {同理由} \(210\le n<4p\) {得} \(p\ge53\){。}

{若} \(h\ge5\){，中国剩余定理将所有禁止类写成同一个整数} \(b\) {在各素数下的余数。长度} \(L\) {区间内避开它们的数目，是某个平移区间内与} \(M\) {互素的整数数目。容斥给出} \[
\#\{1\le k\le L:\gcd(k-b,M)=1\}
>L\rho-(2^h-1).
\tag{C8.3}
\] {这里每个非空因子项的计数与} \(L/D\) {相差严格小于} {1，共有} \(2^h-1\) {项。因此可取} \[
L=\left\lfloor\frac{2^h-1}{\rho}\right\rfloor+1.
\tag{C8.4}
\] {它确实不超过} \(p-1\){。当} \(h=5\) {时，}\(\varphi(M)\ge(2-1)(3-1)(5-1)(7-1)(11-1)=480>2^{h+2}\){；每增加一个素因子，最小可能的} \(\varphi(M)\) {增长超过两倍，故该不等式对所有} \(h\ge5\) {成立。于是} \[
(p-1)\rho>
\left(\frac n4-1\right)\rho
\ge\frac{\varphi(M)}4-\rho
>2^h-1.
\] {结合整数性即得} \(L\le p-1\){。证毕。}

{这是} {Jacobsthal} {短区间问题的一个足够版本。可用一个长度小于} \(p\) {的筛表直接选出} \(k\){，每个素数只标记一个等差数列。这里不依赖一般} {Jacobsthal} {函数的更强上界。}

\paragraph*{\texorpdfstring{{引理} {2：两个} {\ensuremath{p}} {块的可变配比}}{引理 2：两个  块的可变配比}}\label{C8-h48}
\addcontentsline{toc}{paragraph}{{引理} {2：两个} {\ensuremath{p}} {块的可变配比}}

{对两个带标签的} \(p\) {块} \(a^p,b^p\){，分别平均两个不交组} \[
(a^{p-k},b^k),\qquad(a^k,b^{p-k}).
\] {得到} \[
a_k=a+\frac{k}{p}(b-a),\qquad
b_k=b+\frac{k}{p}(a-b).
\tag{C8.5}
\] {对一个固定参照值} \(t\){，若} \(a-t,b-t\) {在某个} \(\ell\) {处不同时为零，则} \(a_k-t\) {或恒非零，或只在一个} \(k\bmod\ell\) {类上为零。因此引理} {1} {允许让} \(a_k-t\) {同时保留这两个差的全部素数支持。}

{关键区别在于，}\(k\) {是分配给真实位置的副本数。它的局部禁止类以互素的素数为模；不再是交换次数的乘法阶同余类。后者可能形成完全覆盖系统，前者在本区间的容量内已由引理} {1} {排除。}

\subsection*{\texorpdfstring{{C8.3} {两块与固定单点：任意多见证的集中}}{C8.3 两块与固定单点：任意多见证的集中}}\label{C8-h49}
\addcontentsline{toc}{subsection}{{C8.3} {两块与固定单点：任意多见证的集中}}

\paragraph*{\texorpdfstring{{引理} {3：固定单点的完整见证}}{引理 3：固定单点的完整见证}}\label{C8-h50}
\addcontentsline{toc}{paragraph}{{引理} {3：固定单点的完整见证}}

{设已有两个} \(p\) {块、若干单点，整体具有完整见证。指定其中一个单点} \(w\){。通过至多} \(3d\) {次平均，可以保持} \(w\) {的值和位置不动，使其中一个块} \(a^p\) {满足} \[
a-w\ne0\pmod\ell\qquad(\ell\in\mathcal P).
\tag{C8.6}
\]

\textbf{{证明。}} {当前已获支持为} \(S=\{\ell:a-w\ne0\}\){。若它未覆盖全部素数，取一个缺失的} \(\ell_0\){。}

{若第二块} \(b\) {在} \(\ell_0\) {处已不同于} \(w\){，直接应用引理} {2} {混合两个块，使第一块与} \(w\) {的差覆盖两个原差的并集。}

{否则} \(a=b=w\pmod{\ell_0}\){。整体合法性保证另有单点} \(x\ne w\pmod{\ell_0}\){。平均} \[
(b^{p-1},x),
\] {使新块值为} \(((p-1)b+x)/p\){，并把未选的一份旧} \(b\) {作为新单点。这是带位置标签的块---单点交换。第一块与} \(w\) {未变，所以全部旧支持保留；第二块现在在} \(\ell_0\) {处不同于} \(w\){。随后再用两次混合平均集中支持。}

{交换本身也不会损失整体见证：其二元变换} \[
(b,x)\mapsto\left(\frac{(p-1)b+x}{p},b\right)
\] {在每个} \(\ell\ne p\) {处可逆，并保持常数向量。因此它把一个整体全同余状态拉回后仍全同余。}

{每个阶段至多三次平均，并严格增加} \(S\){，所以至多} \(d\) {阶段完成。证毕。}

{这里没有周期等待，也没有求某个局部矩阵的正向逆。}

\subsection*{\texorpdfstring{{C8.4} {整个上半带的三块入口}}{C8.4 整个上半带的三块入口}}\label{C8-h51}
\addcontentsline{toc}{subsection}{{C8.4} {整个上半带的三块入口}}

\paragraph*{\texorpdfstring{{第一步：建立两个块}}{第一步：建立两个块}}\label{C8-h52}
\addcontentsline{toc}{paragraph}{{第一步：建立两个块}}

{固定一个位置} \(w\){，对每个} \(\ell\mid n\) {再保护一个与它不同余的位置。总共至多} \(d+1\) {个位置。}

{其余位置足够建立两个不交} \(p\) {块。事实上} \(d\ge p\) {会推出} \(n\ge2^d\ge2^p>4p\){，矛盾；所以} \(d\le p-1\){，特别地} \(n-(d+1)\ge2p\){。在保护集之外平均两个不交} \(p\) {组，耗费两次操作，全部见证由保护集保留。}

{应用引理} {3，得到第一块} \(a^p\) {与固定单点} \(w\) {的完整差见证。}

\paragraph*{\texorpdfstring{{第二步：建立第三块}}{第二步：建立第三块}}\label{C8-h53}
\addcontentsline{toc}{paragraph}{{第二步：建立第三块}}

{两块之外尚有} \(p+r\) {个单点。避开} \(w\){，任选} \(p\) {个进行平均；因为} \(r\ge1\){，容量足够。第一块与} \(w\) {完整保留，故操作安全。得到} \[
(a^p,b^p,c^p,x_1,\ldots,x_r),
\tag{C8.7}
\] {其中} \(w\) {仍为一个单点，且} \(a-w\) {在全部危险素数处非零。}

\paragraph*{\texorpdfstring{{第三步：使三个主块本身保留全部见证}}{第三步：使三个主块本身保留全部见证}}\label{C8-h54}
\addcontentsline{toc}{paragraph}{{第三步：使三个主块本身保留全部见证}}

{称三个主块在} \(\ell\) {处不全相等时，该素数已好。若所有素数已好，无需操作。}

{否则混合} \(b,c\){，选择配比，使新的} \(c_k-a\) {在全部已好素数处非零。引理} {2} {适用，因为这些地方} \(b-a,c-a\) {不同时为零。再交换} \(b_k\) {块与} \(w\){。}

{在旧的好素数处，未动的} \(a,c_k\) {保留见证。在旧的坏素数处，}\(a=b=c\){，而} \(w\ne a\){，所以新} \(b\) {块不再等于} \(a\){。}\textbf{{一次交换同时修复所有原坏素数。}} {此段至多三次平均。}

\paragraph*{\texorpdfstring{{第四步：合并全部单点}}{第四步：合并全部单点}}\label{C8-h55}
\addcontentsline{toc}{paragraph}{{第四步：合并全部单点}}

{平均} \[
(a^{p-r},x_1,\ldots,x_r).
\] {它留下} \(r\) {份旧} \(a\){，以及完整的} \(b,c\) {块。三个旧主块在每个危险素数处不全相等，所以未选位置仍提供完整见证。输出为合法} \[
(a_1^p,b^p,c^p,d_0^r),\qquad d_0=a.
\tag{C8.8}
\]

{至此共至多} \(7+3d\) {次平均，且完全没有使用} \(r\ge d\){。}

\subsection*{\texorpdfstring{{C8.5} {三行折成两行：真正的降秩}}{C8.5 三行折成两行：真正的降秩}}\label{C8-h56}
\addcontentsline{toc}{subsection}{{C8.5} {三行折成两行：真正的降秩}}

\paragraph*{\texorpdfstring{{引理} {4：五次平均制造} {2p} {重块}}{引理 4：五次平均制造 2p 重块}}\label{C8-h57}
\addcontentsline{toc}{paragraph}{{引理} {4：五次平均制造} {2p} {重块}}

{任意合法零和} \((a^p,b^p,c^p,d_0^r)\) {都能在至多五次平均后进入} {(1)，并满足} \[
A-C\ne0\pmod\ell\quad(\ell\mid n).
\tag{C8.9}
\]

\textbf{{证明。}} {由零和关系，在每个} \(\ell\mid n\) {处有} \[
(a-d_0)+(b-d_0)+(c-d_0)=0.
\tag{C8.10}
\] {所以} \(a-d_0,b-d_0\) {不可能同时为零，否则三者全为零，违反合法性。应用引理} {2，先用两次平均使新的} \(a-d_0\) {在所有危险素数处非零。}

{写} \[
p=3m+\varepsilon,\qquad \varepsilon\in\{1,-1\},\qquad m\ge2.
\] {在三个} \(p\) {块中，分别拿出两份具有同样组成的} \(p\) {组，其三种类型的数量均为} \[
(m+\varepsilon,m,m).
\tag{C8.11}
\] {每种类型所需总量不超过} \(p\){：}\(2(m+\varepsilon)\le p\){、}\(2m\le p\){。平均这两个不交组，再平均三个主块的全部剩余位置；剩余也恰有} \(p\) {个。载体块} \(d_0^r\) {始终未动。}

{最终值为} \[
A=\frac{m(a+b+c)+\varepsilon a}{p},\qquad
B=a+b+c-2A,\qquad C=d_0.
\tag{C8.12}
\] {前两组均取值} \(A\){，所以产生} \(2p\) {份相同值。由} {(10)，模每个} \(\ell\mid n\) {有} \[
A-d_0=\frac\varepsilon p(a-d_0)\ne0,\qquad
B-d_0=-2(A-d_0).
\tag{C8.13}
\] {即使} \(\ell=2\){，第一条也保留见证。最终状态合法，从而三次平均之间亦合法。证毕。}

{结合第} {4} {节，得到全文的} \(12+3d\) {次归约定理。它减少了真实独立值参数：四值零和空间为三维；(1)} {为二维。这里没有使用抽象运输矩阵的副本实现定理。}

\subsection*{\texorpdfstring{{C8.6} {归约后的精确算术核心}}{C8.6 归约后的精确算术核心}}\label{C8-h58}
\addcontentsline{toc}{subsection}{{C8.6} {归约后的精确算术核心}}

{取} {(1)} {的本原整数代表。因} \(\gcd(p,r)=1\){，可唯一写成} \[
A=y-pz,\qquad B=-2y+(2p+r)z,\qquad C=-pz,
\tag{C8.14}
\] {其中} \(y,z\in\mathbb Z\){、}\(\gcd(y,z)=1\){。反过来该表示给出全部本原整数核心。利用} \(A-C=y\){、}\(B-C=-2y+nz\){，得到} \[
\boxed{G=\gcd(y,n).}
\tag{C8.15}
\] {所以整个上半带的剩余问题是：对每个} \(\gcd(y,z)=\gcd(y,n)=1\){，从} {(14)} {的真实位置状态到达终点。}

{平方能量在这一核心上恰为} \[
\mathcal E
=p\bigl(6y^2-4nyz+n(2p+r)z^2\bigr)
=\frac p3\bigl(2(3y-nz)^2+rnz^2\bigr).
\tag{C8.16}
\] {括号内二元整二次型的判别式为} \(-8nr\){。因此一般} \(r\) {引入的算术量是} \(nr\){，并不只由} \(n\) {决定。这提供了一个明确的二次格结构，但单有此格或相关} {\ensuremath{\theta }} {级数不产生正向路径。}

{零和} \(p\) {子集也有统一方程。取} \(j\) {份} \(B\){、}\(s\) {份} \(C\){，其余} \(p-j-s\) {份取} \(A\){，则终端条件为} \[
\boxed{(p-3j-s)y+(nj-p^2)z=0,\quad
0\le s\le r,\quad 0\le j\le p-s.}
\tag{C8.17}
\] {一旦命中，真实平均该组就得到} \(p\) {个零；由已有零触发收尾定理完成全系统。}


\section{加权核心及两轮见证修复}\label{construct:core-entry}
\subsection*{\texorpdfstring{{C9.1} {原问题只有一种合法性，两条中间带使用同一个格}}{C9.1 原问题只有一种合法性，两条中间带使用同一个格}}\label{C9-h59}
\addcontentsline{toc}{subsection}{{C9.1} {原问题只有一种合法性，两条中间带使用同一个格}}

{取素数} {\ensuremath{p\ge 5}，令} \[
m\in\{1,2\},\qquad 1\le r<p,\qquad n=(m+1)p+r.
\] {\ensuremath{m=1}} {对应} {\ensuremath{2p<n<3p}；\ensuremath{m=2}} {对应} {\ensuremath{3p<n<4p}。两条带的核心统一写为} \[
\boxed{\mathcal K_{m,p,r}(a,z)=
\bigl(a^{mp},(-ma+rz)^p,(-pz)^r\bigr).}
\tag{C9.1}
\] {记三种块值为} {\ensuremath{A}、\ensuremath{B}、\ensuremath{C}。因为} {\ensuremath{\operatorname{gcd}(p,r)=1}，每个整数三块零和状态都唯一对应整数} {a,z，且} \[
\gcd(A,B,C)=\gcd(a,z).
\] {令} \[
y=a+pz.
\] {两种差为} \[
A-C=y,\qquad A-B=(m+1)y-nz.
\] {本原性给} {\ensuremath{\operatorname{gcd}(y,z)=1}，故有精确公式} \[
\boxed{G=\gcd(y,n).}
\tag{C9.2}
\] {因此下半带的奇数、偶数核心和上半带核心没有不同的原始合法性：统一都是本原整数列} {(a,z)，且} {\ensuremath{y}} {在每个} {\ensuremath{\ell |n}} {处为单位。}

{这与一般输入进入核心的证明分开。上半带已有无素因子数限制的入口；下半带的改进入口见第} {6} {节。核心内产生一次真实零和} {\ensuremath{p}} {元组以后，已有}第\ref{construct:zero-tail}节{统一收尾。}

\subsection*{\texorpdfstring{{C9.3} {所有重复行返回的一套公式}}{C9.3 所有重复行返回的一套公式}}\label{C9-h60}
\addcontentsline{toc}{subsection}{{C9.3} {所有重复行返回的一套公式}}

{从} {(1)} {出发，每个重复的} {\ensuremath{p}} {元组取} \[
i\text{ 份 }A,\quad j\text{ 份 }B,\quad s\text{ 份 }C,
\qquad i+j+s=p.
\] {做} {\ensuremath{m}} {份组成完全相同、位置互不相交的组，再平均剩余的} {\ensuremath{p}} {个位置，另保留} {\ensuremath{r}} {份同值位置作新} {\ensuremath{C}。总计} {m+1} {次原子平均。}

\paragraph*{\texorpdfstring{{保留旧} {\ensuremath{A}}}{保留旧 }}\label{C9-h61}
\addcontentsline{toc}{paragraph}{{保留旧} {\ensuremath{A}}}

{容量条件为} \[
mi\le mp-r,\qquad mj\le p,\qquad ms\le r.
\] {在} {(a,z)} {上的实际矩阵为} {\ensuremath{p^{-1}}} {乘} \[
\boxed{M^A_{j,s}=
\begin{pmatrix}
p-(m+1)j-s&rj-ps\\
-1&0
\end{pmatrix}.}
\tag{C9.5}
\]

\paragraph*{\texorpdfstring{{保留旧} {\ensuremath{B}}}{保留旧 }}\label{C9-h62}
\addcontentsline{toc}{paragraph}{{保留旧} {\ensuremath{B}}}

{容量条件为} \[
mi\le mp,\qquad mj\le p-r,\qquad ms\le r.
\] {实际矩阵为} {\ensuremath{p^{-1}}} {乘} \[
\boxed{M^B_{j,s}=
\begin{pmatrix}
p-(m+1)j-s&rj-ps\\
m&-r
\end{pmatrix}.}
\tag{C9.6}
\]

\paragraph*{\texorpdfstring{{保留旧} {\ensuremath{C}}}{保留旧 }}\label{C9-h63}
\addcontentsline{toc}{paragraph}{{保留旧} {\ensuremath{C}}}

{取} {\ensuremath{s=0}，\ensuremath{0\le mj\le p}，实际矩阵为} \[
F_j=\frac1p
\begin{pmatrix}
p-(m+1)j&rj\\
0&p
\end{pmatrix}.
\tag{C9.7}
\] {在} {(w,z)} {上，这恰为} \[
w'=\frac{p-(m+1)j}{p}w,\qquad z'=z.
\]

{证明只需对每行所选数值求平均；新} {\ensuremath{C}} {等于所保留的旧值，给出新} {\ensuremath{z}。容量条件保证原位置足够，最后一组有且仅有} {\ensuremath{p}} {个位置。一般} {(5)、(6)} {可能奇异或破坏合法性，不能自动视为可逆宏。}

{旧证明的迹零对合分别是} \[
s=p-(m+1)j\quad\text{和}\quad
s=p-r-(m+1)j.
\tag{C9.8}
\] {对这些可行参数，矩阵平方为非零标量，便得到正向射影逆。它们只是同一正向返回集合中容易取逆的元素。参数区间缩为一个点，说明当前的取逆工具不足，不能推出原位置没有其他可用返回。}

{改到始终幺模的} {(z,y)} {坐标，迹零} {\ensuremath{A}、\ensuremath{B}} {返回模} {\ensuremath{n}} {分别为} \[
\begin{pmatrix}p&-1\\0&-p\end{pmatrix},
\qquad
\begin{pmatrix}p&m\\0&-p\end{pmatrix}.
\] {其乘积射影上给平移} \[
\boxed{x\longmapsto x-\frac{m+1}{p},\qquad x=z/y\pmod n.}
\tag{C9.9}
\] {若} {\ensuremath{\ell |(m+1,n)}，这个特定平移在模} {\ensuremath{\ell }} {上失效。这是对该群生成步骤的准确限制，不是对所有真实平均的限制。}

\subsection*{\texorpdfstring{{C9.4} {统一目标是整族零和终端，而非} {\ensuremath{z=0}}}{C9.4 统一目标是整族零和终端，而非 }}\label{C9-h64}
\addcontentsline{toc}{subsection}{{C9.4} {统一目标是整族零和终端，而非} {\ensuremath{z=0}}}

{任意真实} {\ensuremath{p}} {元组取} {\ensuremath{j}} {份} {\ensuremath{B}、\ensuremath{s}} {份} {\ensuremath{C}，其余} {p-j-s} {份} {\ensuremath{A}，其中} \[
0\le s\le r,\qquad0\le j\le p-s.
\] {其和为} \[
[p-(m+1)j-s]a+(rj-ps)z.
\] {令} \[
q=p-(m+1)j-s.
\] {一个整数终端方向是} \[
(a,z)=(ps-rj,q).
\] {在} {(z,y)} {中即为} \[
\boxed{(z,y)=
\bigl(p-(m+1)j-s,\ p^2-nj\bigr).}
\tag{C9.10}
\] {按需约去两坐标的共同因子即可取本原代表。因为} \[
\gcd(p^2-nj,n)=1,
\] {原列非零，所除共同因子与} {\ensuremath{n}} {互素，故这些方向全部合法。}

{式} {(C9.10)} {有两个直接后果。}

{第一，终端的单位坐标} {\ensuremath{y}} {与} {\ensuremath{s}} {无关。固定} {\ensuremath{j}、改变} {\ensuremath{s}，会在另一个坐标上形成长度} {r+1} {的连续区间。原位置容量因此直接提供了一族不同终端，无需坚持最终载体先变成零。}

{第二，对每个} {\ensuremath{\ell |\operatorname{gcd}(m+1,n)}，若} {\ensuremath{r\ge \ell -1}，取} {\ensuremath{j=0}、\ensuremath{s=0,}\ldots,ell-1，则} \[
\frac zy=\frac{p-s}{p^2}\pmod\ell
\] {遍历全部} {\ensuremath{F_\ell }。因此} {(9)} {失效处的第一层合法方向，全部已有真实终端代表：}

\begin{itemize}
\tightlist
\item
  {下半带的模2，\ensuremath{r\ge 1}} {即满足；}
\item
  {上半带的模3，\ensuremath{r\ge 2}} {即满足。}
\end{itemize}

{这不是高阶轨道覆盖证明。不同素数可能需要不同的} {(j,s)，还须找到同一个全局终端并匹配共同尺度。它证明的是：不能把``模} {m+1} {无法移动到} {\ensuremath{z=0}''误当成整个终端族缺失。}

\subsection*{\texorpdfstring{{C9.6} {新的统一入口定理：两轮取代周期与} {\ensuremath{gcd}} {条件}}{C9.6 新的统一入口定理：两轮取代周期与  条件}}\label{C9-h65}
\addcontentsline{toc}{subsection}{{C9.6} {新的统一入口定理：两轮取代周期与} {\ensuremath{gcd}} {条件}}

\textbf{{定理。}} {对素数} {\ensuremath{p\ge 5}、\ensuremath{n=2p+r}、\ensuremath{1\le r<p}，若} \[
\boxed{r\ge\omega(n),}
\tag{C9.13}
\] {则每个合法本原整数零和输入可在至多} {9+3omega(n)} {次原位置} {\ensuremath{p}} {平均后进入合法} {(p,p,r)} {核。}

{不要求旧的周期密度条件，也不要求} {\ensuremath{\operatorname{gcd}(2p-1,n)=1}。\ensuremath{r}} {的奇偶和} {\ensuremath{n}} {的素因子指数均不受限制。}

\paragraph*{\texorpdfstring{{6.1} {已有两个工具}}{6.1 已有两个工具}}\label{C9-h66}
\addcontentsline{toc}{paragraph}{{6.1} {已有两个工具}}

{第一，\ensuremath{r\ge d}:=omega(n)} {时，已有安全分块引理用两次平均建立两个} {\ensuremath{p}} {块及} {\ensuremath{r}} {个单点。分组可在} {pX} {中选取，再把同一位置用于} {\ensuremath{X}，不额外缩放。}

{第二，}第\ref{construct:short-choice}节{仅需两个} {\ensuremath{p}} {块和单点。指定单点} {\ensuremath{w}，至多3d次平均后可令一个块值} {\ensuremath{a}} {满足} \[
a-w\ne0\pmod\ell\qquad(\ell\mid n).
\tag{C9.14}
\] {其中选取短配比的引理只用} {\ensuremath{p\ge 5}、\ensuremath{n<4p}，故在本下半带仍适用：每个} {\ensuremath{\ell |n}} {禁止至多一个类时，区间} {1,\ldots,p-1} {内总有可选参数。}

\paragraph*{\texorpdfstring{{6.2} {一轮操作及原来失败的系数}}{6.2 一轮操作及原来失败的系数}}\label{C9-h67}
\addcontentsline{toc}{paragraph}{{6.2} {一轮操作及原来失败的系数}}

{定义一轮为：先用两个互补} {\ensuremath{p}} {组混合块} {a,b，} \[
a_k=a+\frac{k}{p}(b-a),\qquad
b_k=b+\frac{k}{p}(a-b),
\] {再平均} {\ensuremath{b_{k}}} {的} {p-1} {份和单点} {\ensuremath{w}，留下一份} {\ensuremath{b_{k}}} {作新单点。三步后的角色值为} \[
a_1=a_k,\quad b_1=((p-1)b_k+w)/p,\quad w_1=b_k.
\] {新块差乘} {\ensuremath{p}} {为} \[
D(k)=p(a_1-b_1)
=pa-(p-1)b-w+\frac{2p-1}{p}k(b-a).
\tag{C9.15}
\] {旧证明要求} {\ensuremath{\operatorname{gcd}(2p-1,n)=1}，正是为了让这条斜率不消失。现在分两轮安排每个素数的约束。}

\paragraph*{\texorpdfstring{{6.3} {第一轮预先修复失效素数}}{6.3 第一轮预先修复失效素数}}\label{C9-h68}
\addcontentsline{toc}{paragraph}{{6.3} {第一轮预先修复失效素数}}

{若} {\ensuremath{\ell }} {不整除2p-1，选择} {\ensuremath{k}} {使} {\ensuremath{D(k)\ne 0}} {模} {\ensuremath{\ell }。若} {\ensuremath{a=b}，常数为} {a-w} {非零；否则斜率非零。因此至多禁止一个} {\ensuremath{k}} {类。}

{若} {\ensuremath{\ell |2p-1}，则} {\ensuremath{\ell }} {为奇数，\ensuremath{p=1/2}} {模} {\ensuremath{\ell }。不要强求本轮两块已不同，而要求下一轮的常数} \[
E(k):=p a_1-(p-1)b_1-w_1
=w+\frac{a-3b}{2}+4k(b-a)\pmod\ell
\tag{C9.16}
\] {非零。若} {\ensuremath{a=b}，常数为} {w-a} {非零；否则斜率4(b-a)非零。仍至多禁止一个} {\ensuremath{k}} {类。}

{因此同一个短配比筛选选出} {\ensuremath{k}，同时满足全部素数的本轮或下一轮目标。完整第一轮输出在每个素数处不全相同：普通素数由} {\ensuremath{a_{1}\ne b_{1}}} {保证；失效素数由} {\ensuremath{E(k)\ne 0}} {保证。若输出三角色同为} {\ensuremath{c}，则} {\ensuremath{E(k)=pc-(p-1)c-c=0}，矛盾。}

{全同余状态在} {\ensuremath{p}} {平均下不可能重新产生差异，所以完整输出非全同余也保证本轮每个中间原子安全。}

\paragraph*{\texorpdfstring{{6.4} {第二轮同时得到两个块的完整单位差}}{6.4 第二轮同时得到两个块的完整单位差}}\label{C9-h69}
\addcontentsline{toc}{paragraph}{{6.4} {第二轮同时得到两个块的完整单位差}}

{对} {(\ensuremath{a_{1},b_{1},w_{1})}} {再做一轮，参数} {\ensuremath{h}} {对应的块差乘} {\ensuremath{p}} {为} \[
E(k)+\frac{2p-1}{p}h(b_1-a_1).
\tag{C9.17}
\] {对失效素数，斜率为零，但} {E(k)} {已在第一轮修成非零；对普通素数，第一轮给} {\ensuremath{a_{1}\ne b_{1}}，因此斜率非零，只禁止一个} {\ensuremath{h}} {类。再次短配比筛选，得到全部危险素数处均非零的块差。}

{最后平均第一块的} {p-r} {份和全部} {\ensuremath{r}} {个单点，留下} {\ensuremath{r}} {份旧第一块及完整第二块。未选位置已保护全部素数，输出为合法} {(p,p,r)} {核。次数为} \[
2+3d+3+3+1=9+3d.
\] {定理得证。}

{例如} {\ensuremath{p=5}，已有核心外入口角色} {\ensuremath{a=1,b=3,w=2}，在模3下正是旧斜率消失且常数为零的中点情形。两轮仍可完成；核验保留了对应的12元原位置例子。但这只解决进入核心，不宣称五平均12元全部可达。}

\subsection*{\texorpdfstring{{C9.10} {一般偶数核心的二进深度可以统一降低到一}}{C9.10 一般偶数核心的二进深度可以统一降低到一}}\label{C9-h70}
\addcontentsline{toc}{subsection}{{C9.10} {一般偶数核心的二进深度可以统一降低到一}}

\textbf{{定理。}} {设} {\ensuremath{p\ge 5}} {为素数，\ensuremath{r}} {为偶数且} {\ensuremath{2\le r<p}，\ensuremath{n=2p+r}。对任意合法本原核心} \[
(a^p,(rz-a)^p,(-pz)^r),\qquad
\gcd(a,z)=\gcd(a+pz,n)=1,
\] {若} {\ensuremath{az=0}，已有真实零触发收尾。否则令} \[
\nu=v_2(az)\ge1.
\] {存在至多4(nu-1)次原位置} {\ensuremath{p}} {平均，把核心送到本原参数满足} \[
\boxed{v_2(a'z')=1}
\tag{C9.19}
\] {的合法核心。\ensuremath{r=2}及\ensuremath{r=p-1}同样包含在内。}

\paragraph*{\texorpdfstring{{10.1} {先确定偶数坐标的角色}}{10.1 先确定偶数坐标的角色}}\label{C9-h71}
\addcontentsline{toc}{paragraph}{{10.1} {先确定偶数坐标的角色}}

{合法性在模2下使} {a,z} {奇偶相反。若深度} {\ensuremath{\nu \ge 2}} {且} {\ensuremath{a}} {为奇数，则先执行本构造} {(8)} {的已有对合，取} {\ensuremath{s=r-1}、\ensuremath{j=(p-r+1)/2}。它将} \[
(a,z)\longmapsto p^{-1}(d z,-a),\qquad
d=\frac{pr-n(r-1)}2.
\] {\ensuremath{d}} {为奇数且与} {\ensuremath{n}} {互素。取本原代表后，偶坐标位于} {\ensuremath{a}，深度} {\ensuremath{\nu }} {不变；共花两次原子平均。现在可以假设} \[
4\mid a,\qquad z\text{ 为奇数}.
\tag{C9.20}
\] {这一步不是免费交换参数，而是已经有原位置分组的对合。}

\paragraph*{\texorpdfstring{{10.2} {一个行列式恰含一个二因子的真实返回}}{10.2 一个行列式恰含一个二因子的真实返回}}\label{C9-h72}
\addcontentsline{toc}{paragraph}{{10.2} {一个行列式恰含一个二因子的真实返回}}

{使用保留旧} {\ensuremath{A}} {的一般返回} {(5)，在} {\ensuremath{m=1}} {时为} \[
\frac1pM_{j,s},\qquad
M_{j,s}=
\begin{pmatrix}
p-2j-s&D\\-1&0
\end{pmatrix},\quad D=rj-ps.
\tag{C9.21}
\] {下面统一选择偶数} {\ensuremath{s}，使} {\ensuremath{v2(D)=1}，并且} {\ensuremath{D}} {与} {\ensuremath{n}} {的奇数部分互素。}

{若} {\ensuremath{r=2}} {模4，取} {\ensuremath{j=p}、\ensuremath{s=0}。第一组是整个} {\ensuremath{B}} {块，第二组为} {\ensuremath{A}} {的} {p-r} {份和全部} {\ensuremath{r}} {份} {\ensuremath{C}，另留} {\ensuremath{A}} {的} {\ensuremath{r}} {份。于是} \[
D=rp,\qquad v_2(D)=1.
\] {因为} {\ensuremath{\operatorname{gcd}(r,n)=2}，\ensuremath{D}} {与} {\ensuremath{n}} {的奇数部分互素。}

{若\ensuremath{4|r}，令} {\ensuremath{Q=\operatorname{floor}(n/3)}。选} \[
q=\begin{cases}
Q,&3\nmid n,\\
Q-1,&3\mid n,\quad3\nmid Q-1,\\
Q-2,&3\mid n,\quad3\mid Q-1.
\end{cases}
\tag{C9.22}
\] {有} {\ensuremath{r-1\le q\le p-1}，并且} {\ensuremath{q}} {与} {\ensuremath{n}} {的奇数部分互素。}

{证明后一断言：若} {\ensuremath{n=3Q+1}或3Q+2，\ensuremath{\operatorname{gcd}(Q,n)}只可能整除1或2。若} {\ensuremath{n=3Q}，Q-1与\ensuremath{Q}互素，而Q-2与\ensuremath{Q}的\ensuremath{gcd}只可能整除2；按} {(22)} {又排除了3。容量方面，\ensuremath{3|n}时} {p-r是正的3倍数，故至少为3，\ensuremath{Q-2\ge r-1}；其他情况由} {\ensuremath{r<p}直接得到界。}

{取} {\ensuremath{j=q-1}、\ensuremath{s=2}。此时} \[
i=p-q-1,\quad0\le i\le p-r,\quad
D=2\left(\frac r2(q-1)-p\right),
\] {括号为奇数，所以} {\ensuremath{v2(D)=1}。在任意奇素数} {\ensuremath{\ell |n}} {下，由} {\ensuremath{r=-2p}得} \[
D=-2pq\ne0\pmod\ell.
\] {故所需性质全部成立，原位置容量也由} {\ensuremath{q}} {的范围保证。}

\paragraph*{\texorpdfstring{{10.3} {本原深度严格减一，且不丢失合法性}}{10.3 本原深度严格减一，且不丢失合法性}}\label{C9-h73}
\addcontentsline{toc}{paragraph}{{10.3} {本原深度严格减一，且不丢失合法性}}

{在} {(20)} {下，p-2j-s} {是奇数，\ensuremath{D}} {恰含一个2。整数输出} \[
\widetilde a=(p-2j-s)a+Dz,\qquad
\widetilde z=-a
\] {因此满足} \[
v_2(\widetilde a)=1,\qquad
v_2(\widetilde z)=\nu.
\] {两个输出的共同2因子恰为2。取本原代表只需在二进部分除掉一次2，余下的奇公因子不改变二进赋值，于是} \[
\boxed{v_2(a'z')=\nu-1.}
\tag{C9.23}
\]

{除2后} {\ensuremath{a}'奇、\ensuremath{z}'偶，因此模2合法。对每个奇素数} {\ensuremath{\ell |n}，(21)} {行列式为单位，且} {\ensuremath{y=a+pz}} {的整数输出为} \[
\widetilde y=n jz-(2j+s)y.
\] {两种选择分别有\ensuremath{2j+s=2p}或2q，在这些素数处均为单位，所以奇素数处的合法性保持。本原化不可能除去} {\ensuremath{n}} {的奇素因子，因为此处矩阵可逆且原列本原。}

{每个完整返回的输出本原\ensuremath{G=1}，也保证中间原子没有永久共同非零余数障碍：一旦某个中间本原状态在} {\ensuremath{\ell }处\ensuremath{G}不合法，后续} {\ensuremath{p}平均会保留对应最低赋值上的共同非零余数，不能恢复到\ensuremath{G=1}。程序逐原子按本原\ensuremath{G}检查；它没有把全体暂时被2整除误当成不可达。}

{反复先调整角色再执行} {(21)，每阶段最多四个原子，\ensuremath{\nu }严格减一，最终得到} {(19)。当初始\ensuremath{\nu =1}时使用空词。共同缩放和本原化仅用于计算下一组索引；实际数值始终由原位置平均更新。}


\section{上半带核心的两次转换}\label{construct:common-shape}
\subsection*{\texorpdfstring{{C10.2} {上半带两次平均的转换其实覆盖所有余量}}{C10.2 上半带两次平均的转换其实覆盖所有余量}}\label{C10-h74}
\addcontentsline{toc}{subsection}{{C10.2} {上半带两次平均的转换其实覆盖所有余量}}

{设\ensuremath{p}\ensuremath{\ge}5为素数、1\ensuremath{\le}\ensuremath{s<p}、\ensuremath{n=3p+s}。已有全输入入口得到} \[
(A^{2p},B^p,C^s),\qquad A=a,\ B=-2a+sz,\ C=-pz.
\] {欲转成两个等大\ensuremath{p}块及\ensuremath{r=p+s}份载体，只需选} \[
0\le u\le s,\qquad s\le t+u\le p,\qquad
\gcd(q,n)=1,\quad q=3t+u.
\tag{C10.2.1}
\] {先平均(\ensuremath{A^{p-t-u},B^{t},C^{u})}，产生\ensuremath{p}份\ensuremath{D}。从剩余旧\ensuremath{A}中保留\ensuremath{r}份；它们足够正是t+u\ensuremath{\ge}\ensuremath{s}。再平均全部其余\ensuremath{p}份。两次平均覆盖除载体外的原位置，输出为} \[
((a')^p,(rz'-a')^p,(-pz')^r),
\] {其中} \[
\binom{a'}{z'}=\frac1p
\begin{pmatrix}p-q&st-pu\\-1&0\end{pmatrix}\binom az.
\tag{C10.2.2}
\]

\textbf{{全余量选择规则。}} \[
(t,u)=
\begin{cases}
(s,0),&3\nmid s,\\
((p-u)/3,\ s-3+\rho),&3\mid s,\ \rho=p\bmod3,\ (\rho,s)\ne(1,p-1),\\
((p+1)/2,(p-3)/2),&p\equiv1\pmod3,\ s=p-1.
\end{cases}
\tag{C10.2.3}
\] {第二行先定义\ensuremath{u=s-3+\rho }，再定义\ensuremath{t}，避免把它读成隐式方程。}

{证明：第一行\ensuremath{\operatorname{gcd}(3s,3p+s)=\operatorname{gcd}(3,s)=1}，因为\ensuremath{p>s}且\ensuremath{p}素；容量直接成立。第二行\ensuremath{q=p}，且\ensuremath{u}\ensuremath{\in}{[}1,s−1{]}。条件t+u\ensuremath{\ge}\ensuremath{s}等价于p−s\ensuremath{\ge}6−2ρ。\ensuremath{\rho =2}时p−s\ensuremath{\ge}2；\ensuremath{\rho =1}时p−s\ensuremath{\equiv}1} {mod3，唯一不满足的值为\ensuremath{p-s=1}，恰是第三行。第三行\ensuremath{t+u=p-1=s}、\ensuremath{q=2p}，而\ensuremath{n=4p-1}，所以\ensuremath{\operatorname{gcd}(q,n)=1}，全部容量成立。}

\textbf{{安全性。}} {记\ensuremath{y=a+pz}。式(C10.2.2)给} \[
y'=a'+pz'\equiv-\frac qp y\pmod n,
\qquad
\det\begin{pmatrix}p-q&st-pu\\-1&0\end{pmatrix}
=st-pu\equiv-pq\pmod n.
\tag{C10.2.4}
\] {因此在每个\ensuremath{\ell}\textbar{}\ensuremath{n}处局部可逆，且合法见证仍为单位；清分母／本原化不改变这一点。最终合法又保证两次原子之间不会陷入全体同余。这里不声称转换具有正向逆，只需安全地进入新形状。}

{结合下半带安全入口及三块运输，得到如下共同形状：}

\begin{quote}
{对\ensuremath{p}\ensuremath{\ge}5和整个\ensuremath{2p<n<4p}，合法输入都能进入同一物理形状(p,p,r)，\ensuremath{r=n-2p}。\ensuremath{r=p}时使用三个\ensuremath{p}块的正确本原整数格；其余\ensuremath{r}满足\ensuremath{\operatorname{gcd}(p,r)=1}。}
\end{quote}

{这是共同入口形状定理，不是新的全区间终止证明。旧(2p,p,s)上的控制器不能经单向转换自动变成新形状的控制器；不存在未经证明的逆搬运。}

\subsection*{\texorpdfstring{{C10.3} {固定来源后的全部迹零计数移动}}{C10.3 固定来源后的全部迹零计数移动}}\label{C10-h75}
\addcontentsline{toc}{subsection}{{C10.3} {固定来源后的全部迹零计数移动}}

{现在用共同核心} \[
K_{p,r}(a,z)=(a^p,(rz-a)^p,(-pz)^r),\qquad n=2p+r,
\quad 1\le r<2p,\ r\ne p.
\] {先平均两个不交且组成相同的来源组\ensuremath{c=(i,j,k)}，\ensuremath{i+j+k=p}，产生2p份\ensuremath{D}，再保留\ensuremath{r}份\ensuremath{D}为新载体。令} \[
\alpha=i-j,\qquad \beta=rj-pk.
\] {假设\ensuremath{\alpha }\ensuremath{\ne}0；交换等大块可使\ensuremath{0<\alpha <p}，故\ensuremath{\operatorname{gcd}(\alpha ,p)=1}。剩余库存为} \[
h=(p-2i,p-2j,r-2k,2p-r).
\tag{C10.3.1}
\] {一个输出\ensuremath{o=(x,y,u,v)}要满足0\ensuremath{\le}\ensuremath{o}\ensuremath{\le}\ensuremath{h}、各项和为\ensuremath{p}；其\ensuremath{p}元补组再平均，得到实际矩阵\ensuremath{p}⁻²M(o)，其中} \[
M(o)=\begin{pmatrix}
p(x-y)+\alpha v&pry-p^2u+\beta v\\
-\alpha&-\beta
\end{pmatrix}.
\tag{C10.3.2}
\] {迹零只要求} \[
p(x-y)+\alpha v=\beta.
\tag{C10.3.3}
\]

{本节量词只覆盖这一双来源载体更新模板。保留旧载体的两步返回、不同来源的多阶段返回，以及\ensuremath{\alpha =0}的纯对称配平，不由本节自动分类。第2节统一形状也没有将它们的成功路径全部变成这个模板。}

\textbf{{完整整数核定理。}} {对任何两个满足输出总数与迹零条件的整数输出，其差唯一写为} \[
\boxed{\Delta o=h_0b_2+e b_1,\quad h_0,e\in\mathbb Z,}
\] \[
\boxed{b_1=(-1,-1,2,0),\qquad b_2=(-i,-j,-k,p).}
\tag{C10.3.4}
\] {这里\ensuremath{h_{0}}是整数系数，不是库存向量\ensuremath{h}。}

{证明：令差为(X,Y,U,V)，则\ensuremath{X+Y+U+V=0}、\ensuremath{p(X-Y)+\alpha V=0}。互素性给\ensuremath{V=ph_{0}}及\ensuremath{X-Y=-\alpha h_{0}}。令\ensuremath{e=-Y-jh_{0}}，即得\ensuremath{X=-ih_{0}-e}、\ensuremath{Y=-jh_{0}-e}、\ensuremath{U=-kh_{0}+2e}、\ensuremath{V=ph_{0}}，且唯一。证毕。}

{这正是矩阵} \[
T=\begin{pmatrix}1&1&1&1\\p&-p&0&\alpha\end{pmatrix}
\] {的整数核；其二阶子式的\ensuremath{gcd}为\ensuremath{\operatorname{gcd}(p,\alpha )=1}，Smith不变因子为1、1。因而在无容量上界时，迹零仿射纤维没有额外整除障碍；不能由此忽略非负库存。}

\paragraph*{\texorpdfstring{{3.1} {两种移动有不同的数学作用}}{3.1 两种移动有不同的数学作用}}\label{C10-h76}
\addcontentsline{toc}{paragraph}{{3.1} {两种移动有不同的数学作用}}

{\ensuremath{b}₂把一份来源组合(i,j,k)替换为\ensuremath{p}份\ensuremath{D}。由\ensuremath{pD=iA+jB+kC}，它的均值效果完全为零：} \[
M(o+b_2)=M(o).
\tag{C10.3.5}
\] {\ensuremath{b}₁才是有效方向：} \[
M(o+b_1)=M(o)-pnE_{12}.
\tag{C10.3.6}
\] {因此全部固定来源的迹零返回，去掉不改变作用的重新分组后，恰是一条仿射整数直线，而不是任意二维矩阵空间。}

{写\ensuremath{M(o)=(\beta ,Q};−α,−β)，则由Cayley--Hamilton} \[
M(o)^2=(\beta^2-\alpha Q)I=p f I.
\] {pf被\ensuremath{p}整除由(3.3)立即得出，且} \[
\boxed{f(o+b_1)-f(o)=\alpha n.}
\tag{C10.3.7}
\] {这解释九偏移及载体更新文稿中的相邻返回倍率为何总有等差结构。不是为每个参数猜一个周期；先求关系格，迹零使它成为反射，唯一有效计数方向控制平方标量。}

{安全性也统一。由来源总数得\ensuremath{\beta =nj+p\alpha -p^{2}}；由输出总数与迹零条件得} \[
Q=-p^3+p\beta+n(py+jv).
\] {令\ensuremath{C=(0,1};1,p)把(a,z)改为(z,y\_}\emph{{)=(z,a+pz)，直接相乘得到} \[
CMC^{-1}=\begin{pmatrix}
p^2-nj&-\alpha\\
n[p(y-j)+jv]&-p^2+nj
\end{pmatrix}.
\] {因此} \[
\boxed{\beta^2-\alpha Q\equiv p^4\pmod n,\qquad
p^{-2}CMC^{-1}\equiv
\begin{pmatrix}1&-\alpha/p^2\\0&-1\end{pmatrix}\pmod n.}
\tag{C10.3.8}
\] {由于\ensuremath{\operatorname{gcd}(p,n)=1}，平方标量非零且在所有\ensuremath{\ell}\textbar{}\ensuremath{n}处是单位；真实返回将见证y\_}}{变为−\ensuremath{y_*}模\ensuremath{n}，故保留合法性。重复返回取得共同非零标量，提供真逆。此结论针对容量确实成立的表，不证明任何给定来源有可行输出。}

\paragraph*{\texorpdfstring{{3.2} {容量截面是至多两个整数区间}}{3.2 容量截面是至多两个整数区间}}\label{C10-h77}
\addcontentsline{toc}{paragraph}{{3.2} {容量截面是至多两个整数区间}}

{取\ensuremath{v}₀为\ensuremath{\alpha }⁻¹\ensuremath{\beta }} {modp在{[}0,p−1{]}的代表，令\ensuremath{d}₀=(β−αv₀)/\ensuremath{p}，选一个整数特解} \[
o_0=(d_0,0,p-v_0-d_0,v_0).
\] {所有整数输出为\ensuremath{o=o}₀+\ensuremath{h}₀\ensuremath{b}₂+eb₁。因0\ensuremath{\le}\ensuremath{v}\ensuremath{\le}\ensuremath{2p-r<2p}，只可能\ensuremath{h}₀=0、1；\ensuremath{r>p}时至多\ensuremath{h}₀=0。}

{固定\ensuremath{h}₀，记\ensuremath{a=o}₀+\ensuremath{h}₀\ensuremath{b}₂，则可行整数\ensuremath{e}恰满足} \[
\max\{a_1-h_1,a_2-h_2,\lceil-a_3/2\rceil\}
\le e\le
\min\{a_1,a_2,\lfloor(h_3-a_3)/2\rfloor\},
\tag{C10.3.9}
\] {并检查0\ensuremath{\le}\ensuremath{a}₄\ensuremath{\le}\ensuremath{h}₄。因此固定来源不需要枚举四维输出计数：精确消元后至多两个区间。这里不声称两区间投影一定相接，也不声称区间长度总够供应根或横向作用。}


\section{奇数下半带的统一反射构造}\label{construct:odd-pencil}
{本节设} \(p\ge3\) {为奇素数、}\(1\le r<p\) {为奇数、}\(n=2p+r\){。采用全系统核心} \(K_{p,r}(a,z)\){，并令} \(u=z\){、}\(v=a-rz/2\){，故} \(A=ru/2+v\){、}\(B=ru/2-v\){、}\(C=-pu\){。}

\subsection*{\texorpdfstring{{C11.2} {为什么改变来源，而不是在一张输出纤维里找第二点}}{C11.2 为什么改变来源，而不是在一张输出纤维里找第二点}}\label{C11-h78}
\addcontentsline{toc}{subsection}{{C11.2} {为什么改变来源，而不是在一张输出纤维里找第二点}}

{旧内部奇余量证明用两个旧载体反射，其参数为} \(r,r-2\){；在} \(r=1\) {时第二个参数失效。新构造从两个等大块中先制造新载体，再只从旧载体取零份或一份参加输出。}

{置} \[
\epsilon=\left(\frac{p-r}{2}\bmod2\right)\in\{0,1\},\qquad
d=r-2\epsilon,\qquad
\alpha_0=-r,\quad \alpha_1=4-r.
\tag{C11.2.1}
\] {两张表共用下面的公式。对} \(\alpha=\alpha_0,\alpha_1\){，先平均来源} \[
c_\alpha=(i,j,0),\qquad
i=\frac{p+\alpha}{2},\quad j=\frac{p-\alpha}{2}.
\tag{C11.2.2}
\] {它从} \(A,B\) {各取} \(i,j\) {份，不使用旧} \(C\){。新均值记为} \(D\){，其中保留} \(r\) {份作为最终载体。剩余库存为} \[
\left(\frac{p-\alpha}{2},\frac{p+\alpha}{2},r,p-r\right).
\tag{C11.2.3}
\] {取输出计数} \[
o_\alpha=(x,y,c,v_D)=
\left(
\frac{p+2r-\alpha-2\epsilon}{4},
\frac{p+\alpha-2\epsilon}{4},
\epsilon,
\frac{p-r}{2}
\right).
\tag{C11.2.4}
\] {平均该} \(p\) {元组，再平均其余的} \(p\) {个位置，共三次平均返回同一} \((p,p,r)\) {形状。}

\paragraph*{\texorpdfstring{{2.1} {对所有参数的容量证明}}{2.1 对所有参数的容量证明}}\label{C11-h79}
\addcontentsline{toc}{paragraph}{{2.1} {对所有参数的容量证明}}

{令} \[
s=p-r=4h+2\epsilon,\qquad h\ge0.
\] {由于} \(s\) {是正偶数，}\(\epsilon=0\) {时} \(h\ge1\){。两组来源、输出及库存分别化为：}

{ % do not increment counter
\begin{longtable}[]{@{}
  >{\raggedright\arraybackslash}p{(\linewidth - 4\tabcolsep) * \real{0.3333}}
  >{\raggedright\arraybackslash}p{(\linewidth - 4\tabcolsep) * \real{0.3333}}
  >{\raggedright\arraybackslash}p{(\linewidth - 4\tabcolsep) * \real{0.3333}}@{}}
\toprule\noalign{}
\begin{minipage}[b]{\linewidth}\raggedright
{项}
\end{minipage} & \begin{minipage}[b]{\linewidth}\raggedright
\(\alpha_0=-r\)
\end{minipage} & \begin{minipage}[b]{\linewidth}\raggedright
\(\alpha_1=4-r\)
\end{minipage} \\
\midrule\noalign{}
\endhead
\bottomrule\noalign{}
\endlastfoot
{来源} \((i,j,0)\) & \((2h+\epsilon,r+2h+\epsilon,0)\) & \((2h+\epsilon+2,r+2h+\epsilon-2,0)\) \\
{剩余库存} & \((r+2h+\epsilon,2h+\epsilon,r,4h+2\epsilon)\) & \((r+2h+\epsilon-2,2h+\epsilon+2,r,4h+2\epsilon)\) \\
{输出} & \((r+h,h,\epsilon,2h+\epsilon)\) & \((r+h-1,h+1,\epsilon,2h+\epsilon)\) \\
\end{longtable}
}

{总数都等于} \(p=r+4h+2\epsilon\){。唯一紧的比较是第二张表的第一库存：} \[
r+h-1\le r+2h+\epsilon-2
\quad\Longleftrightarrow\quad h+\epsilon\ge1,
\] {这正由} \(s>0\) {保证。第二来源的} \(j\) {非负也成立：}\(\epsilon=0\) {时} \(h\ge1\){；}\(\epsilon=1,h=0\) {时} \(j=r-1\ge0\){。所有其他比较直接成立。}

{因此这两张表在整个范围有效，连} \((p,r)=(3,1)\) {和} \(r=p-2\) {都无需额外处理。这里的} \(\epsilon\) {只是由整性决定的一位计数，不对应两套算术证明。}

\subsection*{\texorpdfstring{{C11.3} {两张矩阵来自同一反射铅笔}}{C11.3 两张矩阵来自同一反射铅笔}}\label{C11-h80}
\addcontentsline{toc}{subsection}{{C11.3} {两张矩阵来自同一反射铅笔}}

{在} \((a,z)\) {基中记} \(\beta=rj\){，实际返回为} \(J_\alpha/p^2\){，其中} \[
J_\alpha=
\begin{pmatrix}
\beta&pry-p^2\epsilon+\beta v_D\\
-\alpha&-\beta
\end{pmatrix}.
\tag{C11.3.1}
\] {迹零等式是} \(p(x-y)+\alpha v_D=\beta\){。换到} {(1.3)} {后，同一矩阵为} \[
\boxed{
\mathcal J_\alpha=
\begin{pmatrix}
-rp/2&-\alpha\\
pn(r-2\epsilon)/4&rp/2
\end{pmatrix}.}
\tag{C11.3.2}
\] {它的对角和左下项均不随} \(\alpha\) {改变，且} \[
\mathcal J_{\alpha+4}-\mathcal J_\alpha=-4E_{12}.
\tag{C11.3.3}
\] {这是一个沿平方零方向变化的反射族。它来自改变来源和输出的匹配，而非固定来源的同行列式纤维。}

{置} \[
F_\alpha=\frac{pr^2-n(r-2\epsilon)\alpha}{4}.
\] {有} \[
\mathcal J_\alpha^2=pF_\alpha I,\qquad
pF_\alpha\equiv p^4\pmod n.
\tag{C11.3.4}
\] {后一式用} \(r\equiv-2p\pmod n\) {即得。因} \(p\nmid n\){，平方标量非零，且在每个危险素数处为单位。故重复每个三步返回给出其正向射影逆。}

{在} \((z,y_*)=(z,a+pz)\) {基中，} \[
J_\alpha\equiv
\begin{pmatrix}p^2&-\alpha\\0&-p^2\end{pmatrix}\pmod n.
\tag{C11.3.5}
\] {因此真实返回保持合法性} \(\gcd(y_*,n)=1\){。原子中途也不可能制造不可消除的障碍，否则最终合法性不可能恢复。}

\subsection*{\texorpdfstring{{C11.4} {同一倍率环直接给完整上根}}{C11.4 同一倍率环直接给完整上根}}\label{C11-h81}
\addcontentsline{toc}{subsection}{{C11.4} {同一倍率环直接给完整上根}}

{设} \[
c_\epsilon=\begin{cases}r,&\epsilon=0,\\1,&\epsilon=1,\end{cases}
\qquad f_i=F_{\alpha_i}/c_\epsilon,\quad i=0,1.
\tag{C11.4.1}
\] {则} \(f_i\) {是非零整数，且} \[
\gcd(f_0,f_1)=1,\qquad
\gcd(f_0f_1,n)=1,\qquad
r\mid f_0,\qquad f_0-f_1\text{ 为奇数}.
\tag{C11.4.2}
\]

\textbf{{证明。}} {若} \(\epsilon=0\){，} \[
f_0=\frac{r(3p+r)}4,\qquad f_1=f_0-n.
\] {由} \(r,p\) {与} \(n\) {互素，得到} \(\gcd(f_0,n)=1\){，故两数互素。}

{若} \(\epsilon=1\){，令} \(d=r-2\ne0\){，则} \[
f_0=\frac{r[pr+n(r-2)]}{4},\qquad f_1=f_0-nd.
\] {中括号被4整除，因为} \(p-r\equiv2\pmod4\){。又} \[
4f_0\equiv pr^2\pmod n,\qquad
4f_0\equiv pr^2\pmod d.
\] {这里} \(n,d\) {均奇，}\(\gcd(r,d)=1\){，且} \(0<|d|<p\){，所以} \(f_0\) {与} \(nd\) {互素。奇偶性由两种情况下的差均为奇数得到。证毕。}

{令} \[
R=\mathbb Z[1/|f_0f_1|].
\tag{C11.4.3}
\] {由} {(4.2)，}\(2,r\) {都是} \(R\) {单位；没有反演任何危险素数。}

{记} \(S=\operatorname{diag}(1,-1)\) {为两个等大块的角色交换。两条已有真逆的返回给出} \[
T=\mathcal J_{\alpha_1}\mathcal J_{\alpha_0}^{-1}
=\begin{pmatrix}
f_1/f_0&-2r/(c_\epsilon f_0)\\0&1
\end{pmatrix},
\] \[
\boxed{STS^{-1}T^{-1}=U\!\left(\frac{4r}{c_\epsilon f_0}\right).}
\tag{C11.4.4}
\] {上式根参数是} \(R\) {单位。互素性给} \[
\mathbb Z[f_1/f_0,f_0/f_1]=R.
\] {通过} \(T\) {的正反共轭及上根相加，得到全部} \(U(R)\){，包括全部负参数。这里每条形式逆都有先行标量周期，未借助尚未建立的深核。}

\subsection*{\texorpdfstring{{C11.5} {一个两步返回给下根，深核没有循环依赖}}{C11.5 一个两步返回给下根，深核没有循环依赖}}\label{C11-h82}
\addcontentsline{toc}{subsection}{{C11.5} {一个两步返回给下根，深核没有循环依赖}}

{保留} \(r\) {份旧} \(B\){，对其余位置平均两组} \[
\left(A^{(p+r)/2},B^{(p-r)/2}\right),\qquad
\left(A^{(p-r)/2},B^{(p-r)/2},C^r\right).
\tag{C11.5.1}
\] {两个组均长} \(p\){，与保留载体不交。它们给实际矩阵} \(\mathcal K/p\){，其中} \[
\mathcal K=
\begin{pmatrix}-r/2&1\\rn/4&r/2\end{pmatrix},
\qquad
\mathcal K^2=\frac{r(p+r)}2 I.
\tag{C11.5.2}
\] {标量非零且与} \(n\) {互素，故} \(\mathcal K^{-1}\) {正向可执行。由于} \(r,2\in R^\times\){，已有} \(U(2/r)\){，并且} \[
H=\mathcal K U(2/r)
=\begin{pmatrix}-r/2&0\\rn/4&p+r\end{pmatrix},
\qquad
\boxed{SHS^{-1}H^{-1}=L(n).}
\tag{C11.5.3}
\] {这一步没有假设} \(p+r\) {已在} \(R\) {中可逆：}\(H^{-1}\) {来自} \(\mathcal K^{-1}\) {和上根逆。}

{要把整数下根扩为} \(L(nR)\){，必须再给一个真实深核，不能把环系数免费加入。设} \[
P=\begin{pmatrix}1&-2\\0&n\end{pmatrix}.
\] {直接计算得，对每个} \(t\in R\){，} \[
P^{-1}U(t)P=U(nt),
\] \[
P^{-1}\mathcal K
U\!\left(\frac{8(p+r)}r t\right)
\mathcal K^{-1}P=L(nt).
\tag{C11.5.4}
\] {中间上根参数属于} \(R\){，不需要反演} \(p+r\){。故在} \(P\) {基里两条完整} \(nR\) {根已经实际存在。调用既有深同余引理，得到实际群包含} \[
P\Gamma(n^2,R)P^{-1}\supseteq\Gamma(n^3,R).
\tag{C11.5.5}
\] {最后的包含由} \(P\) {和} \(nP^{-1}\) {在} \(R\) {上整直接验证。换基只是分析，没有执行} \(P\){。}

{现在对任意} \(t\in R\){，因全部分母与} \(n\) {互素，可取整数} \(k\) {使} \[
t-k\in n^2R.
\] {于是} \[
L(nt)=L(nk)\,L(n(t-k)).
\] {第一项由} {(5.3)} {的整数幂实现；第二项属于已经实际取得的} \(\Gamma(n^3,R)\){。故完整} \(L(nR)\) {成立。再次应用通用深同余引理，取得} \(\Gamma(n^2,R)\){，证明} {(1.2)。}

{依赖顺序是：两张反射表} {\ensuremath{\to}} \(U(R)\) {\ensuremath{\to}} \(\mathcal K\) {和} \(L(n)\) {\ensuremath{\to}} {换基后的深核} {\ensuremath{\to}} \(L(nR)\) {\ensuremath{\to}} {共同深核。没有用深核反过来制造最初的反射或倍率。}

\subsection*{\texorpdfstring{{C11.6} {接到共同终端，完整覆盖奇数下半区}}{C11.6 接到共同终端，完整覆盖奇数下半区}}\label{C11-h83}
\addcontentsline{toc}{subsection}{{C11.6} {接到共同终端，完整覆盖奇数下半区}}

{换到} \[
(z,y_*)^T=C(u,v)^T,\qquad
C=\begin{pmatrix}1&0\\n/2&1\end{pmatrix}\in\operatorname{SL}_2(R).
\] \(L(nR)\) {与} \(C\) {交换，}\(\Gamma(n^2,R)\) {在该换基下保持。特别已有整数} \(L(n\mathbb Z)\) {和} \(\Gamma(n^2,\mathbb Z)\){。}

{选整数} \(k\) {满足} \[
4\mid k,\qquad k\equiv-p^{-2}\pmod n.
\] {已有上根给真实返回} \[
V=CU(k)C^{-1}
=\begin{pmatrix}
1-kn/2&k\\-kn^2/4&1+kn/2
\end{pmatrix}\in\operatorname{SL}_2(\mathbb Z),
\] \[
V\equiv U(-p^{-2})\pmod n.
\tag{C11.6.1}
\] {它满足共同终端的方向条件。}

{共同终端还要求有限内容对应能够提供标量2。由} \(f_0,f_1\) {奇偶相反，至少一张} {(3.1)} {的整数矩阵有偶行列式；其条目共同} {\ensuremath{gcd}} {为奇数，因为左下项} \(-\alpha_i\) {为奇数。除掉条目} {\ensuremath{gcd}} {后，行列式仍为偶数且与} \(n\) {互素。已有} \(\Gamma(n^2,\mathbb Z)\) {因而允许调用有限内容分支定理，取得所需模} \(n^2\) {的标量2作用。}

{至此}第\ref{sec:finite}节{的全部前提均已落实。其全区间小单位引理对\ensuremath{p}\ensuremath{\ge}5直接适用；\ensuremath{p=3}时这里只有\ensuremath{n=7}，模7单位由−1和2生成，所需补充直接成立。正向互补数字返回供应其余有限尺度，两根列修正与整数主同余纤维完成归零。可固定选终端} \[
t=\left(1,p^2-\frac{n(p-1)}2\right).
\] {在原块值中，它有} \[
A=-\frac{r(p-1)}2,\quad B=\frac{r(p+1)}2,\quad C=-p.
\] \((p+1)/2\) {份} \(A\) {加} \((p-1)/2\) {份} \(B\) {恰为零和} \(p\) {元组。}

{全过程使用原位置与共同尺度；终端有限对应不要求后来数字的有理全局逆。入口和零触发尾部的依赖沿用既有独立证明，故没有把``核心有根''冒充``全输入定理''。}


\section{两个剩余位置的状态相关返回}\label{construct:r-two}
\subsection*{\texorpdfstring{{C12.2} {已有入口、平移和数字返回}}{C12.2 已有入口、平移和数字返回}}\label{C12-h84}
\addcontentsline{toc}{subsection}{{C12.2} {已有入口、平移和数字返回}}

{置\ensuremath{m=p+1}。无条件下半带入口给出} \[
D(u,v)=((u+v)^p,(u-v)^p,(-pu)^2),
\quad \gcd(u,v)=\gcd(v,m)=1.
\tag{C12.2.1}
\] {因此\ensuremath{v}为奇数。共同本原化仅用于分析，不是额外数值操作。}

{整数上平移的具体分组已在正文给出；另一种数字返回的分组及矩阵如下。}

{首先，每个整数\ensuremath{k}的上平移} \[
U(k):(u,v)\longmapsto(u+kv,v)
\tag{C12.2.2}
\] {均可正向执行，且逆也如此。基本U(1)的实际尺度为-1/p；交换两个\ensuremath{p}块的角色给U(-1)。}

{其次，对每个\ensuremath{2\le q\le p}，平均两个不交组} \[
(A^{q-2},B^{p-q},C^2),\qquad
(A^{p-q},B^q),
\] {留下两份旧\ensuremath{A}作载体，给出实际矩阵\ensuremath{p^{-1}}乘} \[
M_q=\begin{pmatrix}-1&-1\\-m&2q-m\end{pmatrix}.
\tag{C12.2.3}
\] {\ensuremath{M_{q}}并非在所有本原状态上都合法。下面每次使用时都先选择\ensuremath{U}平移，并严格证明输出本原及合法；不把它当作免费可逆群元。}

\subsection*{\texorpdfstring{{C12.3} {四次平均给保载体缩放的正向逆}}{C12.3 四次平均给保载体缩放的正向逆}}\label{C12-h85}
\addcontentsline{toc}{subsection}{{C12.3} {四次平均给保载体缩放的正向逆}}

{两个\ensuremath{p}块作互补配比，各取(p+1)/2份一种值、(p-1)/2份另一种值，载体不动，得到} \[
F:(u,v)\longmapsto(u,v/p).
\tag{C12.3.1}
\]

{其逆有一个全局原位置构造。先分别平均} \[
(A^{(p-1)/2},B^{(p-1)/2},C),
\] {使用两个不同的\ensuremath{C}位置和互不相交的旧块位置，产生2p份同值} \[
D=-u/p.
\] {旧\ensuremath{A}、\ensuremath{B}各剩一份。从新\ensuremath{D}保留两份作最终载体，再分别平均一份旧\ensuremath{A}或旧\ensuremath{B}与p-1份\ensuremath{D}，得到} \[
(u,v)\longmapsto(u/p^2,v/p)
=\frac1{p^2}F^{-1}(u,v).
\tag{C12.3.2}
\] {这四次平均覆盖原位置，并共同缩放载体；没有留下未跟随的外部坐标。}

{因此\ensuremath{F}与\ensuremath{F^{-1}}均正向可执行。通过它们对(2.2)的共轭，已有} \[
\boxed{U(\mathbb Z[1/p]).}
\tag{C12.3.3}
\]

\subsection*{\texorpdfstring{{C12.4} {一个统一的三步反射}}{C12.4 一个统一的三步反射}}\label{C12-h86}
\addcontentsline{toc}{subsection}{{C12.4} {一个统一的三步反射}}

{令\ensuremath{t=(p-3)/2}。先平均} \[
(A^{(p+1)/2},B^t,C),
\] {产生\ensuremath{p}份} \[
D=(-u+2v)/p.
\] {保留其中两份作最终载体，余下旧\ensuremath{A}、\ensuremath{B}、\ensuremath{C}、新\ensuremath{D}的数量分别为} \[
(p-1)/2,\quad(p+3)/2,\quad1,\quad p-2.
\]

{令\ensuremath{w=(p-1)/2}，并定义} \[
b=\begin{cases}1,&p=1\pmod4,\\0,&p=3\pmod4,\end{cases}
\quad
x=(p-b-w)/2,\quad y=x+1,\quad z=w-1.
\] {平均一个\ensuremath{p}元组(\ensuremath{A^{x},B^{y},C^{b},D^{z})}，再平均剩余的\ensuremath{p}份。这里是}\textbf{{一个输出组及其补组}}{，不能误写成两个重复组再加第三组。}

{计数对全部\ensuremath{p\ge 5}成立：若\ensuremath{p=4k+1}，(\ensuremath{x,y,b,z)=(k,k+1,1,2k-1)}；若\ensuremath{p=4k+3}，(\ensuremath{x,y,b,z)=(k+1,k+2,0,2k)}。均在上述实际重数内。}

{所得实际矩阵为\ensuremath{p^{-2}}乘} \[
J=\begin{pmatrix}1&-2\\c&-1\end{pmatrix},
\quad
c=\begin{cases}
-(p^2-1)/2,&p=1\pmod4,\\
(p+1)^2/2,&p=3\pmod4.
\end{cases}
\tag{C12.4.1}
\] {于是} \[
J^2=(1-2c)I,\qquad
1-2c=
\begin{cases}p^2,&p=1\pmod4,\\-p(p+2),&p=3\pmod4.\end{cases}
\] {平方标量非零，故\ensuremath{J}的逆也正向可执行。\ensuremath{c}被\ensuremath{m}整除，且平方标量与\ensuremath{m}互素，所以整个返回保持合法性。}

{自由角色交换为\ensuremath{S=\operatorname{diag}(1,-1)}。令\ensuremath{H=J}} {U(2)，则} \[
H=\begin{pmatrix}1&0\\c&2c-1\end{pmatrix},
\qquad
\boxed{[S,H]=L(-2c).}
\tag{C12.4.2}
\] {由\ensuremath{F}的正逆共轭，得到L(2c} {\ensuremath{Z}{[}1/p{]})。置} \[
h=|2c|=
\begin{cases}p^2-1,&p=1\pmod4,\\m^2,&p=3\pmod4,\end{cases}
\quad R=\mathbb Z[1/p],\quad Q=h^2.
\tag{C12.4.3}
\] {已有U(hR)、L(hR)，故深同余引理给} \[
\boxed{\Gamma(Q,R)\text{ 正向可执行}.}
\tag{C12.4.4}
\] {\ensuremath{m}整除\ensuremath{h}、\ensuremath{\operatorname{gcd}(p,h)=1}。第一种情况的\ensuremath{Q}可以包含p-1的额外素因子；下面直接处理它们，不再逐个局部化消除。}

\subsection*{\texorpdfstring{{C12.5} {一次} {\ensuremath{M_{p}}} {把\ensuremath{v}变成整个控制层的单位}}{C12.5 一次  把变成整个控制层的单位}}\label{C12-h87}
\addcontentsline{toc}{subsection}{{C12.5} {一次} {\ensuremath{M_{p}}} {把\ensuremath{v}变成整个控制层的单位}}

{从本原(u,v)出发，先以\ensuremath{U}平移选择某个整数\ensuremath{a}，使} \[
2a=u\pmod v.
\tag{C12.5.1}
\] {这等价于把\ensuremath{u}调整为u\_}\emph{{=2a-v。因为\ensuremath{v}奇，u\_}}{也奇。随后(2.3)的\ensuremath{q=p}返回给} \[
\frac1pM_p(u_*,v)^T
=\frac2p(-a,\ pv-ma)^T.
\tag{C12.5.2}
\] {我们可以选\ensuremath{a}，使这个整数分析列本原，且第二坐标在模\ensuremath{Q}下可逆。}

{具体地，(5.1)保证\ensuremath{\operatorname{gcd}(a,v)=1}。对每个\ensuremath{\ell |Q}：}

\begin{itemize}
\tightlist
\item
  {若\ensuremath{\ell |m}，原\ensuremath{v}为单位，而\ensuremath{pv-ma=pv}} {\ensuremath{mod}} {\ensuremath{\ell }，始终非零。}
\item
  {若\ensuremath{\ell |v}，由原条件知\ensuremath{\ell }不整除\ensuremath{m}；(5.1)使\ensuremath{a}非零，所以\ensuremath{pv-ma=-ma}非零。}
\item
  {其余\ensuremath{\ell }处，\ensuremath{m}、\ensuremath{v}都可逆，可指定\ensuremath{a=(pv-1)/m}} {\ensuremath{mod}} {\ensuremath{\ell }，使新第二坐标等于1。}
\end{itemize}

{这些额外同余的素数都不除\ensuremath{v}，故与(5.1)由CRT同时满足。此外若\ensuremath{p}不整除\ensuremath{v}，另指定\ensuremath{a=1}} {modp；若\ensuremath{p|v}，则(5.1)已保证\ensuremath{p}不整除\ensuremath{a}。因为\ensuremath{p}不整除\ensuremath{Q}，该条件也相容。}

{于是令} \[
(u_1,v_1)=(-a,pv-ma),
\] {有} \[
\boxed{\gcd(u_1,v_1)=1,\quad\gcd(v_1,Q)=1.}
\tag{C12.5.3}
\] {本原性来自\ensuremath{\operatorname{gcd}(a,pv-ma)=\operatorname{gcd}(a,p)=1}。因\ensuremath{m|Q}，输出仍合法。完整返回的输出合法，也保证各原子不会在途中进入不可消除的非零共同余类。}

{这一阶段没有使某个能量必然下降，也不需要；它直接完成局部单位准备。}

\subsection*{\texorpdfstring{{C12.6} {第二次} {\ensuremath{M_{q}}} {同时设定全部终端尺度}}{C12.6 第二次  同时设定全部终端尺度}}\label{C12-h88}
\addcontentsline{toc}{subsection}{{C12.6} {第二次} {\ensuremath{M_{q}}} {同时设定全部终端尺度}}

{取\ensuremath{v_{1}}在模\ensuremath{m}下的逆元\ensuremath{q_{0}}，\ensuremath{1\le q_{0}\le p}。它为奇数。若\ensuremath{q_{0}\ge 2}，令(\ensuremath{q,\epsilon )=(q_{0},1)}；若\ensuremath{q_{0}=1}，则令(\ensuremath{q,\epsilon )=(p,-1)}。因\ensuremath{p=-1}模\ensuremath{m}，总有} \[
2\le q\le p,\quad q\text{ 奇},\quad
qv_1=\epsilon\pmod m,\quad\epsilon\in\{1,-1\}.
\tag{C12.6.1}
\] {这明确避开不可实现的数字\ensuremath{q=1}。}

{置\ensuremath{K=Q/m}。因\ensuremath{v_{1}}与\ensuremath{Q}互素，CRT可以选择整数\ensuremath{b}，使} \[
2b=u_1\pmod{v_1},\qquad
b=\frac{qv_1-\epsilon}{m}\pmod K.
\tag{C12.6.2}
\] {右侧为整数，两个模数互素。还可使\ensuremath{\operatorname{gcd}(b,q)=1}：}

{若\ensuremath{\ell |q}同时整除\ensuremath{v_{1}}，第一条同余已经迫使\ensuremath{b}非零；若\ensuremath{\ell |q}同时整除\ensuremath{K}，因\ensuremath{q}与\ensuremath{m}互素，第二条同余给\ensuremath{mb=-\epsilon }} {\ensuremath{mod}} {\ensuremath{\ell }，也迫使\ensuremath{b}非零。其余\ensuremath{\ell |q}不除两个模数，可再用CRT令\ensuremath{b=1}} {\ensuremath{mod}} {\ensuremath{\ell }。}

{先以\ensuremath{U}平移使第一坐标为\ensuremath{2b-v_{1}}，再执行\ensuremath{M_{q}}。输出的本原整数分析列恰为} \[
(u_2,v_2)=(-b,qv_1-mb),
\quad
\boxed{v_2=\epsilon\pmod Q.}
\tag{C12.6.3}
\] {因为\ensuremath{\operatorname{gcd}(b,v_{1})=\operatorname{gcd}(b,q)=1}，坐标\ensuremath{gcd}严格为1，不会在本原化中破坏已经安排好的尺度。}

{这一步不要求\ensuremath{R}单位模\ensuremath{Q}满射，也不要求\ensuremath{q}在\ensuremath{R}中预先可逆。它是根据输入选取的真实正向返回，直接把尺度设为\ensuremath{R}单位\ensuremath{\epsilon =}±1。}

\subsection*{\texorpdfstring{{C12.7} {精确终端与全输入}}{C12.7 精确终端与全输入}}\label{C12-h89}
\addcontentsline{toc}{subsection}{{C12.7} {精确终端与全输入}}

{用U(k)使\ensuremath{u_{3}=u_{2}+kv_{2}=0}} {modQ。仍有\ensuremath{\operatorname{gcd}(u_{3},v_{2})=1}。取\ensuremath{\alpha }} {\ensuremath{u_{3}+\beta }} {\ensuremath{v_{2}=1}，并调整(alpha,beta)使\ensuremath{\alpha =0}} {modQ。则} \[
P=\begin{pmatrix}
v_2/\epsilon&-u_3/\epsilon\\
\epsilon\alpha&\epsilon\beta
\end{pmatrix}\in\Gamma(Q,\mathbb Z),
\qquad P(u_3,v_2)^T=(0,\epsilon)^T.
\] {由(4.4)，\ensuremath{P}有原位置正向实现。最终状态是非零共同尺度乘} \[
(\epsilon^p,(-\epsilon)^p,0^2),
\] {三次相反值消去即可全零。}

{无条件入口已给全部合法原输入；必要性为原\ensuremath{G}不变量。因此定理完整成立，没有剩下``原位置分解''或``有限尺度类''的额外假设。}


\section{偶数核心的初始双根}\label{construct:even-roots}
\subsection*{\texorpdfstring{{C13.3} {偶数核心的同一整数格}}{C13.3 偶数核心的同一整数格}}\label{C13-h90}
\addcontentsline{toc}{subsection}{{C13.3} {偶数核心的同一整数格}}

{令} {\ensuremath{r=2t}，\ensuremath{n=2p+r}。整数零和核心可以写成} \[
\mathcal K(a,z)=(a^p,(rz-a)^p,(-pz)^r),\qquad \gcd(a,z)=1.
\tag{C13.4}
\] {因为} {\ensuremath{\operatorname{gcd}(p,r)=1}，式} {(C13.4)} {给出全部整数三块零和状态，其本原性恰为} {\ensuremath{\operatorname{gcd}(a,z)=1}。}

{置} {\ensuremath{y=a+pz}。两种差为} \[
a-(-pz)=y,\qquad a-(rz-a)=2y-nz.
\] {由} {\ensuremath{\operatorname{gcd}(y,z)=1}} {得精确等式} \[
\boxed{G=\gcd(y,n).}
\tag{C13.5}
\] {这里没有遗漏模} {2} {条件。旧坐标} {\ensuremath{a=tz+v}} {给} {\ensuremath{y=v+(p+t)z}，模} {2} {恰为} {v+(t+1)z；奇素数} {\ensuremath{\ell |n}} {时} {\ensuremath{y=v}} {模} {\ensuremath{\ell }。旧两条合法条件实际上是} {(5)} {的不同局部写法。}

{沿用}第\ref{construct:even-roots}节{。允许奇数} {\ensuremath{s}} {满足} {\ensuremath{max(1,2r-p)\le s\le r}，令} {\ensuremath{i=(p-s)/2}。平均两个不交组} \[
(a^i,b^i,c^s),\qquad
(a^{p-r-i},b^{p-i},c^{r-s}),
\tag{C13.6}
\] {留下} {\ensuremath{r}} {份原} {\ensuremath{a}。其} {(a,z)} {矩阵为} \[
\frac1pJ_s,\qquad
J_s=\begin{pmatrix}0&d_s\\-1&0\end{pmatrix},\quad
d_s=\frac{pr-ns}{2}.
\tag{C13.7}
\] {因为} {\ensuremath{s}、\ensuremath{p}} {奇，\ensuremath{d_{s}}} {为奇数；对每个} {\ensuremath{\ell |n}，\ensuremath{d_{s}=-p^{2}}} {模} {\ensuremath{\ell }。故} {\ensuremath{d_{s}}} {非零且与} {\ensuremath{n}} {互素，\ensuremath{J_{s}^{2}=-d_{s}}} {\ensuremath{I}，返回及射影逆均有真实正向实现。}

{在} {(z,y)} {整幺模坐标中，同一返回的矩阵为} \[
\frac1pK_i,\qquad
K_i=\begin{pmatrix}p&-1\\ni&-p\end{pmatrix},\quad
K_i^2=(p^2-ni)I.
\tag{C13.8}
\] {交换两个} {\ensuremath{p}} {块的角色分别对应} \[
\Sigma=\begin{pmatrix}-1&r\\0&1\end{pmatrix}
\quad\hbox{和}\quad
S=\begin{pmatrix}1&0\\n&-1\end{pmatrix}.
\] {块角色交换只是更新原位置账本。}

\subsection*{\texorpdfstring{{C13.4} {任意内部偶余量都有深主同余控制}}{C13.4 任意内部偶余量都有深主同余控制}}\label{C13-h91}
\addcontentsline{toc}{subsection}{{C13.4} {任意内部偶余量都有深主同余控制}}

{现设} {\ensuremath{4\le r\le p-3}。\ensuremath{s=r-1}、r-3} {都满足} {(6)} {的容量。记} \[
d=\frac{pr-n(r-1)}2,\qquad e=d+n,\qquad
R=\mathbb Z[1/|de|],\quad h=rn.
\] {有} {\ensuremath{\operatorname{gcd}(d,e)=1}，\ensuremath{d}、\ensuremath{e}} {均奇且与} {\ensuremath{n}} {互素。又} {\textbar{}\ensuremath{de|>1}，所以} {\ensuremath{R}} {有无限多个单位。}

{令} {\ensuremath{J=J_{r-1}}、\ensuremath{Q=J_{r-3}}。直接计算} \[
D=J^{-1}Q=\operatorname{diag}(1,e/d),\qquad
\Sigma D\Sigma D^{-1}=U(rn/e).
\tag{C13.9}
\] {由} {\ensuremath{\operatorname{gcd}(d,e)=1}，} \[
\mathbb Z[e/d,d/e]=\mathbb Z[1/(de)]=R.
\] {用} {\ensuremath{D}} {的共轭、逆共轭及根参数相加、取负，得到} {U(hR)。又} \[
JU(x)J^{-1}=L(-x/d)
\] {给} {L(hR)。全部因子及其逆均有正向实现。}

\textbf{{深同余引理。}} {对} {\ensuremath{R=Z}{[}1/q{]}、\ensuremath{q>1}} {和正整数} {\ensuremath{h}，使用下述外部定理时，} \[
\boxed{\Gamma(h^2,R)\subseteq
E_h:=\langle U(hR),L(hR)\rangle.}
\tag{C13.10}
\]

{证明：Morris} {2007} {定理} {6.1(2)} {应用于有限指数子群} {\ensuremath{\Gamma (h,R)}。其中全部初等矩阵恰为两种根，故} {\ensuremath{E_{h}}} {有限指数。Serre} {对有理数域的} {\ensuremath{S}} {算术} {\ensuremath{\mathrm{SL}_{2}}} {强同余子群性质给某个} {\ensuremath{\Gamma (M,R)}} {包含于} {\ensuremath{E_{h}}；可令} {\ensuremath{h^{2}R}} {整除} {MR。}

{在有限环} {\ensuremath{B=R/MR}} {中取} {\ensuremath{g=(a,b};c,d)} {属于} {\ensuremath{\Gamma (h^{2},R)}} {的像。可用} {U(hk)} {使} {\ensuremath{u=a+hkc}} {在} {\ensuremath{B}} {中为单位：在除} {\ensuremath{h}} {的素数处} {\ensuremath{a=1}；在不除} {\ensuremath{h}} {的素数处，用} {CRT} {选择} {\ensuremath{k}} {避开至多一个坏类。置} {\ensuremath{v=b+hkd}，则} \[
g=U(-hk)L(c/u)\operatorname{diag}(u,u^{-1})U(v/u).
\] {这些非对角参数都属于} {hB，而且} {u-1} {属于} {\ensuremath{h^{2}B}。选择} {\ensuremath{t}} {属于} {hB} {使} {\ensuremath{ht=u-1}（从整数的} {\ensuremath{h^{2}}} {倍代表可选），便有} \[
\operatorname{diag}(u,u^{-1})
=U(h)L(t)U(-h/u)L(-tu).
\tag{C13.11}
\] {所有参数都在} {hB，可分别提升到} {hR。因此} {\ensuremath{E_{h}}} {在模} {\ensuremath{M}} {覆盖} {\ensuremath{\Gamma (h^{2},R)}} {的像，再用已包含的} {\ensuremath{\Gamma (M,R)}，得到} {(10)。}

{这里有限环里的除法仅用于挑选允许根参数的剩余类；没有把任意模} {\ensuremath{M}} {单位当作} {\ensuremath{R}} {单位。也未使用仅在高秩成立的相对初等生成定理。}

{结合} {(9)、(10)，每个内部偶余量都有真实可执行的} \[
\boxed{\Gamma((rn)^2,\mathbb Z[1/|de|]).}
\tag{C13.12}
\] {它把无限轨道问题降为有限的同余轨道和单位尺度问题，但本身不保证所有有限轨道都有终端。\ensuremath{r=2}、\ensuremath{r=p-1}} {各只有一个允许的标准对合，本文不将它们包含在} {(12)} {中。}

{外部来源与}第\ref{construct:odd-pencil}节{一致：Morris,} {NYJM} {13} {(2007),} {定理} {6.1(2)；Serre,} {Annals} {of} {Mathematics} {92} {(1970),} {489--527,} \href{https://doi.org/10.2307/1970630}{DOI}{。此处将} {Serre} {定理用在} {\ensuremath{R=Z}{[}1/65{]}，即} {S=\{无穷处,5,13\}} {的有理数域特例；不是由} {\ensuremath{Z}{[}1/2{]}} {的特例推演得到。Morris} {相关原文已在项目中核对；Serre} {原文全文未取得，这里作为已发表外部定理引用，不声称自行证明或数值验证了该定理。}


\section{偶数下半带内部的物理资源}\label{construct:even-interior}
{本节设} \(p\ge7\) {为奇素数，}\(r\) {为偶数，满足} \[4\le r\le p-3,\qquad n=2p+r.\tag{C14.1}\] {两条外边界另有独立构造。}

\subsection*{\texorpdfstring{{C14.2} {坐标、真实返回与已有深同余控制}}{C14.2 坐标、真实返回与已有深同余控制}}\label{C14-h92}
\addcontentsline{toc}{subsection}{{C14.2} {坐标、真实返回与已有深同余控制}}

{写} \(r=2t\){、}\(m=p+t\){、}\(n=2m\){。核心采用整数幺模坐标} \[
\mathcal D_{p,t}(u,v)=((tu+v)^p,(tu-v)^p,(-pu)^{2t}),
\quad a=tu+v,\ z=u.
\] {本原性为} \(\gcd(u,v)=1\){，差分} {\ensuremath{gcd}} {为} \[
G=\gcd(2v,mu+v).
\] {因此合法性恰为} \[
\begin{cases}
\gcd(v,m)=1,&t\text{ 奇，此时 }m\text{ 偶},\\
\gcd(v,m)=1,\quad u+v\text{ 奇},&t\text{ 偶，此时 }m\text{ 奇}.
\end{cases}\tag{C14.3}
\]

{对每个可行奇数} \(s\){，已有两次原位置平均的对合。其参数范围及矩阵为} \[
\max(1,4t-p)\le s\le2t,\qquad
J_s=\begin{pmatrix}-t&-1\\m(t-s)&t\end{pmatrix},
\] \[
d_s=pt-ms,\qquad J_s^2=-d_sI,\qquad \gcd(d_s,n)=1.
\tag{C14.4}
\] {实际矩阵是} \(J_s/p\){。分组为} \[
(A^i,B^i,C^s),\quad(A^{p-r-i},B^{p-i},C^{r-s}),\quad
i=(p-s)/2,
\] {并留下} \(r\) {份旧} \(A\) {作新载体。换块在} \((u,v)\) {中是} \(S=\operatorname{diag}(1,-1)\){，在} \((a,z)\) {中是} \(\Sigma=(-1,r;0,1)\){。}

{内部范围总有} \(s_0=r-1,s_1=r-3\){。设} \(d=d_{s_0}\){、}\(e=d+n\){，则} \(d,e\) {互素、奇数且与} \(n\) {互素。偶数核心双根构造中的对角倍率与根公式，在} \[
R_0=\mathbb Z[1/|de|],\qquad h=rn
\] {上给出} \(U(hR_0),L(hR_0)\){，并由深同余引理得到} \[
\Gamma(h^2,R_0).\tag{C14.5}
\] {这里根公式首先在} \((a,z)\) {坐标成立；主同余子群在整数幺模换元下保持同一个} {level。}

{若还有第三个返回} \(s_2=r-5\){，则可以把} \(R_0\) {扩为} \(\mathbb Z[1/|d(d+n)(d+2n)|]\){。两个对角倍率的正负幂和根参数相加给出同一根公式；三个奇数两两互素。}

\subsection*{\texorpdfstring{{C14.3} {从有限逆提升到真实逆}}{C14.3 从有限逆提升到真实逆}}\label{C14-h93}
\addcontentsline{toc}{subsection}{{C14.3} {从有限逆提升到真实逆}}

\textbf{{有限逆提升引理。}} {假设在全系统核心上已正向实现} \(\Gamma(Q,R)\){，且一个正向返回的射影矩阵} \(Z\) {属于} \(\mathrm{SL}_2(R)\){。则} \(Z^{-1}\) {也可正向实现。}

{证明：}\(R/QR\) {有限，故某个} \(k>0\) {使} \(Z^k\in\Gamma(Q,R)\){。于是} \[
Z^{-1}=Z^{k-1}(Z^k)^{-1}.
\] {右端全部已正向可执行。这里主同余子群提供真正的有理逆，不能省去这个前提。}

\paragraph*{\texorpdfstring{{3.1} {保载体倍率}}{3.1 保载体倍率}}\label{C14-h94}
\addcontentsline{toc}{paragraph}{{3.1} {保载体倍率}}

{两次互补平均保留全部载体，给} \[
F_c:(u,v)\mapsto(u,cv/p),\qquad c\text{ 奇},\quad |c|\le p.
\] \(F_1\) {始终合法，其整数射影矩阵为} \(\operatorname{diag}(p,1)\){。因为} \(p\) {与} {(5)} {的} {level} {互素，已有矩阵夹乘引理给} \(F_1^{-1}\) {的正向实现。故} \[
C_c=F_cF_1^{-1}=\operatorname{diag}(1,c)\tag{C14.6}
\] {正向可执行。若} \(c\) {为奇数且与} \(n\) {互素，它保持} {(3)。}

{一个直接的夹乘证明如下：若} \(D=\operatorname{diag}(1,c)\){、}\(\gcd(c,Q)=1\){，选择} \(g\in\Gamma(Q,\mathbb Z)\){，使} \(g_{11}\equiv0\pmod c\){。可由模} \(c\) {的旋转矩阵和模} \(Q\) {的恒等矩阵，用CRT及整数初等矩阵提升得到。则} \(H=DgD/c\) {是正向可实现的整数} \(\mathrm{SL}_2\) {矩阵。有限逆提升给} \(H^{-1}\){，而} \(gDH^{-1}=cD^{-1}\) {给所需射影逆。}

\paragraph*{\texorpdfstring{{3.2} {反演} \(t\) {的全部奇因子}}{3.2 反演 t 的全部奇因子}}\label{C14-h95}
\addcontentsline{toc}{paragraph}{{3.2} {反演} \(t\) {的全部奇因子}}

{设奇数} \(c\mid t\){。它可能整除原} {level，不能直接套用上一段的互素条件。但取} {(4)} {中的} \(J=J_{s_0}\){，有} \[
H=\frac1c C_cJC_c
=\begin{pmatrix}-t/c&-1\\m(t-s_0)&tc\end{pmatrix}.
\] \(\det H=\det J=d\) {是} \(R_0\) {单位。因此} \(Z=HJ^{-1}\in\mathrm{SL}_2(R_0)\){，且其射影作用正向可实现。有限逆提升先给} \(Z^{-1}\){，再给} \(H^{-1}\){。恒等式} \[
JC_cH^{-1}=cC_c^{-1}\tag{C14.7}
\] {于是给出} \(C_c^{-1}\){。}

{这没有在物理上除以} \(c\){：}\(H\) {只是完整返回的非零共同尺度规范化。所有} \(n\) {个位置都在同一个核心内。}

{记} \(t=2^a t_0\){，}\(t_0\) {奇。现在可反演倍率} \(p,t_0\){。在} \((a,z)\) {坐标中} \(C_c\) {是上三角矩阵} \[
\begin{pmatrix}c&t(1-c)\\0&1\end{pmatrix};
\] {共轭原上根使参数乘} \(c\){，再用} \(J\) {把上根转成下根。因此，令} \[
R=R_0[1/(pt_0)],\qquad Q=(2^{a+2}m)^2,
\] {可重新应用深同余引理，得到} \[
\boxed{\Gamma(Q,R)\text{ 正向可执行}.}\tag{C14.8}
\] {原} \(r\) {中的奇素数都已局部化；}\(Q\) {只含} \(n\) {的素因子。后续任何} \(|c|\le p\) {的奇数、}\(\gcd(c,n)=1\){，均可通过} {(6)} {及互素夹乘反演并加入} \(R\){。同样的根共轭保持} {(8)} {的} {level。}

\subsection*{\texorpdfstring{{C14.4} {射影循环与精确纤维提升}}{C14.4 射影循环与精确纤维提升}}\label{C14-h96}
\addcontentsline{toc}{subsection}{{C14.4} {射影循环与精确纤维提升}}

{取一个可行} \(s\){，令} \[
G_s=-SJ_s=\begin{pmatrix}t&1\\m(t-s)&t\end{pmatrix},
\quad c_s=\frac{m(t-s)}{t^2}.
\] {其射影作用正向可执行。对奇素数} \(\ell\mid m\){，在} \(x=u/v\) {坐标中，它模} \(\ell\) {是} \[
x\mapsto x+1/t.
\] {若} \(3\mid m\){，选择} \(s\) {避免} \(c_s\equiv-3\pmod9\){。两个相邻的可行奇数} \(s,s-2\) {保证至少一个可用：当} \(v_3(m)=1\) {时，两者的} \(c_s/3\) {模3不同；更深时均不退化。}

\textbf{{局部循环引理。}} {对这样的} \(s\){，该射影作用在每个} \(\mathbb Z/\ell^b\mathbb Z\) {上形成一个完整循环。若} \(t\) {奇，则在所有} \(2^b\) {上也成立。}

{证明：在对应局部环中写} \[
G_s/t=I+N,\qquad N^2=c_s I.
\] \(\ell\ge5\) {时，二项式展开给} \((I+N)^\ell=A I+B N\){，其中} \(A\) {为单位、}\(v_\ell(B)=1\){。}\(\ell=3\) {时} \(B=3+c_s\){，所选条件正保证其赋值为1。若} \(t\) {奇，则} \(m\) {偶且} \(c_s\) {偶；在2处} \((I+N)^2=(1+c_s)I+2N\){，也有单位对角项和精确的一阶非对角项。}

{重复取} \(\ell\) {次幂得到} \(G_s^{\ell^j}\) {的射影非对角系数赋值恰为} \(j\){。对任意} \(x\){，其位移还乘一个单位} \(1-t^2c_sx^2\){。因此} \(G_s^k\) {固定} \(x\bmod\ell^b\) {当且仅当} \(\ell^b\mid k\){，轨道长度正好} \(\ell^b\){。证毕。}

{不同奇素数的循环长度互素，CRT能同时选择同一个正向幂。指数可逐素数幂提升求出，不搜索平均词。}

\textbf{{精确纤维引理。}} {若} \(x,y\) {是} \(R^2\) {中本原列，且} \(x\equiv y\pmod{QR^2}\){，则存在} \(H\in\Gamma(Q,R)\) {使} \(Hx=y\){。证明是把两列补成行列式一矩阵，并调整第二列使两个矩阵模} \(Q\) {相同。对} \(x\equiv\epsilon y\){、}\(\epsilon\in R^\times\) {同样适用。}

{下面每一分支均同时证明方向相等和尺度属于真实} \(R\) {单位，最后才调用这个引理。仅有普通射影方向相同是不够的。}

\subsection*{\texorpdfstring{{C14.5} \(t\) {奇：全部} \(r\equiv2\pmod4\) {的内部核心}}{C14.5 t 奇：全部 r\textbackslash equiv2\textbackslash pmod4 的内部核心}}\label{C14-h97}
\addcontentsline{toc}{subsection}{{C14.5} \(t\) {奇：全部} \(r\equiv2\pmod4\) {的内部核心}}

{此时} \(t\ge3\){、}\(m\) {偶，}\(t\) {已由第3节完全局部化。取} \(Q=(4m)^2\){。合法性使} \(v\) {为模} \(Q\) {单位。第4节的循环同时覆盖全部素数幂，因此某个} \(G_s^k\) {把} \(u/v\) {送到0模} \(Q\){。}

{其矩阵属于} \(\mathrm{GL}_2(R)\){，行列式为已反演的} \(-d_s\){，故输出仍为} \(R\) {本原列。取整数本原代表} \((u_1,v_1)\){，有} \(u_1\equiv0\pmod Q\){、}\(v_1\) {为模} \(Q\) {单位。}

{选} \(v_1^{-1}\pmod m\) {的平衡代表} \(c\){。因} \(m\) {偶，}\(c\) {奇，且} \[
0<|c|\le m/2<p.
\] {先用已实现的} \(U(Qk)\) {使} \(u_1+Qkv_1\) {与} \(c\) {互素。这可按} \(c\) {的素因子用CRT选择；}\(\gcd(c,Q)=1\){，原列本原保证每处可选。随后执行} \(C_c\){。输出} \((u_2,v_2)\) {仍本原，且} \(u_2=0\pmod Q,\ v_2=1\pmod m\){。}

{模} \(Q\) {到模} \(m\) {的单位核，全部由两个已经存在的} \(R\) {单位生成：} \[
\eta_1=-p/t=1-m/t,\qquad \eta_2=e/d=1+2m/d.\tag{C14.9}
\] {若} \(4\mid m\){，LTE给} \(\eta_1\) {在每个局部单位核上的最大阶，各阶互素，所以它单独生成整个核。若} \(m=2\pmod4\){，}\(\eta_2\) {生成模} \(2m\) {的单位核，}\(\eta_1=3\pmod4\) {补齐另一陪集。这些断言包括} \(m\) {的任意素数幂。}

{于是存在实际} \(R\) {单位} \(\epsilon\equiv v_2\pmod Q\){。精确纤维引理把} \((u_2,v_2)\) {送到} \((0,\epsilon)\){，对应} \[
(\epsilon^p,(-\epsilon)^p,0^{2t}),
\] {已有三次相反值平均完成归零。该分支不要求} \(3t\le p\){。}

{当} \(3t\le p\) {时还可直接取} \(s=t\){，得到} \[
SJ_t=-t\,U(1/t).
\] {这是同一机制的较短特例，不是一般定理的容量前提。}

\subsection*{\texorpdfstring{{C14.6} \(t\) {偶且有三个返回：}\(8\le r\le p-5\)}{C14.6 t 偶且有三个返回：8\textbackslash le r\textbackslash le p-5}}\label{C14-h98}
\addcontentsline{toc}{subsection}{{C14.6} \(t\) {偶且有三个返回：}\(8\le r\le p-5\)}

{此时} \(m\) {奇。取} \(s_0=r-1,s_1=r-3,s_2=r-5\){，对应} \(d,d+n,d+2n\){。按第3节反演} \(t_0,p\){，并加入} \[
c_*=(m+\sigma)/2\text{ 为奇数},\quad \sigma\in\{1,-1\},\qquad
\epsilon_2=\sigma/c_*.
\] {有} \(c_*\le p,\ \gcd(c_*,n)=1\){，故它及其逆正向可用，且} \[
\epsilon_2\equiv2\pmod m.
\] {这只是2的剩余类由一个奇数} \(R\) {单位表示，没有反演有理数2。}

{先用}第\ref{construct:core-entry}节{达到} \(v_2(az)=1\){。必要时用} \(J_s\) {交换角色，得到} \(a\) {奇、}\(z=u\) {的赋值为1，因而} \(v=a-tu\) {奇。}

{选} \(v^{-1}\pmod m\) {的平衡代表} \(w=2^h c\){、}\(c\) {奇。仍有} \(|c|<p\){。先用} \(U(Qk)\) {保证施加} \(C_c\) {后本原，再施加} \(C_c\){，使} \(v\bmod m\) {属于} \(R\) {单位像，具体可由} \(\epsilon_2^{-h}\) {表示。这些操作保持二进深度一。}

\paragraph*{\texorpdfstring{{二进方向}}{二进方向}}\label{C14-h99}
\addcontentsline{toc}{paragraph}{{二进方向}}

{两个比值} \[
\frac{d+n}{d},\qquad\frac{d+2n}{d+n}
\] {模8分别为3和7，因为} \(n=2m\){、}\(m\) {奇。选} \(\kappa=3\pmod8\) {的那个。对应实际对角返回在} \((a,z)\) {中使} \(x=z/a\) {乘} \(\kappa\){。而换块使} \[
x\mapsto\frac{x}{-1+rx}\equiv-x\pmod{16}
\quad(v_2(x)=1,\ 4\mid r).
\] {由一次换块及零次或一次} \(\kappa\) {倍率，可覆盖赋值一的全部模16方向。又} \[
v_2(\kappa^2-1)=3,
\] {所以} \(\kappa^2\) {的幂覆盖固定模16方向的全部高阶提升。}

{取终端} \[
u_T=2,\quad a_T=p-r(p-3)/2,\quad v_T=a_T-2t.
\tag{C14.10}
\] {它本原、合法，且} \(v_2(a_Tu_T)=1\){。上面的动作可使} \(z/a\) {在} \(Q\) {的2部分上与该终端相同。}

\paragraph*{\texorpdfstring{{奇素数方向}}{奇素数方向}}\label{C14-h100}
\addcontentsline{toc}{paragraph}{{奇素数方向}}

{选择第4节的非退化} \(G_s\){。它模2是交换矩阵，因此其在任何} \(2^b\) {上的矩阵阶都是2的幂。用CRT把奇素数处所需的射影循环指数，和``在2部分矩阵为恒等''的指数条件合并，即可在保持二进方向的同时，使} \(u/v\) {在} \(m^2\) {上命中} {(10)。}

\paragraph*{\texorpdfstring{{单位尺度}}{单位尺度}}\label{C14-h101}
\addcontentsline{toc}{paragraph}{{单位尺度}}

{到此输出与} {(10)} {模} \(Q\) {只差一个单位} \(\lambda\){。所有方向调整把} \(v\bmod m\) {乘以} \(R\) {单位的剩余类：}\(J_s\) {的倍率} \(t\) {可由} \[
\tau=\epsilon_2^{v_2(t)}t_0\in R^\times
\] {表示，换块用} \(-1\){，对角比值模} \(m\) {为1。而} \(v_T=t^2\pmod m\){。因此} \(\lambda\bmod m\) {也来自} \(R^\times\){。}

{剩余单位核由} \[
\kappa_0=(d+n)/d=1+2m/d,\qquad
\eta=(d+2n)/d=1+4m/d
\] {生成。}\(\eta\) {通过LTE生成模} \(4m\) {的全部单位核；}\(\kappa_0=3\pmod4\) {补另一陪集。故存在} \(\epsilon\in R^\times\) {使输出等于} \(\epsilon(u_T,v_T)\pmod Q\){。}

{精确纤维引理将其真正送到该终端。选取} \((p+1)/2\) {份} \(A\){、}\((p-3)/2\) {份} \(B\) {和一份} \(C\){，和为零。一次平均造出} \(p\) {个零，再用}第\ref{construct:zero-tail}节{完整收尾。}

\subsection*{\texorpdfstring{{C14.7} {两种只需固定二进证书的内部边界}}{C14.7 两种只需固定二进证书的内部边界}}\label{C14-h102}
\addcontentsline{toc}{subsection}{{C14.7} {两种只需固定二进证书的内部边界}}

{还剩} \(r=4\){，以及} \(r=p-3\) {且} \(p=3\pmod4\){。它们的} \(t\) {偶但只有两个标准对合。下面增加一个有实际计数的迹为} \(\pm2\) {的返回。}

{保留} \(r\) {份旧} \(A\){，第一组取} \(i=p-j-s\) {份} \(A\){、}\(j\) {份} \(B\){、}\(s\) {份} \(C\){，第二组取全部剩余位置。矩阵在} \((a,z)\) {中为} \[
M=\begin{pmatrix}a_0&D\\-1&0\end{pmatrix},\quad
a_0=p-2j-s,\quad D=rj-ps.
\tag{C14.11}
\]

{ % do not increment counter
\begin{longtable}[]{@{}lrrrr@{}}
\toprule\noalign{}
{范围} & \(s\) & \(j\) & \(a_0\) & \(D\) \\
\midrule\noalign{}
\endhead
\bottomrule\noalign{}
\endlastfoot
\(r=4,\ p\ge11\) & {1} & \((p-3)/2\) & {2} & \(p-6\) \\
\(r=p-3,\ p=1\pmod3\) & \(p-4\) & {1} & {2} & \(-p^2+5p-3\) \\
\(r=p-3,\ p=2\pmod3\) & \(p-6\) & {4} & {-2} & \(-p^2+10p-12\) \\
\end{longtable}
}

{各行容量非负且} \(\gcd(D,rn)=1\){。第二、三行分别避开} \(p-2\){、}\(p+2\) {与} \(3(p-1)\) {的共同因子3。}\(p=7,r=4\) {使用第二行。}

\paragraph*{\texorpdfstring{{新返回的真逆和统一} {level}}{新返回的真逆和统一 level}}\label{C14-h103}
\addcontentsline{toc}{paragraph}{{新返回的真逆和统一} {level}}

{设已有主同余} {level} {为} \(Q_0\){。取整数} \(k=0\pmod{Q_0}\){、}\(k=a_0\pmod D\){，则} \[
H=M U(k) M/D\in\mathrm{SL}_2(\mathbb Z).
\] {它是正向可实现的完整返回。有限逆提升给} \(H^{-1}\){，再由} \(U(k)MH^{-1}=DM^{-1}\) {得} \(M^{-1}\){。}

\(MJ_s^{-1}\) {是上三角矩阵，比值为} \(D/d_s\){。其正负共轭使根参数局部化到} \(D\){。与} \(\Sigma\) {作交换子给真正的上根} \[
U\!\left(2a_0+r(1-D/d_s)\right).\tag{C14.12}
\] {括号是} \(nR\) {中的元素，二进赋值恰为2。反演} \(t_0\) {后，原根理想是} \(2^{v_2(t)+1}nR\){，与} {(12)} {合并得到} \(2nR\){：清除一个奇数} \(R\) {单位分母后，用奇数与2幂的贝祖关系先取得} \(U(2n)\){，再利用} \(R/2^kR\) {的整数代表扩展到全部参数。下根由} \(J_s\) {共轭取得。故可使用} \[
\Gamma((2n)^2,R).
\] \(m\) {奇，所以2部分只需模16；对} \(r=4\) {下面方便起见使用其子群} \(\Gamma(64m^2,R)\){，多保留两位精度。}

\paragraph*{\texorpdfstring{{归一化二进证书保留尺度}}{归一化二进证书保留尺度}}\label{C14-h104}
\addcontentsline{toc}{paragraph}{{归一化二进证书保留尺度}}

{令} \(\epsilon_2=\sigma/c_*\) {如第6节，}\(\tau\in R^\times\) {表示} \(t\bmod m\){。在} \(r=4\) {时取} \(\tau=\epsilon_2\){；在另一边界取} \(\tau=\epsilon_2^{v_2(t)}t_0\){。}

{对每个实际矩阵，除以一个明确} \(R\) {单位，使其在模} \(m\) {上保持} \(v\) {不变。这是本原列的分析规范化，不是额外的物理除法。证书使用：} \[
-S,\qquad J_{s_0}/\tau,\quad J_{s_1}/\tau,\qquad
M_{uv}/\alpha,\qquad \operatorname{diag}(c^{-1},1).
\tag{C14.13}
\] {其中} \(M_{uv}=(-t,-1;m(t-s),t+a_0)\){，}\(r=4\) {时} \(\alpha=\epsilon_2^2\){，另一边界取} \(\alpha=-D/\tau\){。这些都是} \(R\) {单位，且分别表示} \(M\) {的模} \(m\) {倍率。最后一类来自已取得正向逆的} \(C_c\){：}

\begin{itemize}
\tightlist
\item
  \(r=4\){：使用} \(c=p,p-2,p-6,c_*\){，均为允许的小奇数且与} \(n\) {互素。}
\item
  \(r=p-3\){：使用} \(c=p,\tau,c_*\){；}\(C_\tau\) {是} \(C_{t_0}\){、}\(C_{c_*}\) {与换块的乘积。}
\end{itemize}

{下面的二进列提升引理对全部参数给出同一个证明。对} \(r=4\) {取精度} \(2^6\){，另一边界取} \(2^4\){；其生成元已在本节取得真实逆，故可以直接使用该引理，无须枚举参数或轨道。}

\paragraph*{\texorpdfstring{{从固定二进证书到完整有理终端}}{从固定二进证书到完整有理终端}}\label{C14-h105}
\addcontentsline{toc}{paragraph}{{从固定二进证书到完整有理终端}}

{先选} \(v^{-1}\pmod m\) {的平衡代表} \(2^h c\){，用} \(U(Qk)\) {保证本原后执行} \(C_c\){，再以} \(R\) {单位} \(\epsilon_2^h\) {规范化，使} \(v=1\pmod m\){。}

{用} {(13)} {把二进列精确送到} \((0,1)\){。选择第4节的非退化} \(G_s\){，并将它规范化为} \(G_s/\tau\){，其模} \(m\) {的} \(v\) {倍率为1。它在二进有限环上的阶为2幂，故可选择一个保持二进列的正向幂，同时在全部奇素数幂上把} \(u/v\) {送到0。}

{最终} \(u=0\pmod Q\){，而} \(v=1\) {同时成立于} \(m\) {和二进模数。令} \(\kappa=e/d=1+2m/d\){，取其在二进模数上的阶} \(T\){，则} \(T\) {是2幂。}\(\kappa^T\) {在奇数模数} \(m^2\) {上生成模} \(m\) {的整个单位核，且在二进部分为1。因此} \(v\bmod Q\) {由一个实际} \(R\) {单位} \(\epsilon\) {表示。}

{精确纤维引理现在把输出送到} \((0,\epsilon)\){，三次相反值平均归零。两种内部边界完成。}

\subsection*{\texorpdfstring{{C14.8} {全输入结论}}{C14.8 全输入结论}}\label{C14-h106}
\addcontentsline{toc}{subsection}{{C14.8} {全输入结论}}

{第5节覆盖全部奇} \(t\){；第6节覆盖有三个返回的偶} \(t\){；第7节补齐只有两个返回的} \(r=4\) {和} \(r=p-3\) {偶} \(t\){。因此所有} {(1)} {的合法核心均已完整归零。}

第\ref{construct:entry}节{已把所有合法原输入送到这些核心。必要性仍是非} \(p\) {共同非零余数不变量。定理得证。}

{这不是有限方向覆盖之后再次留下``有理实现''缺口：原位置基本返回、真逆、主同余包含、完整单位尺度和最终精确纤维提升都已分别落实。全长平均词可能很长，没有给出统一高效或最短构造。}


\section{二进本原列的统一有限提升}\label{construct:dyadic}
\subsection*{\texorpdfstring{{C15.4} {全参数二进列传递引理}}{C15.4 全参数二进列传递引理}}\label{C15-h107}
\addcontentsline{toc}{subsection}{{C15.4} {全参数二进列传递引理}}

{记} \(\mathcal O=\mathbb Z_{(2)}\){，即分母为奇数的有理数环。给定} \[
m\text{ 奇},\quad t\text{ 偶},\quad s_0\text{ 奇},\quad s_1=s_0-2,
\] {以及} \(\tau,\alpha\in\mathcal O^\times\){。令} \[
J_s=\begin{pmatrix}-t&-1\\m(t-s)&t\end{pmatrix},
\qquad d_s=m(t-s)-t^2.
\] {从} \(s_*=s_0,s_1\) {中任选一个，取} \(a\in2\mathcal O^\times\){，置} \[
M=J_{s_*}+aE_{22},\qquad
Z_i=J_{s_i}/\tau,\qquad T=M/\alpha.
\]

\textbf{{定理。}} {对每个整数} \(b\ge1\){，由} \(Z_0,Z_1,T\) {及其逆生成的群，在} \[
\Omega_b=\{(u,v)\in(\mathbb Z/2^b\mathbb Z)^2:\ u+v\text{ 奇}\}
\] {上传递。可从任意输入显式构造送到} \((0,1)\) {的有限词；不要求} \(p\) {是素数，不枚举参数类或状态空间。}

{这个定理是有限精度的群作用结论。应用于平均问题时，三个生成元的正向射影实现及真逆必须已经建立；本定理不把模数中的逆免费提升为实际逆。}

\paragraph*{\texorpdfstring{{4.1} {所有矩阵整可逆，且保持正确的模2集合}}{4.1 所有矩阵整可逆，且保持正确的模2集合}}\label{C15-h108}
\addcontentsline{toc}{paragraph}{{4.1} {所有矩阵整可逆，且保持正确的模2集合}}

\(d_s\) {是奇数，并且} \[
J_s^2=-d_sI,\qquad
d_{s_1}=d_{s_0}+2m,\qquad
\det M=d_{s_*}-at\text{ 奇}.
\] {三个规范化矩阵模2均为} \[
\begin{pmatrix}0&1\\1&0\end{pmatrix}.
\] {故它们及其逆都保持} \(\Omega_b\){。}

\paragraph*{\texorpdfstring{{4.2} {从已有返回制造一阶下方向}}{4.2 从已有返回制造一阶下方向}}\label{C15-h109}
\addcontentsline{toc}{paragraph}{{4.2} {从已有返回制造一阶下方向}}

{令} \(Z=Z_0\){，}\(Z_*=J_{s_*}/\tau\){，并令} \[
V=T Z_*^{-1},\qquad \kappa=\tau/\alpha.
\] {因为} \(J_{s_*}^{-1}=-J_{s_*}/d_{s_*}\){，有} \[
V=\kappa B,\qquad
B=\begin{pmatrix}1&0\\c&1+e\end{pmatrix},
\] {其中} \[
c=-\frac{am(t-s_*)}{d_{s_*}}\in2\mathcal O^\times,
\qquad e=-\frac{at}{d_{s_*}}\in4\mathcal O.
\tag{C15.4.1}
\] {因此} \[
B\equiv I+2E_{21}\pmod4.
\] {这不是说已有一个精确有理根} \(L(2)\){；这里只需要它的有限剩余类。}

{由于} \(d_{s_0},d_{s_1}\) {模4相差2，其中恰有一个为1。选该} \(i\){，令} \[
S=Z_i^2=-d_{s_i}\tau^{-2}I.
\] {奇单位平方模4为1，所以} \(S\equiv-I\pmod4\){。}

{取} \(\epsilon\in\{0,1\}\){，使} \((-1)^\epsilon\kappa\equiv1\pmod4\){。则} \[
L_*=S^\epsilon V\equiv I+2E_{21}\pmod4,\qquad
U_*=ZL_*Z^{-1}\equiv I+2E_{12}\pmod4.
\tag{C15.4.2}
\]

\paragraph*{\texorpdfstring{{4.3} {不用表格就能完成模4归一化}}{4.3 不用表格就能完成模4归一化}}\label{C15-h110}
\addcontentsline{toc}{paragraph}{{4.3} {不用表格就能完成模4归一化}}

{给定} \(u+v\) {奇：}

\begin{enumerate}
\def\labelenumi{\arabic{enumi}.}
\tightlist
\item
  {若} \(u\) {奇，先作用} \(Z\){，使} \(u\) {偶、}\(v\) {奇。}
\item
  {若} \(u\equiv2\pmod4\){，作用一次} \(U_*\){，使} \(u\equiv0\pmod4\){。}
\item
  {若} \(v\equiv3\pmod4\){，作用一次} \(S\){，使} \(v\equiv1\pmod4\){。}
\end{enumerate}

{于是模4精确到达} \((0,1)\){，包括尺度。模2情形只需第一步。}

\paragraph*{\texorpdfstring{{4.4} {交换子补出尺度方向}}{4.4 交换子补出尺度方向}}\label{C15-h111}
\addcontentsline{toc}{paragraph}{{4.4} {交换子补出尺度方向}}

{定义} \[
K_{21}=V^2,\qquad K_{12}=ZV^2Z^{-1},\qquad
K_0=[V,ZVZ^{-1}],
\] {交换子约定为} \([A,B]=ABA^{-1}B^{-1}\){。}

{这里下文的} \(XY-YX\) {是矩阵的李括号，与上述群交换子用展开式联系；不会将群交换子本身等同于线性李括号。}

{由} \(\kappa^2\equiv1\pmod8\){、}\(c\in2\mathcal O^\times\){、}\(e\in4\mathcal O\){，} \[
K_{21}\equiv I+4E_{21},\qquad
K_{12}\equiv I+4E_{12}\pmod8.
\tag{C15.4.3}
\] {交换子中的共同标量相消。对} \(B\equiv I+2E_{21}\pmod4\) {及其横向共轭，直接展开得到} \[
K_0\equiv I+4[E_{21},E_{12}]
       \equiv I+4I\pmod8.
\tag{C15.4.4}
\] {最后一步使用特征2中} \[
[E_{21},E_{12}]=\operatorname{diag}(-1,1)\equiv I.
\] {所以这一``中心方向''恰能修正列的尺度。}\(K_0\) {不是作为有理矩阵的标量；它在所需第一层上具有该作用，足以继续提升。}

{补充展开使符号可以直接核对：} \[
B^2=\begin{pmatrix}1&0\\c(2+e)&(1+e)^2\end{pmatrix}.
\] {因为} \(c/2\) {奇、}\(4\mid e\){，下左项模8为4、下右项模8为1。再乘} \(\kappa^2\equiv1\pmod8\){，即得(4.3)。若} \[
A=I+2X+4A_1,\qquad C=I+2Y+4C_1
\] {是在} \(\mathcal O\) {中整可逆的矩阵，展开其逆至模8给} \[
ACA^{-1}C^{-1}\equiv I+4(XY-YX)\pmod8.
\] {取} \(X\bmod2=E_{21}\){、}\(Y\bmod2=E_{12}\){，得(4.4)。模4以上的未知系数完全消去，所以不需要逐参数计算交换子。}

\paragraph*{\texorpdfstring{{4.5} {同一个平方规则提升到任意有限精度}}{4.5 同一个平方规则提升到任意有限精度}}\label{C15-h112}
\addcontentsline{toc}{paragraph}{{4.5} {同一个平方规则提升到任意有限精度}}

{若} \(K\equiv I+4X\pmod8\){，则对所有} \(k\ge2\){，} \[
K^{\,2^{k-2}}\equiv I+2^kX\pmod{2^{k+1}}.
\tag{C15.4.5}
\] {归纳只需平方：从} \(k\ge2\) {起，二次项} \(2^{2k}X^2\) {在下一层恒为零。}

{现在若当前列} \(w\equiv e_2\pmod{2^k}\){，写} \[
w=e_2+2^k(\delta_1,\delta_2)^T
             \pmod{2^{k+1}},\qquad \delta_i\in\mathbb F_2.
\] {因} \(E_{12}e_2=e_1\){、}\(Ie_2=e_2\){，依次作用} \[
K_{12}^{\,\delta_1 2^{k-2}},\qquad
K_0^{\,\delta_2 2^{k-2}}
\tag{C15.4.6}
\] {就把两项误差同时消掉，得到} \(w\equiv e_2\pmod{2^{k+1}}\){。}

{从模4开始迭代至} \(2^b\){，过程只有有限层。每层至多两个指定词的幂，不要求任何无穷} \(2\) {进极限。定理得证。}

\subsection*{\texorpdfstring{{C15.5} {与偶数内部构造的衔接}}{C15.5 与偶数内部构造的衔接}}\label{C15-h113}
\addcontentsline{toc}{subsection}{{C15.5} {与偶数内部构造的衔接}}

{旧}第\ref{construct:even-interior}节{的两种分支恰满足上述前提：}

{ % do not increment counter
\begin{longtable}[]{@{}
  >{\raggedright\arraybackslash}p{(\linewidth - 6\tabcolsep) * \real{0.2500}}
  >{\raggedright\arraybackslash}p{(\linewidth - 6\tabcolsep) * \real{0.2500}}
  >{\raggedright\arraybackslash}p{(\linewidth - 6\tabcolsep) * \real{0.2500}}
  >{\raggedright\arraybackslash}p{(\linewidth - 6\tabcolsep) * \real{0.2500}}@{}}
\toprule\noalign{}
\begin{minipage}[b]{\linewidth}\raggedright
{分支}
\end{minipage} & \begin{minipage}[b]{\linewidth}\raggedright
\(t,m\)
\end{minipage} & \begin{minipage}[b]{\linewidth}\raggedright
\(s_0,s_1\)
\end{minipage} & \begin{minipage}[b]{\linewidth}\raggedright
\(a,s_*\)
\end{minipage} \\
\midrule\noalign{}
\endhead
\bottomrule\noalign{}
\endlastfoot
\(r=4,\ p\ge11\) & \(2,p+2\) & \(3,1\) & \(2,1\) \\
\(r=p-3,\ p\equiv3\bmod4,\ p\equiv1\bmod3\) & \((p-3)/2,(3p-3)/2\) & \(p-4,p-6\) & \(2,p-4\) \\
{同上且} \(p\equiv2\bmod3\) & {同上} & {同上} & \(-2,p-6\) \\
\end{longtable}
}

\(p=7,r=4\) {使用第二行；不能使用第一行缺乏位置容量的计数。}

{这些} \(J_s,M\) {的实际位置计数、合法性、正向真逆和根理想扩环均在旧证明到达二进表之前已证明。具体地，先由已有深同余控制完成} \(M\) {的逆夹乘，再加入其行列式倍率，取得统一level。有限二进表没有参与这些先行论证。本构造只替换后面的二进覆盖步骤，所以没有用待证轨道去制造其自身生成元的循环。}

{旧规范化参数为} \[
\tau=\epsilon_2\quad(r=4),\qquad
\tau=\epsilon_2^{v_2(t)}\,t/2^{v_2(t)}
       \quad(r=p-3);
\] {相应} \(\alpha=\epsilon_2^2\) {或} \(-D/\tau\){。它们全为奇的} \(R\) {单位。}\(J_s/\tau\) {与} \(M/\alpha\) {模} \(m\) {的第二行均为} \((0,1)\){，因此本构造全部组合还保持} \[
v\equiv1\pmod m.
\] {除以这些单位是完整列的分析规范化；记录所得单位即可，不是额外实施一次物理数乘。}

{在旧证明准备好} \(v=1\pmod m\) {后：}

\begin{enumerate}
\def\labelenumi{\arabic{enumi}.}
\tightlist
\item
  {对} \(r=4\) {取} \(b=6\){，另一边界取} \(b=4\){，用本构造算法精确达到} \((0,1)\pmod{2^b}\){，并保持} \(v=1\pmod m\){。}
\item
  {沿用旧奇素数方向调整，以CRT使其不改变二进列。}
\item
  {沿用旧主单位核和精确同余纤维提升，真正到达有理终端并归零。}
\end{enumerate}

{因此64类和256类的枚举在该完整证明接口中均可取消；其它先行位置构造、根理想、单位准备和深核定理仍然使用原有证明。}


\section{上边界 n=3p-1 的重叠反射}\label{construct:r-last}
{本节设} \(p\ge5\) {为奇素数、}\(n=3p-1\){。一般公式取} \(p\ge7\){，}\(p=5\) {的最小容量由末段显式返回处理。}

\subsection*{\texorpdfstring{{C16.2} {核心、标准对合及自由角色交换}}{C16.2 核心、标准对合及自由角色交换}}\label{C16-h114}
\addcontentsline{toc}{subsection}{{C16.2} {核心、标准对合及自由角色交换}}

{置} \[
r=p-1,\quad n=3p-1,\quad d=p^2-3p+1.
\] {无条件下半带入口把任意合法输入送到} \[
K(a,z)=(a^p,(rz-a)^p,(-pz)^r),
\quad \gcd(a,z)=\gcd(a+pz,n)=1.
\tag{C16.2.1}
\] {标准两步对合及等大块角色交换为} \[
A=\begin{pmatrix}0&-d\\-1&0\end{pmatrix},
\quad A^2=dI,\qquad
S=\begin{pmatrix}-1&r\\0&1\end{pmatrix}.
\tag{C16.2.2}
\] {\ensuremath{A}} {的实际尺度为1/p，\ensuremath{S}只更新两个\ensuremath{p}块的角色标签。\ensuremath{d}与pn互素。}

{以下全部矩阵表示整体非零标量意义下的真实核心返回。所有位置共同缩放只用于分析，不作为额外物理操作。}

\subsection*{\texorpdfstring{{C16.3} {三步重叠对合}}{C16.3 三步重叠对合}}\label{C16-h115}
\addcontentsline{toc}{subsection}{{C16.3} {三步重叠对合}}

{给定} {\ensuremath{\sigma }和\ensuremath{h}，先平均} \[
(A^{p-\sigma},C^\sigma),
\] {得到\ensuremath{p}份新值} \[
D=\frac{p-\sigma}{p}a-\sigma z.
\] {保留其中\ensuremath{r=p-1}份作最终载体。其余位置为} {\ensuremath{\sigma }份旧\ensuremath{A}、\ensuremath{p}份旧\ensuremath{B}、r-sigma份旧\ensuremath{C}以及一份\ensuremath{D}。}

{平均} \[
(A^{h-\sigma},B^h,C^{p+\sigma-2h}),
\] {再平均其余\ensuremath{p}份，总计三次原子平均。只需} \[
0\le h-\sigma\le\sigma,\quad h\le p,\quad
0\le p+\sigma-2h\le p-1-\sigma.
\tag{C16.3.1}
\] {实际返回矩阵为\ensuremath{p^{-2}}乘} \[
J_{\sigma,h}=
\begin{pmatrix}
-p\sigma&p[nh-p(p+\sigma)]\\
-(p-\sigma)&p\sigma
\end{pmatrix}.
\] {其平方为} \[
\boxed{J_{\sigma,h}^2=p f_{\sigma,h}I,\qquad
f_{\sigma,h}=p^3-n(p-\sigma)h.}
\tag{C16.3.2}
\] {\ensuremath{f}模\ensuremath{n}为\ensuremath{p^{3}}，故非零且与\ensuremath{n}互素。重复三步返回提供真实正向射影逆。}

{对所有\ensuremath{p\ge 7}，以下四个参数同时满足(3.1)：} \[
(\sigma,h)=(2,3),(2,4),(3,4),(3,5).
\tag{C16.3.3}
\] {最紧的计数是p+2-2}\emph{{\ensuremath{4=p-6}和p+3-2}}{\ensuremath{5=p-7}，均非负。这是整个一般证明唯一的\ensuremath{p\ge 7}容量门槛。}

\subsection*{\texorpdfstring{{C16.4} {自由角色交换把根理想降到2n}}{C16.4 自由角色交换把根理想降到2n}}\label{C16-h116}
\addcontentsline{toc}{subsection}{{C16.4} {自由角色交换把根理想降到2n}}

{固定\ensuremath{\sigma =2}或3，取\ensuremath{h=\sigma +1}，记} \[
f_\sigma=f_{\sigma,h},\quad
g_\sigma=f_{\sigma,h+1}=f_\sigma-n(p-\sigma).
\] {由于\ensuremath{p}不整除\ensuremath{\sigma }，且\ensuremath{f_\sigma }与\ensuremath{n}互素，} \[
\gcd(f_\sigma,g_\sigma)
=\gcd(p^3,n(p-\sigma))=1.
\] {相邻对合的商为} \[
T_\sigma=J_{\sigma,h}^{-1}J_{\sigma,h+1}
=\begin{pmatrix}1&-p\sigma n/f_\sigma\\0&g_\sigma/f_\sigma\end{pmatrix}.
\] {它与角色交换\ensuremath{S}的交换子恰为} \[
\boxed{[S,T_\sigma]=
U\left(\frac{n[n\sigma-pr]}{g_\sigma}\right).}
\tag{C16.4.1}
\] {两个整数系数相差\ensuremath{n}，并且} \[
\gcd(2n-pr,3n-pr)=\gcd(n,pr)=2.
\tag{C16.4.2}
\] {令} \[
R_0=\mathbb Z[1/|f_2g_2f_3g_3|].
\] {两个倍率及其逆生成\ensuremath{R_{0}}，因而上根参数在\ensuremath{R_{0}}上封闭。式(C16.4.1)、(4.2)与Bezout给出\ensuremath{U(2nR_{0})}。}

{不需要先假定\ensuremath{d}可逆：对任意\ensuremath{x}属于\ensuremath{2nR_{0}}，已可执行U(dx)，而} \[
A\,U(dx)\,A^{-1}=L(x).
\] {因此也有\ensuremath{L(2nR_{0})}。深同余引理给出} \[
\boxed{\Gamma(4n^2,R_0)\text{ 正向可执行}.}
\tag{C16.4.3}
\] {2已整除\ensuremath{n}，故控制层没有任何额外素数。这比先产生含\ensuremath{r}或\ensuremath{p}的深根再逐项消除更直接。}

\subsection*{\texorpdfstring{{C16.5} {单位像满时，可以在方向调整后重新本原化}}{C16.5 单位像满时，可以在方向调整后重新本原化}}\label{C16-h117}
\addcontentsline{toc}{subsection}{{C16.5} {单位像满时，可以在方向调整后重新本原化}}

{以下小引理避免一个不必要的行列式局部化要求。}

\textbf{{局部本原化引理。}} {设\ensuremath{R}的反演素数均不除正整数\ensuremath{M}，\ensuremath{\Gamma (M,R)}已可执行，且\ensuremath{R}单位模\ensuremath{M}满射。给定本原整数列\ensuremath{v}，若一条实际正向路径的矩阵属于\ensuremath{\mathrm{GL}_{2}(Z_(M))}，其中\ensuremath{Z_(M)}表示分母与\ensuremath{M}互素的有理数环，并使其方向满足\ensuremath{z/y=0}模\ensuremath{M}、\ensuremath{y}为局部单位，则可以正向精确送到(0,epsilon)，\ensuremath{\epsilon }为某个\ensuremath{R}单位。}

{证明：将路径后的矩阵列清去与\ensuremath{M}互素的分母，得到整数列(\ensuremath{z_{0},y_{0})}。因矩阵在每个\ensuremath{\ell |M}处可逆，该列仍局部本原，故其整数坐标\ensuremath{gcd}与\ensuremath{M}互素。除这个\ensuremath{gcd}后得到本原整数列(z,y)，仍有\ensuremath{z=0}模\ensuremath{M}、\ensuremath{y}为单位。单位像满射给\ensuremath{\epsilon =y}模\ensuremath{M}。}

{取\ensuremath{\alpha }} {z+beta} {\ensuremath{y=1}，调整(alpha,beta)为(alpha+ty,beta-tz)，使\ensuremath{\alpha =0}模\ensuremath{M}。则} \[
P=\begin{pmatrix}
y/\epsilon&-z/\epsilon\\
\epsilon\alpha&\epsilon\beta
\end{pmatrix}
\in\Gamma(M,R),\qquad
P(z,y)^T=(0,\epsilon)^T.
\tag{C16.5.1}
\] {\ensuremath{\Gamma (M,R)}提供真实正向实现。实际状态和分析列之间仅差共同非零尺度，因此同样产生零终端。证毕。}

{引理不允许忽略本原化；它明确追踪被除因子在模\ensuremath{M}下是单位。只有在\ensuremath{R}单位像满时，才可不将路径中每个行列式因子预先反演。若单位像只是一个真子群，原来保留\ensuremath{R}本原性的要求仍然重要。}

\subsection*{\texorpdfstring{{C16.6} {奇数数字与乘3折返给全部模\ensuremath{n}单位}}{C16.6 奇数数字与乘3折返给全部模单位}}\label{C16-h118}
\addcontentsline{toc}{subsection}{{C16.6} {奇数数字与乘3折返给全部模\ensuremath{n}单位}}

{保留旧\ensuremath{C}的两块混合为} \[
F_t=\frac1p
\begin{pmatrix}
t&r(p-t)/2\\0&p
\end{pmatrix},
\qquad t\text{ 奇},\quad -p\le t\le p.
\tag{C16.6.1}
\] {它的原位置两组分别取(p-j,j)和(j,p-j)份\ensuremath{A}、\ensuremath{B}，其中\ensuremath{j=(p-t)/2}。}

{\ensuremath{F_{1}}的整数代表行列式为\ensuremath{p}，与\ensuremath{4n^{2}}互素，故既有安全整数返回取逆引理适用。其上三角倍率反演\ensuremath{p}，先将\ensuremath{p}加入\ensuremath{R}。随后} \[
C_b=F_bF_1^{-1}
=\begin{pmatrix}b&r(1-b)/2\\0&1\end{pmatrix}
\] {对每个奇数\ensuremath{1\le b\le p}、\ensuremath{\operatorname{gcd}(b,n)=1}都实际存在。行列式\ensuremath{b}与控制层互素，故\ensuremath{C_{b}}的逆可执行。上下根随倍率共轭扩张到这些\ensuremath{b}全部反演后的有限局部化环\ensuremath{R}，且\ensuremath{\Gamma (4n^{2},R)}继续成立。}

{因为\ensuremath{n=3p-1}不被3整除，3属于这些安全数字。对任意模\ensuremath{n}单位\ensuremath{u}，取正平衡绝对代表\ensuremath{b<n/2}；\ensuremath{n}偶使\ensuremath{b}必奇。若\ensuremath{b>n/4}，以} \[
b\longmapsto |n-3b|
\tag{C16.6.2}
\] {替代。结果仍是正奇数、与\ensuremath{n}互素，并严格减小；它只改变\ensuremath{u}所乘的符号和3的幂。最终} \[
1\le b\le n/4<p.
\] {这个\ensuremath{b}已是\ensuremath{R}单位。倒推(6.2)，得到\ensuremath{u}的\ensuremath{R}单位提升。因此} \[
R^\times\longrightarrow(\mathbb Z/n\mathbb Z)^\times
\quad\text{满射}.
\tag{C16.6.3}
\] {没有要求3是模\ensuremath{n}的原根，也没有单位代表表。}

\subsection*{\texorpdfstring{{C16.7} {同一条方向循环适用于全部\ensuremath{p\ge 7}}}{C16.7 同一条方向循环适用于全部}}\label{C16-h119}
\addcontentsline{toc}{subsection}{{C16.7} {同一条方向循环适用于全部\ensuremath{p\ge 7}}}

{取} \[
V=J_{3,4}A/p^3,
\] {在整幺模(\ensuremath{z,y)=(z,a+pz)}坐标中有} \[
V=\begin{pmatrix}1&-3/p^2\\0&1\end{pmatrix}\pmod n,
\] \[
\operatorname{tr}V=2-\frac{n(5p-3)}{p^3}=2\pmod4.
\tag{C16.7.1}
\] {平移系数在所有\ensuremath{\ell |n}处均为单位，因为3不整除\ensuremath{n}。迹的模4条件适用于全部\ensuremath{p}，不必选择奇偶分支。}

{已有局部射影循环引理因此适用于所有2幂和所有奇素数幂\ensuremath{\ell |n}。指数CRT用\ensuremath{V}的一次非负幂将任意合法方向送到} \[
z/y=0\pmod{4n^2}.
\tag{C16.7.2}
\] {\ensuremath{A}的行列式可能未在\ensuremath{R}中反演，但它与\ensuremath{n}互素；第5节正好允许在方向调整后处理这个共同内容。}

\subsection*{\texorpdfstring{{C16.8} {从\ensuremath{n}到\ensuremath{4n^{2}}的全部单位尺度}}{C16.8 从到的全部单位尺度}}\label{C16-h120}
\addcontentsline{toc}{subsection}{{C16.8} {从\ensuremath{n}到\ensuremath{4n^{2}}的全部单位尺度}}

{已有\ensuremath{R}单位} \[
\kappa_2=g_2/f_2=1-\frac{n(p-2)}{f_2},\quad
\kappa_3=g_3/f_3=1-\frac{n(p-3)}{f_3}.
\] {对任意奇素数\ensuremath{\ell |n}，有\ensuremath{\operatorname{gcd}(p-3,n)}只含2，因此} \[
v_\ell(\kappa_3-1)=v_\ell(n).
\] {\ensuremath{\kappa _{3}}的幂生成全部奇素数处的主单位核。}

{二进处，p-2为奇数，故} \[
v_2(\kappa_2-1)=v_2(n).
\] {若\ensuremath{4|n}，先用\ensuremath{\kappa _{2}}调好全部二进位；随后用\ensuremath{\kappa _{3}}的完整二进阶次幂修复奇素数处。该阶是2幂，不改变任何奇素数处的LTE起始深度。}

{若\ensuremath{n=2}模4，则\ensuremath{p=1}模4，\ensuremath{v_{2}(p-3)=1}。\ensuremath{\kappa _{2}=3}模4补两类中的一类，\ensuremath{\kappa _{3}}从1模4开始生成余下全部二进主单位。再同样用完整二进阶次幂修复奇素数处。于是主单位核全部生成。}

{接(6.3)，\ensuremath{R}单位模\ensuremath{4n^{2}}满射。第5节把(7.2)精确送到(0,epsilon)。核心成为} \[
(\epsilon^p,(-\epsilon)^p,0^{p-1}).
\] {以(p-1)/2份\ensuremath{\epsilon }、同样多份-\ensuremath{\epsilon }和一个零组成零和\ensuremath{p}元组，即可触发统一零收尾。这完成全部\ensuremath{p\ge 7}。}

\subsection*{\texorpdfstring{{C16.9} {五平均14元：一个显式夹乘补齐最小容量}}{C16.9 五平均14元：一个显式夹乘补齐最小容量}}\label{C16-h121}
\addcontentsline{toc}{subsection}{{C16.9} {五平均14元：一个显式夹乘补齐最小容量}}

{这里\ensuremath{p=5}、\ensuremath{r=4}、\ensuremath{n=14}。只需以下三个已实现变换：} \[
A=\begin{pmatrix}0&-11\\-1&0\end{pmatrix},\quad
S=\begin{pmatrix}-1&4\\0&1\end{pmatrix},\quad
J=J_{1,2}=\begin{pmatrix}-5&-10\\-4&5\end{pmatrix}.
\] {\ensuremath{A^{2}=11I}、\ensuremath{J^{2}=65I}，\ensuremath{J}仍是第3节的三原子返回。}

{令} \[
W=ASASJ=
\begin{pmatrix}-121&330\\0&65\end{pmatrix}.
\] {直接得到} \[
[S,W]=U(84/65).
\tag{C16.9.1}
\] {这个上三角词也可以由有理方向确定：\ensuremath{J}送无穷远到5/4，\ensuremath{S}送到11/4，\ensuremath{A}送到4，\ensuremath{S}送到0，再由\ensuremath{A}送回无穷远。不是平均词搜索的输出。}

{由于121与65互素，\ensuremath{W}的倍率及其逆生成} \[
R_0=\mathbb Z[1/(5\cdot11\cdot13)].
\] {所以\ensuremath{U(84R_{0})}、\ensuremath{L(84R_{0})}及\ensuremath{\Gamma (7056,R_{0})}已可执行。}

{\ensuremath{F_{1}}的整数行列式5与7056互素，故可激活其逆，得到} \[
C_3=F_3F_1^{-1}=\begin{pmatrix}3&-4\\0&1\end{pmatrix}.
\] {虽然3整除旧控制层，下面的夹乘直接提供其逆：} \[
H=\frac13C_3JC_3
\ =\begin{pmatrix}1&-18\\-4&7\end{pmatrix},
\quad \det H=-65=\det J.
\tag{C16.9.2}
\] {\ensuremath{HJ^{-1}}属于\ensuremath{\mathrm{SL}_{2}(R_{0})}，且正向可执行。有限逆提升引理因而给出\ensuremath{H^{-1}}，再用} \[
JC_3H^{-1}=3C_3^{-1}
\] {取得\ensuremath{C_{3}}的真正正向射影逆。}

{共轭上根反演3。令\ensuremath{R=Z}{[}1/(3}\emph{{5}}{11*13){]}，则\ensuremath{84R=28R}，再次用深同余引理得到} \[
\boxed{\Gamma(784,R).}
\tag{C16.9.3}
\]

\paragraph*{\texorpdfstring{{9.1} {不用有限轨道图的方向与尺度}}{9.1 不用有限轨道图的方向与尺度}}\label{C16-h122}
\addcontentsline{toc}{paragraph}{{9.1} {不用有限轨道图的方向与尺度}}

{取\ensuremath{V=JA/125}。其模14平移为-1/25，迹为\ensuremath{2-196/125=2}模4。因此一个\ensuremath{V}的CRT指数覆盖模16和模49的全部合法仿射方向。}

{以下三个实际\ensuremath{R}单位的余数为} \[
\begin{array}{c|cc}
&\bmod16&\bmod7\\ \hline
15&-1&1\\
10725&5&1\\
2145&1&3
\end{array}
\] {因为-1和5生成模16单位群，3生成模7单位群，这三者已覆盖模112的所有单位。它们分别为3}\emph{{5、3}}{\ensuremath{5^{2}}}\emph{{11}}{13、3}\emph{{5}}{11*13。}

{再令\ensuremath{\lambda =-121/65}，有} \[
\lambda^2-1=\frac{10416}{4225},
\quad v_2(10416)=4,\quad v_7(10416)=1.
\] {故\ensuremath{\lambda ^{2}}模16为1，并生成模49到模7的七阶单位核。于是\ensuremath{R}单位模784满射。}

{由(9.3)、第5节和同一个相反值终端，全部合法14元核心及全输入均归零。}


\section{小数字、二进修正及额外二进层的终端}\label{construct:small-digit}
\subsection*{\texorpdfstring{{C17.1} {已有的对称坐标揭示一个乘法群}}{C17.1 已有的对称坐标揭示一个乘法群}}\label{C17-h123}
\addcontentsline{toc}{subsection}{{C17.1} {已有的对称坐标揭示一个乘法群}}

{九偏移核心写作} \[
\mathcal K(a,z)=
\bigl(a^p,(rz-a)^p,(-pz)^r\bigr),
\qquad r=p+s,\ n=2p+r=3p+s.
\] {令} \[
y=a+pz,\qquad \eta=y-\frac n2z=a-\frac r2z.
\tag{C17.1.1}
\] {\ensuremath{\eta }是两个等大\ensuremath{p}块数值之差的一半。\ensuremath{z}记录它们共同的均值方向，\ensuremath{\eta }记录差方向。换块使\ensuremath{\eta }变号而不改\ensuremath{z}。}

{旧两个互补\ensuremath{p}平均给} \[
F_q:
(a,z)\longmapsto
\left(\frac qp\,a+\frac{r(p-q)}{2p}z,\ z\right),
\qquad q\text{奇},\quad1\le q\le p.
\] {在(z,η)坐标中，它恰为\ensuremath{\operatorname{diag}(1,q/p)}。在已有根和深核之后，\ensuremath{F}₁及安全\ensuremath{F_{q}}的正向逆已由旧夹乘定理取得，所以} \[
C_q=F_qF_1^{-1}
\] {给\ensuremath{\operatorname{diag}(1,q)}。以一个共同非零尺度选择射影矩阵代表，令\ensuremath{D_{q}=q}⁻¹\ensuremath{C_{q}}，于是} \[
\boxed{
D_q=\begin{pmatrix}
q^{-1}&0\\
\frac n2(q^{-1}-1)&1
\end{pmatrix}
\quad\text{在}(z,y)\text{基中};
\qquad
(z,\eta)\mapsto(q^{-1}z,\eta).
}
\tag{C17.1.2}
\] {\ensuremath{D_{q}}不需要执行一次额外的\ensuremath{q}⁻¹数乘；实际词为\ensuremath{C_{q}}，\ensuremath{D_{q}}只是同一射影作用的代表。}

{当\ensuremath{n}偶、\ensuremath{q}奇时，n(q⁻¹−1)/2属于nR。因此} \[
\boxed{D_q:\quad z\mapsto q^{-1}z,\quad y\mapsto y\pmod n.}
\tag{C17.1.3}
\] {全部分母来自已真实反演的小数字\ensuremath{q}。换块的负代表} \[
D_{-1}=-S=\begin{pmatrix}-1&0\\-n&1\end{pmatrix}
\] {也在这一乘法族中。}

{这就是此处``分裂环面''的具体含义：两个自然特征方向中的一个缩放，另一个保持。它已存在于原平均操作中，不是引入新的抽象可执行性假设。}

\subsection*{\texorpdfstring{{C17.2} {所需的小数字总存在：一个一般计数引理}}{C17.2 所需的小数字总存在：一个一般计数引理}}\label{C17-h124}
\addcontentsline{toc}{subsection}{{C17.2} {所需的小数字总存在：一个一般计数引理}}

\textbf{{引理。}} {设\ensuremath{p}\ensuremath{\ge}11为整数，\ensuremath{n}为偶数且\ensuremath{n<4p}。则存在奇整数\ensuremath{q}满足} \[
1<q\le p,\qquad q\equiv3\text{或}5\pmod8,\qquad\gcd(q,n)=1.
\tag{C17.2.1}
\]

{证明先检查3、5、11、13。\ensuremath{p=11}或12时，若3、5、11都不与\ensuremath{n}互素，\ensuremath{n}被165整除，与\ensuremath{n<4p}\ensuremath{\le}48矛盾。所以此时已有合适\ensuremath{q}。}

{设\ensuremath{p}\ensuremath{\ge}13。如果四个数字都不与\ensuremath{n}互素，则3·5·11·13整除\ensuremath{n}。记\ensuremath{n}的不同奇素因子数为\ensuremath{d}\ensuremath{\ge}4，令} \[
\rho=\prod_{\ell\mid n,\ \ell\text{奇}}(1-\ell^{-1}).
\] {考察\ensuremath{q=3+8t}、1\ensuremath{\le}\ensuremath{q}\ensuremath{\le}\ensuremath{p}组成的等差数列，项数\ensuremath{L}\ensuremath{\ge}(p−4)/8。每个奇平方自由除数在该数列中的计数与\ensuremath{L}除以该除数相差小于1。容斥故给与\ensuremath{n}互素的项数至少} \[
L\rho-2^d.
\] {因\ensuremath{p>n/4}且\ensuremath{\rho }\ensuremath{\le}1，} \[
L\rho>\frac{n\rho}{32}-\frac12.
\] {而} \[
n\rho\ge2\prod_{\ell\mid n,\ \ell\text{奇}}(\ell-1)
\ge1920\cdot6^{d-4}.
\] {前四个因子来自3、5、11、13；每增加一个不同奇素数，其\ensuremath{\ell}−1至少为6。因此} \[
L\rho>60\cdot6^{d-4}-\frac12>2^d.
\] {至少一个\ensuremath{q}符合要求。证毕。}

{这个引理不假设\ensuremath{p}或\ensuremath{n}为素数，也不要求分解所有候选\ensuremath{q}。算法可先检查四个小数，必要时沿3} {mod8的数列计算\ensuremath{gcd}。有限测试不是该存在性证明的依据。}

\subsection*{\texorpdfstring{{C17.3} {同一个二进单位规则}}{C17.3 同一个二进单位规则}}\label{C17-h125}
\addcontentsline{toc}{subsection}{{C17.3} {同一个二进单位规则}}

{若\ensuremath{q}\ensuremath{\equiv}3或5} {mod8，则} \[
v_2(q^2-1)=3.
\] {因此\ensuremath{q^{2}}的幂在模\ensuremath{2^L}下生成全部1} {mod8的单位，−1和\ensuremath{q}再覆盖模8的四个奇类，得到} \[
\boxed{\langle-1,q\rangle=(\mathbb Z/2^L\mathbb Z)^\times
\qquad(L\ge3).}
\tag{C17.3.1}
\] {这是逐次平方的初等推论：若\ensuremath{u=1+8c}、\ensuremath{c}奇，则\ensuremath{u^{2^{j}}-1}的二进赋值为3+j。先选符号及\ensuremath{q}的指数奇偶匹配模8，再逐位选择\ensuremath{q^{2^{j+1}}}的有无，就能匹配任意奇单位，不需要离散对数黑箱或有限群轨道表。}

{(3.1)只主张这些已可用整数单位的二进像足够，不主张它们单独生成所有奇数模数的单位。后者仍由旧基础单位模\ensuremath{n}满射承担。}

\subsection*{\texorpdfstring{{C17.4} {方向只从原根反射生成}}{C17.4 方向只从原根反射生成}}\label{C17-h126}
\addcontentsline{toc}{subsection}{{C17.4} {方向只从原根反射生成}}

{沿用}第\ref{sec:finite}节{：每条根反射\ensuremath{J}的计数差\ensuremath{\alpha =i-j}给} \[
T_\alpha=S\widehat J
\equiv\begin{pmatrix}1&-\alpha/p^2\\0&1\end{pmatrix}\pmod n.
\] {旧根供应已要求\ensuremath{gcd}} {\ensuremath{\alpha _{i}=1}；整数Bezout给\ensuremath{c_{i}}满足\ensuremath{\Sigma c_{i}\alpha _{i}=1}。设} \[
\boxed{V=\prod_iT_{\alpha_i}^{c_i}
\equiv\begin{pmatrix}1&-p^{-2}\\0&1\end{pmatrix}\pmod n.}
\tag{C17.4.1}
\] {本构造不再对tr} {\ensuremath{V}模4提任何条件。}

\textbf{{环与本原性的账本。}} {\ensuremath{R}取旧根对约分倍率实际生成的环，再加入\ensuremath{p}及所有安全小数字的素因子。旧根供应的核验检查了每条所用根反射的平方系数均为\ensuremath{R}单位，因此}\(\widehat J\){、其逆及\ensuremath{S}都在GL₂(R)，故\ensuremath{V}也在GL₂(R)。这比上一轮使用任意独立direction反射时的前提强，但这里直接由已有根供应得到，不需要新增局部化。}

{对于} \((p,s)=(11,3)\){，使用守恒输出交换构造所列的} \(\alpha=2,3\) {反射及其深核。整数等式} \(-2+3=1\) {给出同一粗略平移；不需要辅助13的局部群计算。}

{因此所有后续矩阵都在同一个GL₂(R)中。当前\ensuremath{R}本原列经它们作用仍\ensuremath{R}本原；共同有理尺度只作分析代表，不额外平均或扩增位置。}\textbf{{无需每次重新取整数本原代表，也无需Smith准备。}}

\subsection*{\texorpdfstring{{C17.5} {共同前提与目标}}{C17.5 共同前提与目标}}\label{C17-h127}
\addcontentsline{toc}{subsection}{{C17.5} {共同前提与目标}}

{令\ensuremath{\delta =1}（\ensuremath{n}奇）或\ensuremath{\delta }\ensuremath{\in}\{2,4\}（\ensuremath{n}偶），\ensuremath{h=\delta n}，\ensuremath{Q=h^{2}}。旧\ensuremath{root_{pairs}}、基础小数字和Morris／Serre给} \[
L(hR),\quad\Gamma(Q,R),\quad
R^\times\twoheadrightarrow(\mathbb Z/n\mathbb Z)^\times.
\] {全输入入口和零触发尾部沿用旧证明。它们均先于本构造终端过程，不需要direction或\ensuremath{binary_{unit}}。}

{显式终端仍为} \[
t_j=(1,p^2-nj),\qquad0\le j<\delta.
\] {选j+1份\ensuremath{A}、\ensuremath{j}份\ensuremath{B}、p−1−2j份\ensuremath{C}就恰好零和。\ensuremath{p}\ensuremath{\ge}11、\ensuremath{r=p+s}保证非负容量；\ensuremath{t_{j}}第一坐标为1，自动本原。}

{从任意合法\ensuremath{R}本原列(z,y)出发，先以\ensuremath{V}的非负幂把它变为} \[
p^2z\equiv y\pmod n.
\] {再取\ensuremath{\epsilon }\ensuremath{\in}\ensuremath{R}×使\ensuremath{\epsilon }\ensuremath{\equiv}\ensuremath{z}} {modn，得到} \[
(z,y)\equiv\epsilon(1,p^2)\pmod n.
\tag{C17.5.1}
\] {此时\ensuremath{z}、\ensuremath{y}在二进处均为单位（\ensuremath{n}偶时）；全程保持\ensuremath{R}本原性。}

\subsection*{\texorpdfstring{{C17.6} {偶数\ensuremath{n}的首坐标：先用乘法自由度完成}}{C17.6 偶数的首坐标：先用乘法自由度完成}}\label{C17-h128}
\addcontentsline{toc}{subsection}{{C17.6} {偶数\ensuremath{n}的首坐标：先用乘法自由度完成}}

{写\ensuremath{n=2}ᵃ\ensuremath{m}，\ensuremath{m}奇，\ensuremath{a}\ensuremath{\ge}1；记\ensuremath{h=2}ᵇ\ensuremath{m}、\ensuremath{v=2}ᵇ、\ensuremath{M=v^{2}}，则\ensuremath{Q=Mm^{2}}且\ensuremath{b}\ensuremath{\ge}2。}

{按第2节取\ensuremath{q}，利用(3.1)求} \[
(-1)^e q^t\equiv z/\epsilon\pmod M,
\qquad e\in\{0,1\},\ t\ge0.
\] {执行\ensuremath{D_{-1}^{e}D_{q}^{t}}。由(1.2)---(1.3)得到} \[
z'\equiv\epsilon\pmod M,\qquad y'\equiv\epsilon p^2\pmod n.
\tag{C17.6.1}
\] {这一步可能改变\ensuremath{z}的奇数模\ensuremath{m}分量，下面用原\ensuremath{V}修正，但冻结整个二进分量。}

{由于\ensuremath{V}} {mod2是非平凡上平移，\ensuremath{V^{2}}\ensuremath{\equiv}\ensuremath{I}} {mod2，逐次平方得} \[
\boxed{V^M\equiv I\pmod M.}
\tag{C17.6.2}
\] {这个结论对任何这样的\ensuremath{V}成立，无需tr} {\ensuremath{V}\ensuremath{\equiv}2} {mod4。\ensuremath{M}是2幂，故模\ensuremath{m}可逆。取} \[
k=M\left[p^2(z'-\epsilon)(y')^{-1}M^{-1}\right]_m.
\] {执行\ensuremath{V^{k}}，得到} \[
\boxed{z''\equiv\epsilon\pmod M,\quad
z''\equiv\epsilon\pmod m,\quad
y''\equiv\epsilon p^2\pmod n.}
\tag{C17.6.3}
\] {因\ensuremath{M}被\ensuremath{v}整除，第一坐标已经匹配模\ensuremath{h}，而且二进最终精度一次到位。}

{这里没有使用\ensuremath{V}在二进方向上形成长循环。\ensuremath{V}只负责奇数处加法，乘法\ensuremath{D_{q}}负责二进处缩放，两者用指数的互素性隔离。}

{\ensuremath{n}奇时\ensuremath{\delta =1}，(5.1)已是模\ensuremath{h}目标，跳过本节。}

\subsection*{\texorpdfstring{{C17.7} {终端选择、两根消元、一次下根}}{C17.7 终端选择、两根消元、一次下根}}\label{C17-h129}
\addcontentsline{toc}{subsection}{{C17.7} {终端选择、两根消元、一次下根}}

{选择} \[
j\equiv\epsilon^{-1}\frac{\epsilon p^2-y''}{n}\pmod\delta,
\quad0\le j<\delta.
\] {于是当前列\ensuremath{w}和目标\ensuremath{w_{T}=\epsilon t_{j}}相符模\ensuremath{h}。}

{令\ensuremath{E=E}₂₁、\ensuremath{X=VEV}⁻¹。在模\ensuremath{m}目标列上，} \[
\det(Ew_T,Xw_T)=2\epsilon^2/p^2
\] {为单位。清除\ensuremath{X}的分母后，按}第\ref{sec:finite}节{解一次二维线性方程；根参数通过CRT选为在二进部分恒等，得到\ensuremath{w}′\ensuremath{\equiv}\ensuremath{w_{T}}} {\ensuremath{modm^{2}}。由于本构造\ensuremath{V}\ensuremath{\in}GL₂(R)，\ensuremath{X}本来就在\ensuremath{M}₂(R)，清分母仅便于用整数系数计算，不扩大环。}

{\ensuremath{n}奇时\ensuremath{Q=m^{2}}，已全部完成。}

{\ensuremath{n}偶时第一坐标保持\ensuremath{\epsilon }} {modM；最后用一个已有L(hr)修正第二坐标模\ensuremath{M}，并选\ensuremath{r}\ensuremath{\equiv}0} {modm保持\ensuremath{m^{2}}部分。这就是上轮第7节的最后一个线性同余式，不再需要规范化循环。得到} \[
w_{\rm out}\equiv w_T\pmod Q.
\] {两列都\ensuremath{R}本原，由旧Γ(Q,R)精确纤维真实到达\ensuremath{w_{T}}；随后零和组平均产生\ensuremath{p}个零，接旧尾部完整归零。}


\section{上半带的大参数共同资源}\label{construct:common-controller}
{本节用于} \(p\ge83\){、}\(p<r<2p\){、}\(n=2p+r\){。第3段的计数也列出其可同时覆盖的部分下半带，但本证明的下半带由前述各构造承担。}

\subsection*{\texorpdfstring{{C18.2} {一组或两组来源：同一实际反射}}{C18.2 一组或两组来源：同一实际反射}}\label{C18-h130}
\addcontentsline{toc}{subsection}{{C18.2} {一组或两组来源：同一实际反射}}

{由}第\ref{construct:common-shape}节{，所有本范围合法输入可以进入} \[
K_{p,r}(a,z)=(A^p,B^p,C^r),\quad
A=a,\quad B=rz-a,\quad C=-pz,
\] \[
\gcd(a,z)=\gcd(a+pz,n)=1.
\] {入口由双块安全入口、上半带三值压缩及两次核心转换组成。}

{选\ensuremath{q}\ensuremath{\in}\{1,2\}，来源\ensuremath{c=(i,j,k)}、\ensuremath{i+j+k=p}。在原位置上平均\ensuremath{q}组不交且组成相同的\ensuremath{p}元组，产生qp份} \[
D=((i-j)a+(rj-pk)z)/p.
\] {保留\ensuremath{r}份\ensuremath{D}作新载体。剩余库存恰为} \[
(p-qi,\ p-qj,\ r-qk,\ qp-r).
\tag{C18.2.1}
\] {选输出\ensuremath{o=(x,y,u,v)}，非负且不超过(2.1)，总数\ensuremath{p}；平均该组及其\ensuremath{p}元补组。共q+2次原子平均，重新得到相同核心形状。}

{令\ensuremath{\alpha =i-j>0}、\ensuremath{\beta =rj-pk}。实际返回是\ensuremath{p}⁻²\ensuremath{J}，其中} \[
J=\begin{pmatrix}
p(x-y)+\alpha v&pry-p^2u+\beta v\\
-\alpha&-\beta
\end{pmatrix}.
\tag{C18.2.2}
\] {若} \[
p(x-y)+\alpha v=\beta,
\tag{C18.2.3}
\] {则\ensuremath{J}迹零且\ensuremath{J^{2}=pfI}。由}第\ref{construct:common-shape}节{，} \[
pf\equiv p^4\pmod n,
\qquad
\widehat J=p^{-2}J\equiv
\begin{pmatrix}1&-\alpha/p^2\\0&-1\end{pmatrix}\pmod n
\tag{C18.2.4}
\] {（第二式在(\ensuremath{z,y_*)=(z,a+pz)}基中）。因此返回非奇异、合法，重复它即给全系统非零标量，提供真实正向射影逆。\ensuremath{q}只改变容量和原子次数，不改变(2.2)---(2.4)。}

{所有以下反射成对选取，第二输出为} \[
o'=o+(-1,-1,2,0).
\] {记相应平方系数为f,g，则\ensuremath{g-f=\alpha n}。}

\subsection*{\texorpdfstring{{C18.3} {四个容量选择覆盖整个定理范围}}{C18.3 四个容量选择覆盖整个定理范围}}\label{C18-h131}
\addcontentsline{toc}{subsection}{{C18.3} {四个容量选择覆盖整个定理范围}}

\paragraph*{\texorpdfstring{{3.1} {随载体余量变化的计数：下半段与上端共用}}{3.1 随载体余量变化的计数：下半段与上端共用}}\label{C18-h132}
\addcontentsline{toc}{paragraph}{{3.1} {随载体余量变化的计数：下半段与上端共用}}

{用于\ensuremath{\operatorname{ceil}(p/3)+4}\ensuremath{\le}\ensuremath{r<p}，以及\ensuremath{r>3p/2}。令} \[
q=\begin{cases}1&r<p,\\2&r>3p/2,\end{cases}
\quad d=qp-r,\quad
a_0=\left\lceil\frac{d+4}{q}\right\rceil.
\] {对\ensuremath{\alpha =a}₀,\ensuremath{a}₀+1，来源取} \[
\boxed{(i,j,k)=(\alpha,0,p-\alpha).}
\tag{C18.3.1}
\] {输出取} \[
\boxed{o_l=(l,p-\alpha+l,\alpha-2l,0).}
\tag{C18.3.2}
\] {总数与迹零自动成立。所有来源和输出容量由} \[
\boxed{
\max\left(0,\left\lceil\frac{d-(q-1)\alpha}{2}\right\rceil\right)
\le l\le
\min\left(\lfloor\alpha/2\rfloor,p-q\alpha\right)
}
\tag{C18.3.3}
\] {保证。设左端为\ensuremath{L}，取\ensuremath{l=L+1,L}；第二张表减一份\ensuremath{A}、一份\ensuremath{B}、增加两份\ensuremath{C}。}

\textbf{{\ensuremath{q=1}的容量。}} {此时\ensuremath{\alpha =d+4}或d+5、\ensuremath{L=\operatorname{ceil}(d/2)}。两种\ensuremath{\alpha }均有\ensuremath{\operatorname{floor}(\alpha /2)}\ensuremath{\ge}L+1。另一上界要求} \[
p\ge d+\lceil d/2\rceil+6.
\] {由\ensuremath{r}\ensuremath{\ge}\ensuremath{\operatorname{ceil}(p/3)+4}，有\ensuremath{d}\ensuremath{\le}\ensuremath{\operatorname{floor}(2p/3)-4}。代入\ensuremath{p=3k+1}或3k+2，右边最大分别为p−1或\ensuremath{p}，故成立。来源剩余\ensuremath{C}为α−d\ensuremath{\ge}4，其他库存也非负。}

\textbf{{\ensuremath{q=2}的容量。}} {此时\ensuremath{0<d<p/2}，\ensuremath{\alpha =\operatorname{ceil}(d/2)+2}或\ensuremath{\operatorname{ceil}(d/2)+3}。因为\ensuremath{d}\ensuremath{\le}2α−4且\ensuremath{\alpha }\ensuremath{\ge}3，(3.3)左端\ensuremath{L}满足L+1\ensuremath{\le}\ensuremath{\operatorname{floor}(\alpha /2)}。又} \[
2\alpha\le d+7,\qquad L+1\le d/4+1,
\] {所以p−2α\ensuremath{\ge}L+1由\ensuremath{p}\ensuremath{\ge}5d/4+8保证。\ensuremath{d}\ensuremath{\le}(p−1)/2且\ensuremath{p}\ensuremath{\ge}83蕴含该条件。来源剩余\ensuremath{C}为2α−d\ensuremath{\ge}4，其余库存同样非负。}

{因此(3.1)---(3.3)共用一个公式覆盖下半带的一整段与靠近4p的整段。每处只有两组相邻\ensuremath{\alpha }、每组两张实际反射表。}

\paragraph*{\texorpdfstring{{3.2} {平衡来源：覆盖上半带中部}}{3.2 平衡来源：覆盖上半带中部}}\label{C18-h133}
\addcontentsline{toc}{paragraph}{{3.2} {平衡来源：覆盖上半带中部}}

{设\ensuremath{r=p+s}，10\ensuremath{\le}\ensuremath{s}\ensuremath{\le}\ensuremath{\operatorname{floor}(p/2)}。取\ensuremath{q=2}、\ensuremath{a}₀=\ensuremath{\operatorname{ceil}(p/6)}，对\ensuremath{\alpha =a}₀,\ensuremath{a}₀+1取} \[
\boxed{(i,j,k)=(2\alpha,\alpha,p-3\alpha),}
\] \[
\boxed{o_e=(\alpha-e,p-3\alpha-e,2\alpha-s+2e,s).}
\tag{C18.3.4}
\] {来源总数、输出总数和迹零恒成立。精确可行区间为} \[
L=\max(5\alpha-p,\lceil s/2\rceil-\alpha),\qquad
U=\min(\alpha,2\alpha+s-\lceil p/2\rceil).
\tag{C18.3.5}
\] {取\ensuremath{e=L,L+1}。}

{证明L+1\ensuremath{\le}\ensuremath{U}，只须四个比较：} \[
4\alpha\le p-1,\quad
3\alpha\le\lfloor p/2\rfloor+s-1,\quad
\lceil s/2\rceil+1\le2\alpha,\quad
\lceil p/2\rceil+1\le3\alpha+\lfloor s/2\rfloor.
\tag{C18.3.6}
\] {由p/6\ensuremath{\le}\ensuremath{\alpha }\ensuremath{\le}(p+11)/6、\ensuremath{p}\ensuremath{\ge}83、10\ensuremath{\le}\ensuremath{s}\ensuremath{\le}(p−1)/2逐项成立：第一式\ensuremath{p}\ensuremath{\ge}25足够；第二式最多差6而s−1\ensuremath{\ge}9；第三式\ensuremath{p}\ensuremath{\ge}15足够；最后一式因3α\ensuremath{\ge}\ensuremath{\operatorname{ceil}(p/2)}、\ensuremath{\operatorname{floor}(s/2)}\ensuremath{\ge}5。4α\ensuremath{\le}\ensuremath{p}还保证来源及被删去的上界p−3α\ensuremath{\ge}\ensuremath{\alpha }；\ensuremath{s}\ensuremath{\le}p/2保证\ensuremath{D}库存，6α\ensuremath{\ge}\ensuremath{p}保证\ensuremath{C}来源容量。}

\paragraph*{\texorpdfstring{{3.3} {紧邻3p：使用已证计数公式，共用同一收尾}}{3.3 紧邻3p：使用已证计数公式，共用同一收尾}}\label{C18-h134}
\addcontentsline{toc}{paragraph}{{3.3} {紧邻3p：使用已证计数公式，共用同一收尾}}

{剩下\ensuremath{r=p+s}、1\ensuremath{\le}\ensuremath{s}\ensuremath{\le}9。继续使用已经书面证明的}第\ref{construct:three-parameters}节{，其\ensuremath{p}\ensuremath{\ge}83界不需要有限反射记录。}

{令\ensuremath{a}₀=\ensuremath{\operatorname{floor}(s/3)+1}，对\ensuremath{\alpha =a}₀,\ensuremath{a}₀+1及} \[
h=\lfloor p/(3\alpha)\rfloor-1,\quad\lfloor p/(3\alpha)\rfloor
\] {取} \[
i=\alpha(h+1),\quad j=\alpha h,\quad k=p-\alpha(2h+1),
\] \[
o_e=(j-e,k-e,(\alpha-s)h+\alpha+2e,sh).
\tag{C18.3.7}
\] {令} \[
e_0=\max\{0,\alpha(3h+2)-p,\lceil((s-\alpha)h-\alpha)/2\rceil\},
\] {使用\ensuremath{e=e}₀,\ensuremath{e}₀+1，\ensuremath{q=2}。旧容量定理保证连续三个\ensuremath{e}，这里只需前两个。不再要求第三个\ensuremath{e}承担二进周期／单位修补；第6节会统一完成。}

\subsection*{\texorpdfstring{{C18.4} {同一根理想与同一深同余群}}{C18.4 同一根理想与同一深同余群}}\label{C18-h135}
\addcontentsline{toc}{subsection}{{C18.4} {同一根理想与同一深同余群}}

{在上述每组来源中，\ensuremath{j}为\ensuremath{\alpha }的倍数，故} \[
f\equiv p^3\pmod\alpha,\quad
g-f=\alpha n,\quad
\gcd(f,n)=\gcd(g,n)=1,
\quad\gcd(f,g)=1.
\] {最后用\ensuremath{\operatorname{gcd}(p,\alpha )=1}。所有计数均使\ensuremath{0<\alpha <p}。}

{令\ensuremath{R}₀为全部f,g素因子的有限局部化环。相邻反射与正向逆给} \[
T=J_f^{-1}J_g=
\begin{pmatrix}1&-\beta n/f\\0&g/f\end{pmatrix}.
\] {其倍率及逆倍率因\ensuremath{\operatorname{gcd}(f,g)=1}恰生成\ensuremath{R}₀。两个等大块的自由角色交换\ensuremath{S=(-1,r};0,1)给} \[
\boxed{[S,T]=U(nH/g),\qquad
H=2\beta+r\alpha=n(2j+\alpha)-2p^2.}
\tag{C18.4.1}
\] {通过倍率共轭与根参数的整数加法，得到U(δnR₀)，其中\ensuremath{\delta }为全部\ensuremath{H}的正\ensuremath{gcd}。}

{四种容量选择的\ensuremath{gcd}均为} \[
\boxed{\delta=\gcd(n,2).}
\tag{C18.4.2}
\] {第3.1节\ensuremath{j=0}，相邻\ensuremath{\alpha }使\ensuremath{H}相差\ensuremath{n}，故\ensuremath{gcd}为\ensuremath{\operatorname{gcd}(n,2p^{2})=\operatorname{gcd}(n,2)}。第3.2节\ensuremath{j=\alpha }，相差3n，且\ensuremath{p}\ensuremath{\ne}3，所以\ensuremath{\operatorname{gcd}(3n,2p^{2})=\operatorname{gcd}(n,2)}。第3.4节\ensuremath{j=(t-1)\alpha }，差为(2t−1)n；而2t−1为正奇数且小于\ensuremath{p}（rt\ensuremath{\le}\ensuremath{p}、\ensuremath{r}\ensuremath{\ge}4t蕴含2t\ensuremath{\le}sqrt(p)），故同样的\ensuremath{gcd}为\ensuremath{\operatorname{gcd}(n,2)}。第3.3节固定\ensuremath{\alpha }的两个\ensuremath{h}使\ensuremath{H}相差2αn，相邻\ensuremath{\alpha }互素使总\ensuremath{gcd}整除2n；再用\ensuremath{H}\ensuremath{\equiv}−\ensuremath{2p^{2}}} {modn，去掉奇因子。\ensuremath{n}奇时奇\ensuremath{\alpha }给奇\ensuremath{H}；\ensuremath{n}偶时偶\ensuremath{\alpha }给\ensuremath{H}\ensuremath{\equiv}2} {mod4，故仍恰为(4.2)。}

{横向反射的\ensuremath{\alpha }含两个相邻数，\ensuremath{gcd}为1。沿用原已证换基level计算：设\ensuremath{b=\delta n}，\ensuremath{P=(e}₁,Je₁)=(1,β;0,−α)，则} \[
P^{-1}U(t)P=U(-\alpha t),\qquad
P^{-1}JU(t)J^{-1}P=L(\alpha t/\det J).
\] {第二式以参数\ensuremath{t=(det}} {J)bx清去分母，不假设\ensuremath{p}已在\ensuremath{R}₀中可逆。Morris／Serre深核定理在\ensuremath{P}基给\ensuremath{\Gamma (\alpha ^{2}b^{2},R}₀)，换回原基包含\ensuremath{\Gamma (\alpha ^{3}b^{2},R}₀)。互素\ensuremath{\alpha }的level合并得到} \[
\boxed{\Gamma(b^2,R_0).}
\tag{C18.4.3}
\] {\ensuremath{R}₀有无限单位：相邻f,g互素且差\ensuremath{\alpha n>2}，不可能全部绝对值为1。原深核定理的适用条件满足，未从有限模像倒推无限群包含。}

{先用旧安全返回夹乘引理取得\ensuremath{F}₁及其逆，反演\ensuremath{p}；再激活所有奇数1\ensuremath{\le}\ensuremath{q}\ensuremath{\le}\ensuremath{p}、\ensuremath{\operatorname{gcd}(q,n)=1}的保载体互补混合\ensuremath{F_{q}}。倍率共轭扩环，并在最终\ensuremath{R}上重新建立同一\ensuremath{\Gamma (b^{2},R)}。\ensuremath{n}奇时一组奇\ensuremath{\alpha }使f,g奇偶相反，2已经在\ensuremath{R}₀中反演。最终\ensuremath{R}单位包含全部安全小整数1\ensuremath{\le}\ensuremath{q}\ensuremath{\le}\ensuremath{p}。全部所用规范反射属于GL₂(R)，完整本原列可在同一环内追踪。}

{在(\ensuremath{z,y_*)}基中，原U(bR)恰变成L(bR)。再由相邻\ensuremath{\alpha }及Bezout，反射乘换块给实际可逆\ensuremath{V}满足} \[
V\equiv\begin{pmatrix}1&-p^{-2}\\0&1\end{pmatrix}\pmod n.
\tag{C18.4.4}
\] {不要求\ensuremath{V}的迹模4或模9满足额外条件。}

\subsection*{\texorpdfstring{{C18.5} {全部\ensuremath{2p<n<4p}的小单位引理}}{C18.5 全部的小单位引理}}\label{C18-h136}
\addcontentsline{toc}{subsection}{{C18.5} {全部\ensuremath{2p<n<4p}的小单位引理}}

\textbf{{引理。}} {\ensuremath{p}\ensuremath{\ge}5为奇素数、\ensuremath{2p<n<4p}、\ensuremath{\operatorname{gcd}(p,n)=1}。由−1和所有1\ensuremath{\le}\ensuremath{q}\ensuremath{\le}\ensuremath{p}、\ensuremath{\operatorname{gcd}(q,n)=1}的剩余类生成的子群\ensuremath{H}，等于(Z/nZ)×。}

\textbf{{\ensuremath{n}奇。}} {若\ensuremath{H}不满，取最小正缺失单位\ensuremath{a}。利用−1可设\ensuremath{p<a<n/2}，又\ensuremath{p>n/4}。若\ensuremath{a<n/3}，取4a−n；若\ensuremath{a>n/3}，取n−2a。它总是严格介于0与\ensuremath{a}之间，且为缺失单位，矛盾。\ensuremath{a=n/3}与\ensuremath{\operatorname{gcd}(a,n)=1}、\ensuremath{n>2p}\ensuremath{\ge}10不相容。2、4都属于\ensuremath{H}，所以折返合法。}

\textbf{{\ensuremath{n=2m}，\ensuremath{m}奇。}} {单位约化到模\ensuremath{m}为同构。选\ensuremath{\epsilon }\ensuremath{\in}\{−1,1\}使\ensuremath{q=(m+\epsilon )/2}为奇数。因\ensuremath{m<2p}，有\ensuremath{q}\ensuremath{\le}\ensuremath{p}且\ensuremath{\operatorname{gcd}(q,n)=1}，故\ensuremath{H}模\ensuremath{m}包含2的逆（差一个符号）。任一单位模\ensuremath{m}取正平衡代表\ensuremath{a}\ensuremath{\le}\ensuremath{m/2<p}，除尽2因子得到奇数\ensuremath{c}\ensuremath{\le}\ensuremath{p}；\ensuremath{c}也与\ensuremath{n}互素，已在\ensuremath{H}中。乘回2的幂即可覆盖原单位。}

\textbf{{4整除\ensuremath{n}。}} {在区间(0,n/2)内，\ensuremath{a}\ensuremath{\mapsto}n/2−a是单位的无固定点对合，故(0,n/4)恰有φ(n)/4个单位。它们都小于\ensuremath{p}，连同负数已有φ(n)/2个\ensuremath{H}元素。\ensuremath{p}本身与\ensuremath{n}互素且\ensuremath{n/4<p<n/2}，故±\ensuremath{p}又给两个不同元素。因此\textbar{}\ensuremath{H|>\phi (n)/2}，\ensuremath{H}只能是满群。}

{证毕。这个引理无九偏移限制，不用有限单位群枚举、原根假设或容斥误差。它与第4节实际小数字激活结合后，给} \[
\boxed{R^\times\twoheadrightarrow(\mathbb Z/n\mathbb Z)^\times.}
\tag{C18.5.1}
\]

\subsection*{\texorpdfstring{{C18.6} {移动终端族，共用一次列输送}}{C18.6 移动终端族，共用一次列输送}}\label{C18-h137}
\addcontentsline{toc}{subsection}{{C18.6} {移动终端族，共用一次列输送}}

{终端列在(\ensuremath{z,y_*)}基中仍取} \[
t_j=(1,p^2-nj).
\tag{C18.6.1}
\] {对应零和\ensuremath{p}元组为(j+1)份\ensuremath{A}、\ensuremath{j}份\ensuremath{B}及p−1−2j份\ensuremath{C}。其精确容量是} \[
\boxed{\max(0,\lceil(p-1-r)/2\rceil)\le j\le(p-1)/2.}
\tag{C18.6.2}
\] {旧九偏移取\ensuremath{j=0}或1。\ensuremath{r<p}时只需把选择区间移到(6.2)，不必改变列算法。该区间含至少两个连续整数（本范围\ensuremath{r}\ensuremath{\ge}12），所以可以选出任意\ensuremath{j}} {modδ，\ensuremath{\delta }\ensuremath{\in}\{1,2\}。所有\ensuremath{t_{j}}本原，且\ensuremath{p^{2}-nj}与\ensuremath{n}互素。}

{现在已有L(bR)、\ensuremath{\Gamma (b^{2},R)}、(4.4)、(5.1)，以及安全小数字的可逆乘法。按}第\ref{construct:small-digit}节{：}

\begin{enumerate}
\def\labelenumi{\arabic{enumi}.}
\tightlist
\item
  {用\ensuremath{V}的非负幂和一个\ensuremath{R}单位\ensuremath{\epsilon }，使当前列\ensuremath{w}\ensuremath{\equiv}\ensuremath{\epsilon (1,p^{2})}} {modn。}
\item
  {\ensuremath{n}偶时，用安全\ensuremath{q}\ensuremath{\equiv}3或5} {mod8与符号组成的分裂乘法，匹配首坐标的完整二进精度；再用\ensuremath{V}的一个2幂倍次数恢复奇数部分。该小\ensuremath{q}对\ensuremath{p}\ensuremath{\ge}11、\ensuremath{n<4p}始终存在。}
\item
  {由当前第二坐标确定\ensuremath{j}} {modδ，在(6.2)选取\ensuremath{j}，使\ensuremath{w}\ensuremath{\equiv}\ensuremath{\epsilon t_{j}}} {modb。与旧证明相比，只改变这次整数代表选择。}
\item
  {在奇数部分，两个根方向\ensuremath{E}₂₁与VE₂₁\ensuremath{V}⁻¹在\ensuremath{\epsilon t_{j}}上作用的行列式为\ensuremath{2\epsilon ^{2}/p^{2}}} {modn，是单位；解一次二维线性同余修正到奇数平方精度。根参数用CRT保持二进精度。}
\item
  {\ensuremath{n}偶时再用一次L(bx)补第二坐标的二进精度；最后得到\ensuremath{w_{out}}\ensuremath{\equiv}\ensuremath{\epsilon t_{j}}} {\ensuremath{modb^{2}}。}
\item
  {两列均\ensuremath{R}本原，\ensuremath{\Gamma (b^{2},R)}的精确纤维作用将其真正送到\ensuremath{\epsilon t_{j}}。执行(6.2)对应的零和\ensuremath{p}平均，再接独立已证的零触发尾部，全系统归零。}
\end{enumerate}

{这些步骤对整个(1.1)使用相同算术条件。零触发依赖的是先行端点含零尾部、3p、4p与宽尾部，不循环调用本文新合并的区间证明。单个反射的正向逆、扩环与完整列尺度已在前面准备好，未把模逆当成有理逆。}


\section{九偏移的连续参数容量证明}\label{construct:three-parameters}
{本节设} \(p\ge83\){、}\(1\le s\le9\){、}\(n=3p+s\){，用于共同上半带的最小九种偏移。}

\subsection*{\texorpdfstring{{C19.6} {连续三个反射已够，公式分界降至83}}{C19.6 连续三个反射已够，公式分界降至83}}\label{C19-h138}
\addcontentsline{toc}{subsection}{{C19.6} {连续三个反射已够，公式分界降至83}}

{原统一区间公式保留连续四个\ensuremath{e}，是为了同时选择二进奇偶、三进好类，以及二进隔一单位。现在三进要求删除，连续三个\ensuremath{e}即可：两个相邻值负责迹的奇偶，首尾相差2负责二进主单位。}

{沿用}第\ref{construct:capacity-interval}节{的\ensuremath{a=}\ensuremath{\lfloor}s/3\ensuremath{\rfloor}+1、\ensuremath{\alpha }\ensuremath{\in}\{a,a+1\}、\ensuremath{h=}\ensuremath{\lfloor}p/(3α)\ensuremath{\rfloor}−1或\ensuremath{\lfloor}p/(3α)\ensuremath{\rfloor}，以及\ensuremath{e}的精确下界\ensuremath{L}。只需\ensuremath{L}\ensuremath{\le}U−2。相同的逐项比较给出充分条件} \[
\begin{gathered}
h\ge4,\quad(3\alpha-s)h\ge s,\quad(3\alpha-s)h\ge\alpha+4,\\
sh\ge5\alpha+3-s,\quad sh\ge3\alpha+4-s,\\
\alpha h\ge4\alpha+4-s,\quad\alpha h\ge3\alpha+2.
\end{gathered}
\tag{C19.6.1}
\] {其证明与旧文一致，只把所需的额外间隔3改成2。例如候选\ensuremath{\operatorname{ceil}(((s-\alpha )h-\alpha )/2)}与最后一个上界比较，用2ceil(A/2)\ensuremath{\le}A+1，得到αh\ensuremath{\ge}t+5−2α−s；再用\ensuremath{t}\ensuremath{\le}6α−1即为倒数第二条。}

{ % do not increment counter
\begin{longtable}[]{@{}rlrr@{}}
\toprule\noalign{}
{\ensuremath{\alpha }} & {可能\ensuremath{s}} & {足够\ensuremath{h}} & {\ensuremath{p}\ensuremath{\ge}83时较小\ensuremath{h}} \\
\midrule\noalign{}
\endhead
\bottomrule\noalign{}
\endlastfoot
{1} & {1、2} & {7} & {26} \\
{2} & {1至5} & {12} & {12} \\
{3} & {3至8} & {8} & {8} \\
{4} & {6至9} & {4} & {5} \\
{5} & {9} & {4} & {4} \\
\end{longtable}
}

{每个范围代入(6.1)即可。因此\ensuremath{p}\ensuremath{\ge}83全部由统一三连反射公式承担，余下素数11\ensuremath{\le}\ensuremath{p<83}共18个，所需有限计数证书由180组缩为162组。原两项特殊补充仍保留。}


\section{反射输出的完整容量区间}\label{construct:capacity-interval}
\subsection*{\texorpdfstring{{C20.2} {保留自由参数，容量成为一个完整整数区间}}{C20.2 保留自由参数，容量成为一个完整整数区间}}\label{C20-h139}
\addcontentsline{toc}{subsection}{{C20.2} {保留自由参数，容量成为一个完整整数区间}}

{继续用原反射计数与矩阵。给定正整数} {α,h，令} \[
\begin{aligned}
j&=\alpha h,& i&=\alpha(h+1),& k&=p-2\alpha h-\alpha,\\
v&=sh,&x&=\alpha h-e,&y&=k-e,\\
u&=(\alpha-s)h+\alpha+2e.
\end{aligned}
\tag{C20.2.1}
\] {组成和为} {\ensuremath{p}，迹零条件} {\ensuremath{p(x-y)+\alpha v=\beta }} {恒成立。旧公式仅用} {\ensuremath{e=0,1,2,3}；现在允许这四个连续整数从任意合法起点开始。}

{写} {\ensuremath{p=3\alpha h+t}。输出非负及上界精确要求} \[
L\le e\le U,
\quad
L=\max\left(0,2\alpha-t,
\left\lceil\frac{(s-\alpha)h-\alpha}{2}\right\rceil\right),
\] \[
U=\min\left(\alpha h,\ \alpha h+t-\alpha,
\left\lfloor\frac{sh+\alpha+s-t}{2}\right\rfloor\right).
\tag{C20.2.2}
\] {另须满足来源容量} {2i\ensuremath{\le}\ensuremath{p}、2k\ensuremath{\le}p+s} {和} {\ensuremath{v}\ensuremath{\le}p−s；2j\ensuremath{\le}\ensuremath{p}} {由2i\ensuremath{\le}\ensuremath{p}得到。}

{因此``连续四个反射''恰可取} {\ensuremath{e=L,L+1,L+2,L+3}，只须证明} {\ensuremath{L}\ensuremath{\le}U−3。取整没有额外格空洞：\ensuremath{e}} {的可行集合就是闭区间内的全部整数。}


\section{小参数的五返回公式}\label{construct:five-returns}
\subsection*{\texorpdfstring{{C21.1} {三条返回承担三个不同的任务}}{C21.1 三条返回承担三个不同的任务}}\label{C21-h140}
\addcontentsline{toc}{subsection}{{C21.1} {三条返回承担三个不同的任务}}

{设\ensuremath{p}\ensuremath{\ge}2、1\ensuremath{\le}\ensuremath{s<p}、\ensuremath{r=p+s}、\ensuremath{n=3p+s}。核心为} \[
(A^p,B^p,C^r),\qquad A=a,\ B=rz-a,\ C=-pz.
\]

{对一个来源参数\ensuremath{\alpha }，以下三条返回分别承担：}

\begin{enumerate}
\def\labelenumi{\arabic{enumi}.}
\tightlist
\item
  {\ensuremath{J}：迹零，重复一次就提供真实正向射影逆；}
\item
  {\ensuremath{G}：同来源的相邻迹零输出，\ensuremath{J}⁻¹\ensuremath{G}提供可逆三角倍率；}
\item
  {\ensuremath{M}：由\ensuremath{J}的输出加守恒移动，直接满足\ensuremath{M=U(-n)J}。}
\end{enumerate}

{因为新根U(−n)已经有准确参数，第二个横向参数只需两条相邻反射。它不必再完成一遍守恒交换。这个任务划分正是五条返回而非六条或十二条的原因。}

\subsection*{\texorpdfstring{{C21.2} {固定\ensuremath{\alpha }的显式计数}}{C21.2 固定的显式计数}}\label{C21-h141}
\addcontentsline{toc}{subsection}{{C21.2} {固定\ensuremath{\alpha }的显式计数}}

{设\ensuremath{\alpha }\ensuremath{\ge}1、1\ensuremath{\le}\ensuremath{s}\ensuremath{\le}2α，写} \[
p=3\alpha d+\rho,\qquad 0\le\rho<3\alpha.
\] {来源取} \[
j=\alpha d,\qquad i=j+\alpha,\qquad
k=\alpha d+\rho-\alpha.
\tag{C21.2.1}
\] {输出取} \[
o_e=(x,y,u,v)=
\bigl(\alpha d-e,\ k-e,\ (\alpha-s)d+\alpha+2e,\ sd\bigr).
\tag{C21.2.2}
\] {这个参数化恒满足来源与输出总数\ensuremath{p}、迹零方程。因此每条容量合法的\ensuremath{o_{e}}都给一个实际迹零返回\ensuremath{J_{e}}。}

{守恒移动为} \[
\Delta_3=(3\alpha-\rho,-\rho,2\rho-3\alpha-s,s).
\] {所以改变后的输出为} \[
o_e+\Delta_3=
\bigl(\alpha d+3\alpha-\rho-e,\ \alpha d-\alpha-e,\
(\alpha-s)d-2\alpha+2\rho-s+2e,\ s(d+1)\bigr).
\tag{C21.2.3}
\]

{我们同时需要\ensuremath{o_{e}}、\ensuremath{o_{e+1}}及(2.3)可行。它们各自四个位置容量为} \[
p-2i,\quad p-2j,\quad r-2k,\quad p-s.
\tag{C21.2.4}
\]

\textbf{{精确区间。}} {来源容量与固定\ensuremath{v}容量成立时，这三条输出的所有约束等价于整数\ensuremath{e}满足} \[
L\le e\le U,
\] {其中} \[
\boxed{
L=\max\left(
0,\ 5\alpha-2\rho,\
\left\lceil\frac{(s-\alpha)d-\alpha}{2}\right\rceil,\
\left\lceil\frac{(s-\alpha)d+2\alpha+s-2\rho}{2}\right\rceil
\right),
}
\tag{C21.2.5}
\] \[
\boxed{
U=\min\left(
\alpha d-\alpha-\mathbf1_{\rho=0},\
\left\lfloor\frac{sd+\alpha+s-\rho-2}{2}\right\rfloor,\
\left\lfloor\frac{sd+4\alpha+2s-3\rho}{2}\right\rfloor
\right).
}
\tag{C21.2.6}
\]

{证明是逐项把(2.2)、(2.3)及e+1的非负和库存上界移项。这里\ensuremath{s}\ensuremath{\le}2α用于删去被更强条件蕴含的界；例如输出\ensuremath{o_{e}}与\ensuremath{o_{e+1}}的\ensuremath{u}上界合并为式(C21.2.6)第二项。}

\subsection*{\texorpdfstring{{C21.3} {不用参数扫描的一般容量充分条件}}{C21.3 不用参数扫描的一般容量充分条件}}\label{C21-h142}
\addcontentsline{toc}{subsection}{{C21.3} {不用参数扫描的一般容量充分条件}}

\textbf{{容量定理。}} {若} \[
\boxed{
d\ge7,\qquad sd\ge9\alpha+2-s,\qquad
(3\alpha-s)d\ge4\alpha+s+2,
}
\tag{C21.3.1}
\] {则来源及全部输出容量均成立，且\ensuremath{L}\ensuremath{\le}\ensuremath{U}。故可以直接选\ensuremath{e=L}。}

\textbf{{证明。}} {由\ensuremath{s}\ensuremath{\le}2α，} \[
p-2i=\alpha d+\rho-2\alpha\ge0,
\] \[
r-2k=\alpha d-\rho+2\alpha+s
\ge\alpha d-\alpha+s+1>0.
\] {固定\ensuremath{D}输出的最强上界是s(d+1)\ensuremath{\le}p−s，等价于} \[
(3\alpha-s)d+\rho\ge2s,
\] {由第三条充分条件及4α+s+2\ensuremath{\ge}2s得到。其余来源及固定输出容量更弱。}

{为证明\ensuremath{L}\ensuremath{\le}\ensuremath{U}，记} \[
A=(s-\alpha)d-\alpha,\quad
B=(s-\alpha)d+2\alpha+s-2\rho,
\] \[
C=sd+\alpha+s-\rho-2,\quad
D=sd+4\alpha+2s-3\rho,\quad
U_1=\alpha d-\alpha-\mathbf1_{\rho=0}.
\] {比较四个下界与三个上界。对整数下界\ensuremath{l}，\ensuremath{l}\ensuremath{\le}\ensuremath{\operatorname{floor}(C/2)}等价于2l\ensuremath{\le}\ensuremath{C}；对\ensuremath{\operatorname{ceil}(A/2)}，与整数\ensuremath{U}₁比较只需\ensuremath{A}\ensuremath{\le}2U₁，与\ensuremath{\operatorname{floor}(C/2)}比较用充分条件A+1\ensuremath{\le}\ensuremath{C}。}

{十二个比较中最紧的三个分别归结为} \[
5\alpha\le\alpha d-\alpha-1,
\] \[
10\alpha\le sd+\alpha+s-2,
\] \[
(s-\alpha)d+2\alpha+s\le2(\alpha d-\alpha-1).
\] {它们分别由(3.1)三条保证。其余比较的余量可以直接写出：}

\[
C-2(5\alpha-2\rho)=sd-9\alpha+s+3\rho-2\ge0,
\] \[
D-2(5\alpha-2\rho)=sd-6\alpha+2s+\rho>0,
\] \[
2U_1-A=(3\alpha-s)d-\alpha-2\mathbf1_{\rho=0}\ge0,
\] \[
C-A=\alpha d+2\alpha+s-\rho-2\ge1,
\] \[
D-A=\alpha d+5\alpha+2s-3\rho\ge1,
\] \[
2U_1-B=(3\alpha-s)d-4\alpha-s+2\rho-2\mathbf1_{\rho=0}\ge0,
\] \[
C-B=\alpha d-\alpha+\rho-2\ge1,\qquad
D-B=\alpha d+2\alpha+s-\rho\ge1.
\] {零下界的比较由同样三条立即成立。每个使用\ensuremath{\rho }的下界只取0\ensuremath{\le}\ensuremath{\rho }\ensuremath{\le}3α−1，因此量词覆盖全部余数，不是从有限\ensuremath{p}试验外推。证毕。}

{定义} \[
H(s,\alpha)=\max\left(
7,\left\lceil\frac{9\alpha+2-s}{s}\right\rceil,
\left\lceil\frac{4\alpha+s+2}{3\alpha-s}\right\rceil
\right).
\tag{C21.3.2}
\] {只要\ensuremath{p}\ensuremath{\ge}3αH(s,α)，上述三个实际返回就由一个闭式规则取得。}

\subsection*{\texorpdfstring{{C21.4} {第二个横向参数只需两条反射}}{C21.4 第二个横向参数只需两条反射}}\label{C21-h143}
\addcontentsline{toc}{subsection}{{C21.4} {第二个横向参数只需两条反射}}

{对第二个\ensuremath{\alpha }，同样取(2.1)---(2.2)，但只要求\ensuremath{o_{e}}和\ensuremath{o_{e+1}}，不要求守恒平移输出。此时可取} \[
L_{\rm pair}=\max\left(
0,2\alpha-\rho,\left\lceil\frac{(s-\alpha)d-\alpha}{2}\right\rceil
\right),
\] \[
U_{\rm pair}=\min\left(
\alpha d-1,\ \alpha d+\rho-\alpha-1,\
\left\lfloor\frac{sd+\alpha+s-\rho}{2}\right\rfloor-1
\right).
\tag{C21.4.1}
\] {以下两个条件足够：} \[
\boxed{d\ge4,\qquad sd\ge3\alpha+2-s.}
\tag{C21.4.2}
\] {来源容量及\ensuremath{v}\ensuremath{\le}p−s由\ensuremath{s}\ensuremath{\le}2α、\ensuremath{d}\ensuremath{\ge}4成立。对0与2α−ρ，逐项比较显然，唯一紧式是} \[
sd-3\alpha+s+\rho-2\ge0.
\] {对\ensuremath{A=(s-\alpha )d-\alpha }，两个整数上界比较由(3α−s)d−α−2\ensuremath{\ge}0保证；半整数比较由} \[
sd+\alpha+s-\rho-2-A
=\alpha d+2\alpha+s-\rho-2\ge1
\] {保证。故\ensuremath{L_{pair}}\ensuremath{\le}\ensuremath{U_{pair}}。}

{令} \[
H_{\rm pair}(s,\alpha)=
\max\left(4,\left\lceil\frac{3\alpha+2-s}{s}\right\rceil\right).
\]

{现在统一选} \[
\boxed{a=\lceil s/2\rceil,\qquad \alpha=a,\ a+1,}
\tag{C21.4.3}
\] {对\ensuremath{a}使用三返回，对a+1只使用两反射。总充分界为} \[
\boxed{p\ge T_s:=
\max\bigl(3aH(s,a),\ 3(a+1)H_{\rm pair}(s,a+1)\bigr).}
\tag{C21.4.4}
\] {这是对任意\ensuremath{s}\ensuremath{\ge}1的容量公式；\ensuremath{p}的素数性只在下一节的算术资源中使用。}

\subsection*{\texorpdfstring{{C21.5} {五条返回已经接出完整根与方向接口}}{C21.5 五条返回已经接出完整根与方向接口}}\label{C21-h144}
\addcontentsline{toc}{subsection}{{C21.5} {五条返回已经接出完整根与方向接口}}

{设\ensuremath{p}为奇素数、\ensuremath{p>s}并满足(4.4)，\ensuremath{n=3p+s}。两个\ensuremath{\alpha }均小于\ensuremath{p}，所以\ensuremath{\operatorname{gcd}(p,\alpha n)=1}。}

{对固定\ensuremath{\alpha }，令\ensuremath{J=J_{e}}、\ensuremath{G=J_{e+1}}，平方系数分别为pf、pg。原恒等式给} \[
g-f=\alpha n,\quad
f\equiv p^3\pmod\alpha,\quad
pf\equiv p^4\pmod n.
\] {所以} \[
\gcd(f,g)=\gcd(f,n)=\gcd(g,n)=1.
\] {\ensuremath{J}、\ensuremath{G}均非零且由重复四步提供真逆。\ensuremath{J}⁻¹\ensuremath{G}的三角倍率g/f及其逆，将\ensuremath{f}、\ensuremath{g}全部素因子局部化。}

{对第一组\ensuremath{\alpha =a}，守恒返回\ensuremath{M}满足\ensuremath{M=U(-n)J}，所以U(±n)真实可执行。用两组相邻反射的倍率共轭扩环，得到} \[
U(nR_0).
\] {注意只需第一组供应这个\ensuremath{n}根，第二组的倍率也能共轭同一根，不必再次构造守恒输出。}

{两组横向反射的左下系数为−\ensuremath{a}、−(a+1)。由旧横向根、Morris／Serre及换基level引理，先得} \[
\Gamma(a^3n^2,R_0),\quad
\Gamma((a+1)^3n^2,R_0),
\] {再由\ensuremath{\operatorname{gcd}(a,a+1)=1}合并为} \[
\boxed{\Gamma(n^2,R_0).}
\tag{C21.5.1}
\]

{旧安全行列式夹乘激活\ensuremath{p}与小数字，再在扩大的\ensuremath{R}上重建同样根和深核。两个反射乘换块，由−\ensuremath{a+(a+1)=1}直接给模\ensuremath{n}单位步长\ensuremath{V}；不使用任何独立direction或迹条件。所有参与返回在最终GL₂(R)中，故接}第\ref{construct:small-digit}节{即可完成核心到零和终端。}

{在1\ensuremath{\le}\ensuremath{s}\ensuremath{\le}9时，基础单位模\ensuremath{n}满射、全输入入口和零触发尾部均已有一般证明，因此得到旧完整定理的一条新统一构造证明。对更大的可变\ensuremath{s}，(4.4)和根供应仍成立，但全输入结论必须另核对相应入口、基础单位和尾部，不能直接宣称全部区间已被同一方法覆盖。}

\subsection*{\texorpdfstring{{C21.6} {九偏移的实际替代范围}}{C21.6 九偏移的实际替代范围}}\label{C21-h145}
\addcontentsline{toc}{subsection}{{C21.6} {九偏移的实际替代范围}}

{将\ensuremath{s=1,}\ldots,9代入(4.4)：}

{ % do not increment counter
\begin{longtable}[]{@{}rrrrr@{}}
\toprule\noalign{}
{\ensuremath{s}} & {\ensuremath{a}} & {H(s,a)} & {\ensuremath{H_{pair}(s,a+1)}} & {\ensuremath{T_{s}}} \\
\midrule\noalign{}
\endhead
\bottomrule\noalign{}
\endlastfoot
{1} & {1} & {10} & {7} & {42} \\
{2} & {1} & {8} & {4} & {24} \\
{3} & {2} & {7} & {4} & {42} \\
{4} & {2} & {7} & {4} & {42} \\
{5} & {3} & {7} & {4} & {63} \\
{6} & {3} & {7} & {4} & {63} \\
{7} & {4} & {7} & {4} & {84} \\
{8} & {4} & {7} & {4} & {84} \\
{9} & {5} & {7} & {4} & {105} \\
\end{longtable}
}

\begin{itemize}
\tightlist
\item
  {所有\ensuremath{p}\ensuremath{\ge}83继续允许使用旧已证根供应公式；}
\item
  {\ensuremath{p<83}时，只要\ensuremath{p}\ensuremath{\ge}\ensuremath{T_{s}}，直接用本构造全参数五返回公式；}
\item
  {其他小参数保留原根记录，(11,3)按上一轮守恒修补完成。}
\end{itemize}

{旧162组中恰有48组满足新增条件，故需保留的逐参数根记录为114组。按\ensuremath{s}计数，新增替代数为} \[
(8,12,8,8,6,6,0,0,0).
\] {48组的替代不是从它们的试验成功归纳：它们只是一般定理范围中的有限成员。完整映射和剩余根记录保存在}第\ref{sec:certificate}节{。该清单不是独立的全部证明，必须连同一般公式、(11,3)修补及已证算术收尾阅读。}


\section{守恒输出交换与辅助素数的消去}\label{construct:conservation}
\subsection*{\texorpdfstring{{C22.1} {一个不依赖素数性的守恒公式}}{C22.1 一个不依赖素数性的守恒公式}}\label{C22-h146}
\addcontentsline{toc}{subsection}{{C22.1} {一个不依赖素数性的守恒公式}}

{记} \[
\mathcal K(a,z)=(A^p,B^p,C^r),\qquad
A=a,\quad B=rz-a,\quad C=-pz,\qquad n=2p+r.
\] {因此} \[
pA+pB+rC=0.
\tag{C22.1.1}
\]

{固定来源组成(i,j,k)，\ensuremath{i+j+k=p}，平均两个不交且组成相同的\ensuremath{p}元组，产生2p份} \[
D=\frac{iA+jB+kC}{p}.
\] {保留\ensuremath{r}份\ensuremath{D}作新载体；其余位置为} \[
A^{p-2i},\ B^{p-2j},\ C^{r-2k},\ D^{2p-r}.
\] {选一个\ensuremath{p}元输出组，计数\ensuremath{o=(x,y,u,v)}，平均为\ensuremath{a}′，再平均其补组。新载体值\ensuremath{D=-pz}′，即} \[
z'=-D/p.
\]

{对任意整数\ensuremath{t}，定义计数移动} \[
\boxed{
\Delta_t=(ti-p,\ tj-p,\ tk-r,\ n-tp).
}
\tag{C22.1.2}
\] {四个分量之和为0，所以它保持输出组大小。}

\textbf{{守恒交换定理。}} {若\ensuremath{o}与\ensuremath{o+\Delta _{t}}均满足上述四类剩余位置的非负容量约束，则两条实际四步返回\ensuremath{W}、\ensuremath{W}′满足} \[
\boxed{W'=U(-n)W\quad\text{在}(a,z)\text{坐标中}.}
\tag{C22.1.3}
\] {若\ensuremath{W}已有真实正向射影逆，则U(−n)也正向可执行；两个等大块的换角色给U(n)。}

\textbf{{证明。}} {两个输出组的总和值之差为} \[
(ti-p)A+(tj-p)B+(tk-r)C+(n-tp)D.
\] {由来源关系\ensuremath{iA+jB+kC=pD}及(1.1)，上式恰为} \[
tpD-0+(n-tp)D=nD.
\] {故新第一块值增加\ensuremath{nD/p=-nz}′，新载体\ensuremath{z}′不变，即} \[
(a',z')\longmapsto(a'-nz',z').
\] {这就是(1.3)，而且等式在整个核心上成立。证毕。}

{这里没有假设\ensuremath{p}为素数，没有为新返回求迹零条件，也没有预先要求\ensuremath{W}′的逆。需要的只有两个实际可行的分组和旧\ensuremath{W}的真实逆。}

{九偏移\ensuremath{r=p+s}、\ensuremath{0<s<p}时，最自然的选择是\ensuremath{t=3}：} \[
\boxed{\Delta_3=(3i-p,\ 3j-p,\ 3k-p-s,\ s).}
\tag{C22.1.4}
\] {它的意义是：减去全部旧库存的总和零关系，加入三份来源均值关系，再补\ensuremath{s}份\ensuremath{D}。这个操作只是计数恒等式；是否有足够实际位置仍由容量不等式判定。}

\subsection*{\texorpdfstring{{C22.2} {行列式不变是同一守恒关系的矩阵形式}}{C22.2 行列式不变是同一守恒关系的矩阵形式}}\label{C22-h147}
\addcontentsline{toc}{subsection}{{C22.2} {行列式不变是同一守恒关系的矩阵形式}}

{不要求输出本身迹零。令} \[
\alpha=i-j,\qquad\beta=rj-pk.
\] {一般四步实际矩阵是\ensuremath{p}⁻²乘} \[
M(o)=
\begin{pmatrix}
p(x-y)+\alpha v & pry-p^2u+\beta v\\
-\alpha&-\beta
\end{pmatrix}.
\tag{C22.2.1}
\] {固定来源后，两条返回的底行相同。沿(1.2)，第一行变化恰为} \[
(\alpha n,\beta n)=-n(-\alpha,-\beta),
\] {所以} \[
M(o+\Delta_t)=U(-n)M(o),\qquad
\det M(o+\Delta_t)=\det M(o).
\tag{C22.2.2}
\]

{一般来说，若两个可逆2×2矩阵底行相同且行列式相同，则第一行之差与底行平行，两者之商必为上剪切。式(C22.1.2)更进一步：从实际守恒关系直接指定了剪切参数为−\ensuremath{n}。}

{对任意输出移动(Δx,Δy,Δu,Δv)、四项之和0，有} \[
\Delta(\det M)
=p[-\beta\Delta x+(\beta+\alpha r)\Delta y-\alpha p\Delta u].
\tag{C22.2.3}
\] {这给其他行列式不变移动的完整线性条件；本构造不假定该条件的所有整数解都满足位置容量。}

\subsection*{\texorpdfstring{{C22.3} {(11,3)的明确计数：原位置直接得到U(−36)}}{C22.3 (11,3)的明确计数：原位置直接得到U(−36)}}\label{C22-h148}
\addcontentsline{toc}{subsection}{{C22.3} {(11,3)的明确计数：原位置直接得到U(−36)}}

{取\ensuremath{p=11}、\ensuremath{r=14}、\ensuremath{n=36}。沿用旧根对的来源} \[
(i,j,k)=(4,2,5).
\] {两个来源组用去8份\ensuremath{A}、4份\ensuremath{B}、10份\ensuremath{C}，保留14份\ensuremath{D}后，剩余库存为} \[
(A,B,C,D)=(3,7,4,8).
\]

{原迹零返回使用} \[
o=(2,5,1,3).
\] {守恒移动为} \[
\Delta_3=(1,-5,1,3),
\] {所以} \[
\boxed{o'=o+\Delta_3=(3,0,2,6).}
\] {两组计数均非负、和为11，且不超过(3,7,4,8)。第四个输出组取各自补集，因此所有步骤都在原36个位置完成。}

{它们的整数矩阵分别为} \[
J=\begin{pmatrix}-27&568\\-2&27\end{pmatrix},
\qquad
M=\begin{pmatrix}45&-404\\-2&27\end{pmatrix}.
\] {有} \[
\det J=\det M=407=11\cdot37,\qquad
J^2=-407I,\qquad
\boxed{MJ^{-1}=U(-36).}
\tag{C22.3.1}
\]

{\ensuremath{W=J/11^{2}}重复一次提供射影逆。按时间先执行旧\ensuremath{J}返回、再执行新\ensuremath{M}返回，共8次原子平均，得到} \[
\frac1{11^4}MJ
=-\frac{37}{11^3}U(-36).
\] {整个系统只差共同非零尺度；换块共轭后同样8步得到U(36)。因此两个符号和全部整数倍U(36Z)都真实正向可执行。无需在此之前调用任何同余群来制造这个逆。}

\textbf{{安全性。}} {在(a,z)坐标的危险见证\ensuremath{y=a+11z}下，U(±36)满足\ensuremath{y}′=\ensuremath{y}} {mod36。共同尺度−\ensuremath{37/11^{3}}在2、3处为单位。更强地，正负根词合成是危险素数处单位的标量闭路；若中间进入全体相同的危险余数，后续平均无法恢复，故原位置中间步骤同样保持合法性。}

\subsection*{\texorpdfstring{{C22.4} {原倍率把直接\ensuremath{n}根扩为整族}}{C22.4 原倍率把直接根扩为整族}}\label{C22-h149}
\addcontentsline{toc}{subsection}{{C22.4} {原倍率把直接\ensuremath{n}根扩为整族}}

{旧相邻反射的另一个计数是} \[
G_{\rm row}=(4,2,5;1,4,3,3),\qquad
G=\begin{pmatrix}-27&172\\-2&27\end{pmatrix}.
\] {其平方为385I，已有真实逆。令} \[
T=J^{-1}G=
\begin{pmatrix}1&-972/37\\0&-35/37\end{pmatrix}.
\] {其三角倍率及逆实际可用；由\ensuremath{\operatorname{gcd}(35,37)=1}，} \[
\mathbb Z[-35/37,-37/35]=
R_0:=\mathbb Z[1/(5\cdot7\cdot37)].
\]

{\ensuremath{T}及其逆共轭U(36)，使参数乘−35/37或−37/35；再相加、取负和整数Bezout，得到} \[
\boxed{U(36R_0)\text{全部真实可执行}.}
\tag{C22.4.1}
\] {根理想已直接含\ensuremath{n}，不再是26n；13从未作为待消除的根因子进入这条论证。}

\subsection*{\texorpdfstring{{C22.5} {两个横向系数按一般引理取得深核}}{C22.5 两个横向系数按一般引理取得深核}}\label{C22-h150}
\addcontentsline{toc}{subsection}{{C22.5} {两个横向系数按一般引理取得深核}}

{另一个旧已计数的反射为} \[
K=\begin{pmatrix}-52&905\\-3&52\end{pmatrix},
\qquad K^2=-11I,
\] {位置计数(4,1,6;2,7,1,1)。它有真实正向逆，不必先假设11已经在\ensuremath{R}₀中可逆。}

{\ensuremath{J}、\ensuremath{K}的左下项分别为−2、−3。用}第\ref{sec:arithmetic}节{，结合(4.1)，分别得到} \[
\Gamma(2^3n^2,R_0),\qquad
\Gamma(3^3n^2,R_0)
\] {的真实正向实现。该引理只要求\ensuremath{J}、\ensuremath{K}及其逆实际存在；其行列式可通过根参数乘\ensuremath{det}} {\ensuremath{J}或\ensuremath{det}} {\ensuremath{K}清除，不要求先在\ensuremath{R}₀反演它们。}

{由同一通用引理的\ensuremath{gcd}} {level合并，} \[
\gcd(8n^2,27n^2)=n^2,
\] {故} \[
\boxed{\Gamma(n^2,R_0)\text{真实可执行}.}
\tag{C22.5.1}
\] {具体\ensuremath{n=36}时，两个level为10368、34992，合并为1296。这三个level都不含13。}

{Morris／Serre仍用于从完整双向根取得深主同余群；取消的是原来特别针对13的有限群与局部分离计算，而非全部外部群论依赖。}

\subsection*{\texorpdfstring{{C22.6} {小数字扩环和完整终端仍按统一路线}}{C22.6 小数字扩环和完整终端仍按统一路线}}\label{C22-h151}
\addcontentsline{toc}{subsection}{{C22.6} {小数字扩环和完整终端仍按统一路线}}

{有(5.1)以后，\ensuremath{F}₁和所有安全小数字\ensuremath{F_{q}}的逆由原安全整数行列式夹乘取得。将\ensuremath{p=11}及安全小数字加入\ensuremath{R}，按已有根共轭和横向引理重建同一个\ensuremath{n^{2}}深核。这里\ensuremath{p=11}及新单位都与36互素，没有循环依赖。}

{在(z,y)坐标，U(nR)就是L(nR)。为直接复用}第\ref{construct:small-digit}节{，可使用较弱资源L(2nR)、\ensuremath{\Gamma (4n^{2},R)}，取\ensuremath{\delta =2}；两者已包含于更强的(4.1)、(5.1)中。}

{方向返回由\ensuremath{J}和\ensuremath{K}的\ensuremath{\alpha =2}、3通过−\ensuremath{2+3=1}组合生成。在最终\ensuremath{R=Z}{[}1/(5·7·11·37){]}中，两者及其逆都整可逆，所以全程保持\ensuremath{R}本原列。再用已有互补混合的乘法自由度和两个终端} \[
(1,121),\qquad(1,85),
\] {即接统一的全输入核心到终端论证。}

{因此(11,3)已完成到末端的完整替代：} \[
\text{守恒输出交换}
\to U(nR_0)
\to \text{两横向系数的gcd深核}
\to \text{基础小数字}
\to \text{统一分裂乘法终端}.
\] {不调用SL₂(\ensuremath{F}₁₃)生成、13进高层提升、完美局部群／可解商分离，也不调用原\ensuremath{\alpha =1}、平方系数含107的额外反射。}


\section{九偏移剩余资源的精确接口}\label{construct:nine-offsets}
\subsection*{\texorpdfstring{{C23.3} {四原子大载体反射的一般计数}}{C23.3 四原子大载体反射的一般计数}}\label{C23-h152}
\addcontentsline{toc}{subsection}{{C23.3} {四原子大载体反射的一般计数}}

{从} \(\mathcal L\) {先平均两个不交、组成相同的组} \[
(A^i,B^j,C^k),\qquad i+j+k=p.
\] {记} \[
\alpha=i-j>0,\quad \beta=rj-pk=nj-p(p-\alpha).
\] {得到} {2p} {份} {\ensuremath{D=(\alpha }} {a+beta} {z)/p。保留} {\ensuremath{r=p+s}} {份} {\ensuremath{D}；剩余旧} {A,B,C} {和新} {\ensuremath{D}} {数目为} \[
p-2i,\quad p-2j,\quad r-2k,\quad p-s.
\] {平均一组} {(\ensuremath{A^{x},B^{y},C^{u},D^{v})}，再平均其} {\ensuremath{p}} {元补组。注意这里只有}\textbf{{一个输出组及其补组}}{，不是上半带的两条重复输出行。}

{完整容量是} \[
\begin{gathered}
2i\le p,\quad2j\le p,\quad2k\le r,\qquad x+y+u+v=p,\\
0\le x\le p-2i,\quad0\le y\le p-2j,\quad
0\le u\le r-2k,\quad0\le v\le p-s.
\end{gathered}
\tag{C23.3.1}
\] {当} \[
p(x-y)+\alpha v=\beta
\tag{C23.3.2}
\] {时，总返回恰为} {\ensuremath{p^{-2}}} {乘} \[
J=\begin{pmatrix}\beta&Q\\-\alpha&-\beta\end{pmatrix},
\quad Q=pry-p^2u+v\beta.
\tag{C23.3.3}
\] {定义} {\ensuremath{pf=\beta ^{2}-\alpha }} {\ensuremath{Q}，则} \[
J^2=pfI,\qquad pf=p^4\pmod n.
\tag{C23.3.4}
\] {这里} {\ensuremath{f}} {是整数：由} {\ensuremath{\beta =\alpha }} {\ensuremath{v}} {\ensuremath{mod}} {\ensuremath{p}，有} {\ensuremath{Q=v}} {\ensuremath{\beta }} {\ensuremath{mod}} {\ensuremath{p}，故} {\ensuremath{p}} {整除} {\ensuremath{\beta ^{2}-\alpha }} {\ensuremath{Q}。又} {\ensuremath{\operatorname{gcd}(p,n)=1}，故} {\ensuremath{f}} {非零且与} {\ensuremath{n}} {互素。}

{每个反射因此有真正的正向射影逆：重复同样四次平均。换到始终幺模的} {(\ensuremath{z,y)=(z,a+pz)}} {基，实际模} {\ensuremath{n}} {作用为} \[
p^{-2}J\equiv
\begin{pmatrix}1&-\alpha/p^2\\0&-1\end{pmatrix}.
\tag{C23.3.5}
\] {这也直接证明它保持合法性。}

\subsection*{\texorpdfstring{{C23.5} {换块交换子、双向根和主同余控制}}{C23.5 换块交换子、双向根和主同余控制}}\label{C23-h153}
\addcontentsline{toc}{subsection}{{C23.5} {换块交换子、双向根和主同余控制}}

{先给适用于一般计数证书的接口。若同一个} {alpha,j,v} {有两条反射，后一条取} {(x-1,y-1,u+2,v)，记平方系数为} {f,g，则} \[
g=f+\alpha n,\qquad
T=J_f^{-1}J_g=
\begin{pmatrix}1&-\beta n/f\\0&g/f\end{pmatrix}.
\tag{C23.5.1}
\] {与角色交换} {\ensuremath{S}} {的交换子为} \[
\boxed{[S,T]=U\!\left(\frac{nH}{g}\right),\quad
H=2\beta+r\alpha=n(2j+\alpha)-2p^2.}
\tag{C23.5.2}
\] {这里} {\ensuremath{U(t)=(1,t};0,1)。所有逆都已由实际反射和角色交换提供。}

{令} {\ensuremath{R0}} {是由各倍率} {g/f} {及} {f/g} {生成的整数环。对采用} {\ensuremath{j=\alpha }} {\ensuremath{h}} {的计数公式，\ensuremath{\operatorname{gcd}(f,g)=1}：因为} {\ensuremath{\operatorname{gcd}(p,\alpha )=1}、\ensuremath{j=\alpha }} {\ensuremath{h}，(3.4)} {给} {\ensuremath{f=p^{3}}} {\ensuremath{mod}} {\ensuremath{\alpha }，再用} {\ensuremath{\operatorname{gcd}(f,n)=1}。因此} {\ensuremath{R0}} {恰为相应} {f,g} {的素因子局部化。小参数证书独立核对所有} {f,g} {的素因子确实被这些}\textbf{{约分后}}{的倍率反演；没有假定约掉的公因子免费可逆。}

{各} {\ensuremath{T}} {共轭} {\ensuremath{U}} {将根参数乘} {f/g。由共轭、求逆、根参数相加和整数} {Bezout，得到} \[
U(\delta nR_0),\qquad\delta=\gcd\{H\}.
\tag{C23.5.3}
\] {大} {\ensuremath{p}} {公式中，相邻} {\ensuremath{h}} {的} {\ensuremath{H}} {相差2alpha} {\ensuremath{n}，而两个} {\ensuremath{\alpha }} {互素。故} {\ensuremath{\delta |2n}；又} {\ensuremath{H=-2p^{2}}} {\ensuremath{mod}} {\ensuremath{n}，\ensuremath{\delta }} {没有奇素因子。\ensuremath{n}} {奇时有奇} {\ensuremath{\alpha }} {给奇} {\ensuremath{H}，所以} {\ensuremath{\delta =1}；\ensuremath{n}} {偶时一条偶} {\ensuremath{\alpha }} {排除额外的4，故} {\ensuremath{\delta =2}。}

{解释横向根时无需把换基自身当成可执行操作。对任一反射} {\ensuremath{J}，设} \[
P=(e_1,Je_1)=\begin{pmatrix}1&\beta\\0&-\alpha\end{pmatrix}.
\] {精确有} \[
P^{-1}U(t)P=U(-\alpha t),\qquad
P^{-1}J U(t)J^{-1}P=L(\alpha t/\det J).
\tag{C23.5.4}
\] {置} {\ensuremath{h}}\emph{{=\ensuremath{\delta }} {\ensuremath{n}，第二式取} {\ensuremath{t=(det}} {J)h}}{\ensuremath{x}，清去分母；于是} {\ensuremath{P}} {基下已有} {U(alpha} {h*} {R0)、L(alpha} {h*} {R0)。深同余引理给} {\ensuremath{\Gamma (\alpha ^{2}}} {\ensuremath{h}}\emph{{²,R0)。换回原基，因} {\ensuremath{\alpha }} {\ensuremath{P^{-1}}} {和} {\ensuremath{P}} {都整于} {\ensuremath{R0}，有} \[
\Gamma(\alpha^3h_*^2,R_0)\text{ 正向可执行}.
\tag{C23.5.5}
\] {不同} {\ensuremath{\alpha }} {的} {\ensuremath{gcd}} {为1时，这些群生成包含} {\ensuremath{\Gamma (h}}}{²,R0)} {的群。理由是先模各} {level} {的公倍数：各素数幂下的主同余条件是嵌套的，CRT} {和} {\ensuremath{\mathrm{SL}_{2}}} {的有限环满射使它们的乘积给} {\ensuremath{gcd}} {level，再接已知深层。这一步不要求任何} {\ensuremath{\alpha }} {是} {\ensuremath{R0}} {单位。}

{故大} {\ensuremath{p}} {一律已有} {\ensuremath{\Gamma (n^{2},R0)}} {或} {\ensuremath{\Gamma (4n^{2},R0)}。小参数证书除一个补充外给} {\ensuremath{\operatorname{gcd}(\alpha )=1}、\ensuremath{\delta }} {属于\{1,2,4\}，且} {\ensuremath{\delta }} {的素数都整除} {\ensuremath{n}，故同样得到仅支持于} {\ensuremath{n}} {的控制层。例外} \((11,3)\) {由守恒输出交换构造处理。}

{下文统一使用更深但方便的} \[
\mathfrak M=\begin{cases}n^4,&n\text{ 奇},\\64n^4,&n\text{ 偶}.
\end{cases}
\tag{C23.5.6}
\] {上述控制层均包含} {\ensuremath{\Gamma (mathfrak}} {M,R0)。}


\section{两种载体的三行返回计数}\label{construct:three-row-transport}
{本节取} \(n=3p+r\){、}\(1\le r<p\){，在} \((y,z)\) {坐标中写核心为} \(((y-pz)^{2p},(-2y+(2p+r)z)^p,(-pz)^r)\){。以下矩阵作用于} \((y,z)^T\){；改到其他构造的} \((a,z)\) {坐标时使用} \(a=y-pz\){。}

\subsection*{\texorpdfstring{{C24.2} {一次不交划分的完整通用返回公式}}{C24.2 一次不交划分的完整通用返回公式}}\label{C24-h154}
\addcontentsline{toc}{subsection}{{C24.2} {一次不交划分的完整通用返回公式}}

{记核心三个实际值为} \(A,B,C\){。把原位置分成两个组成相同的} \(p\) {组、第三个} \(p\) {组，以及保留的} \(r\) {个载体位置。两个相同组各含} \(i\) {份} \(A\){、}\(j\) {份} \(B\){、}\(s\) {份} \(C\){，其中} \(i+j+s=p\){。}

{若保留旧} \(C\) {作载体，}\(s=0\){，有} \[
D_j=\frac1p\begin{pmatrix}p-3j&nj\\0&p\end{pmatrix},
\qquad 0\le j\le\lfloor p/2\rfloor.
\tag{C24.3}
\] {若保留旧} \(A\){，有} \[
T^A_{j,s}=\frac1p\begin{pmatrix}-(3j+s)&nj\\-1&p\end{pmatrix},
\tag{C24.4}
\] {且} \(2i\le2p-r,\ 2j\le p,\ 2s\le r\){。若保留旧} \(B\){，则} \[
T^B_{j,s}=\frac1p\begin{pmatrix}3p-3j-s&n(j-p)\\2&-(2p+r)\end{pmatrix},
\tag{C24.5}
\] {且} \(2i\le2p,\ 2j\le p-r,\ 2s\le r\){。}

{三式均直接来自真实副本数，并不只包含迹零参数。一般参数的合法性必须在输出取本原代表后检查；共同因子的消去可能改变结论，不能一律要求未归一化输出在每个危险素数处都不全相等。}

{对于要求在整个二维核心上恒成立的这类不交返回，两个新} \(p\) {组具有相同值就必须具有相同组成。因为行差若在加权零和平面上恒为零，必为重数向量} \((2p,p,r)\) {的倍数；两行长度相同又迫使该倍数为零。}

{这不排除某个特殊输入发生偶然相等，也不包括重叠平均或多轮操作。}


\section{加权上半带的四返回根公式}\label{construct:upper-four}
\subsection*{\texorpdfstring{{C25.2} {上半带的整数坐标与真实对合}}{C25.2 上半带的整数坐标与真实对合}}\label{C25-h155}
\addcontentsline{toc}{subsection}{{C25.2} {上半带的整数坐标与真实对合}}

{设} {\ensuremath{p\ge 5}} {为素数、\ensuremath{1\le r<p}、\ensuremath{n=3p+r}。由}第\ref{construct:short-choice}节{，只需处理} \[
\mathcal K(a,z)=\bigl(a^{2p},(-2a+rz)^p,(-pz)^r\bigr),
\qquad \gcd(a,z)=\gcd(a+pz,n)=1.
\tag{C25.1}
\] {这是旧坐标} {\ensuremath{y=a+pz}} {的整幺模变换。记三种实际值为} {\ensuremath{A}、\ensuremath{B}、\ensuremath{C}。}

{使用}第\ref{construct:three-row-transport}节{的原位置构造：平均两个组成相同的} {\ensuremath{p}} {组，再平均其余} {\ensuremath{p}} {组，保留} {\ensuremath{r}} {份旧} {\ensuremath{A}} {或旧} {\ensuremath{B}} {作新载体。每个返回恰有三个原子平均。}

{保留} {\ensuremath{A}} {时取} {\ensuremath{s=p-3j}，两个相同组各取} {(\ensuremath{A^{2j},B^{j},C^{s})}。在} {(a,z)} {上的整数矩阵及实际尺度为} \[
\frac1p A_j,\qquad
A_j=\begin{pmatrix}0&-d_j\\-1&0\end{pmatrix},
\quad d_j=p^2-nj.
\tag{C25.2}
\] {允许} {\ensuremath{j}} {的连续整数区间是} \[
I_A=\left[
\left\lceil\frac{2p-r}{6}\right\rceil,
\min\left(\left\lfloor\frac{2p-r}{4}\right\rfloor,
\left\lfloor\frac p3\right\rfloor\right)\right]\cap\mathbb Z.
\tag{C25.3}
\]

{保留} {\ensuremath{B}} {时取} {\ensuremath{s=p-r-3k}，两个相同组各取} {(\ensuremath{A^{2k+r},B^{k},C^{s})}。对应} \[
\frac1p B_k,\qquad
B_k=\begin{pmatrix}r&b_k\\2&-r\end{pmatrix},
\quad b_k=nk-p^2+rp,
\tag{C25.4}
\] \[
I_B=\left[
\max\left(0,\left\lceil\frac{2p-3r}{6}\right\rceil\right),
\left\lfloor\frac{p-r}{3}\right\rfloor\right]\cap\mathbb Z.
\tag{C25.5}
\] {直接有} \[
A_j^2=d_j I,\qquad B_k^2=\mu_k I,
\quad \mu_k=r^2+2b_k.
\] {\ensuremath{d_{j}}、\ensuremath{\mu _{k}}} {模} {\ensuremath{n}} {都等于} {\ensuremath{p^{2}}，故非零且与} {\ensuremath{n}} {互素。这些返回和射影逆都有真实正向实现，重复一次就是六原子的全局非零标量词。返回在所有} {\ensuremath{\ell |n}} {处保持合法性。}

\subsection*{\texorpdfstring{{C25.3} {两个对角方向之间的显式交换子}}{C25.3 两个对角方向之间的显式交换子}}\label{C25-h156}
\addcontentsline{toc}{subsection}{{C25.3} {两个对角方向之间的显式交换子}}

{假设} {\ensuremath{I_{A}}、\ensuremath{I_{B}}} {各含两个相邻整数} {j,j+1} {和} {k,k+1。记} \[
d=d_j,\quad e=d-n,\qquad
\mu=\mu_k,\quad\nu=\mu+2n,
\quad R=\mathbb Z[1/|de\mu\nu|].
\] {四个对合已经提供两个正向可逆的上三角变换：} \[
D=A_j^{-1}A_{j+1}
=\begin{pmatrix}1&0\\0&e/d\end{pmatrix},
\quad
T=B_k^{-1}B_{k+1}
=\begin{pmatrix}1&rn/\mu\\0&\nu/\mu\end{pmatrix}.
\tag{C25.6}
\] {它们的交换子为} \[
\boxed{DTD^{-1}T^{-1}=U\left(\frac{rn^2}{e\nu}\right).}
\tag{C25.7}
\] {这是一条对所有满足容量条件的} {p,r} {恒成立的公式，不是逐参数的周期搜索。}

{下面核对共同局部化环，避免将任意分母视为免费操作。有} {\ensuremath{\operatorname{gcd}(d,e)=1}；又} \[
\gcd(\mu,\nu)=
\begin{cases}1,&r\text{ 奇},\\2,&r\text{ 偶}.\end{cases}
\] {这是因为} {\ensuremath{\operatorname{gcd}(\mu ,n)=1}} {且} {\ensuremath{\mu }} {的奇偶性与} {\ensuremath{r}} {相同。\ensuremath{r}} {偶时} {\ensuremath{n}} {奇，\ensuremath{d}、\ensuremath{e}} {奇偶相反，所以} {\ensuremath{Z}{[}e/d,d/e{]}} {已含} {1/2。\ensuremath{r}} {奇时} {\ensuremath{\mu }、\ensuremath{\nu }} {已互素。两种情形都得到} \[
\mathbb Z[e/d,d/e,\nu/\mu,\mu/\nu]=R.
\tag{C25.8}
\]

{对任意} {\ensuremath{x}，\ensuremath{D}} {共轭} {U(x)} {将其参数乘} {d/e；\ensuremath{T}} {共轭} {U(x)} {将参数乘} {mu/nu，\ensuremath{T}} {的上平移部分在此消去。结合逆共轭、根参数的加法和取负，(7)、(8)} {给出完整} {\ensuremath{U(rn^{2}R)}。由} \[
A_jU(x)A_j^{-1}=L(x/d)
\] {再得到完整} {\ensuremath{L(rn^{2}R)}。所有操作和逆均已有原位置实现。}

{因此令} {\ensuremath{h=rn^{2}}，调用}第\ref{construct:even-roots}节{，即有} \[
\boxed{\Gamma(h^2,R)=\Gamma(r^2n^4,R)
\text{ 的作用正向可执行}.}
\tag{C25.9}
\] {\ensuremath{R}} {有无限多个单位：相邻} {\ensuremath{d}、\ensuremath{e}} {互素且相差} {\ensuremath{n>2}，不可能同时绝对值为1。(9)} {使用} {Morris} {2007} {定理6.1(2)、Serre} {对有理数域} {\ensuremath{S}} {整数} {\ensuremath{\mathrm{SL}_{2}}} {的强同余子群性质以及已证任意有限环消元。外部来源与原文访问边界沿用所引深同余文档，不由有限样本证明这两项外部定理。}

{与下半带相比，这里没有交换} {2p} {块与} {\ensuremath{p}} {块的免费角色对称。两组载体返回分别给出} {(6)，通过它们的交换子获得根，所以自然多出一个} {\ensuremath{n}} {因子。不能将下半带的换块矩阵直接搬来。}

\subsection*{\texorpdfstring{{C25.4} {一个所有} {\ensuremath{p}} {都适用的内部区间}}{C25.4 一个所有  都适用的内部区间}}\label{C25-h157}
\addcontentsline{toc}{subsection}{{C25.4} {一个所有} {\ensuremath{p}} {都适用的内部区间}}

\textbf{{容量定理。}} {对所有素数} {\ensuremath{p\ge 13}，若} \[
10\le r\le p-3,
\quad\text{即}\quad
\boxed{3p+10\le n\le4p-3,}
\tag{C25.10}
\] {则} {\ensuremath{I_{A}}、\ensuremath{I_{B}}} {均至少含两个整数，故第} {3} {节全部结论成立。}

{证明} {\ensuremath{I_{B}}：允许} {\ensuremath{s}} {的区间是} {{[}0,min(floor(r/2),p-r){]}，且} {\ensuremath{s=p-r}} {模3。\ensuremath{r\ge 10}} {使} {\ensuremath{\operatorname{floor}(r/2)\ge 5}。若} {\ensuremath{p-r\ge 5}，则} {0,\ldots,5} {至少给每个模3类两个代表；若} {\ensuremath{p-r=3}，则选} {\ensuremath{s=0,3}；若} {\ensuremath{p-r=4}，则选} {\ensuremath{s=1,4}。}

{证明} {\ensuremath{I_{A}}，先设} {\ensuremath{p\ge 23}。若} {\ensuremath{r\le 2p/3}，则} {(3)} {的上界为} {\ensuremath{\operatorname{floor}(p/3)}，选} {0,\ldots,5} {中与} {\ensuremath{p}} {同余模3的两个} {\ensuremath{s}} {即可，其他容量由} {\ensuremath{r\le 2p/3}} {保证。若} {\ensuremath{r>2p/3}，则} {\ensuremath{j}} {的实区间为} {{[}(2p-r)/6,(2p-r)/4{]}，长度} \[
\frac{2p-r}{12}\ge\frac{p+3}{12}>2,
\] {所以至少有两个整数。\ensuremath{p=13}} {只需} {\ensuremath{r=10}，此时} {\ensuremath{I_{A}={3,4}}；\ensuremath{p=17}} {的} {\ensuremath{r=10,}\ldots,14} {均有} {\ensuremath{I_{A}={4,5}}；\ensuremath{p=19}} {的} {\ensuremath{r=10,}\ldots,13} {有} {\{5,6\}，\ensuremath{r=14}} {有} {\{4,5,6\}，\ensuremath{r=15,16}} {有} {\{4,5\}。证毕。}

{(3)、(5)} {是更精确的适用条件，个别} {\ensuremath{r<10}} {也可覆盖，例如} {\ensuremath{p=13,r=9}。式} {(C25.10)} {取一个无需逐} {\ensuremath{p}} {分类的统一范围。}


\section{安全扩环、局部方向与精确尺度}\label{construct:upper-local}
{本节设} \(p\ge13\) {为奇素数、}\(n=3p+r\){、}\(10\le r\le p-3\){，并使用加权上半带四返回的容量及根。}

\subsection*{\texorpdfstring{{C26.2} {共同核心与已有资源}}{C26.2 共同核心与已有资源}}\label{C26-h158}
\addcontentsline{toc}{subsection}{{C26.2} {共同核心与已有资源}}

{置} {\ensuremath{r=n-3p}，故} {\ensuremath{10\le r\le p-3}。已有原位置入口将任意合法输入送到} \[
K(a,z)=(a^{2p},(-2a+rz)^p,(-pz)^r),\qquad
\gcd(a,z)=\gcd(a+pz,n)=1.
\] {令} {\ensuremath{y=a+pz}，使用整幺模坐标} {(z,y)。合法性恰为} {\ensuremath{y}} {在模} {\ensuremath{n}} {下可逆。}

{四返回文档给出相邻参数} {\ensuremath{j_{0},j_{0}+1}} {和} {\ensuremath{k_{0},k_{0}+1}，以及} \[
A=\begin{pmatrix}0&-d\\-1&0\end{pmatrix},\quad
B=\begin{pmatrix}r&b\\2&-r\end{pmatrix},
\] \[
d=p^2-nj_0,\quad e=d-n,\quad
b=nk_0-p^2+rp,\quad \mu=r^2+2b,\quad\nu=\mu+2n.
\] {相邻} {\ensuremath{A}、\ensuremath{B}} {分别记为} {\ensuremath{A}',\ensuremath{B}'。每个矩阵实际乘以} {1/p，恰由三个原位置平均实现；平方分别为} {dI,eI,mu} {I,nu} {\ensuremath{I}，故射影逆也正向可执行。}

{令} {\ensuremath{R_*}} {=} {\ensuremath{Z}{[}\ensuremath{1/|de}} {\ensuremath{\mu }} {\ensuremath{\nu |}{]}。已有} \[
U(rn^2R_*),\quad L(rn^2R_*),\quad
\Gamma(r^2n^4,R_*).
\tag{C26.2.1}
\] {反演素数均不除} {\ensuremath{n}；\ensuremath{n}} {奇时} {2} {已在} {\ensuremath{R_*}} {中可逆。还可执行} \[
D_0=A^{-1}A'=\operatorname{diag}(1,e/d),\quad
T=B^{-1}B'=
\begin{pmatrix}1&rn/\mu\\0&\nu/\mu\end{pmatrix}.
\] {两者的共轭倍率及逆倍率生成整个} {\ensuremath{R_*}。在本文每次扩张} {\ensuremath{R}} {后，同一组倍率继续使根参数构成相应的} {\ensuremath{R}} {模。}

{以下``可执行矩阵''均指：一个有限原位置平均词在整个核心上实现该矩阵的某个非零共同标量倍。共同标量不改变归零目标；不得用于单独缩放外部位置。}

\subsection*{\texorpdfstring{{C26.3} {两个取逆引理}}{C26.3 两个取逆引理}}\label{C26-h159}
\addcontentsline{toc}{subsection}{{C26.3} {两个取逆引理}}

\textbf{{有限逆提升。}} {若} {\ensuremath{\Gamma (Q,R)}} {已可执行，\ensuremath{Z}} {属于} {\ensuremath{\mathrm{SL}_{2}(R)}} {且} {\ensuremath{Z}} {正向可执行，则} {\ensuremath{Z}} {的射影逆也正向可执行。}

{因为} {R/QR} {有限，存在} {\ensuremath{k\ge 1}} {使} {\ensuremath{Z^{k}}} {属于} {\ensuremath{\Gamma (Q,R)}，而} \[
Z^{-1}=Z^{k-1}(Z^k)^{-1}.
\tag{C26.3.1}
\] {这里最后一因子是已知主同余群中的真正逆，故不只是模} {\ensuremath{Q}} {的逆。}

\textbf{{安全整数返回取逆。}} {若} {\ensuremath{\Gamma (Q,R)}} {已可执行，整数非奇异返回} {\ensuremath{M}} {真实存在，且} {\ensuremath{\operatorname{gcd}(det}} {\ensuremath{M,Q)=1}，则} {\ensuremath{M}} {的射影逆也可执行。}

{证明使用}引理\ref{lem:smith}{的} {Smith--CRT} {夹乘，那里} {\ensuremath{n}} {可逐字替换为} {\ensuremath{Q}：取正偶次} {\ensuremath{P=M^{k}=I}} {\ensuremath{mod}} {\ensuremath{Q}，令} {\ensuremath{\Delta =det}} {\ensuremath{P}。以整数行列式一的行列变换对角化} {\ensuremath{P}，构造} {\ensuremath{g_*}} {使} {Pg\_}\emph{{\ensuremath{P}} {被} {\ensuremath{\Delta }} {整除，再用} {CRT} {取} {g=g\_}} {\ensuremath{mod}} {\ensuremath{\Delta }、\ensuremath{g=I}} {\ensuremath{mod}} {\ensuremath{Q}。于是} \[
h=\Delta^{-1}PgP\in\Gamma(Q,\mathbb Z),\qquad
gPh^{-1}=\Delta P^{-1}.
\] {\ensuremath{g}、\ensuremath{h}} {的逆由主同余控制提供。证毕。}

{对本文常用的伴随形矩阵，还有无需展开巨大} {\ensuremath{M^{k}}} {的恒等式。若} \[
M=\begin{pmatrix}\alpha&D\\-1&0\end{pmatrix},
\quad \gcd(D,Q)=1,
\] {取整数} {\ensuremath{t}} {满足} {\ensuremath{t=\alpha }} {\ensuremath{mod}} {\ensuremath{D}、\ensuremath{t=0}} {\ensuremath{mod}} {\ensuremath{Q}，则} \[
H=D^{-1}MU(t)M\in SL_2(\mathbb Z),\qquad
U(t)MH^{-1}=DM^{-1}.
\tag{C26.3.2}
\] {\ensuremath{H}} {本身正向可执行，故} {(3.1)} {提供} {\ensuremath{H^{-1}}。当} {\textbar{}\ensuremath{D|=1}} {时模} {\ensuremath{D}} {条件为空。}

{所有应用均先指出真实返回，再激活其逆；两个引理不将任意形式矩阵变为平均操作。}

\subsection*{\texorpdfstring{{C26.4} {消除来自} {\ensuremath{r}} {的额外素因子}}{C26.4 消除来自  的额外素因子}}\label{C26-h160}
\addcontentsline{toc}{subsection}{{C26.4} {消除来自} {\ensuremath{r}} {的额外素因子}}

{目标是在某个有限局部化环} {\ensuremath{R}} {中获得} \[
U(3^a n^2R),\quad L(3^a n^2R),\quad
\Gamma(M,R),\qquad
a=v_3(r),\quad M=3^{2a}n^4.
\tag{C26.4.1}
\] {\ensuremath{R}} {的反演素数全部不除} {\ensuremath{n}；因此} {\ensuremath{M}} {的素数支持恰为} {\ensuremath{n}} {的素数支持。}

\paragraph*{\texorpdfstring{{4.1} {先激活} {\ensuremath{p}} {和小素因子}}{4.1 先激活  和小素因子}}\label{C26-h161}
\addcontentsline{toc}{paragraph}{{4.1} {先激活} {\ensuremath{p}} {和小素因子}}

{保留} {\ensuremath{r}} {份原载体} {\ensuremath{C}，将三个} {\ensuremath{p}} {组分别平均，给出} \[
F_t=\begin{pmatrix}t/p&r(p-t)/(3p)\\0&1\end{pmatrix},
\quad t=p-3j,\quad0\le j\le(p-1)/2.
\tag{C26.4.2}
\] {两个重复组各取} {p-j} {份} {\ensuremath{A}、\ensuremath{j}} {份} {\ensuremath{B}。它们的互补组恰有} {\ensuremath{p}} {份，故容量成立。}

{取} {\ensuremath{\epsilon =1}} {或} {-1，使} {\ensuremath{p=\epsilon }} {mod3。\ensuremath{F_\epsilon }} {的整数代表行列式为} {\ensuremath{\epsilon }} {\ensuremath{p}，与} {\ensuremath{r^{2}n^{4}}} {互素；由第3节取得逆。其上三角倍率为} {epsilon/p，通过根的正逆共轭，先将} {\ensuremath{p}} {加入} {\ensuremath{R}} {的反演集合。}

{对} {\ensuremath{q=1}} {mod3，只要} {\ensuremath{\epsilon }} {\ensuremath{q}} {是} {(4.2)} {的允许数字，便有} \[
C_q=F_{\epsilon q}F_\epsilon^{-1}
=\begin{pmatrix}q&r(1-q)/3\\0&1\end{pmatrix}.
\tag{C26.4.3}
\] {设素数} {\ensuremath{\ell |r}、\ensuremath{\ell \ne 3}、\ensuremath{\ell \le (p-1)/2}，选} {\ensuremath{q=\ell }} {或} {-\ensuremath{\ell }} {使} {\ensuremath{q=1}} {mod3。若} {\ensuremath{\ell \le (p-3)/2}，数字容量立即成立。若} {\ensuremath{\ell =(p-1)/2}，则} {\ensuremath{p=2\ell +1>7}} {为素数迫使} {\ensuremath{\ell =p=2}} {mod3，取} {\ensuremath{q=-\ell }、\ensuremath{\epsilon =-1}，所需数字为正} {\ensuremath{\ell }，容量仍成立。}

{(4.3)} {的右上项被} {\ensuremath{q}} {整除，故} \[
H=q^{-1}C_qAC_q
\] {为整数矩阵，\ensuremath{det}} {\ensuremath{H=det}} {\ensuremath{A=-d}。于是} {\ensuremath{HA^{-1}}} {属于} {\ensuremath{\mathrm{SL}_{2}(R)}，且正向可执行。由} {(3.1)} {得} {\ensuremath{H^{-1}}，从} \[
AC_qH^{-1}=qC_q^{-1}
\tag{C26.4.4}
\] {得到} {\ensuremath{C_{q}}} {的真正射影逆。共轭上根将} {\ensuremath{\ell }} {局部化，再以} {\ensuremath{A}} {交换上下根。}

{这里允许} {\ensuremath{\ell }} {整除旧主同余层：式} {(C26.4.4)} {正是绕过``行列式与旧层互素''限制的必要步骤。}

\paragraph*{\texorpdfstring{{4.2} {唯一可能的大素因子}}{4.2 唯一可能的大素因子}}\label{C26-h162}
\addcontentsline{toc}{paragraph}{{4.2} {唯一可能的大素因子}}

{若} {\ensuremath{r}} {有一个大于} {p/2} {的素因子，则} {\ensuremath{r}} {本身就是该素数。此时} {\ensuremath{r}} {奇、\ensuremath{n}} {偶，且} {3} {不整除} {\ensuremath{n}。}

{若} {\ensuremath{p=r}} {mod3，保留} {\ensuremath{B}} {的迹零返回} {\ensuremath{k_*=(p-r)/3}、\ensuremath{s=0}} {可执行，其平方标量} \[
\mu_*=\frac{r(2p+r)}3
\] {被} {\ensuremath{r}} {整除且与} {\ensuremath{n}} {互素。它与原} {\ensuremath{B}} {的商提供上三角倍率，直接将} {\ensuremath{r}} {局部化。}

{余下假设} {\ensuremath{p\ne r}} {mod3。选择一个奇数} {\ensuremath{q}} {满足} \[
r<q<p,\qquad2q\le3r,\qquad\gcd(q,n)=1.
\tag{C26.4.5}
\] {以下闭式选择证明其存在，不需逐素数搜索。}

\begin{itemize}
\tightlist
\item
  {若} {\ensuremath{5r\ge 3p}，取严格小于} {n/4} {的最大奇数。此时} {n-4q} {属于} {\{2,4,6,8\}。因} {\ensuremath{q}} {奇且} {3} {不整除} {\ensuremath{n}，\ensuremath{q}} {与} {\ensuremath{n}} {互素。\ensuremath{p-r\ge 4}} {保证} {\ensuremath{q>r}，其余容量由} {\ensuremath{3p\le 5r}} {得到。}
\item
  {若} {\ensuremath{p/2<r<3p/5}，取离} {n/5} {最近的奇数。距离至多1；n-5q} {为} {±1、±3} {或} {±5。只有} {\ensuremath{5|n}} {时须在} {n/5±1} {中选不被5整除的一项。n/5} {是偶数，至少一项可用。由} {\ensuremath{r>p/2}、\ensuremath{r<3p/5}} {和} {\ensuremath{r\ge (p+1)/2}，可直接得} {\ensuremath{r<q<p}} {及} {\ensuremath{q\le 3r/2}。}
\end{itemize}

{置} {\ensuremath{s=q-r}，为正偶数，且} {\ensuremath{2s\le r}。保留} {\ensuremath{B}、两个重复行各取} {p-s} {份} {\ensuremath{A}} {和} {\ensuremath{s}} {份} {\ensuremath{C}，给} \[
N_s=\begin{pmatrix}p-s&-ps\\2&-r\end{pmatrix}.
\] {其行列式} {\ensuremath{\Delta =s(2p+r)-pr}} {模} {\ensuremath{n}} {为} {-pq，模} {\ensuremath{r}} {为} {2ps，故与} {rn} {互素。第3节已可给出} {\ensuremath{N_{s}^{-1}}。}

{令} {\ensuremath{H=N_{s}B^{-1}=(\rho ,\beta };0,1)。其右上项满足} \[
\beta=(p-s)/2\pmod r,
\] {非零：\ensuremath{s=q-r<p-r}，故} {\ensuremath{r<p-s<p<2r}。于是} \[
[D_0,H]=U(x),\qquad x=\beta n/e,
\tag{C26.4.6}
\] {且} {\ensuremath{x}} {是模} {\ensuremath{r}} {单位。选整数} {\ensuremath{t}} {使} {\ensuremath{tx=1}} {\ensuremath{mod}} {\ensuremath{r}，则} {\ensuremath{n^{2}(1-tx)}} {属于} {\ensuremath{rn^{2}R}。根参数相加给出} {\ensuremath{U(n^{2})}。再将任意} {\ensuremath{R}} {元素取整数模} {\ensuremath{r}} {代表，可得} {\ensuremath{U(n^{2}R)}，\ensuremath{A}} {共轭给下根。由深同余引理得到} {\ensuremath{\Gamma (n^{4},R)}。}

{最后} {\ensuremath{N_{0}=(p,0};2,-r)} {行列式} {-pr} {已与新层互素，故} {\ensuremath{N_{0}^{-1}}} {可执行。\ensuremath{N_{0}B^{-1}}} {的倍率将} {\ensuremath{r}} {局部化。若中途为扩大根参数环加入} {\ensuremath{det}} {\ensuremath{N_{s}}，仍只反演不除} {rn} {的素数。}

{上述操作消掉} {\ensuremath{r}} {的全部非3素因子。若} {\ensuremath{3|r}，则它的所有非3素因子均小于} {p/2，已属于第4.1节。因此} {(4.1)} {在全部参数下成立。}

\subsection*{\texorpdfstring{{C26.5} {可执行数字与模} {\ensuremath{n}} {的全部单位}}{C26.5 可执行数字与模  的全部单位}}\label{C26-h163}
\addcontentsline{toc}{subsection}{{C26.5} {可执行数字与模} {\ensuremath{n}} {的全部单位}}

{保留} {\ensuremath{A}} {的一般返回为} \[
M_A(j,s)=
\begin{pmatrix}p-3j-s&rj-ps\\-1&0\end{pmatrix}.
\tag{C26.5.1}
\] {两个重复} {\ensuremath{p}} {组各取} {p-j-s} {份} {\ensuremath{A}、\ensuremath{j}} {份} {\ensuremath{B}、\ensuremath{s}} {份} {\ensuremath{C}，留下} {\ensuremath{r}} {份旧} {\ensuremath{A}，再平均其余} {\ensuremath{p}} {份。其容量恰为} \[
0\le j\le\lfloor p/2\rfloor,\quad
0\le s\le\lfloor r/2\rfloor,\quad
j+s\ge\lceil r/2\rceil.
\tag{C26.5.2}
\] {置} {\ensuremath{q=3j+s}、\ensuremath{\alpha =p-q}、\ensuremath{D=rj-ps=nj-pq}。若} {\ensuremath{\operatorname{gcd}(q,n)=1}，则} {\ensuremath{D}} {与} {\ensuremath{M}} {互素，(3.2)} {激活正向逆。且} \[
M_AA^{-1}=\begin{pmatrix}-D/d&-\alpha\\0&1\end{pmatrix}.
\] {通过共轭将} {\ensuremath{D}} {局部化。由于} {\ensuremath{D=-pq}} {\ensuremath{mod}} {\ensuremath{n}、\ensuremath{p}} {已可逆，\ensuremath{R}} {单位的模} {\ensuremath{n}} {像便包含} {\ensuremath{q}。}

{可以一次将全部满足} {(5.2)、\ensuremath{\operatorname{gcd}(q,n)=1}} {的有限多返回的} {\ensuremath{D}} {加入} {\ensuremath{R}。每次扩大} {\ensuremath{R}} {后，根和} {(4.1)} {继续成立。所有反演素数仍不除} {\ensuremath{n}。}

\textbf{{容量区间。}} {令} {\ensuremath{h=\operatorname{floor}(r/2)}、\ensuremath{J=\operatorname{floor}(p/2)}，则每个整数} \[
L\le q\le U,\qquad
U=3J+h,\quad
L=\begin{cases}h+2,&r\text{ 偶},\\h+5,&r\text{ 奇},\end{cases}
\tag{C26.5.3}
\] {均可实现。可取} \[
\max(0,\lceil(q-h)/3\rceil)\le j\le
\min(J,\lfloor q/3\rfloor,\lfloor(q-\lceil r/2\rceil)/2\rfloor),
\quad s=q-3j.
\tag{C26.5.4}
\] {证明也可看每个固定} {\ensuremath{j}} {的数字区间：} \[
[3j+\max(0,\lceil r/2\rceil-j),\,3j+h].
\] {偶} {\ensuremath{r}} {从} {\ensuremath{j=1}} {开始，奇} {\ensuremath{r}} {从} {\ensuremath{j=2}} {开始；\ensuremath{r\ge 10}} {保证这些整数区间相邻或交叠，起点和终点正是} {(5.3)。}

\textbf{{单位覆盖命题。}} {最终有限局部化环} {\ensuremath{R}} {满足} \[
R^\times\longrightarrow(\mathbb Z/n\mathbb Z)^\times
\quad\text{满射}.
\tag{C26.5.5}
\]

{若} {3} {不整除} {\ensuremath{n}，第4节已使} {\ensuremath{r}} {成为实际} {\ensuremath{R}} {单位，因此} {-r/p} {是模} {\ensuremath{n}} {的单位3的一个} {\ensuremath{R}} {单位提升。对任意模} {\ensuremath{n}} {单位} {\ensuremath{u}，先取正平衡绝对代表} {\ensuremath{b\le n/2}。若} {\ensuremath{b>n/4}，以} {\textbar{}\ensuremath{n-3b|}} {代替} {\ensuremath{b}。这是非零正整数严格下降，且只改变符号和3的乘法幂。重复至} {\ensuremath{b\le n/4}。若} {\ensuremath{b<L}，逐次乘3至首次达到} {\ensuremath{L}，所得} {\ensuremath{q\le 3L-3\le U}。原} {\ensuremath{b}} {若已\textgreater=\ensuremath{L}，则} {\ensuremath{b\le n/4\le U}。最后} {\ensuremath{q}} {是} {(5.3)} {的安全数字，所以} {\ensuremath{q}} {属于} {\ensuremath{R}} {单位像；倒推得到} {\ensuremath{u}} {也属于该像。\ensuremath{r}} {奇时} {\ensuremath{r\le p-4}、\ensuremath{p\ge 17}；这正保证上述} {\ensuremath{3L-3\le U}，偶} {\ensuremath{r}} {情形更直接。}

{若} {\ensuremath{3|n}，置} {\ensuremath{k=n/3}。\ensuremath{k}±1} {都在} {(5.3)} {中。模} {\ensuremath{n}} {到模} {\ensuremath{k}} {的单位群核由以下安全数字提供：}

\begin{itemize}
\tightlist
\item
  {\ensuremath{3|k}} {时，1+k} {生成阶3的核；}
\item
  {\ensuremath{k=1}} {mod3} {时，1+k} {是阶2核的非平凡元；}
\item
  {\ensuremath{k=2}} {mod3} {时，1-k} {是阶2核的非平凡元。}
\end{itemize}

{这里取负是} {\ensuremath{R}} {单位} {-1，不要求负数数字本身满足} {(5.3)。}

{对任意单位} {\ensuremath{u}，取它模} {\ensuremath{k}} {的正平衡绝对代表} {\ensuremath{1\le b\le k/2}。考虑} {k-b、k+b、\ensuremath{b}；三者模} {\ensuremath{k}} {都等于} {±\ensuremath{u}。若} {k-b} {不被3整除，取} {\ensuremath{q=k-b}；否则} {\ensuremath{k}} {不被3整除，而} {k+b} {和} {\ensuremath{b}} {均不被3整除，取落在} {{[}L,U{]}} {的一个。}

{容量说明如下。总有} {\ensuremath{k/2\ge L}；k-b} {从而位于} {{[}L,U{]}。若} {\ensuremath{k+b>U}，则} \[
b>U-k\ge k/2-2.
\] {除} {(\ensuremath{p,r)=(19,15)}} {外，所有} {\ensuremath{3|r}} {的允许参数均有} {\ensuremath{k/2-2\ge L}：偶} {\ensuremath{r}} {由} {\ensuremath{3p-2r\ge 24}，奇} {\ensuremath{r}} {由} {\ensuremath{3p-2r\ge 39}；\ensuremath{p\ge 31}} {立即成立，较小} {\ensuremath{p}} {的允许} {\ensuremath{3|r}} {项直接满足这两不等式，唯一例外正是} {(19,15)。在这个例外中} {\ensuremath{k=24}} {被3整除，所以第一选择} {k-b} {始终可用，不会进入溢出分支。于是总有一个安全} {\ensuremath{q}} {可用。q/u} {在模} {\ensuremath{k}} {下等于} {±1，前述核元素和} {-1} {将其提升；因} {\ensuremath{q}} {已由} {(5.1)} {激活，得到} {\ensuremath{u}} {的} {\ensuremath{R}} {单位提升。证毕。}

\subsection*{\texorpdfstring{{C26.6} {一条局部射影循环引理}}{C26.6 一条局部射影循环引理}}\label{C26-h164}
\addcontentsline{toc}{subsection}{{C26.6} {一条局部射影循环引理}}

\textbf{{引理。}} {设} {\ensuremath{\ell }} {为素数，\ensuremath{T}} {属于} {\ensuremath{\mathrm{GL}_{2}(Z_\ell )}，其模} {\ensuremath{\ell }} {作用为} \[
T\sim\begin{pmatrix}1&b\\0&1\end{pmatrix},\qquad b\ne0.
\] {作用在仿射坐标} {\ensuremath{x=z/y}、\ensuremath{y}} {为单位的图上。}

\begin{enumerate}
\def\labelenumi{\arabic{enumi}.}
\tightlist
\item
  {\ensuremath{\ell \ge 5}} {时，\ensuremath{T}} {在每个} {\ensuremath{Z/\ell ^{e}Z}} {上都是单个长度} {\ensuremath{\ell ^{e}}} {的循环。}
\item
  {\ensuremath{\ell =3}} {时，令} {\ensuremath{T}} {的迹半数归一化后为} {I+N、\ensuremath{N^{2}=cI}；若} {\ensuremath{c=0}} {mod3} {且} {\ensuremath{c\ne 6}} {mod9，同样成立。}
\item
  {\ensuremath{\ell =2}} {时，若} {tr} {\ensuremath{T=2}} {mod4，也同样成立。}
\end{enumerate}

{证明：奇} {\ensuremath{\ell }} {可将} {\ensuremath{T}} {除以迹半数，得到} {I+N。模} {\ensuremath{\ell }} {的} {\ensuremath{N}} {非零幂零，故} {\ensuremath{c=0}} {\ensuremath{mod}} {\ensuremath{\ell }。写} \[
(I+N)^\ell=A_\ell I+B_\ell N,\qquad
B_\ell=\sum_{j\ge0}\binom{\ell}{2j+1}c^j.
\] {\ensuremath{\ell \ge 5}} {时} {\ensuremath{v_\ell (B_\ell )=1}、\ensuremath{A_\ell }} {为单位；\ensuremath{\ell =3}} {时} {\ensuremath{B_{3}=3+c}，条件恰保证同样的赋值。此后重复} {\ensuremath{\ell }} {次幂，非标量系数的赋值每次恰增加1。因为} {\ensuremath{N}} {在模} {\ensuremath{\ell }} {下为非零上根，\ensuremath{T^{\ell ^{h}}}} {在} {\ensuremath{x}} {上的首个非零位移总具有精确赋值} {\ensuremath{h}。因此一个循环每提升一级都分裂为长度增加} {\ensuremath{\ell }} {倍的单循环。}

{\ensuremath{\ell =2}} {时} {Cayley--Hamilton} {给} {\ensuremath{T^{2}=(tr}} {T)T-(det} {T)I。非标量系数恰有赋值1，而} {\ensuremath{T}} {模2的上平移系数为单位。并且} \[
\operatorname{tr}(T^2)=(\operatorname{tr}T)^2-2\det T=2\pmod4.
\] {继续平方得到逐层精确位移，证明同上。标量归一化不改变射影作用。}

{引理不仅说某个方向能动，而且提供逐位计算到指定方向的指数；无需枚举} {\ensuremath{\ell ^{e}}} {个方向。}

\subsection*{\texorpdfstring{{C26.7} {同时消除全部素数处的方向}}{C26.7 同时消除全部素数处的方向}}\label{C26-h165}
\addcontentsline{toc}{subsection}{{C26.7} {同时消除全部素数处的方向}}

{在} {(z,y)} {坐标中取} {\ensuremath{V=BA/p^{2}}，直接有} \[
V=\begin{pmatrix}1&-3/p\\0&1\end{pmatrix}\pmod n.
\tag{C26.7.1}
\] {因此每个} {\ensuremath{\ell \ge 5}、\ensuremath{\ell |n}} {处均可用第6节。}

{若} {\ensuremath{n}} {偶，整数} {BA} {的迹为} \[
\operatorname{tr}(BA)=n(2j_0-k_0)-p(p+r).
\tag{C26.7.2}
\] {\ensuremath{n=0}} {mod4} {时此迹自动为2} {mod4；\ensuremath{n=2}} {mod4} {时把} {\ensuremath{k_{0}}} {换成} {\ensuremath{k_{0}+1}} {恰改变迹} {\ensuremath{n}，所以相邻两个} {\ensuremath{B}} {中总有一个使迹为2} {mod4。除以} {\ensuremath{p^{2}}} {不改变模4的迹，故这个} {\ensuremath{V}} {也在全部二进层形成完整循环。}

{若} {\ensuremath{3|n}，(7.1)} {在模3下是单位矩阵，须先用另一条真实} {\ensuremath{A}} {返回处理3方向。置} {\ensuremath{k=n/3}，从} {\ensuremath{\delta }} {属于} {\{-1,1,3\}} {中选} \[
q=k+\delta=-p\pmod3.
\] {具体地，-p-k} {模3分别为1、2、0时取} {\ensuremath{\delta =1}、-1、3。总有} {\ensuremath{\operatorname{gcd}(q,n)=1}。可行} {\ensuremath{j}} {区间} {(5.4)} {含两个相邻整数。为说明这不是经验容量：\ensuremath{q}} {位于} {{[}p+r/3-1,p+r/3+3{]}，将} {(5.4)} {的三个实上界减去} {(q-h)/3，每个差均至少2；\ensuremath{p\ge 19}} {的条件直接给出，唯一更小的允许} {\ensuremath{3|r}} {参数} {(17,12)} {也成立。}

{若} {\ensuremath{v_{3}(n)=1}，选其中一个} {\ensuremath{j}} {使} {\ensuremath{kj\ne 1}} {mod3；若} {\ensuremath{v_{3}(n)\ge 2}，任取。实际} {(z,y)} {矩阵为} \[
W=\frac1p\begin{pmatrix}p&-1\\nj&-q\end{pmatrix}.
\] {它模3是非平凡上平移，且迹半数归一化后的参数为} \[
c=\frac{(p+q)^2-4nj}{(p-q)^2}.
\tag{C26.7.3}
\] {若} {\ensuremath{9|n}，则} {\ensuremath{c=0}} {mod9。若} {\ensuremath{v_{3}(n)=1}，则} {\ensuremath{c/3=-kj}} {mod3，所选} {\ensuremath{j}} {保证} {\ensuremath{c\ne 6}} {mod9。因此第6节提供每个三进层的完整方向循环。}

{现在给定任意合法} {\ensuremath{R}} {本原列} {(z,y)。先用} {\ensuremath{W}} {的某个非负次幂使} {\ensuremath{z/y=0}} {\ensuremath{mod}} {\ensuremath{3^{v_{3}(M)}}。\ensuremath{V}} {在模3下为} {\ensuremath{I}，其模任意3幂的}\textbf{{完整矩阵阶}}{是3的幂；令后续} {\ensuremath{V}} {的指数被该阶整除，就保留刚调好的3数据。其他} {\ensuremath{\ell |n}} {的方向循环长度分别是} {\ensuremath{\ell ^{v_\ell (M)}}，彼此及上述3幂互素。指数} {CRT} {同时完成所有方向。因此得到} \[
z=0\pmod M,\qquad y\in(R/MR)^\times.
\tag{C26.7.4}
\] {若} {3} {不整除} {\ensuremath{n}，省去} {\ensuremath{W}} {即可。}

{这里明确使用} {\ensuremath{V}} {的完整3进矩阵阶来保护已完成数据；不能拿局部射影逆在不同素数处随意拼接。}

\subsection*{\texorpdfstring{{C26.8} {单位尺度从} {\ensuremath{n}} {提升到} {\ensuremath{M}}}{C26.8 单位尺度从  提升到 }}\label{C26-h166}
\addcontentsline{toc}{subsection}{{C26.8} {单位尺度从} {\ensuremath{n}} {提升到} {\ensuremath{M}}}

{第5节已给模} {\ensuremath{n}} {的全部单位像。还需证明这个像覆盖模} {\ensuremath{M}} {的主单位核，不能只证明射影方向。}

{已有实际} {\ensuremath{R}} {单位} \[
\kappa_A=e/d=1-n/d,\qquad
\kappa_B=\nu/\mu=1+2n/\mu.
\tag{C26.8.1}
\] {对任意奇素数} {\ensuremath{\ell |n}，} \[
v_\ell(\kappa_A-1)=v_\ell(n).
\] {LTE} {表明} {\ensuremath{\kappa _{A}}} {在模} {\ensuremath{\ell ^{v_\ell (M)}}} {下生成所有等于1模} {\ensuremath{\ell ^{v_\ell (n)}}} {的单位。不同素数处这些循环的阶为互素素数幂，故一条} {\ensuremath{\kappa _{A}}} {的幂已通过} {CRT} {覆盖其乘积。}

{若} {\ensuremath{4|n}，二进} {LTE} {同样适用于} {\ensuremath{\kappa _{A}}。若} {\ensuremath{n=2}} {mod4，则} {\ensuremath{\kappa _{A}=3}} {mod4，而} \[
v_2(\kappa_B-1)=2.
\] {\ensuremath{\kappa _{B}}} {在奇素数处仍恰从} {\ensuremath{v_\ell (n)}} {层开始，所以其幂覆盖``等于1模} {2n''的全部主单位；再用} {\ensuremath{\kappa _{A}}} {补另一个模4类。故无论} {\ensuremath{n}} {的局部类型如何，} \[
\boxed{R^\times\longrightarrow(R/MR)^\times
\text{ 满射}.}
\tag{C26.8.2}
\]

{注意本文并未要求实际反演2或3；它们若整除} {\ensuremath{n}，就始终不在反演集合。利用的是不除} {\ensuremath{n}} {的真实单位在有限商中的像。}

\subsection*{\texorpdfstring{{C26.9} {从有限同余精确到终端}}{C26.9 从有限同余精确到终端}}\label{C26-h167}
\addcontentsline{toc}{subsection}{{C26.9} {从有限同余精确到终端}}

{全部使用的规范返回均属于} {\ensuremath{\mathrm{GL}_{2}(R)}：\ensuremath{p}、四对合标量、所用安全返回的行列式已经逐步局部化。因此它们保持} \textbf{{\ensuremath{R}} {本原性}}{。即使整数代表后来出现共同因子，该因子只能是} {\ensuremath{R}} {单位；不能忽略这一点而直接对模方向作结论。}

{由} {(7.4)、(8.2)，选} {\ensuremath{\epsilon }} {属于} {\ensuremath{R^\times}，使} {\ensuremath{\epsilon =y}} {\ensuremath{mod}} {\ensuremath{M}。\ensuremath{R}} {为} {\ensuremath{Z}} {的有限局部化，故是} {PID；(z,y)} {与} {(0,epsilon)} {都是} {\ensuremath{R}} {本原列且模} {\ensuremath{M}} {相同。以下显式主同余输送补上精确性。}

{取} {alpha,beta} {属于} {\ensuremath{R}，使} {\ensuremath{\alpha }} {z+beta} {\ensuremath{y=1}。通过} \[
(\alpha,\beta)\mapsto(\alpha+ty,\beta-tz)
\] {选取} {\ensuremath{t}，使} {\ensuremath{\alpha =0}} {\ensuremath{mod}} {\ensuremath{M}；因为} {\ensuremath{y}} {模} {\ensuremath{M}} {为单位，总能做到。此时} {\ensuremath{\beta =\epsilon ^{-1}}} {\ensuremath{mod}} {\ensuremath{M}。于是} \[
P=\begin{pmatrix}
y/\epsilon&-z/\epsilon\\
\epsilon\alpha&\epsilon\beta
\end{pmatrix}
\in\Gamma(M,R),
\quad
P\binom zy=\binom0\epsilon.
\tag{C26.9.1}
\] {(4.1)} {保证} {\ensuremath{P}} {有真实正向实现。因此有限同余终端提升为精确核心} \[
(\epsilon^{2p},(-2\epsilon)^p,0^r).
\] {实际平均词可能带另一个非零共同尺度，它仍保持同一个零和} {\ensuremath{p}} {组。}

{令} {\ensuremath{j=\operatorname{floor}(p/3)}，取} {2j} {份} {\ensuremath{\epsilon }、\ensuremath{j}} {份} {-2epsilon} {和} {p-3j} {份零；数目恰为} {\ensuremath{p}，和为零，且} {p-3j} {属于} {\{1,2\}\textless=r。执行一次平均产生} {\ensuremath{p}} {个零。调用}第\ref{construct:zero-tail}节{完成全系统。必要性仍由既有} {\ensuremath{G}} {不变量给出，定理得证。}


\section{上端重叠构造及辅助因子修正}\label{construct:upper-overlap}
\subsection*{\texorpdfstring{{C27.2} {新的四步原位置对合}}{C27.2 新的四步原位置对合}}\label{C27-h168}
\addcontentsline{toc}{subsection}{{C27.2} {新的四步原位置对合}}

{沿用共同核心} \[
K(a,z)=(a^{2p},(-2a+rz)^p,(-pz)^r).
\] {记三种实际值为} {\ensuremath{A}、\ensuremath{B}、\ensuremath{C}。先平均} \[
(A^{p-1},C),
\] {产生} {\ensuremath{p}} {份} \[
D=\frac{p-1}{p}a-z.
\] {现在的重数为} \[
(A^{p+1},B^p,C^{r-1},D^p).
\] {从新} {\ensuremath{D}} {块保留} {\ensuremath{r}} {份作最终载体。对整数} {\ensuremath{h}，另平均两个互不相交、组成相同的} {\ensuremath{p}} {组} \[
(A^{2h-1},B^h,C^{p+1-3h}),
\tag{C27.2.1}
\] {再平均其余} {\ensuremath{p}} {份。这最后一组含新} {\ensuremath{D}} {的其余} {p-r} {份，所以整个词确实发生重叠。}

{容量等价于} \[
\boxed{
h\in I_J=
\left[
\left\lceil\frac{2p-r+3}{6}\right\rceil,
\left\lfloor\frac{p+3}{4}\right\rfloor
\right]\cap\mathbb Z.}
\tag{C27.2.2}
\] {在本文参数下} {\ensuremath{h\ge 1}；上界也保证} {\ensuremath{p+1-3h\ge 0}、\ensuremath{2h\le p}。其余两个约束分别是} \[
2(2h-1)\le p+1,\qquad2(p+1-3h)\le r-1.
\]

{所得} {(a,z)} {返回矩阵为} \[
\frac1{p^2}J_h,\qquad
J_h=
\begin{pmatrix}
-p&p[nh-p(p+1)]\\
-(p-1)&p
\end{pmatrix}.
\tag{C27.2.3}
\] {直接计算} \[
\boxed{J_h^2=p f_h I,\qquad
f_h=p^3-n(p-1)h.}
\tag{C27.2.4}
\] {\ensuremath{f_{h}=p^{3}}} {\ensuremath{mod}} {\ensuremath{n}，故非零且与} {\ensuremath{n}} {互素。重复该四步返回，实际共同尺度为} {\ensuremath{f_{h}/p^{3}}。于是每个} {\ensuremath{J_{h}}} {都有真正正向射影逆。}

{在整幺模} {(\ensuremath{z,y)=(z,a+pz)}} {坐标中，实际矩阵更简单：} \[
\frac1{p^2}J_h\sim
\begin{pmatrix}
1&-(p-1)/p^2\\
nh/p&-1
\end{pmatrix}.
\tag{C27.2.5}
\] {它在每个} {\ensuremath{\ell |n}} {处可逆并将} {\ensuremath{y}} {变为} {-\ensuremath{y}} {\ensuremath{mod}} {\ensuremath{\ell }，故保持合法域。中间原子也安全：一旦全同余，后续} {\ensuremath{p}} {平均不能重新产生不同余，和最终合法性矛盾。}

\paragraph*{\texorpdfstring{{2.1} {同一矩阵适用于一般加权核心}}{2.1 同一矩阵适用于一般加权核心}}\label{C27-h169}
\addcontentsline{toc}{paragraph}{{2.1} {同一矩阵适用于一般加权核心}}

{这一恒等式不依赖} {\ensuremath{m=2}。对任意整数} {\ensuremath{m\ge 1}，考虑} \[
K_m(a,z)=(a^{mp},(-ma+rz)^p,(-pz)^r),\qquad n=(m+1)p+r.
\] {先平均} {(\ensuremath{A^{p-1},C)}，保留新} {\ensuremath{D}} {的} {\ensuremath{r}} {份，再平均} {\ensuremath{m}} {个组成相同的组} \[
(A^{mh-1},B^h,C^{p+1-(m+1)h}),
\] {最后平均其余} {\ensuremath{p}} {份。只要所有计数非负、} \[
m(mh-1)\le(m-1)p+1,\quad mh\le p,\quad
m[p+1-(m+1)h]\le r-1,
\tag{C27.2.6}
\] {这就是原位置} {m+2} {原子返回，其矩阵仍然恰为} {(2.3)，没有额外} {\ensuremath{m}} {系数。因此} {(2.4)、(2.5)} {和后文的单位平移} {-\ensuremath{1/p^{2}}} {也不依赖} {\ensuremath{m}。}

{当} {\ensuremath{p\ge 2m+1}} {时，容量区间简化为} \[
\left\lceil\frac{mp+m+1-r}{m(m+1)}\right\rceil
\le h\le
\left\lfloor\frac{(m-1)p+m+1}{m^2}\right\rfloor.
\tag{C27.2.7}
\] {其实端点差为} {(\ensuremath{mr-p+m+1)/(m^{2}(m+1))}。这同时说明优点和限制：代数控制可跨块权重统一，但同一个模板仍在} {\ensuremath{r}} {约小于} {p/m} {时容量不足。没有因此证明任意} {\ensuremath{m}、任意} {\ensuremath{r}} {的全输入定理。}

\subsection*{\texorpdfstring{{C27.3} {容量一次覆盖上端}}{C27.3 容量一次覆盖上端}}\label{C27-h170}
\addcontentsline{toc}{subsection}{{C27.3} {容量一次覆盖上端}}

{旧} {\ensuremath{A}} {对合为} \[
A_j=\begin{pmatrix}0&-d_j\\-1&0\end{pmatrix},
\quad d_j=p^2-nj,
\] {实际矩阵} {\ensuremath{A_{j}/p}，参数区间} \[
I_A=
\left[
\left\lceil\frac{2p-r}{6}\right\rceil,
\min\left(\left\lfloor\frac{2p-r}{4}\right\rfloor,
\left\lfloor\frac p3\right\rfloor\right)
\right]\cap\mathbb Z.
\]

{在} {(1.1)} {下，\ensuremath{I_{J}}} {的实端点差为} \[
\frac{2r-p+3}{12}\ge2,
\] {故包含两个相邻整数} {h,h+1。}

{\ensuremath{I_{A}}} {至少含三个相邻整数} {j,j+1,j+2：若} {\ensuremath{r\le 2p/3}，其实区间长度为} {\ensuremath{r/6>3}；若} {\ensuremath{r>2p/3}，长度为} {(\ensuremath{2p-r)/12>3}，因为} {\ensuremath{p\ge 37}、\ensuremath{r\le p-1}。取整后仍至少有三个整数。}

{第三个} {\ensuremath{A}} {只在后面的二进单位尺度补齐中使用。以上容量同时适用于} {\ensuremath{r=p-2}、p-1，不需要两个不同证明。}

\subsection*{\texorpdfstring{{C27.4} {根参数中的} {\ensuremath{r}} {消失}}{C27.4 根参数中的  消失}}\label{C27-h171}
\addcontentsline{toc}{subsection}{{C27.4} {根参数中的} {\ensuremath{r}} {消失}}

{记} \[
d=d_j,\quad e=d_{j+1},\quad d_2=d_{j+2},\qquad
f=f_h,\quad g=f_{h+1}.
\] {有} \[
e=d-n,\qquad g=f-n(p-1),
\] {且} {\ensuremath{\operatorname{gcd}(d,e)=\operatorname{gcd}(f,g)=1}。后一个等式来自} \[
\gcd(f,n(p-1))=\gcd(p^3,n(p-1))=1.
\] {由于} {\ensuremath{0<j,j+1,j+2,h,h+1<p}，\ensuremath{d,e,d_{2},f,g}} {均与} {\ensuremath{p}} {互素，也均与} {\ensuremath{n}} {互素。}

{两个相邻商为} \[
D_0=A_j^{-1}A_{j+1}=\operatorname{diag}(1,e/d),
\quad
T_0=J_h^{-1}J_{h+1}
=\begin{pmatrix}1&-pn/f\\0&g/f\end{pmatrix}.
\] {因此} \[
\boxed{[D_0,T_0]=U\!\left(-\frac{pn^2}{eg}\right).}
\tag{C27.4.1}
\] {令} {\ensuremath{R_{0}=Z}{[}\ensuremath{1/|defg|}{]}。两个倍率及其逆生成} {\ensuremath{R_{0}}；通过共轭、逆共轭和根参数相加，得到} \[
U(pn^2R_0).
\] {\ensuremath{A_{j}}} {将上根交换为下根，得到} {\ensuremath{L(pn^{2}R_{0})}。全部因子均有原位置实现和正向逆。}

{既有深同余引理给出} \[
\Gamma(p^2n^4,R_0).
\tag{C27.4.2}
\] {不允许直接删去这里的} {\ensuremath{p}。下一节提供删除它的完整论证。}

\subsection*{\texorpdfstring{{C27.5} {用两个幂零方向消去唯一辅助素数} {\ensuremath{p}}}{C27.5 用两个幂零方向消去唯一辅助素数 }}\label{C27-h172}
\addcontentsline{toc}{subsection}{{C27.5} {用两个幂零方向消去唯一辅助素数} {\ensuremath{p}}}

{置} \[
N=\begin{pmatrix}
p(p-1)&-p^2\\
(p-1)^2&-p(p-1)
\end{pmatrix},
\qquad N^2=0.
\] {由} {(2.3)、(2.4)} {得} \[
J_hU(pn^2t)J_h^{-1}=I+\frac{n^2t}{f}N.
\tag{C27.5.1}
\] {右侧属于} {\ensuremath{\mathrm{SL}_{2}(R_{0})}，原位置正向可逆。}

{将} {\ensuremath{t}} {换成} {\ensuremath{n^{2}t}，则这些矩阵模} {\ensuremath{n^{4}}} {恒为} {\ensuremath{I}，而模} {\ensuremath{p^{2}}} {的参数} {\ensuremath{n^{4}t/f}} {遍历全部} {\ensuremath{Z/p^{2}Z}。它们的非零幂零方向由本原列} \[
v=\binom p{p-1}
\] {给出。\ensuremath{A_{j}}} {共轭再给出方向} {\ensuremath{A_{j}}} {\ensuremath{v}。两列的行列式} \[
\det(v,A_jv)=d(p-1)^2-p^2
\tag{C27.5.2}
\] {为模} {\ensuremath{p}} {单位。}

{更明确地，令} {\ensuremath{P=(v,A_{jv})}、\ensuremath{\Delta =det}} {\ensuremath{P}，则} \[
P^{-1}NP=\begin{pmatrix}0&-\Delta\\0&0\end{pmatrix},
\quad
P^{-1}A_jNA_j^{-1}P=
\begin{pmatrix}0&0\\-\Delta/d&0\end{pmatrix}.
\tag{C27.5.3}
\] {\ensuremath{P}} {在} {\ensuremath{Z/p^{2}Z}} {上可逆，两个系数也都是单位。因此} {(5.1)} {及其} {\ensuremath{A_{j}}} {共轭在模} {\ensuremath{p^{2}}} {下生成整个} {\ensuremath{\mathrm{SL}_{2}(Z/p^{2}Z)}，同时在模} {\ensuremath{n^{4}}} {下恒为} {\ensuremath{I}。}

{结合} {(4.2)} {与} {CRT，得到真正的包含} \[
\boxed{\Gamma(n^4,R_0)\text{ 正向可执行}.}
\tag{C27.5.4}
\] {证明中只有固定深度} {\ensuremath{p^{2}}，且用显式初等根生成；不存在逐} {\ensuremath{p}} {的有限群搜索。}

\subsection*{\texorpdfstring{{C27.6} {环扩张和全部单位像}}{C27.6 环扩张和全部单位像}}\label{C27-h173}
\addcontentsline{toc}{subsection}{{C27.6} {环扩张和全部单位像}}

{先用}第\ref{construct:upper-local}节{激活保载体返回} \[
F_\epsilon=
\begin{pmatrix}
\epsilon/p&r(p-\epsilon)/(3p)\\
0&1
\end{pmatrix},
\quad\epsilon\in\{1,-1\},\quad\epsilon=p\pmod3.
\] {其整数代表行列式} {\ensuremath{\epsilon }} {\ensuremath{p}} {与} {\ensuremath{n^{4}}} {互素。上三角倍率将} {\ensuremath{p}} {局部化；原有根于是扩展为} {\ensuremath{U(n^{2}R)}、\ensuremath{L(n^{2}R)}。再用深同余引理，\ensuremath{\Gamma (n^{4},R)}} {在扩大的环上也成立。}

{第三个} {\ensuremath{A}} {的商提供} {\ensuremath{d_{2}/d}} {及其逆，可将} {\ensuremath{d_{2}}} {的素因子加入} {\ensuremath{R}。所有反演素数均不除} {\ensuremath{n}。}

{若} {3} {不整除} {\ensuremath{n}，再用真实保留} {\ensuremath{B}} {的返回} \[
N_0=\begin{pmatrix}p&0\\2&-r\end{pmatrix}.
\] {实际矩阵} {\ensuremath{N_{0}/p}} {由三个原子平均实现：两个} {\ensuremath{p}} {组分别全取原} {\ensuremath{A}，留下} {\ensuremath{r}} {份原} {\ensuremath{B}，最后平均其余} {\ensuremath{p}} {份。\ensuremath{det}} {\ensuremath{N_{0}=-pr}} {与} {\ensuremath{n}} {互素，所以逆已由} {(5.4)} {提供。它对下根的共轭倍率为} {-r/p，因而将} {\ensuremath{r}} {局部化。-r/p} {是模} {\ensuremath{n}} {的单位3的实际} {\ensuremath{R}} {单位提升。}

{现在沿用前一证明的一般安全} {\ensuremath{A}} {返回} \[
M_{q,j}=\begin{pmatrix}p-q&nj-pq\\-1&0\end{pmatrix},
\] {其中} {\ensuremath{q=3j+s}} {满足原位置容量。每个} {\ensuremath{\operatorname{gcd}(q,n)=1}} {的返回都由} {(5.4)} {激活逆，其与} {\ensuremath{A_{j}}} {的上三角商将行列式局部化，使} {\ensuremath{q}} {进入} {\ensuremath{R}} {单位的模} {\ensuremath{n}} {像。只需有限多个} {\ensuremath{q}，所以} {\ensuremath{R}} {仍为有限局部化。}

{为避免外推旧参数范围，重新核对单位覆盖用到的区间：} \[
L\le q\le U,\quad
L=\begin{cases}r/2+2,&r\text{ 偶},\\(r+9)/2,&r\text{ 奇},\end{cases}
\quad U=3\lfloor p/2\rfloor+\lfloor r/2\rfloor.
\tag{C27.6.1}
\] {它的容量证明对} {\ensuremath{r\le p-1}、\ensuremath{r\ge 10}} {原样成立。}

\begin{itemize}
\tightlist
\item
  {3} {不整除} {\ensuremath{n}} {时，正平衡代表通过} {\ensuremath{b->|n-3b|}} {降至} {\ensuremath{b\le n/4}，再乘3进入} {(6.1)。所需} {\ensuremath{3L-3\le U}：偶} {\ensuremath{r}} {归为} {\ensuremath{2r+9\le 3p}，奇} {\ensuremath{r}} {归为} {\ensuremath{2r+25\le 3p}，均由} {\ensuremath{p\ge 37}、\ensuremath{r\le p-1}} {成立。}
\item
  {\ensuremath{3|n}} {时，置} {\ensuremath{k=n/3}，使用} {\ensuremath{k}±1} {生成单位约化到模} {\ensuremath{k}} {的核；对正平衡代表} {\ensuremath{b\le k/2}，从} {k-b、k+b、\ensuremath{b}} {中选安全代表。所需} {\ensuremath{k/2-2\ge L}：偶} {\ensuremath{r}} {归为} {\ensuremath{3p-2r\ge 24}，奇} {\ensuremath{r}} {归为} {\ensuremath{3p-2r\ge 39}。后者的最坏情形为} {\ensuremath{r=p-2}，此时左侧} {\ensuremath{p+4\ge 41}，所以两条边界均无小参数例外。}
\end{itemize}

{具体代表构造和核生成证明见前文第5节，由这些不等式保证全部可移用。因此} \[
R^\times\longrightarrow(\mathbb Z/n\mathbb Z)^\times
\quad\text{满射}.
\tag{C27.6.2}
\] {这里不要求3在} {\ensuremath{n}} {可被3整除时成为} {\ensuremath{R}} {单位。}

\subsection*{\texorpdfstring{{C27.7} {同一个返回积覆盖所有局部方向}}{C27.7 同一个返回积覆盖所有局部方向}}\label{C27-h174}
\addcontentsline{toc}{subsection}{{C27.7} {同一个返回积覆盖所有局部方向}}

{取} {\ensuremath{V=J_{h}}} {\ensuremath{A_{j}/p^{3}}。在} {(z,y)} {坐标中，} \[
\boxed{V=\begin{pmatrix}1&-1/p^2\\0&1\end{pmatrix}\pmod n.}
\tag{C27.7.1}
\] {其平移系数在每个} {\ensuremath{\ell |n}} {处均为单位。与旧} {BA} {平移} {-3/p} {相比，模3的失效已经消失。}

{其迹为} \[
\operatorname{tr}V=
2-\frac{n[ph+(p-1)j]}{p^3}.
\tag{C27.7.2}
\] {若} {\ensuremath{4|n}，迹自动为2} {mod4；若} {\ensuremath{n=2}} {mod4，从相邻} {h,h+1} {中选偶数} {\ensuremath{h}，亦得到迹为2} {mod4。因此前文局部循环引理在全部二进层适用。}

{若} {\ensuremath{3|n}，令} {\ensuremath{c=1-4det(V)/(tr}} {\ensuremath{V)^{2}}。若} {\ensuremath{9|n}，有} {\ensuremath{c=0}} {mod9。若} {\ensuremath{v_{3}(n)=1}，直接展开得} \[
\frac c3=\frac{n/3}{p^3}(j-h)\pmod3.
\tag{C27.7.3}
\] {从相邻} {j,j+1} {中选一个，使右侧不等于2，即满足} {\ensuremath{c\ne 6}} {mod9。这一选择不影响二进条件，因为} {(7.2)} {中} {\ensuremath{j}} {的系数} {p-1} {为偶数。}

{所有} {\ensuremath{\ell \ge 5}} {直接适用同一局部循环引理。因此选定的一条} {\ensuremath{V}，在每个} {\ensuremath{Z/\ell ^{e}Z}} {的合法仿射图上都形成完整单循环。指数} {CRT} {使任意合法列满足} \[
z=0\pmod{n^4},\qquad y\text{ 为模 }n^4\text{ 单位}.
\tag{C27.7.4}
\] {不再需要先用额外三进返回、再让其他返回保持三进数据的两阶段拼接。}

\subsection*{\texorpdfstring{{C27.8} {共同尺度与精确终端}}{C27.8 共同尺度与精确终端}}\label{C27-h175}
\addcontentsline{toc}{subsection}{{C27.8} {共同尺度与精确终端}}

{已知} {\ensuremath{R}} {单位} \[
\kappa_1=e/d=1-n/d,\qquad
\kappa_2=d_2/d=1-2n/d.
\] {\ensuremath{n}} {奇或} {\ensuremath{4|n}} {时，\ensuremath{\kappa _{1}}} {的幂借助} {LTE、CRT} {生成模} {\ensuremath{n^{4}}} {到模} {\ensuremath{n}} {的全部主单位核。}

{\ensuremath{n=2}} {mod4} {时，\ensuremath{\kappa _{1}=3}} {mod4，而} {\ensuremath{\kappa _{2}}} {满足} \[
v_2(\kappa_2-1)=2;
\] {在每个奇} {\ensuremath{\ell |n}} {处，\ensuremath{v_\ell (\kappa _{2}-1)=v_\ell (n)}。因此} {\ensuremath{\kappa _{2}}} {生成等于1} {\ensuremath{mod}} {2n} {的全部单位核，再由} {\ensuremath{\kappa _{1}}} {补另一个模4类。配合} {(6.2)，\ensuremath{R}} {单位模} {\ensuremath{n^{4}}} {满射。}

{全部所用返回的规范代表在最终} {\ensuremath{R}} {中属于} {\ensuremath{\mathrm{GL}_{2}(R)}，保持} {\ensuremath{R}} {本原性；尤其} {\ensuremath{J_{h}}} {的行列式} {\ensuremath{p}} {\ensuremath{f_{h}}} {中的} {\ensuremath{p}} {已在第6节局部化。于是} {(7.4)} {可匹配一个} {\ensuremath{\epsilon }} {属于} {\ensuremath{R^\times}。正文的精确本原列纤维引理把状态精确送到} \[
(z,y)=(0,\epsilon).
\] {核心成为} \[
(\epsilon^{2p},(-2\epsilon)^p,0^r).
\] {以} {\ensuremath{\operatorname{floor}(p/3)}} {份} {-2epsilon、其两倍份} {\ensuremath{\epsilon }} {和} {\ensuremath{p}} {mod3} {份零组成零和} {\ensuremath{p}} {组，平均后产生} {\ensuremath{p}} {个零，再接统一零触发尾部。全输入入口已由既有上半带归约保证，完成定理。}


\section{上端小素数的有限容量表}\label{construct:upper-small}
\subsection*{\texorpdfstring{{C28.3} {两个} {\ensuremath{A}、两个重叠} {\ensuremath{J}} {与单位密度}}{C28.3 两个 、两个重叠  与单位密度}}\label{C28-h176}
\addcontentsline{toc}{subsection}{{C28.3} {两个} {\ensuremath{A}、两个重叠} {\ensuremath{J}} {与单位密度}}

{在核心} \[
K(a,z)=(a^{2p},(-2a+rz)^p,(-pz)^r),\quad n=3p+r,
\] {使用}第\ref{construct:upper-overlap}节{的} \[
A_j=\begin{pmatrix}0&-d_j\\-1&0\end{pmatrix},
\quad d_j=p^2-nj,
\] {以及} \[
J_h=\begin{pmatrix}-p&p[nh-p(p+1)]\\-(p-1)&p\end{pmatrix},
\quad J_h^2=p f_h I,\quad f_h=p^3-n(p-1)h.
\] {实际尺度分别为1/p、\ensuremath{1/p^{2}}，均已有原位置实现和正向射影逆。}

{若} {\ensuremath{A_{j},A_{j+1}}} {和} {\ensuremath{J_{h},J_{h+1}}} {的容量成立，原证明第4、5节即给} \[
U(pn^2R),\ L(pn^2R)
\longrightarrow\Gamma(p^2n^4,R)
\longrightarrow\Gamma(n^4,R).
\tag{C28.3.1}
\] {第三个} {\ensuremath{A}} {不参与这一控制层的建立；此前它仅用于构造一个二进单位。}

{以下八组全部满足这些容量：}

{ % do not increment counter
\begin{longtable}[]{@{}lllllll@{}}
\toprule\noalign{}
{\ensuremath{p}} & {\ensuremath{n}} & {\ensuremath{A}起始\ensuremath{j}} & {\ensuremath{J}起始\ensuremath{h}} & {安全数字区间\ensuremath{I}} & {\ensuremath{I}中单位数} & {phi(n)} \\
\midrule\noalign{}
\endhead
\bottomrule\noalign{}
\endlastfoot
{19} & {74} & {4} & {4} & {{[}13,35{]}} & {12} & {36} \\
{19} & {75} & {4} & {4} & {{[}11,36{]}} & {14} & {40} \\
{23} & {90} & {5} & {5} & {{[}15,43{]}} & {8} & {24} \\
{23} & {91} & {4} & {5} & {{[}13,44{]}} & {24} & {72} \\
{29} & {114} & {6} & {6} & {{[}18,55{]}} & {12} & {36} \\
{29} & {115} & {5} & {6} & {{[}16,56{]}} & {31} & {88} \\
{31} & {122} & {6} & {6} & {{[}19,59{]}} & {21} & {60} \\
{31} & {123} & {6} & {6} & {{[}17,60{]}} & {28} & {80} \\
\end{longtable}
}

{区间统一为} \[
I=\left[
\left\lfloor r/2\right\rfloor+
\begin{cases}2,&r\text{ 偶}\\5,&r\text{ 奇}\end{cases},
\ 3\lfloor p/2\rfloor+\lfloor r/2\rfloor
\right].
\] {其中每个} {\ensuremath{q=3j+s}} {都有真实} {\ensuremath{A}} {型返回} \[
M_{q,j}=\begin{pmatrix}p-q&nj-pq\\-1&0\end{pmatrix}.
\tag{C28.3.2}
\] {若} {\ensuremath{\operatorname{gcd}(q,n)=1}，行列式与控制层互素，故逆可执行。其与旧} {\ensuremath{A}} {的商局部化行列式，提供} {\ensuremath{q}} {\ensuremath{mod}} {\ensuremath{n}} {的} {\ensuremath{R}} {单位提升；\ensuremath{p}} {的局部化仍由} {\ensuremath{F_\eta }、\ensuremath{\eta =p}} {mod3} {给出。}

{表中每行均有} {\ensuremath{4|I}∩\ensuremath{units|>\phi (n)}，且} {\ensuremath{I}} {位于} {(0,n/2)。所以正负单位代表超过单位群的一半，每个模} {\ensuremath{n}} {单位都是其中两个元素之比。这个抽屉原理证明} {\ensuremath{R}} {单位模} {\ensuremath{n}} {满射，替代此前依赖大} {\ensuremath{p}} {不等式的小代表构造。}

\subsection*{\texorpdfstring{{C28.4} {二进尺度不需要第三个迹零} {\ensuremath{A}}}{C28.4 二进尺度不需要第三个迹零 }}\label{C28-h177}
\addcontentsline{toc}{subsection}{{C28.4} {二进尺度不需要第三个迹零} {\ensuremath{A}}}

{\ensuremath{n}} {为奇数时，相邻} {\ensuremath{A}} {的单位} {\ensuremath{d_{j+1}/d_{j}=1-n/d_{j}}} {已借} {LTE、CRT} {生成全部主单位核。}

{对表中偶数} {\ensuremath{n=4p-2}，统一取} {\ensuremath{q=p+6}，并使用} {(3.2)} {的两个相距2的可行} {\ensuremath{j}：}

{ % do not increment counter
\begin{longtable}[]{@{}lllll@{}}
\toprule\noalign{}
{\ensuremath{p}} & {\ensuremath{q}} & {j,j+2} & {\ensuremath{D=nj-pq}} & {D+2n} \\
\midrule\noalign{}
\endhead
\bottomrule\noalign{}
\endlastfoot
{19} & {25} & {6,8} & {-31} & {117} \\
{23} & {29} & {7,9} & {-37} & {143} \\
{29} & {35} & {8,10} & {-103} & {125} \\
{31} & {37} & {8,10} & {-171} & {73} \\
\end{longtable}
}

{两个返回均安全，逆由} {(3.1)} {提供。它们的行列式已成为} {\ensuremath{R}} {单位，于是} \[
\kappa=(D+2n)/D=1+2n/D
\] {满足} {\ensuremath{v_{2}(\kappa -1)=2}，并在每个奇} {\ensuremath{\ell |n}} {处满足} {\ensuremath{v_\ell (\kappa -1)=v_\ell (n)}。再配合} {\ensuremath{d_{j+1}/d_{j}=3}} {mod4，正好生成全部主单位核。}

{方向仍使用} {\ensuremath{J_{h}}} {\ensuremath{A_{j}/p^{3}}，其模} {\ensuremath{n}} {平移为} {-\ensuremath{1/p^{2}}；相邻} {\ensuremath{h}} {调整二进迹，相邻} {\ensuremath{j}} {排除唯一三进退化类。原局部循环与精确纤维论证完全适用，所以八组均为全输入定理。}

{接} {\ensuremath{p\ge 37}} {的旧定理，已经得到所有} {\ensuremath{p\ge 19}} {的两条上端边界。}

\subsection*{\texorpdfstring{{C28.5} {两种重叠返回给统一仿射行列式}}{C28.5 两种重叠返回给统一仿射行列式}}\label{C28-h178}
\addcontentsline{toc}{subsection}{{C28.5} {两种重叠返回给统一仿射行列式}}

{余下只需} {(13,50)、(13,51)、(17,66)。它们各只有一个所需的旧} {\ensuremath{A}，但可以把首次平均改为} \[
(A^{p-\sigma},C^\sigma),\qquad \sigma=1,3.
\] {保留新} {\ensuremath{D}} {的} {\ensuremath{r}} {份，再平均两个相同组} \[
(A^{2h-\sigma},B^h,C^{p+\sigma-3h}),
\] {最后平均其余} {\ensuremath{p}} {份。所得四步对合为} \[
J_{\sigma,h}=
\begin{pmatrix}
-p\sigma&p[nh-p(p+\sigma)]\\
-(p-\sigma)&p\sigma
\end{pmatrix},
\quad
J_{\sigma,h}^2=p f_{\sigma,h}I,
\] \[
f_{\sigma,h}=p^3-n(p-\sigma)h.
\tag{C28.5.1}
\]

{以下是同一模板的全部参数，\ensuremath{h}和h+1均使用：}

{ % do not increment counter
\begin{longtable}[]{@{}llllll@{}}
\toprule\noalign{}
{\ensuremath{p}} & {\ensuremath{n}} & {唯一\ensuremath{A}参数\ensuremath{j}} & {\ensuremath{d=p^{2}-nj}} & {\ensuremath{\sigma =1}的\ensuremath{h}} & {\ensuremath{\sigma =3}的\ensuremath{h}} \\
\midrule\noalign{}
\endhead
\bottomrule\noalign{}
\endlastfoot
{13} & {50} & {3} & {19} & {3} & {4} \\
{13} & {51} & {3} & {16} & {3} & {4} \\
{17} & {66} & {4} & {25} & {4} & {5} \\
\end{longtable}
}

{每行的四个} {\ensuremath{f}} {两两按相邻族互素，均为奇数，且与} {pn} {互素。记由四个} {\ensuremath{f}} {局部化的环为} {\ensuremath{R_{0}}。相邻} {\ensuremath{J}} {的商为} \[
T_\sigma=
\begin{pmatrix}
1&-p\sigma n/f_{\sigma,h}\\
0&f_{\sigma,h+1}/f_{\sigma,h}
\end{pmatrix}.
\] {两族固定点不同，其交换子为} \[
\boxed{
[T_1,T_3]
=U\left(-\frac{2p^2n^2}{f_{1,h_1+1}f_{3,h_3+1}}\right).
}
\tag{C28.5.2}
\] {倍率及逆倍率生成} {\ensuremath{R_{0}}，因此有} {\ensuremath{U(kR_{0})}，\ensuremath{k=2p^{2}n^{2}}。即使} {\ensuremath{d}} {尚不是} {\ensuremath{R_{0}}} {单位，也有} \[
A\,U(dx)\,A^{-1}=L(x),\qquad x\in kR_0,
\] {所以} {\ensuremath{L(kR_{0})}} {同样可执行。深同余引理给} \[
\Gamma(4p^4n^4,R_0).
\tag{C28.5.3}
\]

\subsection*{\texorpdfstring{{C28.6} {一个共同的两轮局部提升}}{C28.6 一个共同的两轮局部提升}}\label{C28-h179}
\addcontentsline{toc}{subsection}{{C28.6} {一个共同的两轮局部提升}}

{置} {\ensuremath{M=4n^{4}}，\ensuremath{f=f_{1,h_{1}}}，} \[
N=\begin{pmatrix}p(p-1)&-p^2\\(p-1)^2&-p(p-1)\end{pmatrix},
\quad N^2=0.
\] {从} {(5.2)} {已有的上根出发，} \[
E(t)=J_{1,h_1}U(p^2Md\,t)J_{1,h_1}^{-1}
=I+\frac{pMd\,t}{f}N.
\tag{C28.6.1}
\] {E(t)} {与} {\ensuremath{AE(t)A^{-1}}} {都属于} {\ensuremath{\mathrm{SL}_{2}(R_{0})}，且模} {\ensuremath{M}} {为} {\ensuremath{I}。乘} {\ensuremath{d}} {正是为了清除后一共轭的分母，不要求} {\ensuremath{d}} {先可逆。}

{令} {\ensuremath{v=(p,p-1)^T}，\ensuremath{P=(v,Av)}，\ensuremath{\Delta =det}} {\ensuremath{P=d(p-1)^{2}-p^{2}}。\ensuremath{\Delta }} {为模} {\ensuremath{p}} {单位。精确恒等式为} \[
P^{-1}E(t)P=U(-pMd\Delta t/f),\quad
P^{-1}AE(t)A^{-1}P=L(-pM\Delta t/f).
\tag{C28.6.2}
\] {这里} {\ensuremath{P}} {只作有限环中的坐标变换，不宣称它本身正向可执行。}

{在} {\ensuremath{B=Z/p^{4}Z}} {中，两参数遍历} {pB。U(pB)、L(pB)} {包含} {\ensuremath{\Gamma (p^{2},B)}：对角\ensuremath{1+p^{2}a}用一次} {Gauss} {恒等式分解，其余由上下根消去。这也是既有有限环深同余证明的直接特例。结合} {(5.3)，得到} \[
\Gamma(p^2M,R_0).
\] {而} {\ensuremath{AE(t)A^{-1}}} {模} {\ensuremath{p^{2}}为} {U(pM} {\ensuremath{d^{2}t/f)}。选} {\ensuremath{t=f/d^{2}}} {\ensuremath{mod}} {\ensuremath{p}，就与} {U(pM)} {同余；二者也模} {\ensuremath{M}} {同余。用刚取得的主同余层纠正，得到}\textbf{{精确可执行的} {U(pM)}}{。}

{再共轭} {\ensuremath{J_{1,h_{1}}}：} \[
J_{1,h_1}U(pMd\,t)J_{1,h_1}^{-1}
=I+\frac{Md\,t}{f}N.
\] {它及其} {\ensuremath{A}} {共轭在模} {\ensuremath{p^{4}}} {下变成参数为单位倍的完整上下根，因此生成整个} {\ensuremath{\mathrm{SL}_{2}(B)}，同时模} {\ensuremath{M}} {为} {\ensuremath{I}。接原深层即得} \[
\boxed{\Gamma(M,R_0).}
\tag{C28.6.3}
\]

{这两轮完全相同地处理三行参数；没有枚举一个} {\ensuremath{p^{4}}} {模群。随后激活} {\ensuremath{F_\eta }} {局部化} {\ensuremath{p}，已有根成为} {\ensuremath{U(2n^{2}R)}、\ensuremath{L(2n^{2}R)}。}

\subsection*{\texorpdfstring{{C28.7} {唯一的辅助2也由同一根原理消除}}{C28.7 唯一的辅助2也由同一根原理消除}}\label{C28-h180}
\addcontentsline{toc}{subsection}{{C28.7} {唯一的辅助2也由同一根原理消除}}

{\ensuremath{n=50}、66} {时} {\ensuremath{M}} {的全部素数已整除} {\ensuremath{n}，不存在额外局部支持。}

{\ensuremath{n=51}} {时} {\ensuremath{M}} {多一个辅助2，但唯一} {\ensuremath{A}} {的} {\ensuremath{d=16}} {给出} \[
A\,U(16n^4t)\,A^{-1}=L(n^4t).
\] {这里} {\ensuremath{16n^{4}t}} {已属于现有} {\ensuremath{U(2n^{2}R)}。\ensuremath{T_{1}}} {的右上项} {-pn/f} {在模2下为单位。因此} {\ensuremath{L(n^{4}t)}} {与} {\ensuremath{T_{1}}} {\ensuremath{L(n^{4}t)T_{1}^{-1}}} {在模4下具有两个独立幂零方向，参数遍历全部} {Z/4Z，生成} {\ensuremath{\mathrm{SL}_{2}(Z/4Z)}，同时模} {\ensuremath{n^{4}}为} {\ensuremath{I}。}

{由} {(6.3)} {得} {\ensuremath{\Gamma (n^{4},R)}。所以三行最终控制层统一记为} \[
Q=\begin{cases}4n^4,&n\text{ 偶},\\n^4,&n=51.\end{cases}
\tag{C28.7.1}
\] {它的素数支持均仅来自} {\ensuremath{n}。}

\subsection*{\texorpdfstring{{C28.8} {三行共用的尺度与方向接口}}{C28.8 三行共用的尺度与方向接口}}\label{C28-h181}
\addcontentsline{toc}{subsection}{{C28.8} {三行共用的尺度与方向接口}}

{选择同一个安全数字} {\ensuremath{q}} {的两个相邻参数} {j,j+1：}

{ % do not increment counter
\begin{longtable}[]{@{}llllll@{}}
\toprule\noalign{}
{\ensuremath{p}} & {\ensuremath{n}} & {\ensuremath{q}} & {j,j+1} & {\ensuremath{D_{0}=nj-pq}} & {\ensuremath{D_{1}=D_{0}+n}} \\
\midrule\noalign{}
\endhead
\bottomrule\noalign{}
\endlastfoot
{13} & {50} & {17} & {4,5} & {-21} & {29} \\
{13} & {51} & {14} & {3,4} & {-29} & {22} \\
{17} & {66} & {25} & {6,7} & {-29} & {37} \\
\end{longtable}
}

{它们均有原位置容量，行列式与} {\ensuremath{n}} {互素，故逆可执行。\ensuremath{\operatorname{gcd}(D_{0},d)=1}，其与旧} {\ensuremath{A}} {的上仿射商局部化} {\ensuremath{d}} {和} {\ensuremath{D_{0}}；第二个商再局部化} {\ensuremath{D_{1}}。因此旧} {\ensuremath{A}} {的规范代表也属于最终} {\ensuremath{\mathrm{GL}_{2}(R)}。}

{安全数字区间仍用第3节的同一公式：}

{ % do not increment counter
\begin{longtable}[]{@{}llll@{}}
\toprule\noalign{}
{p,n} & {\ensuremath{I}} & {\ensuremath{I}中单位数} & {phi(n)} \\
\midrule\noalign{}
\endhead
\bottomrule\noalign{}
\endlastfoot
{13,50} & {{[}10,23{]}} & {6} & {20} \\
{13,51} & {{[}8,24{]}} & {10} & {32} \\
{17,66} & {{[}12,31{]}} & {7} & {20} \\
\end{longtable}
}

{每行均满足正负代表超过单位群一半，故每个模} {\ensuremath{n}} {单位都有} {\ensuremath{R}} {单位提升。}

{令} {\ensuremath{\delta =D_{1}/D_{0}}，则} {\ensuremath{\delta -1=n/D_{0}}。在所有奇} {\ensuremath{\ell |n}} {处，它恰从} {\ensuremath{v_\ell (n)}} {层生成主单位；\ensuremath{n}偶时} {\ensuremath{\delta =3}} {mod4。另取} {\ensuremath{\kappa =f_{3,h_{3}+1}/f_{3,h_{3}}}，得到：}

{ % do not increment counter
\begin{longtable}[]{@{}lll@{}}
\toprule\noalign{}
{p,n} & {\ensuremath{\delta }} & {\ensuremath{\kappa }} \\
\midrule\noalign{}
\endhead
\bottomrule\noalign{}
\endlastfoot
{13,50} & {-29/21} & {-303/197} \\
{13,51} & {-22/29} & {-353/157} \\
{17,66} & {-37/29} & {-631/293} \\
\end{longtable}
}

{两条偶数行均有} {\ensuremath{v_{2}(\kappa -1)=2}。先以} {\ensuremath{\delta }} {补模4的类，再用} {\ensuremath{\kappa }} {调好全部二进位；最后以} {\ensuremath{\delta }} {的完整二进阶次幂纠正奇素数处的主单位，不改变二进结果。该阶是2幂，在所有奇素数处仍为单位，所以} {LTE} {保留所需起始深度。\ensuremath{n=51}} {只用} {\ensuremath{\delta }} {即可。}

{方向使用} {\ensuremath{J_{1,h}A/p^{3}}，分别取} {\ensuremath{h=4}、3、4。模} {\ensuremath{n}} {的平移均为} {-\ensuremath{1/p^{2}}；偶数行迹为2} {mod4；被3整除的两行} {\ensuremath{j=h}，故三进参数} {\ensuremath{c=0}} {mod9。已有局部单循环引理与} {CRT} {遂完成全部方向。}

{所有规范返回在最终} {\ensuremath{R}} {中都可逆，保持} {\ensuremath{R}} {本原性。方向及单位尺度匹配后，\ensuremath{\Gamma (Q,R)}} {的精确纤维输送到} {(\ensuremath{z,y)=(0,\epsilon )}，再由原位置零和} {\ensuremath{p}} {组触发尾部。因既有入口覆盖任意合法原输入，三行均是完整维数定理。}


\section{全部 n 大于等于 4p 的补全}\label{construct:wide-tail}
\subsection*{\texorpdfstring{{C29.2} {为什么旧缺口反而提供足够资源}}{C29.2 为什么旧缺口反而提供足够资源}}\label{C29-h182}
\addcontentsline{toc}{subsection}{{C29.2} {为什么旧缺口反而提供足够资源}}

{已有定理覆盖} \[
n\ge4p-2+d_p(n),\qquad N(p)\le4p+\lceil\log_2p\rceil.
\] {因此只须考虑} \[
n=4p+s,\quad1\le s<\lceil\log_2p\rceil<p.
\] {此时} {\ensuremath{p}} {不整除} {\ensuremath{n}，令} {\ensuremath{d=\omega (n)}。旧不等式失效恰为} \[
\boxed{d\ge s+3.}
\tag{C29.2.1}
\] {\ensuremath{n=4p}} {本身满足旧条件，已经完成。}

{若} {\ensuremath{n}} {奇，\ensuremath{s}} {奇且} {\ensuremath{d\ge 4}；\ensuremath{n}} {至少含} {\ensuremath{d}} {个不同奇素数。若} {\ensuremath{n}} {偶，\ensuremath{s}} {为正偶数，故} {\ensuremath{d\ge 5}。由最小素数乘积，得到两个统一估计：} \[
\boxed{p\ge16s+21,\qquad
\varphi(n)>8\cdot2^d+7s+36.}
\tag{C29.2.2}
\]

{证明如下。奇数情形的基点} {\ensuremath{d=4}} {给} \[
n\ge3\cdot5\cdot7\cdot11=1155,\quad\varphi(n)\ge480;
\] {偶数情形的基点} {\ensuremath{d=5}} {给} \[
n\ge2\cdot3\cdot5\cdot7\cdot11=2310,\quad\varphi(n)\ge480.
\] {增加一个不同素因子时，最小乘积至少乘13，最小} {\ensuremath{\phi }} {至少乘12。因此归纳得到} \[
n\ge65d-111,\qquad\varphi(n)>8\cdot2^d+7d+15.
\] {基点分别比较1155与149、2310与214，以及480与171、480与306；后续乘积增长超过右侧的增长。重复素因子只会增大这两个下界。}

{将} {\ensuremath{s\le d-3}} {代入，便得} {\ensuremath{n\ge 65s+84}，即} {\ensuremath{p=(n-s)/4\ge 16s+21}，以及} {(2.2)} {的} {\ensuremath{\phi }} {估计。}

{所以不能仅把旧缺口视为``任意大小的余量} {\ensuremath{s}''。它强制} {\ensuremath{p}} {相对} {\ensuremath{s}} {很大，并提供很强的单位密度；这两项正用于下面的容量和尺度证明。}

\subsection*{\texorpdfstring{{C29.3} {短参数筛选仍然成立}}{C29.3 短参数筛选仍然成立}}\label{C29-h183}
\addcontentsline{toc}{subsection}{{C29.3} {短参数筛选仍然成立}}

{记} {\ensuremath{\rho =\phi (n)/n}。每个} {\ensuremath{\ell |n}} {给定至多一个禁止余数类，CRT} {将它们合成一个平移。长度} {\ensuremath{L}} {的整数区间中，避开这些类的个数大于} \[
L\rho-(2^d-1).
\] {取} {\ensuremath{L=p-1}，由} {(2.2)，} \[
(p-1)\rho
=\frac{n-s-4}{4}\rho
\ge\frac{\varphi(n)-s-4}{4}>2^d-1.
\tag{C29.3.1}
\] {因此} {{[}1,p-1{]}} {内总有同时避开所有禁止类的参数。}

{这重新证明了尾部所需的短配比引理，而不是将旧} {\ensuremath{n<4p}} {的引理直接用于} {\ensuremath{n>4p}。具体算法只须对各素数标记一个等差数列；没有平均词搜索。}

\subsection*{\texorpdfstring{{C29.4} {全输入进入大载体核心}}{C29.4 全输入进入大载体核心}}\label{C29-h184}
\addcontentsline{toc}{subsection}{{C29.4} {全输入进入大载体核心}}

{我们保持同一个} {\ensuremath{m=2}} {核心，允许载体大于} {\ensuremath{p}：} \[
\boxed{
K(a,z)=(a^{2p},(-2a+rz)^p,(-pz)^r),\quad
r=p+s,\quad n=3p+r.}
\tag{C29.4.1}
\] {因为} {\ensuremath{\operatorname{gcd}(p,r)=1}，本原整数核心仍满足} \[
\gcd(a,z)=1,\qquad G=\gcd(a+pz,n).
\] {以下证明任意合法输入都能进入} {(4.1)。}

\paragraph*{\texorpdfstring{{4.1} {四个} {\ensuremath{p}} {块与} {\ensuremath{s}} {个单点}}{4.1 四个  块与  个单点}}\label{C29-h185}
\addcontentsline{toc}{paragraph}{{4.1} {四个} {\ensuremath{p}} {块与} {\ensuremath{s}} {个单点}}

{固定一个单点} {\ensuremath{w}，保护它和每个} {\ensuremath{\ell |n}} {处的一个不同余见证，在其余位置平均出两个} {\ensuremath{p}} {块。容量足够：\ensuremath{d<p}} {由} {\ensuremath{2^{d}\le n<5p}、\ensuremath{p\ge 37}} {得到，而} {\ensuremath{n-(d+1)\ge 2p}。}

{用既有块--单点交换和第3节配比筛选，至多3d步使第一块} {\ensuremath{a}} {与} {\ensuremath{w}} {在全部} {\ensuremath{\ell |n}} {处不同余。避开} {\ensuremath{w}，再平均出第三、第四个} {\ensuremath{p}} {块，得到} \[
(a^p,b^p,c^p,d_0^p,x_1,\ldots,x_s).
\] {沿用三个主块的见证修复：先混合} {b,c，使} {c-a} {在原已好素数处非零，再交换} {\ensuremath{b}} {块与} {\ensuremath{w}。原坏素数处} {\ensuremath{a=b=c}、\ensuremath{w\ne a}，因此这一次交换修复全部原坏素数。现在前三块在每个} {\ensuremath{\ell |n}} {处不全相同。}

{这些论证只使用第3节的短配比，以及每个交换真实存在，不需要} {\ensuremath{s\ge d}。}

\paragraph*{\texorpdfstring{{4.2} {第一次三行折叠}}{4.2 第一次三行折叠}}\label{C29-h186}
\addcontentsline{toc}{paragraph}{{4.2} {第一次三行折叠}}

{令} {\ensuremath{p=3m+\epsilon }，\ensuremath{\epsilon }} {属于} {\{1,-1\}。先混合} {b,c，使} {b-a} {在所有前三块非全同余的素数处非零；再混合} {a,b。}

{第二次混合保持} {\ensuremath{S=a+b+c}} {不变，而} \[
mS+\epsilon a-pd_0
\] {对混合参数的斜率在所需素数处非零。因此第3节可使} \[
\frac{mS+\epsilon a}{p}-d_0
\] {在全部} {\ensuremath{\ell |n}} {处非零。随后平均两个组成均为} {(m+epsilon,m,m)} {的} {\ensuremath{p}} {组，再平均余下} {\ensuremath{p}} {份，得到} \[
(A^{2p},B^p,d_0^p,x_1,\ldots,x_s),
\qquad A-d_0\text{ 在模 }n\text{ 下为单位}.
\]

\paragraph*{\texorpdfstring{{4.3} {增大载体，再折叠一次}}{4.3 增大载体，再折叠一次}}\label{C29-h187}
\addcontentsline{toc}{paragraph}{{4.3} {增大载体，再折叠一次}}

{平均} {p-s} {份} {\ensuremath{A}} {和全部} {\ensuremath{s}} {个单点，留下} {p+s} {份旧} {\ensuremath{A}} {作载体。其余是三个} {\ensuremath{p}} {块；其中原} {\ensuremath{d_{0}}} {块与载体有完整见证，故操作安全。}

{相对载体} {\ensuremath{A}，三个} {\ensuremath{p}} {块的差之和在每个} {\ensuremath{\ell |n}} {处为零，这是全局零和与} {\ensuremath{n=0}} {\ensuremath{mod}} {\ensuremath{\ell }} {的直接结果。因此它们任意两个差不可能同时为零。先混合两个块取得完整差见证，再用上一段的} {(m+epsilon,m,m)} {三行折叠，得到} {(4.1)。}

{全过程至多20+3d次平均。载体在最后折叠中保持不动，其大小可以超过} {\ensuremath{p}；没有调用一个} {p+s} {维平均求解器。}

\subsection*{\texorpdfstring{{C29.5} {两次制造新载体的五步对合}}{C29.5 两次制造新载体的五步对合}}\label{C29-h188}
\addcontentsline{toc}{subsection}{{C29.5} {两次制造新载体的五步对合}}

{置} {\ensuremath{\epsilon =1}} {或} {-1，使} {\ensuremath{p=\epsilon }} {mod4，并令} \[
t_0=\frac{p-\epsilon}{4},\qquad
t_1=t_0+s.
\] {以下分别取} {\ensuremath{t=t_{0},t_{1}}，定义} \[
b=p-s-3t,\qquad
\beta=rt-pb=nt-p^2+ps.
\tag{C29.5.1}
\]

{先平均两个互不相交、组成相同的} {\ensuremath{p}} {组} \[
(A^{p-b-t},B^t,C^b).
\] {它们产生} {2p} {份} \[
D=\frac{s a+\beta z}{p}.
\tag{C29.5.2}
\] {这里使用} {\ensuremath{p-b-3t=s}。原} {\ensuremath{A}、\ensuremath{B}、\ensuremath{C}} {的剩余数量分别为} {2(b+t)、p-2t、r-2b。}

{从新} {\ensuremath{D}} {中保留} {\ensuremath{r=p+s}} {份作最终载体。对整数} {\ensuremath{h}，再平均两个组成相同的} {\ensuremath{p}} {组} \[
(A^x,B^h,C^u,D^t),\qquad
x=2h+t-b,\quad u=p-2t+b-3h.
\tag{C29.5.3}
\] {最后平均其余} {\ensuremath{p}} {份。总计五次原子平均。}

{可取全部整数} \[
\boxed{
\max\left(0,\left\lceil\frac{5p-5s-16t}{6}\right\rceil\right)
\le h\le b.}
\tag{C29.5.4}
\] {第2节的} {\ensuremath{p\ge 16s+21}} {保证此区间至少含两个相邻整数。}

{具体容量核对：两个初始重复组需要} {\ensuremath{2t\le p}、\ensuremath{2b\le r}、\ensuremath{p-b-t\ge 0}。最后重复组要求} \[
0\le x\le b+t,\quad
0\le2h\le p-2t,\quad
0\le2u\le r-2b,\quad
2t\le p-s.
\] {将} {\ensuremath{t=t_{0}+\delta }、\ensuremath{\delta =0}或\ensuremath{s}} {代入，有} {\ensuremath{b-t=\epsilon -s-4\delta \le 0}，且} {\ensuremath{b\le (p-2t)/2}，所以} {(5.4)} {保证以上全部条件。其未截断实端点差为} \[
\frac{p+\epsilon-2s-4\delta}{12}>2,
\] {而} {\ensuremath{b\ge 4}，截断下界为0不会破坏两个相邻整数的存在。}

{所得返回为} {\ensuremath{p^{-2}}} {\ensuremath{J_{t,h}}，其中} \[
J_{t,h}=
\begin{pmatrix}\beta&Q\\-s&-\beta\end{pmatrix},
\quad
Q=prh-p^2u+t\beta.
\] {直接得到} \[
\boxed{
J_{t,h}^2=p f_{t,h}I,\qquad
f_{t,h}=p^3+n[4t^2+(s-2p)t-sh].}
\tag{C29.5.5}
\] {\ensuremath{f=p^{3}}} {\ensuremath{mod}} {\ensuremath{n}，所以非零且与} {\ensuremath{n}} {互素。每个五步返回的射影逆都可通过重复它正向实现。}

{在} {(\ensuremath{z,y)=(z,a+pz)}} {基下，整数矩阵为} \[
\begin{pmatrix}
p^2-nt&-s\\
n[ph-t(p-t)]&-p^2+nt
\end{pmatrix}.
\tag{C29.5.6}
\] {故其实际模} {\ensuremath{n}} {作用为} {(\ensuremath{1,-s/p^{2}};0,-1)，保持合法域。}

\subsection*{\texorpdfstring{{C29.6} {四个五步返回统一产生深主同余控制}}{C29.6 四个五步返回统一产生深主同余控制}}\label{C29-h189}
\addcontentsline{toc}{subsection}{{C29.6} {四个五步返回统一产生深主同余控制}}

{旧} {\ensuremath{A}} {对合仍适用于} {\ensuremath{r=p+s}：} \[
A_j=\begin{pmatrix}0&-d_j\\-1&0\end{pmatrix},
\quad d_j=p^2-nj,
\quad
\left\lceil\frac{p-s}{6}\right\rceil\le j\le
\left\lfloor\frac{p-s}{4}\right\rfloor.
\] {因为} {\ensuremath{p-s\ge 36}，区间至少有三个相邻整数} {j,j+1,j+2。分别记其标量为} {\ensuremath{d_{0},d_{1},d_{2}}。}

{固定} {\ensuremath{t}，取} {(5.4)} {的两个相邻} {h,h+1，记} {\ensuremath{f=f_{t,h}}、\ensuremath{g_{t}=f_{t,h+1}}。有} \[
g_t=f-sn,\qquad \gcd(f,g_t)=1.
\tag{C29.6.1}
\] {因为} {\ensuremath{\operatorname{gcd}(f,n)=1}，且} \[
f=p(p-4t)^2\pmod s;
\] {\ensuremath{t=t_{0}}} {或} {\ensuremath{t_{0}+s}} {时} {p-4t} {分别为} {\ensuremath{\epsilon }、epsilon-4s，与} {\ensuremath{s}} {互素，而} {\ensuremath{\operatorname{gcd}(p,s)=1}。}

{相邻商为} \[
T_t=J_{t,h}^{-1}J_{t,h+1}
=\begin{pmatrix}1&\beta_t n/f\\0&g_t/f\end{pmatrix}.
\] {与} {\ensuremath{A_{j}^{-1}A_{j+1}=\operatorname{diag}(1,d_{1}/d_{0})}} {取交换子，得到} \[
U\left(\frac{\beta_t n^2}{d_1g_t}\right).
\tag{C29.6.2}
\] {两个} {\ensuremath{t}} {提供的} {\ensuremath{\beta }} {满足} \[
\beta_{t_1}-\beta_{t_0}=ns,\quad
\gcd(\beta_{t_0},s)=1,\quad
\boxed{\gcd(\beta_{t_0},\beta_{t_1})=\gcd(n,5)=:g.}
\tag{C29.6.3}
\] {最后一步用} {\ensuremath{\beta _{t}=}} {p(s-p)} {\ensuremath{mod}} {\ensuremath{n}、\ensuremath{\operatorname{gcd}(p,n)=1}。}

{令} {\ensuremath{R_{0}}} {局部化所有} {\ensuremath{d_{0},d_{1}}} {及两对} {\ensuremath{f,g_{t}}。相邻商的倍率和逆倍率生成这个环，所有反演素数均不除} {\ensuremath{n}。由} {(6.2)、整数} {Bezout} {和共轭，获得} \[
U(gn^2R_0),\qquad L(gn^2R_0).
\] {于是既有深同余引理给} \[
\boxed{\Gamma(M,R_0),\qquad M=g^2n^4.}
\tag{C29.6.4}
\] {\ensuremath{g}} {为1或5且整除} {\ensuremath{n}，所以控制层没有任何} {\ensuremath{n}} {以外的辅助素数。}

{再取} {\ensuremath{\eta }} {属于} {\{1,-1\}、\ensuremath{\eta =p}} {mod3，由安全返回取逆引理激活保载体} {\ensuremath{F_\eta }，局部化} {\ensuremath{p}；由第三个} {\ensuremath{A}} {的相邻商局部化} {\ensuremath{d_{2}}。这里} {\ensuremath{\eta }} {与第5节按模4选择的} {\ensuremath{\epsilon }} {分别定义。下面每次加入安全返回行列式，都通过真实上三角倍率扩张根参数环，再由深同余引理取得} {\ensuremath{\Gamma (M,R)}。不存在先假定任意返回可逆的循环。}

\subsection*{\texorpdfstring{{C29.7} {单位尺度用一次计数统一覆盖}}{C29.7 单位尺度用一次计数统一覆盖}}\label{C29-h190}
\addcontentsline{toc}{subsection}{{C29.7} {单位尺度用一次计数统一覆盖}}

{一般保留} {\ensuremath{A}} {的真实返回仍是} \[
M_{q,j}=
\begin{pmatrix}p-q&nj-pq\\-1&0\end{pmatrix},\quad q=3j+k.
\] {现在须显式保留} {\ensuremath{j+k\le p}，因为} {\ensuremath{r>p}。完整容量为} \[
0\le j\le\lfloor p/2\rfloor,\quad
0\le k\le\lfloor r/2\rfloor,\quad
\lceil r/2\rceil\le j+k\le p.
\]

{下列整个整数区间都可实现：} \[
I=[L,2p-1],\qquad
L=\begin{cases}
r/2+2,&r\text{ 偶},\\
(r+9)/2,&r\text{ 奇}.
\end{cases}
\tag{C29.7.1}
\] {证明：固定} {\ensuremath{j}} {时，可实现的} {\ensuremath{q}} {区间为} \[
[2j+\lceil r/2\rceil,\ \min(3j+\lfloor r/2\rfloor,2j+p)].
\] {\ensuremath{r}} {偶从} {\ensuremath{j=1}} {开始，\ensuremath{r}} {奇从} {\ensuremath{j=2}} {开始，连续区间相邻或交叠，最终上界为2p-1。}

{每个} {\ensuremath{\operatorname{gcd}(q,n)=1}} {的返回行列式} {\ensuremath{D=nj-pq}} {都与} {\ensuremath{M}} {互素，故正向逆已可执行。其与旧} {\ensuremath{A}} {的商将} {\ensuremath{D}} {局部化，而} {-D/p} {是} {\ensuremath{q}} {\ensuremath{mod}} {\ensuremath{n}} {的实际} {\ensuremath{R}} {单位提升。只须加入有限多个} {\ensuremath{q}} {属于} {\ensuremath{I}。}

{记} {\ensuremath{C}} {为} {\ensuremath{I}} {中模} {\ensuremath{n}} {单位的个数。由容斥，} \[
C>|I|\rho-(2^d-1),\quad
|I|\ge\frac{3n-7s-36}{8}.
\] {利用} {(2.2)，得到} \[
\boxed{C>\frac{\varphi(n)}4.}
\tag{C29.7.2}
\] {又} {\ensuremath{I}} {完全位于} {(0,n/2)，所以} {\ensuremath{I}} {中单位及其负数组成集合} {\ensuremath{S}，大小大于} {phi(n)/2。对任意模} {\ensuremath{n}} {单位} {\ensuremath{u}，\ensuremath{S}} {与} {uS} {必相交，因而} {\ensuremath{u}} {是} {\ensuremath{S}} {中两个元素之比。所有这些元素均已有} {\ensuremath{R}} {单位提升，故} \[
\boxed{R^\times\longrightarrow(\mathbb Z/n\mathbb Z)^\times
\text{ 满射}.}
\tag{C29.7.3}
\]

{这避免了寻找模} {\ensuremath{n}} {原根，也不需要逐类反演} {\ensuremath{r}} {或构造二三单位的小代表。尾部本身的多素因子条件使容斥误差小于主项所留余量。}

\subsection*{\texorpdfstring{{C29.8} {所有局部方向，包括模5}}{C29.8 所有局部方向，包括模5}}\label{C29-h191}
\addcontentsline{toc}{subsection}{{C29.8} {所有局部方向，包括模5}}

{取某个} {\ensuremath{J_{t,h}}} {与} {\ensuremath{A_{j}}} {的积} {\ensuremath{V=J_{t,h}A_{j}/p^{3}}。在} {(z,y)} {基下，} \[
V=\begin{pmatrix}1&-5/p\\0&1\end{pmatrix}\pmod n.
\tag{C29.8.1}
\] {这在全部} {\ensuremath{\ell |n}、\ensuremath{\ell \ne 5}} {处是非平凡上平移。}

{其迹为} \[
\operatorname{tr}V=
2-\frac n{p^3}[pt+t^2+ph+sj].
\tag{C29.8.2}
\] {\ensuremath{n=0}} {mod4} {时迹自动为2} {mod4；\ensuremath{n=2}} {mod4} {时} {\ensuremath{s}} {偶，\ensuremath{pt+t^{2}}} {与} {sj} {均偶，选相邻} {\ensuremath{h}} {中的偶数即可。因此二进完整循环成立。}

{若} {\ensuremath{3|n}，将} {\ensuremath{V}} {除以迹半数后写为} {I+N、\ensuremath{N^{2}=cI}。\ensuremath{9|n}} {时} {\ensuremath{c=0}} {mod9；\ensuremath{v_{3}(n)=1}} {时} \[
\frac c3=
\frac{n/3}{p^3}
\bigl[(p-s)(j+t-h)-5t^2\bigr]\pmod3.
\tag{C29.8.3}
\] {因为} {\ensuremath{p-s=2p}} {mod3} {为单位，从相邻} {j,j+1} {中选一个便可避开} {\ensuremath{c=6}} {mod9。这不改变二进条件。所有} {\ensuremath{\ell \ge 7}} {直接适用已证局部射影循环引理。}

{若} {\ensuremath{5|n}，先选一个} {\ensuremath{q}} {属于} {\ensuremath{I}，使} \[
q=-p\pmod5,\qquad\gcd(q,n)=1.
\tag{C29.8.4}
\] {这样的} {\ensuremath{q}} {确实存在：固定非零模5类后，容斥主项为} {\textbar{}\ensuremath{I|\rho /4}，误差小于\ensuremath{2^{d-1}}；(2.2)} {还给} {\textbar{}\ensuremath{I|\rho >2^{d+1}}，所以剩余数严格为正。}

{该安全} {\ensuremath{A}} {返回在模5下为} {(1,-1/p;0,1)，在每个5幂上具有完整射影循环。先用其非负次幂使} {\ensuremath{z/y=0}} {\ensuremath{mod}} {\ensuremath{5^{v_{5}(M)}}。\ensuremath{V}} {模5为} {\ensuremath{I}，其在模5幂下的完整矩阵阶为5幂。随后令} {\ensuremath{V}} {的指数被该阶整除，再用} {CRT} {同时满足所有其他素数处的完整循环指数。}

{于是任意合法} {\ensuremath{R}} {本原列均能到达} \[
z=0\pmod M,\qquad y\in(R/MR)^\times.
\tag{C29.8.5}
\] {这一步保留完整5进矩阵阶，不以局部形式逆代替真实路径。}

\subsection*{\texorpdfstring{{C29.9} {从同余目标到精确零终端}}{C29.9 从同余目标到精确零终端}}\label{C29-h192}
\addcontentsline{toc}{subsection}{{C29.9} {从同余目标到精确零终端}}

{两个已知实际} {\ensuremath{R}} {单位} \[
\kappa_1=d_1/d_0=1-n/d_0,\qquad
\kappa_2=d_2/d_0=1-2n/d_0
\] {通过} {LTE} {和} {CRT} {生成模} {\ensuremath{M}} {到模} {\ensuremath{n}} {的全部主单位核。\ensuremath{n=2}} {mod4} {时，用} {\ensuremath{\kappa _{2}}} {生成等于1模2n的部分，再以} {\ensuremath{\kappa _{1}=3}} {mod4} {补另一个类；其余情形} {\ensuremath{\kappa _{1}}} {即可。}

{结合} {(7.3)，\ensuremath{R}} {单位模} {\ensuremath{M}} {满射。全部所用规范返回的行列式在最终} {\ensuremath{R}} {中已可逆，所以它们保持} {\ensuremath{R}} {本原性。选} {\ensuremath{\epsilon }} {属于} {\ensuremath{R^\times}} {匹配} {(8.5)} {的} {\ensuremath{y}，由已证主同余纤维引理，将状态}\textbf{{精确}}{送到} {(\ensuremath{z,y)=(0,\epsilon )}。}

{最终核心为} \[
(\epsilon^{2p},(-2\epsilon)^p,0^{p+s}).
\] {选取全部3p个非零位置和\ensuremath{p}个零，调用旧不等式已经独立证明的4p维判据；剩余\ensuremath{s}个零不动。这个子块自身零和且含零，\ensuremath{G=1}，因此调用合法。这也说明收尾没有循环使用本文的上方区间定理。结论承接第4节任意原输入入口，因此不是仅有核心群或有限模像。}

{必要性沿用} {\ensuremath{G}} {不变量。旧定理覆盖的} {\ensuremath{n}} {与本文补齐的} {\ensuremath{n}} {合并，即得定理} {\ensuremath{A}；再接} {3p+10} {至4p-1} {的已有完整区间，得定理} {\ensuremath{B}。}


\section{五平均的十六至十八元基例}\label{construct:five-small}
\subsection*{\texorpdfstring{{C30.3} {固定整数格上的旧不变量}}{C30.3 固定整数格上的旧不变量}}\label{C30-h193}
\addcontentsline{toc}{subsection}{{C30.3} {固定整数格上的旧不变量}}

{对} {\ensuremath{n=16}、17、18，令} {\ensuremath{Q}} {为} {\ensuremath{n}} {的非五素因子集合，分别为}

\[
\{2\},\qquad\{17\},\qquad\{2,3\}.
\]

{称整数零和状态合法，是指在每个} {\ensuremath{q}} {in} {\ensuremath{Q}} {下坐标不全同余。} {本原输入的合法性与} {\ensuremath{G}} {为五幂等价；此处三个} {\ensuremath{n}} {都不被} {5} {整除，因此本原} {\ensuremath{G}} {实际为} {1。}

{从} {5X} {出发，已有初始化在} {\ensuremath{n\ge 15+|Q|}} {时，至多两次安全整数平均建立}

\[
\mathcal I:\quad\text{至少两个不同的五重值，且保持合法整数零和。}
\]

{三个} {\ensuremath{n}} {都满足该容量条件。} {以后只需给每个非终端} {\ensuremath{I}} {状态一个安全非恒等整数平均，继续保持} {\ensuremath{I}。} {平方能量严格下降，便保证有限终止；缩回} {\ensuremath{X}} {后所有值属于} {(1/5)Z。}

\subsection*{\texorpdfstring{{C30.4} {已有分支与十六元三重值边界}}{C30.4 已有分支与十六元三重值边界}}\label{C30-h194}
\addcontentsline{toc}{subsection}{{C30.4} {已有分支与十六元三重值边界}}

{沿用已证的顺序、背景值与共同预留引理：}

\begin{itemize}
\tightlist
\item
  {至少四个重值；}
\item
  {三个重值，但存在低重数值、某个极端多于五份，或两个重值同余} {\ensuremath{mod}} {5；}
\item
  {一个重值至少九份；}
\item
  {两个重值同余} {\ensuremath{mod}} {5；}
\item
  {两个不同余重值总重数至少十三。}
\end{itemize}

{它们都有安全整数闭包或明确的相反值终端。} {\ensuremath{n=17}、18} {剩余三重值情形由此前的两份中间值修复处理。}

{只需补充} {\ensuremath{n=16}} {的三个实际值形状}

\[
(a^5,b^6,c^5),\quad a<b<c,
\quad a,b,c\text{ 模 5 两两不同}.
\]

{取唯一整数调节} {(\ensuremath{a^{i},c^{4-i},b)}。\ensuremath{b}} {还剩五份，均值不等于} {\ensuremath{b}} {时已经完成。} {若均值等于} {\ensuremath{b}，则} {\ensuremath{i=1}、2、3：}

\begin{itemize}
\tightlist
\item
  {\ensuremath{i=1}} {时，\ensuremath{a+3c=4b}，零和给} {\ensuremath{13b-5c=0}，所以} {\ensuremath{b=c}} {\ensuremath{mod}} {2，进而} {\ensuremath{a=c}} {\ensuremath{mod}} {2，违反合法性。}
\item
  {\ensuremath{i=3}} {对称。}
\item
  {\ensuremath{i=2}} {时，\ensuremath{a+c=2b}，零和给} {\ensuremath{16b=0}，故} {\ensuremath{b=0}、\ensuremath{a=-c}，是可逐对消去的终端。}
\end{itemize}

{所以三重值分支也已闭合。}

\subsection*{\texorpdfstring{{C30.5} {最小保护集合与总和障碍}}{C30.5 最小保护集合与总和障碍}}\label{C30-h195}
\addcontentsline{toc}{subsection}{{C30.5} {最小保护集合与总和障碍}}

{以下只剩两个不同余重值} {\ensuremath{a}、\ensuremath{b}。设}

\[
f(a)=5+\alpha,\quad f(b)=5+\beta,
\quad\alpha\ge\beta\ge0,\quad\alpha+\beta\le2.
\tag{C30.2}
\]

{因此} {(alpha,beta)} {仅为} {(0,0)、(1,0)、(1,1)、(2,0)。}

{对每个} {\ensuremath{q}} {in} {\ensuremath{Q}，若某余数类大小至少} {n-5，它是唯一的危险主类；其例外支持} {\ensuremath{F_{q}}} {有} {2} {至} {5} {个位置。} {只有当} {\ensuremath{a}、\ensuremath{b}} {都属于主类时，才需要从低重数位置中额外保护见证；否则始终保留的锚点副本已作见证。} {称前一种} {\ensuremath{F_{q}}} {为需要保护的支持集。}

{先每个支持集保护一个位置，随后删除所有冗余保护，得到包含意义下最小的横截集} {\ensuremath{T}，\textbar{}\ensuremath{T|=h\le |Q|}。} {对每个} {\ensuremath{v}} {in} {\ensuremath{T}，必存在} {\ensuremath{q_{v}}} {使}

\[
F_{q_v}\cap T=\{v\}.
\tag{C30.3}
\]

{否则} {\ensuremath{v}} {可以删掉，违反最小性。} {同时} {\ensuremath{a=b}} {\ensuremath{mod}} {\ensuremath{q_{v}}，\ensuremath{v}} {不同于} {\ensuremath{b}} {\ensuremath{mod}} {\ensuremath{q_{v}}，\ensuremath{T}} {中其他位置都等于} {\ensuremath{b}} {\ensuremath{mod}} {\ensuremath{q_{v}}。}

\textbf{{保护总和引理。}} {假设} {\ensuremath{h>0}，某个} {\ensuremath{v}} {in} {\ensuremath{T}} {的私人素数} {\ensuremath{q_{v}}} {满足}

\[
\sum_{i\notin T}(x_i-b)=0\pmod{q_v}.
\tag{C30.4}
\]

{则状态不可能零和。}

{证明：因} {\ensuremath{q_{v}|n}，零和要求} {\ensuremath{sum_{i}(x_{i}-b)=0}} {\ensuremath{mod}} {\ensuremath{q_{v}}。} {\ensuremath{T}} {内除} {\ensuremath{v}} {外的差全为零，\ensuremath{v}} {的差非零；再加上} {(4)，得到矛盾。}

{注意} {(4)} {只要求}\textbf{{总和}}{模} {\ensuremath{q_{v}}} {为零，不要求每个未保护位置都同余。} {这比上一轮设想的逐值仿射格刚性更弱，也更容易用线性证书证明。}

\subsection*{\texorpdfstring{{C30.6} {全部碰撞的方程}}{C30.6 全部碰撞的方程}}\label{C30-h196}
\addcontentsline{toc}{subsection}{{C30.6} {全部碰撞的方程}}

{令} {\ensuremath{\delta =a-b}，它在模} {5} {下可逆。对每个未保护的低重数位置} {\ensuremath{x_{i}}，定义}

\[
t_i=\frac{x_i-b}{\delta}\in\mathbb Q,
\qquad z_i=t_i\pmod5.
\]

{这里分母} {\ensuremath{\delta }} {不含} {5，因此} {\ensuremath{z_{i}}} {有定义。} {两组重值仍完整包含在未保护部分中。}

{先执行任何可用的旧锚点同余修复。若仍未完成：}

\begin{itemize}
\tightlist
\item
  {\ensuremath{z_{i}=0}} {强迫} {\ensuremath{t_{i}=5}；这样的低重数位置至多一个。}
\item
  {\ensuremath{z_{i}=1}} {强迫} {\ensuremath{t_{i}=-4}；这样的低重数位置至多一个。}
\end{itemize}

{例如} {\ensuremath{z_{i}=0}} {时，(\ensuremath{b^{4},x_{i})}} {的均值若不等于} {\ensuremath{a}} {就已经成功；碰撞给} {\ensuremath{x_{i}=5a-4b}。} {如果这种值至少两份，(\ensuremath{b^{3},x_{i}^{2})}} {的均值为} {2a-b，安全且不碰撞。}

{未保护池中某个模} {5} {余数若有五份，直接平均它们即可保留} {\ensuremath{a}、\ensuremath{b}} {两个重值。} {该操作非恒等，因为低重数的每个实际值不足五份。} {因此只需所有余数重数至多四。}

{现在假设所有下面的非恒等整数候选都不保持双重值。} {对任意非空低重数位置子集} {\ensuremath{S}，记} {\ensuremath{k=|S|\le 5}、r=sum\_\{i} {in} {S\}z\_i} {\ensuremath{mod}} {5。}

{若} {\ensuremath{i=-r}} {\ensuremath{mod}} {5} {的} {{[}0,4{]}} {代表满足} {\ensuremath{i\le min(\alpha ,5-k)}，可选} {\ensuremath{i}} {份} {\ensuremath{a}、\ensuremath{S}，以及} {5-k-i} {份} {\ensuremath{b}。} {它保留五份} {\ensuremath{a}，因此失败必须满足}

\[
\sum_{j\in S}t_j=5-i.
\tag{C30.5}
\]

{如果同时} {\ensuremath{5-k-i\le \beta }，则两个旧重值都保留，该余数模式直接不可能失败。}

{从另一锚点出发，若} {\ensuremath{j=r-k}} {\ensuremath{mod}} {5} {满足} {\ensuremath{j\le min(\beta ,5-k)}，失败必须满足}

\[
\sum_{i\in S}t_i=k+j-5.
\tag{C30.6}
\]

{所有这些候选都是在原位置上的真实五平均，且保留锚点副本和} {\ensuremath{T}，所以安全。} {这些方程中的负号与非整数系数属于分析，不是新增的平均操作。}

\subsection*{\texorpdfstring{{C30.7} {有限线性证书只需给出的总和结论}}{C30.7 有限线性证书只需给出的总和结论}}\label{C30-h197}
\addcontentsline{toc}{subsection}{{C30.7} {有限线性证书只需给出的总和结论}}

{设}

\[
c=n-10-h,\qquad L=c-\alpha-\beta.
\]

{\ensuremath{L}} {是未保护低重数位置数，\ensuremath{c}} {属于} {\{5,6,7,8\}。} {将} {\ensuremath{z_{1}\le }\ldots\textless=\ensuremath{z_{L}}} {排序，仅消除位置置换的重复；每个实际} {\ensuremath{t_{i}}} {仍是独立未知数，没有假设同余值实际相等。}

{在第} {6} {节的余数重数限制下，四种} {(alpha,beta)} {与四个} {\ensuremath{c}} {共给出} {903} {个无序模式。} {逐个建立} {(5)、(6)} {及} {\ensuremath{z=0}、1} {的方程，精确有理消元得到：}

\textbf{{证书引理。}}

\begin{enumerate}
\def\labelenumi{\arabic{enumi}.}
\tightlist
\item
  {\ensuremath{c=7}、8} {时，全部碰撞方程均不相容。}
\item
  {\ensuremath{c=6}} {时，只有两个相容模式，且都满足} \[
  D:=5+\alpha+\sum_i t_i=15.
  \tag{C30.7}
  \]
\item
  {\ensuremath{c=5}} {时，只有二十个相容模式，它们的} {\ensuremath{D}} {的最简分母都为奇数。}
\end{enumerate}

{每个相容模式的有理解唯一；这比本证明所需的总和结论更强。} {下面列出全部相容模式，便于人工审阅；重复上标只表示余数重数。}

{ % do not increment counter
\begin{longtable}[]{@{}rllr@{}}
\toprule\noalign{}
{\ensuremath{c}} & {alpha,beta} & {余数多重集} & {\ensuremath{D}} \\
\midrule\noalign{}
\endhead
\bottomrule\noalign{}
\endlastfoot
{5} & {0,0} & {\ensuremath{0,2^{4}}} & {-22} \\
{5} & {0,0} & {\ensuremath{0,3^{4}}} & {2} \\
{5} & {0,0} & {\ensuremath{0,4^{4}}} & {-4} \\
{5} & {0,0} & {\ensuremath{1,2^{4}}} & {19} \\
{5} & {0,0} & {\ensuremath{1,3^{4}}} & {13} \\
{5} & {0,0} & {\ensuremath{1,4^{4}}} & {37} \\
{5} & {0,0} & {\ensuremath{2^{4},3}} & {1} \\
{5} & {0,0} & {\ensuremath{2^{4},4}} & {17} \\
{5} & {0,0} & {\ensuremath{2,3^{4}}} & {-11} \\
{5} & {0,0} & {\ensuremath{2,4^{4}}} & {-2} \\
{5} & {0,0} & {\ensuremath{3^{4},4}} & {26} \\
{5} & {0,0} & {\ensuremath{3,4^{4}}} & {14} \\
{5} & {1,0} & {\ensuremath{2^{4}}} & {14} \\
{5} & {1,0} & {\ensuremath{3^{4}}} & {34/3} \\
{5} & {1,0} & {\ensuremath{3,4^{3}}} & {11} \\
{5} & {1,0} & {\ensuremath{4^{4}}} & {22} \\
{5} & {1,1} & {\ensuremath{0,3^{2}}} & {-3} \\
{5} & {1,1} & {\ensuremath{1,3^{2}}} & {18} \\
{5} & {2,0} & {\ensuremath{2^{3}}} & {13} \\
{5} & {2,0} & {\ensuremath{2^{2},3}} & {14} \\
{6} & {1,0} & {\ensuremath{3,4^{4}}} & {15} \\
{6} & {2,0} & {\ensuremath{2^{4}}} & {15} \\
\end{longtable}
}

{\ensuremath{c=6}} {两行的实际归一化低重数值分别为} {(-7,4,4,4,4)} {与} {(2,2,2,2)。} {其他二十行的完整唯一解固定保存在文末完整有限核验程序的} {\ensuremath{FIVE_{T}ABLE}} {中，与全部系统消元结果逐项相等。}

{为什么这承担全部输入而不只是有限样本：任意有理输入必有上述某个模} {5} {模式；假设全部候选碰撞就必满足该模式的有理线性方程。} {消元决定的是这些方程的}\textbf{{整个解集}}{，与输入整数的高度、分母或操作长度无关。}

\subsection*{\texorpdfstring{{C30.8} {最小保护与总和证书排除全部例外}}{C30.8 最小保护与总和证书排除全部例外}}\label{C30-h198}
\addcontentsline{toc}{subsection}{{C30.8} {最小保护与总和证书排除全部例外}}

{未保护部分相对} {\ensuremath{b}} {的总和为}

\[
\sum_{i\notin T}(x_i-b)=\delta D.
\tag{C30.8}
\]

\paragraph*{\texorpdfstring{{\ensuremath{n=17}}}{}}\label{C30-h199}
\addcontentsline{toc}{paragraph}{{\ensuremath{n=17}}}

{\ensuremath{h=0}} {时} {\ensuremath{c=7}，没有相容模式。} {\ensuremath{h=1}} {时} {\ensuremath{c=6}，\ensuremath{D=15}。保护位置的私人素数是} {17，且} {\ensuremath{17|\delta }，所以} {(8)} {模} {17} {为零，违反保护总和引理。}

\paragraph*{\texorpdfstring{{\ensuremath{n=18}}}{}}\label{C30-h200}
\addcontentsline{toc}{paragraph}{{\ensuremath{n=18}}}

{\ensuremath{h=0}、1} {时} {\ensuremath{c=8}、7，没有相容模式。} {\ensuremath{h=2}} {时} {\ensuremath{c=6}，\ensuremath{D=15}。任取一个保护位置，其私人素数为} {2} {或} {3，均整除} {\ensuremath{\delta }，因此} {(8)} {模该素数为零，再次矛盾。} {这里另一保护位置在该私人素数下属于主类，正是包含最小性保证的关键。}

\paragraph*{\texorpdfstring{{\ensuremath{n=16}}}{}}\label{C30-h201}
\addcontentsline{toc}{paragraph}{{\ensuremath{n=16}}}

{\ensuremath{h=0}} {时} {\ensuremath{c=6}，\ensuremath{D=15}，由零和得}

\[
16b+15\delta=0.
\]

{因此} {\ensuremath{\delta }} {为偶数；\ensuremath{c=6}} {两个模式的} {\ensuremath{t_{i}}} {都是整数，全部实际值均与} {\ensuremath{b}} {同奇偶，违反合法性。}

{\ensuremath{h=1}} {时} {\ensuremath{c=5}，私人素数为} {2，\ensuremath{\delta }} {为偶数。} {证书保证} {\ensuremath{D}} {的分母为奇数，所以} {\ensuremath{\delta }} {\ensuremath{D}} {在} {2} {处为零。} {由保护总和引理矛盾。}

{注意少数} {\ensuremath{t_{i}}} {可以有分母} {2，但它们相加后的} {\ensuremath{D}} {没有分母} {2。} {这正是只证明总和比逐值仿射格更有效的地方。}

{所有情形均矛盾，所以至少有一个安全非恒等整数候选保持双重值。} {结合第} {3、4} {节及平方能量下降，此处三个维数全部解决。} {再接} {\ensuremath{n=15}} {与} {\ensuremath{n\ge 19}，得到主定理} {(1)。}


\section{七平均的二十二元基例}\label{construct:seven-small}
\subsection*{\texorpdfstring{{C31.2} {原位置约定与正向逆的取得}}{C31.2 原位置约定与正向逆的取得}}\label{C31-h202}
\addcontentsline{toc}{subsection}{{C31.2} {原位置约定与正向逆的取得}}

{继续使用统一核心} \[
\mathcal K_m(a,z)=
\bigl(a^{7m},(-ma+rz)^7,(-7z)^r\bigr),
\qquad n=7(m+1)+r.
\] {二十元为} {\ensuremath{m=1,r=6}；二十二元为} {\ensuremath{m=2,r=1}。令} {\ensuremath{y=a+7z}，则本原核心合法性均为} \[
\gcd(a,z)=\gcd(y,n)=1.
\tag{C31.2}
\]

{记三块值为} {\ensuremath{A}、\ensuremath{B}、\ensuremath{C}。宏} {A(j,s)} {的含义是：保留} {\ensuremath{r}} {份旧} {\ensuremath{A}；另做} {\ensuremath{m}} {份组成完全相同的七元组，每组取} {7-j-s} {份} {\ensuremath{A}、\ensuremath{j}} {份} {\ensuremath{B}、\ensuremath{s}} {份} {\ensuremath{C}；最后平均其余七个位置。宏} {B(j,s)} {同理保留旧} {\ensuremath{B}，宏} {C(j,0)} {保留旧} {\ensuremath{C}。第3、5节表中每个计数都满足原位置容量。}

{以下表列整数矩阵，实际宏在} {(a,z)} {上的作用为矩阵乘} {1/7。\ensuremath{m=1}} {时每宏有两个原子平均；\ensuremath{m=2}} {时每宏有三个。角色交换} {\ensuremath{S}} {不做数值操作，只交换两个等大的七块标签。}

{若正向宏组成一个闭路} \[
M_\ell\cdots M_1=cI,\qquad c\ne0,
\tag{C31.3}
\] {则闭路中每个宏都有正向射影逆：把闭路循环移位，使该宏最先执行，再执行剩余宏。循环移位后的乘积仍为} {cI；实际尺度则另乘每个宏自身的} {1/7。此论证对负标量也成立，不能把负的全局尺度与不可执行的逐坐标取负混淆。}

\subsection*{\texorpdfstring{{C31.5} {二十二元：原来只有一个对合的边界也能取逆}}{C31.5 二十二元：原来只有一个对合的边界也能取逆}}\label{C31-h203}
\addcontentsline{toc}{subsection}{{C31.5} {二十二元：原来只有一个对合的边界也能取逆}}

{二十二元核心为} \[
(a^{14},(-2a+z)^7,-7z).
\] {采用三个宏：}

{ % do not increment counter
\begin{longtable}[]{@{}lll@{}}
\toprule\noalign{}
{名称} & {原位置类型} & {整数矩阵} \\
\midrule\noalign{}
\endhead
\bottomrule\noalign{}
\endlastfoot
{\ensuremath{A}} & {A(1,0)} & \(\begin{pmatrix}4&1\\-1&0\end{pmatrix}\) \\
{\ensuremath{F}} & {C(2,0)} & \(\begin{pmatrix}1&2\\0&7\end{pmatrix}\) \\
{\ensuremath{J}} & {B(2,0)} & \(\begin{pmatrix}1&2\\2&-1\end{pmatrix}\) \\
\end{longtable}
}

{宏} {\ensuremath{A}} {保留一份旧} {\ensuremath{A}，两重复组各为} {(\ensuremath{A^{6},B)}，最后为} {(\ensuremath{A,B^{5},C)}。宏} {\ensuremath{F}} {保留} {\ensuremath{C}，两重复组各为} {(\ensuremath{A^{5},B^{2})}，最后为} {(\ensuremath{A^{4},B^{3})}。宏} {\ensuremath{J}} {保留一份旧} {\ensuremath{B}，两重复组各为} {(\ensuremath{A^{5},B^{2})}，最后为} {(\ensuremath{A^{4},B^{2},C)}。均恰好使用原22个位置。}

{直接有} \[
FA^2=\begin{pmatrix}7&2\\-28&-7\end{pmatrix},
\qquad
\boxed{(FA^2)^2=-7I,\qquad J^2=5I.}
\tag{C31.8}
\] {因此} {\ensuremath{A}、\ensuremath{F}、\ensuremath{J}} {及其逆都能正向执行。实际} {A,A,F,A,A,F} {共18次七平均，尺度为-\ensuremath{7/7^{6}}。}

{这一步绕过``\ensuremath{r=1}只存在一个迹零对合''的旧容量限制：\ensuremath{A}、\ensuremath{F}} {本身不必迹零，联合闭路提供正向逆。}

\paragraph*{\texorpdfstring{{根与深主同余}}{根与深主同余}}\label{C31-h204}
\addcontentsline{toc}{paragraph}{{根与深主同余}}

{令} {\ensuremath{H=AFJA^{-1}}，则} \[
H=\begin{pmatrix}-7&-62\\0&5\end{pmatrix},
\qquad
\boxed{FHF^{-1}H^{-1}=U(396/35).}
\tag{C31.9}
\] {\ensuremath{F}、\ensuremath{H}} {对上根参数的共轭倍率分别为1/7、-7/5；连同逆共轭，生成} {\ensuremath{R_{22}=Z}{[}1/35{]}，故有} {\ensuremath{U(396R_{22})}。}

{取} \[
C=\begin{pmatrix}1&4\\0&-1\end{pmatrix}=C^{-1}.
\] {对任意} {\ensuremath{t}，} \[
CU(t)C=U(-t),\qquad
CAU(t)A^{-1}C=L(t).
\] {所以在这个整数幺模坐标中已有双向396倍根。深同余引理给出} \[
\Gamma(396^2,R_{22})\subseteq\Gamma^+.
\tag{C31.10}
\] {由于} {\ensuremath{C}} {整幺模，结论换回原坐标仍成立。}

\paragraph*{\texorpdfstring{{删去辅助的三进深层}}{删去辅助的三进深层}}\label{C31-h205}
\addcontentsline{toc}{paragraph}{{删去辅助的三进深层}}

{\ensuremath{396^{2}=1936}·81。要避免在辅助素数3处无谓扩大有限商，下面严格证明} \[
\boxed{\Gamma(1936,R_{22})\subseteq\Gamma^+.}
\tag{C31.11}
\]

{置} {\ensuremath{X=F^{2}/7}，它是} {\ensuremath{\mathrm{SL}_{2}(R_{22})}} {中已有正向实现的规范代表。整数模计算给} \[
A^{440}=X^{440}=I\pmod {1936}.
\] {将这两个矩阵约化模9，所得生成元生成全部} {\ensuremath{\mathrm{SL}_{2}(Z/9Z)}，共648个元素。固定证书保存全部父边，独立检查每个矩阵行列式为1；群阶公式} \[
|\mathrm{SL}_2(\mathbb Z/9\mathbb Z)|=3^3|\mathrm{SL}_2(\mathbb F_3)|=648
\] {证明这是完整群，而非一个抽样子集。}

{从模9满射，可取三个真实词在模9下分别等于} {I+3E12、I+3E21、I+3diag(1,-1)。这些词都由} {\ensuremath{A^{440}}、\ensuremath{X^{440}}} {组成，因此在模1936下恒为} {\ensuremath{I}。对任意这样的词} {\ensuremath{W=I+3T}，奇素数3的二项式展开给} \[
W^{3^k}=I+3^{k+1}T\pmod {3^{k+2}}.
\] {三方向逐层补齐全部高阶核，所以这个模1936恒等的子群在模81上仍满射到} {\ensuremath{\mathrm{SL}_{2}(Z/81Z)}。}

{给定} {\ensuremath{g}} {属于} {\ensuremath{\Gamma (1936,R_{22})}，选该子群中的} {\ensuremath{w}} {使} {\ensuremath{w=g}} {模81。两者模1936均为} {\ensuremath{I}，故} {\ensuremath{w^{-1}g}} {属于} {\ensuremath{\Gamma (1936}·\ensuremath{81,R_{22})}，已由} {(10)} {可执行。这证明} {(11)，没有从一个小模像直接推断无限群包含。}

\subsection*{\texorpdfstring{{C31.6} {二十二元必须保留共同尺度}}{C31.6 二十二元必须保留共同尺度}}\label{C31-h206}
\addcontentsline{toc}{subsection}{{C31.6} {二十二元必须保留共同尺度}}

{现在取} {\ensuremath{M=1936=2^{4}}·\ensuremath{11^{2}}。\ensuremath{R_{22}}} {单位} {+-\ensuremath{5^{i}7^{j}}} {的模} {\ensuremath{M}} {像只有440个，而全部模} {\ensuremath{M}} {单位共有880个；这一次单位像不满。}

{固定证书给} \[
U=\overline{R_{22}^{\times}},\qquad
(\mathbb Z/M\mathbb Z)^\times=U\ \sqcup\ 3U.
\tag{C31.12}
\] {对合法本原列，\ensuremath{y=a+7z}} {是模} {\ensuremath{M}} {单位。正确的有限状态因此为} \[
\boxed{(x,\epsilon),\quad x=z/y\pmod M,\quad
\epsilon=
\begin{cases}
0,&y\in U,\\
1,&y\in3U.
\end{cases}}
\tag{C31.13}
\] {共1936·\ensuremath{2=3872}个。取本原整数代表时可能除去的5、7因子及共同符号属于} {\ensuremath{R_{22}}} {单位，所以} {(13)} {与代表选择相容。}

{在} {(z,y)} {坐标中，三个宏为} \[
A=\begin{pmatrix}7&-1\\22&-3\end{pmatrix},\quad
F=\begin{pmatrix}7&0\\44&1\end{pmatrix},\quad
J=\begin{pmatrix}-15&2\\-110&15\end{pmatrix}.
\tag{C31.14}
\] {若矩阵写成} {(}\ensuremath{\alpha}{,}\ensuremath{\beta}{;}\ensuremath{\gamma}{,}\ensuremath{\delta}{)，其作用必须同时更新方向和尺度：} \[
x'=(\alpha x+\beta)(\gamma x+\delta)^{-1}\pmod M,
\] \[
\epsilon'=\epsilon+\chi(\gamma x+\delta)\pmod2,
\tag{C31.15}
\] {其中} {chi} {是单位群对指数2子群} {\ensuremath{U}} {的商映射。所有分母均为模} {\ensuremath{M}} {单位。}

{完整证书用} {\ensuremath{A}、\ensuremath{A^{-1}}、\ensuremath{F}、\ensuremath{F^{-1}}、\ensuremath{J}} {的五条边，证明全部3872个状态均在} {(0,0)} {的单一轨道；最大父边深度为8。证书不只检查投影到} {\ensuremath{x}} {的1936个方向，也没有假设不同素数的局部结果自动拼接。单位类和方向在整个模} {\ensuremath{M}} {上一起验证。}

{因此任意合法核心沿父边回到} {(0,0)} {后，本原代表满足} {\ensuremath{z}'=0} {模} {\ensuremath{M}，\ensuremath{a}'属于} {\ensuremath{U}。由} {(12)} {可选择} {\ensuremath{s}} {属于} {\ensuremath{R_{22}}} {单位，使} {\ensuremath{s=a}'} {模} {\ensuremath{M}。用} {(11)} {和同余本原列的补列引理，精确送到} {(s,0)。}

{精确终端为} \[
(s^{14},(-2s)^7,0).
\] {选四份} {\ensuremath{s}、两份} {-2s} {和一个零，真实七平均产生七个零，再接已证零触发尾部。}

{所以全部合法二十二元核心可达。}第\ref{construct:short-choice}节{已覆盖其全部原输入，因此二十二元完整} {\ensuremath{G=1}} {判据成立。}

\subsection*{\texorpdfstring{{C31.7} {有限证书怎样承担全部输入}}{C31.7 有限证书怎样承担全部输入}}\label{C31-h207}
\addcontentsline{toc}{subsection}{{C31.7} {有限证书怎样承担全部输入}}

{两种证明都遵循同一个严格顺序：}

\begin{enumerate}
\def\labelenumi{\arabic{enumi}.}
\tightlist
\item
  {原位置计数表给宏；标量闭路给正向逆。}
\item
  {精确交换子给双向根。}
\item
  {外部有限指数、强同余子群定理和任意有限环消元给明确的主同余包含。}
\item
  {在这个已经确定的模数上穷尽方向与可实现单位尺度，得到有限父边树。}
\item
  {对任意无界输入，用父边词回到终端剩余类，再用主同余补列矩阵到精确终端。}
\end{enumerate}

{第五步仍处理任意大的有理输入，不要求把所有输入值穷举。有限图只承担已证明的有限商命题。程序另外为图中每个状态选取一个整数本原提升，核对其逆词和精确} {Bezout} {目标矩阵；它没有展开任意主同余元素的完整原子词。}

{二十元全部单位已可实现；二十二元单位像缺少一半，但宏本身可在两种尺度类之间移动。后者展示了为什么``单位像不满''不意味着不可达：必须检查式} {(C31.15)} {的联合轨道。}


\section{七平均的二十三至二十六元碰撞总和}\label{construct:seven-collisions}
\subsection*{\texorpdfstring{{C32.2} {一般有效的对称压缩}}{C32.2 一般有效的对称压缩}}\label{C32-h208}
\addcontentsline{toc}{subsection}{{C32.2} {一般有效的对称压缩}}

{沿用}第\ref{construct:five-small}节{的双重值剩余形状：}

\[
f(a)=p+\alpha,\quad f(b)=p+\beta,
\quad\alpha\ge\beta\ge0,\quad\alpha+\beta\le p-3,
\quad p\nmid a-b.
\]

{保护集合} {\ensuremath{T}} {是危险低重数例外集的包含极小横截集，记} {\ensuremath{h=|T|}。} {令} {\ensuremath{\delta =a-b}，未保护低重数位置的参数为} {\ensuremath{t_{i}=(x_{i}-b)/\delta }，余数} {\ensuremath{z_{i}=t_{i}}} {\ensuremath{mod}} {\ensuremath{p}。} {假设所有保护后候选都碰撞，就有此前的子集关联方程。}

\textbf{{对称压缩引理。}} {将一个余数类的实际参数取算术平均，记其重数} {\ensuremath{f_{r}}、平均值} {\ensuremath{y_{r}}。} {如果全部个体位置方程相容，那么压缩变量} {\ensuremath{y_{r}}} {满足}

\[
\sum_r k_r y_r=p-i
\]

{或另一锚点的}

\[
\sum_r k_r y_r=\sum_r k_r+j-p,
\]

{其中} {\ensuremath{k_{r}}} {是实际可选的位置数，\ensuremath{0\le k_{r}\le f_{r}}，且总选取数及锚点富余数合法。} {总和保持为} {\ensuremath{sum_{i}}} {\ensuremath{t_{i}=sum_{r}}} {\ensuremath{f_{r}}} {\ensuremath{y_{r}}。}

{证明：固定每类选取数量} {\ensuremath{k_{r}}，对该类型的全部位置选择取方程平均。} {一个位置出现在} {\ensuremath{k_{r}/f_{r}}} {比例的选择中，所以它们的贡献恰为} {\ensuremath{k_{r}}} {\ensuremath{y_{r}}。} {这只平均分析方程，不对实际状态新增任何操作。}

{若某条候选方程在一个余数类中选取} {\ensuremath{0<k_{r}<f_{r}}} {个位置，则交换一个选中与未选的位置，两个方程相减，迫使这两个位置的实际参数相等。} {对所有位置重复交换，该余数类便实际常值。}

{因此，求总和时可以从一开始压缩到``每个余数一个变量''，不必先假设同余值实际相等。} {这个引理对任意} {\ensuremath{p}} {成立。}

\subsection*{\texorpdfstring{{C32.3} {七平均剩余模式的完整分类}}{C32.3 七平均剩余模式的完整分类}}\label{C32-h209}
\addcontentsline{toc}{subsection}{{C32.3} {七平均剩余模式的完整分类}}

{此前的顺序、同余双重值、单背景和共同预留分支全部沿用。} {只剩} {\ensuremath{p=7}、\ensuremath{\alpha +\beta \le 4}} {的不同余双重值情况。}

{锚点同余修复表明：余数} {0} {和} {1} {各至多有一个低重数位置，其实际参数分别为} {7、-6。} {任一余数若有七个未保护位置，直接平均它们就保留两个旧重值。因此所有余数重数至多六。}

{令}

\[
n-h=21+e,\qquad L=7+e-\alpha-\beta.
\]

{对} {\ensuremath{n=23}、24、25、26} {及其} {\ensuremath{h\le d_{7}(n)}，总有} {\ensuremath{e}} {in} {\{1,2,3,4,5\}。} {逐项列出全部无序余数多重集，并建立第} {2} {节的精确压缩方程，共有} {57026} {个系统。} {仅下列} {11} {个模式相容：}

{ % do not increment counter
\begin{longtable}[]{@{}rlllr@{}}
\toprule\noalign{}
{\ensuremath{e}} & {alpha,beta} & {余数与重数} & {压缩实际参数} & {\ensuremath{D=7+\alpha +sum}} {\ensuremath{f_{r}}} {\ensuremath{y_{r}}} \\
\midrule\noalign{}
\endhead
\bottomrule\noalign{}
\endlastfoot
{1} & {2,0} & {\ensuremath{2^{6}}} & {2} & {21} \\
{1} & {3,0} & {\ensuremath{2^{5}}} & {2} & {20} \\
{1} & {3,0} & {\ensuremath{2^{4},3}} & {2,3} & {21} \\
{1} & {4,0} & {\ensuremath{2^{4}}} & {2} & {19} \\
{1} & {4,0} & {\ensuremath{2^{3},3}} & {2,3} & {20} \\
{1} & {4,0} & {\ensuremath{2^{3},4}} & {2,4} & {21} \\
{1} & {4,0} & {\ensuremath{2^{2},3^{2}}} & {2,3} & {21} \\
{2} & {3,0} & {\ensuremath{2^{6}}} & {2} & {22} \\
{2} & {4,0} & {\ensuremath{2^{5}}} & {2} & {21} \\
{2} & {4,0} & {\ensuremath{2^{4},3}} & {2,3} & {22} \\
{3} & {4,0} & {\ensuremath{2^{6}}} & {2} & {23} \\
\end{longtable}
}

{\ensuremath{e=4}、5} {没有相容模式。} {每个相容系统都满秩，表中是唯一有理解；每个重数大于一的余数类都出现在至少一条部分选取方程中，因此个体实际参数也都等于表中给出的整数。}

{这个有限证书给出的关键结果是}

\[
\boxed{D\in\mathbb Z,\qquad0<D<21+e.}
\tag{C32.1}
\]

{一般原位置参数是任意有理数，不受高度界限制；有限枚举只枚举其模} {7} {模式。}

\subsection*{\texorpdfstring{{C32.4} {接到最小保护总和矛盾}}{C32.4 接到最小保护总和矛盾}}\label{C32-h210}
\addcontentsline{toc}{subsection}{{C32.4} {接到最小保护总和矛盾}}

{若} {\ensuremath{h>0}，选择某个保护位置的私人素数} {\ensuremath{q}。} {由保护构造，\ensuremath{q|(a-b)=\delta }，其他保护位置在模} {\ensuremath{q}} {下都属于主类。} {未保护部分相对} {\ensuremath{b}} {的总和为} {\ensuremath{\delta }} {\ensuremath{D}，(1)} {保证它模} {\ensuremath{q}} {为零。} {这使整个状态在模} {\ensuremath{q}} {下只剩该保护位置一项非零差，违反} {\ensuremath{q|n}} {及全局零和。}

{若} {\ensuremath{h=0}，全部低重数实际参数都是整数，所以所有实际值都属于} {b+delta} {\ensuremath{Z}。} {合法性于是迫使} {\ensuremath{\operatorname{gcd}(\delta ,n)=1}。} {零和关系}

\[
nb+\delta D=0
\]

{要求} {\ensuremath{n|D}，但} {(1)} {给} {\ensuremath{0<D<n}，矛盾。}

{因此全部碰撞不可能发生，必有安全非恒等整数七平均保持双重值。}

\subsection*{\texorpdfstring{{C32.5} {初始化、其他分支和终止}}{C32.5 初始化、其他分支和终止}}\label{C32-h211}
\addcontentsline{toc}{subsection}{{C32.5} {初始化、其他分支和终止}}

{\ensuremath{n=23}、24、25、26} {都满足旧初始化条件} {\ensuremath{n\ge 21+d_{7}(n)}。} {从} {7X} {出发，至多两次整数平均进入双重值不变量。} {三重值分支适用已证的} {\ensuremath{n\ge 3p+2}} {引理，其余旧分支的前提也都满足。} {第} {4} {节补齐唯一剩余的低总重数不同余双重值分支。}

{每一步均安全、整数、非恒等，平方能量严格下降；相反值终端可直接消去。} {所以有限终止。缩回原输入，所有中间值属于} {(1/7)Z。} {接已有} {\ensuremath{n\ge 27}，主定理成立。}


\section{三平均的双三重值下降}\label{construct:triple-tail}
{本节平均元数固定为} \(3\){，位置数} \(n\ge11\){；}\(d\) {表示} \(n\) {的不同非三素因子个数。}

\subsection*{\texorpdfstring{{C33.2} {不变量和三个推广要点}}{C33.2 不变量和三个推广要点}}\label{C33-h212}
\addcontentsline{toc}{subsection}{{C33.2} {不变量和三个推广要点}}

{记}

\[
\mathcal P_n=\{q:q\mid n,\ q\text{ 为素数},\ q\ne3\},
\qquad d=|\mathcal P_n|.
\]

{对整数配置} {\ensuremath{E}，记实际值} {\ensuremath{a}} {的重数为} {f(a)，三重值集合为} {H=\{a:f(a)\textgreater=3\}。保持}

\[
\mathcal I_n(E):\quad \sum E=0,\quad E\in\mathbb Z^n,
\quad \forall q\in\mathcal P_n\ E\text{ 模 }q\text{ 不全同余},
\quad |H|\ge2.
\]

{与素数证明相比，只需三项修正：同时保护至多} {\ensuremath{d}} {个危险集；在使用} {(a,x,x)} {时保留它已满足的仿射关系，从而覆盖} {\ensuremath{q=2}；无同余邻居时按三重值的剩余份数选择池内操作或三个余类各取一个。}

\subsection*{\texorpdfstring{{C33.3} {一般安全引理}}{C33.3 一般安全引理}}\label{C33-h213}
\addcontentsline{toc}{subsection}{{C33.3} {一般安全引理}}

{称一次整数三平均为安全，如果对所有} {\ensuremath{q}} {属于} {\ensuremath{P_{n}}} {都保持不全同余。}

\paragraph*{\texorpdfstring{{3.1} {危险支持集}}{3.1 危险支持集}}\label{C33-h214}
\addcontentsline{toc}{paragraph}{{3.1} {危险支持集}}

{对固定} {\ensuremath{q}，若有至少} {n-3} {个位置属于同一模} {\ensuremath{q}} {余数类，记其补集为} {\ensuremath{F_{q}}。由于} {\ensuremath{n\ge 11}，主类唯一。由不全同余及} {\ensuremath{q|n}、sum} {\ensuremath{E=0}，有}

\[
2\le |F_q|\le3.
\]

{大小} {1} {不可能：若主类余数为} {\ensuremath{c}、唯一例外为} {c+t，则总和模} {\ensuremath{q}} {等于} {\ensuremath{nc+t=t}，迫使} {\ensuremath{t=0}。}

{一个三元组只要不包含任何完整} {\ensuremath{F_{q}}，未选位置便对每个} {\ensuremath{q}} {仍不全同余，因此操作安全。没有主类的} {\ensuremath{q}} {自动安全。}

\paragraph*{\texorpdfstring{{3.2} {多素数保护池}}{3.2 多素数保护池}}\label{C33-h215}
\addcontentsline{toc}{paragraph}{{3.2} {多素数保护池}}

{给定一个位置池，逐个处理危险} {\ensuremath{q}。若当前可用池仍包含整个} {\ensuremath{F_{q}}，留出其中一个位置；否则无需动作。至多留出} {\ensuremath{d}} {个位置，之后选出的任何三元组都安全。}

{若池中所有值同余模} {3，且每个实际值至多出现两次，则：}

\begin{itemize}
\tightlist
\item
  {池大小至少} {d+3} {时，有安全、非恒等的整数三元组。}
\item
  {池大小至少} {d+4} {时，还可要求均值不等于某个指定值} {\ensuremath{a}。}
\end{itemize}

{第二条只需在留出位置后取四个位置。如果它们的每个三元组和都为} {3a，比较两个仅替换一个位置的三元组，便得四个值全相等，违反重数至多二。}

{这里有统一容量余量}

\[
\boxed{n\ge d+10\qquad(n\ge11).}
\tag{C33.1}
\]

{\ensuremath{d=0}、1} {时显然；\ensuremath{d=2}} {时} {\ensuremath{n}} {至少为} {14；\ensuremath{d\ge 3}} {时} {\ensuremath{n\ge 2*5^{d-1}\ge d+10}。因此不存在随不同素因子个数增长的新维数边界。}

\paragraph*{\texorpdfstring{{3.3} {重复值和仿射见证}}{3.3 重复值和仿射见证}}\label{C33-h216}
\addcontentsline{toc}{paragraph}{{3.3} {重复值和仿射见证}}

{下列整数操作对包括} {\ensuremath{q=2}} {在内的全部} {\ensuremath{q}} {都安全：}

\begin{enumerate}
\def\labelenumi{\arabic{enumi}.}
\tightlist
\item
  {(a,a,x)，其中} {\ensuremath{f(a)\ge 3}。未选位置留下} {\ensuremath{a}；若输出全同余为} {\ensuremath{c}，便有} {\ensuremath{a=c}、新均值=\ensuremath{c}，故} {\ensuremath{x=c}，矛盾。}
\item
  {(a,b,x)，其中不同} {a,b} {各至少有两份。未选位置留下} {\ensuremath{a}、\ensuremath{b}，同样迫使} {\ensuremath{x=c}。}
\item
  {未选位置留下} {\ensuremath{a}、\ensuremath{b}，被选值中至少两个是} {\ensuremath{a}、\ensuremath{b}} {的整数仿射组合。若输出全同余为} {\ensuremath{c}，这些被选值也都等于} {\ensuremath{c}} {\ensuremath{mod}} {\ensuremath{q}，剩下的一个由三元组和强制等于} {\ensuremath{c}。}
\end{enumerate}

{第三条特别覆盖} {\ensuremath{x=3b-2a}、\ensuremath{y=3a-2b}} {时的} {(x,y,z)、(x,b,z)。还覆盖} {(a,x,x)：在实际用到它的分支中，\ensuremath{a}、\ensuremath{b}} {都留下副本，且} {\ensuremath{x=3b-2a}。于是不能出现模} {2} {例外。}

{必须限定最后一条的前提。单独的``(a,x,x)} {且} {\ensuremath{a}、\ensuremath{x}} {各重复两次''在模} {2} {下不安全。例如}

\[
(3^3,0^2,1^6,-15)
\]

{是} {12} {元零和、模} {2} {不全同余的状态，但平均} {(3,0,0)} {后全部为奇数。推广证明使用留下的} {\ensuremath{a}、\ensuremath{b}} {和} {\ensuremath{x}} {的仿射关系，完全不需要除以} {2。}

\paragraph*{\texorpdfstring{{3.4} {两值排除只需要不同余模} {3}}{3.4 两值排除只需要不同余模 3}}\label{C33-h217}
\addcontentsline{toc}{paragraph}{{3.4} {两值排除只需要不同余模} {3}}

\textbf{{引理。}} {满足上述所有不全同余条件的整数零和配置，不可能只有两个实际值} {\ensuremath{a}、\ensuremath{b}} {且} {\ensuremath{a\ne b}} {\ensuremath{mod}} {3。}

{若} {\ensuremath{a}} {出现} {\ensuremath{r}} {次，则}

\[
r(a-b)=-nb.
\]

{对每个} {\ensuremath{q|n}、\ensuremath{q\ne 3}，不全同余意味着} {\ensuremath{q}} {不整除} {a-b；对} {\ensuremath{q=3}，假设也排除了整除。因此} {\ensuremath{\operatorname{gcd}(a-b,n)=1}，得到} {\ensuremath{n|r}，与} {\ensuremath{0<r<n}} {矛盾。}

{素数证明中的所有两值排除恰好发生在} {\ensuremath{a}、\ensuremath{b}} {不同余模} {3} {的分支。因此，不需要在一般维数中排除所有两值状态。}

\subsection*{\texorpdfstring{{C33.4} {至多两步建立不变量}}{C33.4 至多两步建立不变量}}\label{C33-h218}
\addcontentsline{toc}{subsection}{{C33.4} {至多两步建立不变量}}

{研究} {\ensuremath{E=3X}。所有初值被} {3} {整除，对} {\ensuremath{P_{n}}} {的不全同余保持。}

{若已有两个三重值，无需操作。}

{若只有一个三重值} {\ensuremath{a}：\ensuremath{f(a)\ge 5}} {时，取不同值} {\ensuremath{b}，平均} {(a,a,b)，留下三份} {\ensuremath{a}} {并新建另一三重值；\ensuremath{f(a)=3}、4} {时，\ensuremath{a}} {之外的池至少} {\ensuremath{n-4\ge d+6}} {个位置，每值至多两份。用保护池引理选择均值不等于} {\ensuremath{a}} {的三元组。}

{若没有三重值，先在整个初态中用保护池引理制造三重值} {\ensuremath{a}。它的重数在} {3} {到} {5} {之间。\ensuremath{a}} {之外至少} {\ensuremath{n-5\ge d+5}} {个位置，都是尚未改变的初值，仍被} {3} {整除且每值至多两份。再用保护池引理制造不同于} {\ensuremath{a}} {的三重值。}

{第二步不使用新生} {\ensuremath{a}，因此无需再增加三进精度。}

\subsection*{\texorpdfstring{{C33.5} {至少三个三重值}}{C33.5 至少三个三重值}}\label{C33-h219}
\addcontentsline{toc}{subsection}{{C33.5} {至少三个三重值}}

{若} {\textbar{}\ensuremath{H|\ge 4}，有两个不同三重值} {\ensuremath{a}、\ensuremath{b}} {同余模} {3。平均} {(a,a,b)} {安全、非恒等且为整数；至少两个未参与的三重值保留。}

{若} {\textbar{}\ensuremath{H|=3}} {且有同余的} {\ensuremath{a}、\ensuremath{b}，设第三个为} {\ensuremath{c}。(a,a,b)} {与} {(b,b,a)} {的均值不同，因此至少一个均值不同于} {\ensuremath{c}，保留} {\ensuremath{c}} {和新均值两个三重值。}

{若} {\textbar{}\ensuremath{H|=3}} {且三个值各占一个模} {3} {余类，有非三重值} {\ensuremath{x}} {时，将它与其余类中的三重值} {\ensuremath{a}} {组成} {(a,a,x)，其他两个三重值保留。}

{否则只有三个实际值} {\ensuremath{a}、\ensuremath{b}、\ensuremath{c}，平均} {(a,b,c)。若至少两个重数\textgreater=4，操作后它们仍为三重值。否则重数只能是} {3、3、n-6，设} {\ensuremath{c}} {的重数为} {n-6。新均值若不等于} {\ensuremath{c}，不变量保持；若等于} {\ensuremath{c}，则} {\ensuremath{a+b=2c}，零和又给} {\ensuremath{nc=0}，得到显式终端}

\[
T_n(a)=(a^3,(-a)^3,0^{n-6}),\qquad a\ne0.
\tag{C33.2}
\]

{三次} {(a,-a,0)} {立即归零。}

\subsection*{\texorpdfstring{{C33.6} {两个同余模} {3} {的三重值}}{C33.6 两个同余模 3 的三重值}}\label{C33-h220}
\addcontentsline{toc}{subsection}{{C33.6} {两个同余模} {3} {的三重值}}

{设三重值恰为} {\ensuremath{a}、\ensuremath{b}。若} {\ensuremath{f(a)\ge 4}，平均} {(b,b,a)；若} {\ensuremath{f(b)\ge 4}，作对称操作。剩余至少一个旧三重值，新均值与它不同。}

{只剩} {\ensuremath{f(a)=f(b)=3}。其余池} {\ensuremath{S}} {大小} {\ensuremath{n-6\ge d+4}，每个实际值最多出现两次。}

{若} {\ensuremath{S}} {有} {\ensuremath{x}} {与} {\ensuremath{a}、\ensuremath{b}} {同余模} {3，考虑} {(a,a,x)、(b,b,x)。二者安全；同时破坏双三重值条件会要求}

\[
2a+x=3b,\qquad 2b+x=3a,
\]

{从而} {\ensuremath{5(a-b)=0}，矛盾。}

{否则} {\ensuremath{S}} {只占另外两个模} {3} {余类。若只占一类，直接使用保护池引理。}

{若两类} {\ensuremath{S_{1}}、\ensuremath{S_{2}}} {都非空，对跨类位置对} {\ensuremath{c}、\ensuremath{e}} {考虑} {(a,c,e)、(b,c,e)。它们不能安全的充分障碍只能是某个} {\ensuremath{F_{q}}} {恰为} {\{c,e\}：选一次} {\ensuremath{a}} {或} {\ensuremath{b}} {不可能包含该三重值的全部例外位置。因此至多} {\ensuremath{d}} {个跨类位置对被阻塞。}

{实际跨类位置对数量至少}

\[
|S_1||S_2|\ge |S|-1=n-7>d.
\]

{选一个未被阻塞的对，两种操作都安全。如果二者均破坏双三重值条件，则} {\ensuremath{a+c+e=3b}、\ensuremath{b+c+e=3a}，推出} {\ensuremath{4(a-b)=0}，矛盾。}

\subsection*{\texorpdfstring{{C33.7} {两个不同余模} {3} {的三重值}}{C33.7 两个不同余模 3 的三重值}}\label{C33-h221}
\addcontentsline{toc}{subsection}{{C33.7} {两个不同余模} {3} {的三重值}}

{设三重值恰为} {\ensuremath{a}、\ensuremath{b}。先考虑任何不同于} {\ensuremath{a}、但与} {\ensuremath{a}} {同余模} {3} {的值} {\ensuremath{x}。操作} {(a,a,x)} {安全；只有} {\ensuremath{f(a)\le 4}} {且} {\ensuremath{x=3b-2a}} {时，它才不能保持双三重值。}

{若这个} {\ensuremath{x}} {至少出现两次，改用} {(a,x,x)，新均值为} {2b-a，与} {\ensuremath{a}、\ensuremath{b}} {都不同，\ensuremath{b}} {未动。不变量保持；包括} {\ensuremath{q=2}} {在内的安全性由第} {3.3} {节的仿射见证保证。}

{因此，在尚未找到保持操作时，\ensuremath{a}} {的余类中最多有一个额外单点} {\ensuremath{x=3b-2a}，且} {\ensuremath{f(a)\le 4}。对} {\ensuremath{b}} {对称，最多有一个额外单点} {\ensuremath{y=3a-2b}，且} {\ensuremath{f(b)\le 4}。其他位置全在第三个模} {3} {余类中，记其池为} {\ensuremath{Z}。}

\paragraph*{\texorpdfstring{{7.1} {\ensuremath{x}、\ensuremath{y}} {都存在}}{7.1 、 都存在}}\label{C33-h222}
\addcontentsline{toc}{paragraph}{{7.1} {\ensuremath{x}、\ensuremath{y}} {都存在}}

{因为} {\ensuremath{f(a)+f(b)+2\le 10<n}，\ensuremath{Z}} {非空。取} {\ensuremath{z}} {属于} {\ensuremath{Z}，平均} {(x,y,z)。\ensuremath{a}、\ensuremath{b}} {的重数都保留，操作由仿射见证安全，且} {\ensuremath{x-y=5(b-a)\ne 0}。}

\paragraph*{\texorpdfstring{{7.2} {恰有一个额外单点}}{7.2 恰有一个额外单点}}\label{C33-h223}
\addcontentsline{toc}{paragraph}{{7.2} {恰有一个额外单点}}

{不妨只有} {\ensuremath{x}} {存在。若} {\ensuremath{Z}} {为空，安全操作} {(a,a,x)} {将产生只有} {\ensuremath{a}、\ensuremath{b}} {两值且不同余模} {3} {的状态，与第} {3.4} {节矛盾。因此} {\ensuremath{Z}} {非空。}

{对} {\ensuremath{z}} {属于} {\ensuremath{Z}，(x,b,z)} {是安全非恒等整数操作。若} {\ensuremath{f(b)\ge 4}，\ensuremath{a}、\ensuremath{b}} {的三重性保留。若} {\ensuremath{f(b)=3}，由} {\ensuremath{f(a)\le 4}} {得} {\textbar{}\ensuremath{Z|\ge n-8\ge 3}，每值最多两份，因此可以选} {\ensuremath{z\ne 5a-4b}，使新均值不等于} {\ensuremath{a}，保留} {\ensuremath{a}} {与新均值两个三重值。}

{只有} {\ensuremath{y}} {的情形对称。}

\paragraph*{\texorpdfstring{{7.3} {没有额外单点：按重数选择操作}}{7.3 没有额外单点：按重数选择操作}}\label{C33-h224}
\addcontentsline{toc}{paragraph}{{7.3} {没有额外单点：按重数选择操作}}

{此时} {\ensuremath{a}、\ensuremath{b}} {之外的所有值都在} {\ensuremath{Z}。\ensuremath{Z}} {不能为空，否则违反第} {3.4} {节。}

\begin{itemize}
\tightlist
\item
  {若} {\ensuremath{f(a)=f(b)=3}，\textbar{}\ensuremath{Z|=n-6\ge d+4}，使用保护池引理，\ensuremath{a}、\ensuremath{b}} {完整保留。}
\item
  {若} {\ensuremath{f(a),f(b)\ge 4}，对任意} {\ensuremath{z}} {属于} {\ensuremath{Z}，(a,b,z)} {安全且两者仍至少有三份。}
\item
  {否则不妨} {\ensuremath{f(a)=3}、\ensuremath{f(b)\ge 4}。选择} {\ensuremath{z\ne 2b-a}，则} {(a,b,z)} {的新均值不同于保留的} {\ensuremath{b}，不变量保持。}
\end{itemize}

{最后一种选择必然存在。否则} {\ensuremath{Z}} {全部等于} {\ensuremath{z=2b-a}。每个非三重值最多两份，故} {\ensuremath{r=|Z|}} {是} {1} {或} {2。总和为零给出}

\[
(3-r)(a-b)+nb=0.
\tag{C33.3}
\]

{由于全部值是} {\ensuremath{a}、\ensuremath{b}、2b-a，各个} {\ensuremath{q}} {属于} {\ensuremath{P_{n}}} {的不全同余条件保证} {\ensuremath{q}} {不整除} {a-b；而} {\ensuremath{a}、\ensuremath{b}} {不同余模} {3。所以} {\ensuremath{\operatorname{gcd}(a-b,n)=1}，式} {(C33.3)} {强迫} {\ensuremath{n|(3-r)}，与} {3-r} {属于} {\{1,2\}} {矛盾。}

{这一步同时解决多个危险素数下的小池问题和偶数维数边界，不使用对} {2} {的模逆元。}

\subsection*{\texorpdfstring{{C33.8} {有限终止和全维数结论}}{C33.8 有限终止和全维数结论}}\label{C33-h225}
\addcontentsline{toc}{subsection}{{C33.8} {有限终止和全维数结论}}

{除显式终端} {\ensuremath{T_{n}}} {外，第} {5--7} {节给每个满足} {\ensuremath{I_{n}}} {的状态一个安全、非恒等、保持} {\ensuremath{I_{n}}} {的整数操作。平方能量}

\[
\Phi(E)=\sum_i e_i^2\in\mathbb Z_{\ge0}
\]

{每次严格下降，因此必在有限步进入} {\ensuremath{T_{n}}，随后三步归零。初始为} {3X，所有中间值都为整数；在原} {\ensuremath{X}} {上执行同样的位置序列，所有中间值都属于} {(1/3)Z。}

{对本原整数零和的非零状态，\ensuremath{G}} {为三幂保证对} {\ensuremath{P_{n}}} {中所有素数不全同余，故定理适用。必要性沿用非三模同余障碍。7、8、10} {元由已有独立基例，9} {元由三幂网络完成。因此} {\ensuremath{n\ge 7}} {的充要判据不再需要一般素数到合数的归约。}

{本证明的深层条件是重复值提供可保留的同余见证、保护池容量足够，以及有限整数能量下降；素数性并非不变量法本身所必需。}


\section{三平均的八元欧几里得下降}\label{construct:triple-small}
{本节取三平均八元核心} \((u^4,v^3,-4u-3v)\){，其中} \((u,v)\) {本原且} \(u-v\) {为奇数。任意合法八元输入的安全归约由正文八元安全入口引理给出。}

\subsection*{\texorpdfstring{{C34.3} {八元族的三进欧几里得算法}}{C34.3 八元族的三进欧几里得算法}}\label{C34-h226}
\addcontentsline{toc}{subsection}{{C34.3} {八元族的三进欧几里得算法}}

{八元也令}

\[
x=u,\qquad y=u+v.
\]

{这次平方能量出现更简单的对角型：}

\[
\sum_iX_i^2=4Q_8(x,y),\qquad
Q_8(x,y)=2x^2+3y^2.
\]

{而数论条件变成}

\[
u-v\equiv u+v=y\equiv1\pmod2.
\]

{所以八元族恰对应本原整数对} \((x,y)\) {且} \(y\) {为奇数。}

\paragraph*{\texorpdfstring{{3.1} {三个有限平均宏}}{3.1 三个有限平均宏}}\label{C34-h227}
\addcontentsline{toc}{paragraph}{{3.1} {三个有限平均宏}}

{存在以下三个真实平均宏。等号只差一个非零全局比例；全局比例不影响可达性。}

{ % do not increment counter
\begin{longtable}[]{@{}lrll@{}}
\toprule\noalign{}
{宏} & {步数} & {在} \((u,v)\) {上的输出} & {在} \((x,y)\) {上的射影作用} \\
\midrule\noalign{}
\endhead
\bottomrule\noalign{}
\endlastfoot
\(R\) & {4} & \(((2u+v)/3,-u-2v/3)\) & \((x+y,-y)\) \\
\(S\) & {6} & \((v/9,u/9)\) & \((y-x,y)\) \\
\(F\) & {3} & \(((2u+v)/3,-(8u+7v)/9)\) & \((3x+3y,2x-4y)\) \\
\end{longtable}
}

{这些不是形式矩阵假设，而是实际三平均序列。令} \(w=-4u-3v\){。三个宏依次选择的值如下：}

\[
\begin{aligned}
R:\quad &(v,v,u),\ (v,u,u),\\
&\left(w,\frac{u+2v}{3},\frac{2u+v}{3}\right),\
\left(\frac{u+2v}{3},\frac{2u+v}{3},u\right);\\[3pt]
S:\quad &(w,v,u),\ \left(-u-\frac{2v}{3},-u-\frac{2v}{3},u\right),\\
&\left(-u-\frac{2v}{3},u,u\right),\\
&\left(-\frac u3-\frac{4v}{9},v,\frac u3-\frac{2v}{9}\right)
\quad\text{重复两次},\\
&\left(\frac v9,\frac v9,\frac u3-\frac{2v}{9}\right);\\[3pt]
F:\quad &(v,u,u),\
\left(w,\frac{2u+v}{3},\frac{2u+v}{3}\right),\
(v,u,u).
\end{aligned}
\]

{逐行平均后，所列值在当时状态中确实存在；最终重数重新成为} \(4+3+1\){，并得到表中的参数。}

\paragraph*{\texorpdfstring{{3.2} {平移与严格下降}}{3.2 平移与严格下降}}\label{C34-h228}
\addcontentsline{toc}{paragraph}{{3.2} {平移与严格下降}}

{两个反射的合成给出}

\[
S\circ R:(x,y)\longmapsto(x+2y,y),
\]

\[
R\circ S:(x,y)\longmapsto(x-2y,y).
\]

{所以可以用有限次实际平均，把} \(x\) {替换成任意} \(x+2ky\){，同时保持} \(y\) {不变。}

{现在设} \(|y|>1\){。}

{若} \(x\) {为偶数，先施行} \(S\){，把它变成} \(y-x\){。因为} \(y\) {为奇数，新第一坐标为奇数，}\(|y|\) {不变。于是可假设} \(x,y\) {都是奇数。}

{利用上述平移，选择一个代表} \(x_0\equiv x\pmod{2y}\){，使}

\[
|x_0-2y|<|y|.
\]

{严格不等式总能做到：最小绝对值剩余至多} \(|y|\){；若等号成立，则} \(y\mid x\){，与} \(\gcd(x,y)=1\) {一起迫使} \(|y|=1\){。}

{对} \((x_0,y)\) {施行宏} \(F\){，清除无关比例后得到}

\[
(3(x_0+y),2(x_0-2y)).
\]

{因为} \(x_0,y\) {都是奇数，两个坐标都有公因子} \(2\){。先除以} \(2\) {后，第二坐标} \(x_0-2y\) {为奇数，所以之后约去的公因子也是奇数。最终本原第二坐标} \(y'\) {满足}

\[
y'\text{ 为奇数},\qquad |y'|\le |x_0-2y|<|y|.
\]

{这就给出了对正整数} \(|y|\) {的严格归纳下降。}

\paragraph*{\texorpdfstring{{3.3} {终止情形}}{3.3 终止情形}}\label{C34-h229}
\addcontentsline{toc}{paragraph}{{3.3} {终止情形}}

{当} \(|y|=1\) {时，使用} \(x\mapsto x+2ky\) {可把} \(x\) {化成} \(0\) {或} \(\pm1\){。反射} \(R,S\) {又把} \((\pm1,\pm1)\) {送到某个} \((0,\pm1)\){。只需处理}

\[
\mathcal B_8(0,1)=(-3,0,0,0,0,1,1,1).
\]

{它有显式四步收尾：先平均} \((-3,0,0)\){，得到三个} \(-1\){；然后连续三次平均} \((-1,0,1)\){。每次消去一对} \(-1,1\){，最终全零。}

{因此所有合法的八元二参数状态均可达。与第} {1} {节归约结合，八元主定理得证。}


\section{三平均的十元基例}\label{construct:triple-ten}
\subsection*{\texorpdfstring{{C35.2} {十元安全整数下降}}{C35.2 十元安全整数下降}}\label{C35-h230}
\addcontentsline{toc}{subsection}{{C35.2} {十元安全整数下降}}

{称三个坐标构成一个}\textbf{{整三元组}}{，如果它们不全相等且和被} \(3\) {整除。平均这样的三元组得到一个非恒等整数状态。}

{若未选的七个位置中同时保留：}

\begin{itemize}
\tightlist
\item
  {一奇一偶；}
\item
  {两个不同的模} \(5\) {余数，}
\end{itemize}

{则该操作后的本原状态仍有} \(G=1\){。事实上，未选位置的差异说明原始输出的全体公因子不含} \(2\) {或} \(5\){；约去公因子后，这两个差异仍然存在。另一方面，本原十元零和状态的} \(G\) {只能含素因子} \(2,5\){。}

{先记录一个模} \(3\) {结构事实。}

\begin{quote}
\textbf{{模三引理。}} {若一个整数多重集没有整三元组，则它至多使用两个模} \(3\) {余数类；每个至少含三个位置的余数类只能由同一个实际数值组成。}
\end{quote}

{若三个余数类都出现，各取一个即成整三元组；若同一余数类中至少三个数不全相等，任取一个非恒等三元组也满足条件。这就证明了引理。}

\begin{quote}
\textbf{{安全下降引理。}} {设} \(X\) {是本原零和十元整数状态且} \(G(X)=1\){。若没有任何实际数值出现六次，则存在一次保持} \(G=1\) {的非恒等整数三平均。}
\end{quote}

{证明分成四种余数分布。}

\paragraph*{\texorpdfstring{{2.1} {两种局部障碍都不集中}}{2.1 两种局部障碍都不集中}}\label{C35-h231}
\addcontentsline{toc}{paragraph}{{2.1} {两种局部障碍都不集中}}

{一个三元组拿走后，补集在奇偶上单色，只可能发生在全局奇偶重数为} \(2+8\) {且三元组拿走两个少数位置时。零和保证奇数位置数为偶数，所以不存在} \(1+9\) {或} \(3+7\) {的危险情形；全部同奇偶又与} \(G=1\) {矛盾。}

{类似地，补集模} \(5\) {单色只可能发生在某个模} \(5\) {类至少出现七次时。若奇偶不是} \(2+8\){，且每个模} \(5\) {类至多六个，则任意整三元组都安全。}

{整个十元集必有整三元组。否则由模三引理和``无六重值''，两个模} \(3\) {类只能各含五个相等值，状态为} \((a^5,b^5)\){。零和给出} \(b=-a\){；本原化后} \(G=2\){，与} \(G=1\) {矛盾。}

\paragraph*{\texorpdfstring{{2.2} {只有奇偶集中}}{2.2 只有奇偶集中}}\label{C35-h232}
\addcontentsline{toc}{paragraph}{{2.2} {只有奇偶集中}}

{设多数奇偶类为} \(M\){，}\(|M|=8\){，少数类为} \(P\){，}\(|P|=2\){。若没有安全操作，则} \(M\) {内没有整三元组，而且一个} \(P\) {元素与两个} \(M\) {元素也不能组成整三元组。}

{由模三引理，若} \(M\) {的一个模} \(3\) {类含至少六个位置，已经出现六重值。否则} \(M\) {的两个模} \(3\) {类都至少含两个位置。设其余数为} \(r,s\){。两个} \(M\) {元素的余数和可以取}

\[
2r,\quad r+s,\quad2s,
\]

{这恰是全部三个模} \(3\) {余数。于是对任一} \(z\in P\){，都能选两个} \(M\) {元素使三者之和被} \(3\) {整除；奇偶不同又保证三者不全相等。该三元组只拿走一个少数位置，故安全，矛盾。}

\paragraph*{\texorpdfstring{{2.3} {只有模五集中}}{2.3 只有模五集中}}\label{C35-h233}
\addcontentsline{toc}{paragraph}{{2.3} {只有模五集中}}

{设模} \(5\) {主类为} \(C\){，其外位置集为} \(F\){，其中} \(|F|=b\in\{1,2,3\}\){。}

{若} \(b=2\) {或} \(3\){，对任一} \(z\in F\){，三元组``}\(z\) {加两个} \(C\) {元素''没有拿走全部} \(F\){。在假设不存在安全操作时，}\(C\) {内因而没有整三元组，可以应用模三引理。若} \(C\) {的两个模} \(3\) {类都至少含两个位置，其两两余数和覆盖全部模} \(3\) {类，于是得到安全三元组；否则} \(C\) {中至少六个位置属于同一个实际值。}

{若} \(b=1\){，则} \(|C|=9\){。在无六重值的假设下，模三引理迫使} \(C\) {为} \((r^4,s^5)\){，其中} \(r,s\) {是两个实际值。二者模} \(5\) {相同，记为} \(q\){。零和模} \(5\) {给出唯一的外部值也同余于}

\[
-9q\equiv q\pmod5,
\]

{与它位于} \(C\) {外矛盾。}

\paragraph*{\texorpdfstring{{2.4} {奇偶与模五同时集中}}{2.4 奇偶与模五同时集中}}\label{C35-h234}
\addcontentsline{toc}{paragraph}{{2.4} {奇偶与模五同时集中}}

{令} \(P\) {为两个少数奇偶位置，}\(F\) {为模} \(5\) {主类外的} \(b\le3\) {个位置，并令}

\[
K=\{1,\ldots,10\}\setminus(P\cup F),\qquad
|K|=8-b+|P\cap F|.
\]

\(K\) {内所有数模} \(10\) {相同。任何} \(K\) {内整三元组都安全，所以} \(K\) {满足模三引理。}

{若} \(b=1\){，总能在} \(P\setminus F\) {中取一个位置} \(z\){。三元组``}\(z\) {加两个} \(K\) {元素''若为整三元组便安全。若} \(K\) {有两个模} \(3\) {类且各至少两个位置，两两余数和覆盖全部模} \(3\) {类；否则} \(|K|\ge7\) {立即给出六重值。}

{以下设} \(b\ge2\){。对每个} \(z\notin K\){，``}\(z\) {加两个} \(K\) {元素''只拿走至多一个} \(P\) {位置和至多一个} \(F\) {位置，所以只要和被} \(3\) {整除就安全。因此，无六重值时只剩} \(|K|=5\) {或} \(6\){。}

{若} \(|K|=6\){，模三引理和两两余数和条件迫使} \(K\) {的模} \(3\) {重数为} \(5+1\){。记五重类为} \(r\){、单点类为} \(s\){。每个} \(K\) {外位置都必须属于余数} \(s\){；否则它可与两个} \(K\) {元素组成安全三元组。此时或有} \(b=2,P\cap F=\varnothing\){，或有} \(b=3,|P\cap F|=1\){；两种情形都能分别从} \(P\setminus F\) {与} \(F\setminus P\) {各取一个位置。把它们与} \(K\) {中的单点组成三元组，三个余数都是} \(s\){，补集仍保留奇偶与模} \(5\) {的见证，得到安全三元组。}

{若} \(|K|=5\){，必有} \(b=3\) {且} \(P\cap F=\varnothing\){。}\(K\) {或为单一模} \(3\) {类} \(r^5\){，或为} \(r^4+s\){。第二种情形与上一段相同：所有外部位置都为余数} \(s\){，用} \(K\) {单点、一个} \(P\) {位置和一个} \(F\) {位置组成安全三元组。}

{在} \(K=r^5\) {的情形，外部位置只能使用另两个模} \(3\) {余数} \(s,t\){。若存在余数不同的一个} \(P\) {位置和一个} \(F\) {位置，则它们与任一} \(K\) {元素组成余数型} \((r,s,t)\) {的安全三元组。否则所有五个外部位置同属一个模} \(3\) {类；选一个} \(P\) {位置和两个} \(F\) {位置也给出安全三元组。矛盾。}

{四种情形全部完成，安全下降引理得证。}

{每次安全操作后取本原代表。非恒等整数平均严格降低平方能量}

\[
E(X)=\sum_iX_i^2,
\]

{而约去公因子只会继续降低它。因此反复使用安全下降引理，有限步后必定到达零，或者到达一个含六重值的状态。}

\subsection*{\texorpdfstring{{C35.3} {六重状态的四步桥}}{C35.3 六重状态的四步桥}}\label{C35-h235}
\addcontentsline{toc}{subsection}{{C35.3} {六重状态的四步桥}}

{考虑}

\[
X=(u,u,u,u,u,u,a,b,c,d),\qquad 6u+a+b+c+d=0,
\]

{并记四个例外差值}

\[
e_a=a-u,\quad e_b=b-u,\quad e_c=c-u,\quad e_d=d-u.
\]

{它们满足}

\[
e_a+e_b+e_c+e_d=-10u.
\tag{C35.3.1}
\]

\(G(X)=1\) {等价于：至少一个} \(e_i\) {为奇数，且至少一个} \(e_i\) {不被} \(5\) {整除。}

{若某个} \(e_i\) {与} \(10\) {互素，平均其余三个例外值。留下的例外记为} \(t\){，则一步后得到}

\[
\mathcal B_{10}\left(u,\frac{-6u-t}{3}\right).
\]

{清除分母后的两个参数差为}

\[
3u-(-6u-t)=9u+t=10u+(t-u)\equiv e_t\pmod{10},
\]

{所以这是合法终端族。}

{只需处理没有单个} \(e_i\) {与} \(10\) {互素的情形。此时每个奇} \(e_i\) {都被} \(5\) {整除，每个不被} \(5\) {整除的} \(e_i\) {都是偶数。由式} {(C35.3.1)，奇数} \(e_i\) {的个数为偶数；它不可能是四个，否则所有} \(e_i\) {都被} \(5\) {整除。因此四个差值模} \(10\) {恰由以下两类组成：}

\begin{itemize}
\tightlist
\item
  {两个} \(5\){；}
\item
  {两个互为相反数的非零偶数。}
\end{itemize}

{任意跨类配对的和都是模} \(10\) {单位。给四个例外任意编号为} \(a,b,c,d\){；在配对} \((a,b)\) {与} \((a,c)\) {中至少一个是跨类配对。}

{设} \(e_a+e_b\) {与} \(10\) {互素。下面四步是一个统一桥。令}

\[
q=\frac{-5u-a-b}{3},\qquad
r=\frac{-5u+2a+2b}{9},\qquad
s=\frac{-11u-a-b}{27}.
\]

{依次平均}

\[
(u,c,d),\qquad(q,a,b),\qquad(u,q,r),\qquad(u,q,r).
\]

{各步所列的值在当时状态中确实存在，重数变化为}

\[
u^6abcd
\longrightarrow u^5abq^3
\longrightarrow u^5q^2r^3
\longrightarrow u^4qr^2s^3
\longrightarrow u^3rs^6.
\]

{最终状态是} \(\mathcal B_{10}(s,u)\){，其单独坐标为} \(r=-6s-3u\){。参数差满足}

\[
s-u=-\frac{38u+a+b}{27}
=-\frac{40u+e_a+e_b}{27}.
\]

{因为} \(27\) {是模} \(10\) {单位，}\(e_a+e_b\) {是模} \(10\) {单位，所以最终终端族合法。若单位配对是} \((a,c)\){，交换} \(b,c\) {即得同一构造。}

{这证明：}

\begin{quote}
\textbf{{六重桥定理。}} {每个本原零和、}\(G=1\) {的十元六重状态，都在至多四步后进入合法的} \(\mathcal B_{10}\){。}
\end{quote}

\subsection*{\texorpdfstring{{C35.4} {十元二参数核}}{C35.4 十元二参数核}}\label{C35-h236}
\addcontentsline{toc}{subsection}{{C35.4} {十元二参数核}}

{定义}

\[
\mathcal B_{10}(u,v)
=(u,u,u,u,u,u,v,v,v,-6u-3v),
\]

{其中} \((u,v)\) {取本原整数代表。一般公式给出}

\[
G(\mathcal B_{10}(u,v))=\gcd(u-v,10).
\]

{因此合法条件是} \(\gcd(u-v,10)=1\){。}

{取射影坐标}

\[
X=u,\qquad Y=3u+2v.
\]

{每个合法本原状态唯一落入以下两种局部类型之一：}

\begin{enumerate}
\def\labelenumi{\arabic{enumi}.}
\tightlist
\item
  \(X,Y\) {都是奇数且} \(X+Y\equiv0\pmod4\){，尺度} \(h=1\){；}
\item
  \(X,Y\) {奇偶相反，尺度} \(h=2\){。}
\end{enumerate}

{两种类型都要求} \(5\nmid Y\){。本原能量可统一写成}

\[
W(X,Y)=h^2(5X^2+Y^2)=\frac13E(\mathcal B_{10}).
\tag{C35.4.1}
\]

\paragraph*{\texorpdfstring{{4.1} {七个真实返回宏}}{4.1 七个真实返回宏}}\label{C35-h237}
\addcontentsline{toc}{paragraph}{{4.1} {七个真实返回宏}}

{以下七个矩阵作用于列向量} \((X,Y)^T\){。每个矩阵都由一个真实三平均序列实现，并在最后返回同一个} \(6+3+1\) {重数族：}

\[
\begin{aligned}
M_1&=\begin{pmatrix}1&1\\-5&-1\end{pmatrix},&
M_2&=\begin{pmatrix}1&0\\-5&2\end{pmatrix},&
M_3&=\begin{pmatrix}1&-2\\0&3\end{pmatrix},\\
M_4&=\begin{pmatrix}1&1\\5&-3\end{pmatrix},&
M_5&=\begin{pmatrix}4&1\\0&-3\end{pmatrix},&
M_6&=\begin{pmatrix}1&-2\\5&2\end{pmatrix},\\
M_7&=\begin{pmatrix}1&0\\-10&-1\end{pmatrix}.&&
\end{aligned}
\]

{为使宏可以独立复核，令}

\[
U=u,\quad V=v,\quad W=-6u-3v,
\]

{以及}

\[
Q=\frac{-5u-2v}{3},\quad
R=\frac{u+2v}{3},\quad
S=\frac{2u+v}{3},\quad
T=-2u-\frac v3.
\]

{按行依次平均所列三个当前值：}

{ % do not increment counter
\begin{longtable}[]{@{}lrl@{}}
\toprule\noalign{}
{矩阵} & {步数} & {实际平均序列} \\
\midrule\noalign{}
\endhead
\bottomrule\noalign{}
\endlastfoot
\(M_1\) & {3} & \((W,V,U),(V,U,U),(V,U,U)\) \\
\(M_2\) & {5} & \((W,V,U),(Q,U,U),(V,V,U),(Q,R,U),(Q,R,U)\) \\
\(M_3\) & {4} & \((W,U,U),(V,U,U),(V,S,U),(V,S,U)\) \\
\(M_4\) & {3} & \((W,U,U),(V,U,U),(V,U,U)\) \\
\(M_5\) & {4} & \((W,V,U),(V,V,U),(R,U,U),(R,U,U)\) \\
\(M_6\) & {4} & \((V,U,U),(W,S,U),(V,S,U),(V,S,U)\) \\
\(M_7\) & {5} & \((W,V,V),(V,U,U),(T,S,U),(T,S,U),(S,U,U)\) \\
\end{longtable}
}

{核验程序逐步检查每个三元组当时存在，并从最终状态重新读取} \((u',v')\){，得到上列矩阵。七个行列式为}

\[
4,\ 2,\ 3,\ -8,\ -12,\ 12,\ -1.
\]

{它们只含素因子} \(2,3\){。每个矩阵第二行的第一项被} \(5\) {整除、第二项模} \(5\) {非零，所以自动保持} \(5\nmid Y\){。}

\paragraph*{\texorpdfstring{{4.2} {十二种局部签名}}{4.2 十二种局部签名}}\label{C35-h238}
\addcontentsline{toc}{paragraph}{{4.2} {十二种局部签名}}

{对本原合法方向} \(z=(X,Y)\){，令}

\[
d_i(z)=\gcd\bigl((M_iz)_1,(M_iz)_2\bigr).
\]

{因为} \(z\) {本原，}\(d_i\mid\det M_i\){。所以} \(d_i\) {只含} \(2,3\){。约去} \(d_i\) {后的输出若属于合法奇偶类型，记其尺度为} \(h_i'\){；否则记为短横线。}

{模数}

\[
96=2^5\cdot3
\]

{足以同时决定全部} \(d_i\) {与输出模} \(4\) {类型。5,120} {个本原合法局部类恰压缩成下表} {12} {种签名。每个表项是} \((d_i,h_i')\){：}

{ % do not increment counter
\begin{longtable}[]{@{}
  >{\raggedleft\arraybackslash}p{(\linewidth - 18\tabcolsep) * \real{0.1212}}
  >{\raggedleft\arraybackslash}p{(\linewidth - 18\tabcolsep) * \real{0.1212}}
  >{\raggedleft\arraybackslash}p{(\linewidth - 18\tabcolsep) * \real{0.1212}}
  >{\raggedright\arraybackslash}p{(\linewidth - 18\tabcolsep) * \real{0.0909}}
  >{\raggedright\arraybackslash}p{(\linewidth - 18\tabcolsep) * \real{0.0909}}
  >{\raggedright\arraybackslash}p{(\linewidth - 18\tabcolsep) * \real{0.0909}}
  >{\raggedright\arraybackslash}p{(\linewidth - 18\tabcolsep) * \real{0.0909}}
  >{\raggedright\arraybackslash}p{(\linewidth - 18\tabcolsep) * \real{0.0909}}
  >{\raggedright\arraybackslash}p{(\linewidth - 18\tabcolsep) * \real{0.0909}}
  >{\raggedright\arraybackslash}p{(\linewidth - 18\tabcolsep) * \real{0.0909}}@{}}
\toprule\noalign{}
\begin{minipage}[b]{\linewidth}\raggedleft
{类型}
\end{minipage} & \begin{minipage}[b]{\linewidth}\raggedleft
{输入} \(h\)
\end{minipage} & \begin{minipage}[b]{\linewidth}\raggedleft
{类数}
\end{minipage} & \begin{minipage}[b]{\linewidth}\raggedright
\(M_1\)
\end{minipage} & \begin{minipage}[b]{\linewidth}\raggedright
\(M_2\)
\end{minipage} & \begin{minipage}[b]{\linewidth}\raggedright
\(M_3\)
\end{minipage} & \begin{minipage}[b]{\linewidth}\raggedright
\(M_4\)
\end{minipage} & \begin{minipage}[b]{\linewidth}\raggedright
\(M_5\)
\end{minipage} & \begin{minipage}[b]{\linewidth}\raggedright
\(M_6\)
\end{minipage} & \begin{minipage}[b]{\linewidth}\raggedright
\(M_7\)
\end{minipage} \\
\midrule\noalign{}
\endhead
\bottomrule\noalign{}
\endlastfoot
{1} & {1} & {384} & \((4,2)\) & {--} & \((1,1)\) & \((4,1)\) & {--} & {--} & \((1,1)\) \\
{2} & {1} & {384} & \((4,2)\) & {--} & \((1,1)\) & \((8,2)\) & {--} & {--} & \((1,1)\) \\
{3} & {1} & {128} & \((4,2)\) & {--} & \((3,1)\) & \((4,1)\) & {--} & {--} & \((1,1)\) \\
{4} & {1} & {128} & \((4,2)\) & {--} & \((3,1)\) & \((8,2)\) & {--} & {--} & \((1,1)\) \\
{5} & {2} & {768} & \((1,1)\) & \((1,1)\) & \((1,2)\) & {--} & \((2,1)\) & {--} & \((1,2)\) \\
{6} & {2} & {768} & \((1,1)\) & \((1,1)\) & \((1,2)\) & {--} & \((4,2)\) & {--} & \((1,2)\) \\
{7} & {2} & {256} & \((1,1)\) & \((1,1)\) & \((3,2)\) & {--} & \((6,1)\) & {--} & \((1,2)\) \\
{8} & {2} & {256} & \((1,1)\) & \((1,1)\) & \((3,2)\) & {--} & \((12,2)\) & {--} & \((1,2)\) \\
{9} & {2} & {768} & \((1,1)\) & \((2,2)\) & \((1,2)\) & {--} & {--} & \((2,1)\) & \((1,2)\) \\
{10} & {2} & {768} & \((1,1)\) & \((2,2)\) & \((1,2)\) & {--} & {--} & \((4,2)\) & \((1,2)\) \\
{11} & {2} & {256} & \((1,1)\) & \((2,2)\) & \((3,2)\) & {--} & {--} & \((6,1)\) & \((1,2)\) \\
{12} & {2} & {256} & \((1,1)\) & \((2,2)\) & \((3,2)\) & {--} & {--} & \((12,2)\) & \((1,2)\) \\
\end{longtable}
}

\paragraph*{\texorpdfstring{{4.3} {精确实射影覆盖}}{4.3 精确实射影覆盖}}\label{C35-h239}
\addcontentsline{toc}{paragraph}{{4.3} {精确实射影覆盖}}

{固定上表的一种局部签名。若} \(M_i=(a\ b;c\ d)\) {可用，严格下降条件}

\[
h^2(5X^2+Y^2)
>
(h_i')^2\frac{5(aX+bY)^2+(cX+dY)^2}{d_i^2}
\]

{等价于一个齐次整系数二次不等式。令} \(t=Y/X\){，它成为}

\[
D_i(1,t)>0.
\]

{对} {12} {种签名分别取全部可用宏。精确核验程序完成以下有限步骤：}

\begin{enumerate}
\def\labelenumi{\arabic{enumi}.}
\tightlist
\item
  {枚举每个} \(D_i(1,t)\) {的有理根和不可约二次根；}
\item
  {用有理区间隔离全部不同实根；}
\item
  {在每个开区间取有理点检查符号；}
\item
  {单独检查每个根与无穷方向} \(X=0\){；}
\item
  {再核对未覆盖方向是否与当前模} \(96\) {签名相容。}
\end{enumerate}

{结果是：每个合法方向都有某个宏严格降低} \(W\){，唯一可能的相容未覆盖方向是}

\[
X=0.
\]

{此时} \(u=0\){，状态} \(\mathcal B_{10}(0,v)\) {已含零坐标。固定一个零位置，对其余九个位置施行标准} \(3^2\) {三进平均网络即可归零。}

{所以合法} \(\mathcal B_{10}\) {的正整数势能} \(W\) {要么严格下降，要么已经到达可解终点。无限严格下降不可能，故所有合法终端状态均可达零。}

\subsection*{\texorpdfstring{{C35.5} {主定理的合成}}{C35.5 主定理的合成}}\label{C35-h240}
\addcontentsline{toc}{subsection}{{C35.5} {主定理的合成}}

{从任意本原十元} \(G=1\) {状态出发：}

\begin{enumerate}
\def\labelenumi{\arabic{enumi}.}
\tightlist
\item
  {若没有六重值，使用安全下降引理降低本原平方能量；}
\item
  {由于能量是正整数，有限步后到达零或六重状态；}
\item
  {六重桥定理在至多四步后把后者送入合法} \(\mathcal B_{10}\){；}
\item
  {第} {4} {节的有限下降证书把该终端状态归零。}
\end{enumerate}

{这证明了} \(G=1\) {的充分性，并与第} {1} {节的必要性合成十元主定理。}
